% concatenated from Derivation-of-WKE-11.tex


\documentclass[UTF-8,reqno]{amsart}
%\usepackage[utf8]{inputenc}
%\usepackage[T1]{fontenc}
\usepackage{enumerate}
\usepackage{mhequ}
%\linespread{1.2}
\setlength{\topmargin}{-0.3cm}
%\setlength{\oddsidemargin}{0.5cm}\setlength{\evensidemargin}{0.5cm}
%\setlength{\textwidth}{15truecm}
\setlength{\textheight}{21.5truecm}
\usepackage{amsthm,amsmath,amssymb,url,color, booktabs,nccmath}
\usepackage[left=3.2cm,right=3.2cm,top=4cm,bottom=4cm]{geometry}
%\usepackage{showkeys}
\usepackage{mathrsfs}
\usepackage{enumitem,dsfont}
%\usepackage{refcheck}
%\usepackage{appendix}
\usepackage{subfigure}
\usepackage[graphicx]{realboxes}
\usepackage{tikz}


%Graphics Packages 
\usepackage{pinlabel}
\usepackage{graphicx}	
\usepackage{subcaption}
\captionsetup[subfigure]{labelfont=rm}
\usepackage{tikz}
\usepackage{tikz-cd}
%------------------------------------------------------------------------------------------------
\usepackage{color}
\usepackage[colorlinks=true]{hyperref}
\hypersetup{
    %colorlinks=true,       % false: boxed links; true: colored links
    linkcolor=blue,          % color of internal links
    citecolor=red,        % color of links to bibliography
    filecolor=blue,      % color of file links
    urlcolor=cyan
}
\newcommand{\blue}{\color{blue}}
\newcommand{\red}{\color{red}}


\definecolor{darkergreen}{rgb}{0.0, 0.5, 0.0}



%------------------------------------------------------------------------------------------------

\setlength{\parskip}{2pt}

\numberwithin{equation}{section}
\def\theequation{\arabic{section}.\arabic{equation}}
\newcommand{\be}{\begin{eqnarray}}
\newcommand{\ee}{\end{eqnarray}}
\newcommand{\ce}{\begin{eqnarray*}}
\newcommand{\de}{\end{eqnarray*}}
\newtheorem{theorem}{Theorem}[section]
\newtheorem{lemma}[theorem]{Lemma}
\newtheorem{proposition}[theorem]{Proposition}
\newtheorem{Examples}[theorem]{Example}
\newtheorem{corollary}[theorem]{Corollary}
\newtheorem{assumption}[theorem]{Assumption}
\newtheorem{definition}[theorem]{Definition}
\theoremstyle{definition}
\newtheorem{remark}[theorem]{Remark}


\newcommand{\rmb}[1]{\textcolor{blue}{#1}}
\newcommand{\rmm}[1]{\textcolor{magenta}{#1}}
\newcommand{\rmg}[1]{\textcolor[rgb]{0.50,0.25,0.00}{#1}}
\newcommand{\rmk}[1]{\textcolor{red}{#1}}
\newcommand{\VV}{\mathscr{V}}
\newcommand{\LL}{\mathscr{L}}
\newcommand{\UU}{\mathscr{U}}
\def\nor{|\mspace{-1mu}|\mspace{-1mu}|}
\newcommand{\assign}{:=}
\newcommand{\cdummy}{\cdot}
\newcommand{\mathd}{\mathrm{d}}
\newcommand{\tmcolor}[2]{{\color{#1}{#2}}}
\newcommand{\tmmathbf}[1]{\ensuremath{\mathbf{#1}}}
\newcommand{\tmop}[1]{\ensuremath{\operatorname{#1}}}
\newcommand{\tmtextit}[1]{{\itshape{#1}}}



\def\Wick#1{\,\colon\!\! #1 \!\colon}
\DeclareMathOperator{\supp}{supp}
\def\bX{\overline{X}}
\def\bZ{\overline{Z}}
\def\PPhi{\mathbf{\Phi}}
\def\Re{{\mathrm{Re}}}
\def\T{\mathbb{T}^{3}}
\def\Im{{\mathrm{Im}}}
\def\var{{\mathrm{var}}}
\def\HS{{\mathrm{\tiny HS}}}
\def\eps{\varepsilon}
\def\t{\tau}
\def\e{\mathrm{e}}
\def\vr{\varrho}
\def\th{\theta^{i}}
\def\a{\alpha}
\def\om{\omega}
\def\Om{\Omega}
\def\v{\mathrm{v}}
\def\u{\mathbf{u}}
\def\w{\mathrm{w}}
\def\p{\partial}
\def\d{\hh}
\def\g{\gamma}
\def\l{\lambda}
\def\la{\langle}
\def\ra{\rangle}

\def\<{{\langle}}
\def\>{{\rangle}}
\def\bz{{\mathbf{z}}}
\def\by{{\mathbf{y}}}
\def\bx{{\mathbf{x}}}
\def\tr{\mathrm {tr}}
\def\W{{\mathcal W}}
\def\Ric{{\rm Ricci}}
\def\osc{{\rm osc}}
\def\Cap{\mbox{\rm Cap}}
\def\sgn{\mbox{\rm sgn}}
\def\mathcalV{{\mathcal V}}
\def\Law{{\mathord{{\rm Law}}}}
\def\dif{{\mathord{{\rm d}}}}
\def\dis{{\mathord{{\rm \bf d}}}}
\def\Hess{{\mathord{{\rm Hess}}}}
\def\min{{\mathord{{\rm min}}}}
\def\Vol{\mathord{{\rm Vol}}}
\def\bbbn{{\rm I\!N}}
\def\no{\nonumber}
\def\={&\!\!=\!\!&}
\DeclareMathOperator*{\esssup}{esssup}
\newcommand{\norm}[1]{\left\|#1\right\|}
 \newcommand{\eqdef}{\stackrel{\mbox{\tiny def}}{=}}
\def\bB{{\mathbf B}}
\def\bC{{\mathbf C}}

\def\cA{{\mathcal A}}
\def\cB{{\mathcal B}}
\def\cC{{\mathcal C}}
\def\cD{{\mathcal D}}
\def\cE{{\mathcal E}}
\def\cF{{\mathcal F}}
\def\cG{{\mathcal G}}
\def\cH{{\mathcal H}}
\def\cI{{\mathcal I}}
\def\cJ{{\mathcal J}}
\def\cK{{\mathcal K}}
\def\cL{{\mathcal L}}
\def\cM{{\mathcal M}}
\def\cN{{\mathcal N}}
\def\cO{{\mathcal O}}
\def\cP{{\mathcal P}}
\def\cQ{{\mathcal Q}}
\def\cR{{\mathcal R}}
\def\cS{{\mathcal S}}
\def\cT{{\mathcal T}}
\def\cU{{\mathcal U}}
\def\cV{{\mathcal V}}
\def\cW{{\mathcal W}}
\def\cX{{\mathcal X}}
\def\cY{{\mathcal Y}}
\def\cZ{{\mathcal Z}}

\def\mA{{\mathbb A}}
\def\mB{{\mathbb B}}
\def\mC{{\mathbb C}}
\def\mD{{\mathbb D}}
\def\mE{{\mathbb E}}
\def\mF{{\mathbb F}}
\def\mG{{\mathbb G}}
\def\mH{{\mathbb H}}
\def\mI{{\mathbb I}}
\def\mJ{{\mathbb J}}
\def\mK{{\mathbb K}}
\def\mL{{\mathbb L}}
\def\mM{{\mathbb M}}
\def\mN{{\mathbb N}}
\def\mO{{\mathbb O}}
\def\mP{{\mathbb P}}
\def\mQ{{\mathbb Q}}
\def\mR{{\mathbb R}}
\def\mS{{\mathbb S}}
\def\mT{{\mathbb T}}
\def\mU{{\mathbb U}}
\def\mV{{\mathbb V}}
\def\mW{{\mathbb W}}
\def\mX{{\mathbb X}}
\def\mY{{\mathbb Y}}
\def\mZ{{\mathbb Z}}

\def\bB{{\mathbf B}}
\def\bP{{\mathbf P}}
\def\bQ{{\mathbf Q}}
\def\bE{{\mathbf E}}
\def\1{{\mathbf{1}}}

\def\sA{{\mathscr A}}
\def\sB{{\mathscr B}}
\def\sC{{\mathscr C}}
\def\sD{{\mathscr D}}
\def\sE{{\mathscr E}}
\def\sF{{\mathscr F}}
\def\sG{{\mathscr G}}
\def\sH{{\mathscr H}}
\def\sI{{\mathscr I}}
\def\sJ{{\mathscr J}}
\def\sK{{\mathscr K}}
\def\sL{{\mathscr L}}
\def\sM{{\mathscr M}}
\def\sN{{\mathscr N}}
\def\sO{{\mathscr O}}
\def\sP{{\mathscr P}}
\def\sQ{{\mathscr Q}}
\def\sR{{\mathscr R}}
\def\sS{{\mathscr S}}
\def\sT{{\mathscr T}}
\def\sU{{\mathscr U}}
\def\sV{{\mathscr V}}
\def\sW{{\mathscr W}}
\def\sX{{\mathscr X}}
\def\sY{{\mathscr Y}}
\def\sZ{{\mathscr Z}}
\def\E{\mathbf E}


\def\fM{{\mathfrak M}}
\def\fA{{\mathfrak A}}

\def\geq{\geqslant}
\def\leq{\leqslant}
\def\ge{\geqslant}
\def\le{\leqslant}

\def\c{\mathord{{\bf c}}}
\def\div{\mathord{{\rm div}}}
\def\iint{\int\!\!\!\int}
\def\Re{{\mathrm{Re}}}
\def\Im{{\mathrm{Im}}}
\def\var{{\mathrm{var}}}
\def\HS{{\mathrm{\tiny HS}}}
\def\eps{\varepsilon}
\def\t{\tau}
\def\e{\mathrm{e}}
\def\vr{\varrho}
\def\th{\theta^{i}}
\def\a{\alpha}
\def\om{\omega}
\def\Om{\Omega}
\def\v{\mathrm{v}}
\def\u{\mathbf{u}}
\def\w{\mathrm{w}}
\def\p{\partial}
\def\d{\hh}
\def\g{\gamma}
\def\l{\lambda}
\def\la{\langle}
\def\ra{\rangle}

\def\<{{\langle}}
\def\>{{\rangle}}
\def\bz{{\mathbf{z}}}
\def\by{{\mathbf{y}}}
\def\bx{{\mathbf{x}}}
\def\tr{\mathrm {tr}}
\def\W{{\mathcal W}}
\def\Ric{{\rm Ricci}}
\def\osc{{\rm osc}}
\def\Cap{\mbox{\rm Cap}}
\def\sgn{\mbox{\rm sgn}}
\def\mathcalV{{\mathcal V}}
\def\Law{{\mathord{{\rm Law}}}}
\def\dif{{\mathord{{\rm d}}}}
\def\dis{{\mathord{{\rm \bf d}}}}
\def\Hess{{\mathord{{\rm Hess}}}}
\def\min{{\mathord{{\rm min}}}}
\def\Vol{\mathord{{\rm Vol}}}
\def\bbbn{{\rm I\!N}}
\def\no{\nonumber}
\def\={&\!\!=\!\!&}
\def\bt{\begin{theorem}}
\def\et{\end{theorem}}
\def\bl{\begin{lemma}}
\def\el{\end{lemma}}
\def\br{\begin{remark}}
\def\er{\end{remark}}
\def\bx{\begin{Examples}}
\def\ex{\end{Examples}}
\def\bd{\begin{definition}}
\def\ed{\end{definition}}
\def\bp{\begin{proposition}}
\def\ep{\end{proposition}}
\def\bc{\begin{corollary}}
\def\ec{\end{corollary}}

\def\iso{{\bB^\beta\!\!_{\!{\rm iso}}}}
\def\ani{{\bB^{\beta;\alpha}\!\!_{\!\!{\rm ani}}}}
\def\mix{{\bB^{\beta,\gamma}\!\!_{\!{\rm mix}}}}

\def\fM{{\mathfrak M}}
\def\fA{{\mathfrak A}}

\def\geq{\geqslant}
\def\leq{\leqslant}
\def\ge{\geqslant}
\def\le{\leqslant}

\def\c{\mathord{{\bf c}}}
\def\div{\mathord{{\rm div}}}
\def\iint{\int\!\!\!\int}
\def\Id{\textrm{Id}}

\def\bH{{\mathbf H}}
\def\bW{{\mathbf W}}
\def\bP{{\mathbf P}}
\def\bA{{\mathbf A}}\def\si{\phi}
\def\bT{{\mathbf T}} \def\R{\mathbb R}\def\nn{\nabla}
\def\ff{\frac} \def\R{\mathbb R}  \def\ff{\frac} \def\ss{\sqrt} \def\B{\scr B}
\def\N{\mathbb N} \def\kk{ \zeta} \def\m{{\bf m}}
\def\dd{\hh} \def\DD{\Dd} \def\vv{\varepsilon} \def\rr{\rho^{i}}
\def\<{\langle} \def\>{\rangle}
\def\x{{\bf x}}
\def\y{{\bf y}}
\def\dd{{a}}
\def\hh{{h}}
\newcommand{\Prec}{\prec\!\!\!\prec}
\allowdisplaybreaks

\setcounter{tocdepth}{1}


\begin{document}

\title[The Wave  Kinetic Theory for Quasi MMT]{The Wave  Kinetic Theory for Quasilinear MMT %Majda-McLaughlin-Tabak
Equation}

\author{Huaxiang L\"u}
\address[H. L\"u]{Academy of Mathematics and Systems Science, Chinese Academy of Sciences, Beijing 100190, China, and Fakult\"at f\"ur Mathematik, Bielefeld Universit\"at, D 33615 Bielefeld, Germany}\email{lvhuaxiang22@mails.ucas.ac.cn }

 

\thanks{%Research  supported   by National Key R\&D Program of China (No. 2022YFA1006300, 2020YFA0712700) and the NSFC (No. 12426205, 12090014, 12288201). The financial support by the DFG through the CRC 1283 ``Taming uncertainty and profiting from randomness and low regularity in analysis, stochastics and their applications'' is greatly acknowledged.
Funded by the Deutsche Forschungsgemeinschaft (DFG,German Research Foundation)-Project-ID 317210226-SFB 1283"
}


\begin{abstract}

We study the one-dimensional quasilinear Majda--McLaughlin--Tabak (MMT) equation on a large torus $[0,L]$:
\begin{align*}
    i \partial_t u +2\pi|\nabla|^\sigma u +\lambda^{2}|\nabla|^\beta\left[ \left||\nabla|^\beta u\right|^{2} |\nabla|^\beta u\right]=0.
\end{align*}
Our focus is on the well-posedness of its dynamics and the emergence of kinetic behavior where the domain size $L$ tends to infinity and the nonlinearity $\alpha=\lambda^2L^{-1}$ vanishes. In contrast to semilinear dispersive models, the quasilinear structure leads to unavoidable derivative loss, which prevents the construction of solutions via iteration of the Duhamel formula.  


Our results exhibit a dichotomy depending on the dispersion exponent $\sigma$. For $\sigma\in(1,2]$, we prove that, with high probability, solutions exist up to time scales $T_0 \sim \alpha^{-\frac54+} \wedge \alpha^{-\frac1{1-\beta}+}$, and that only trivial resonances occur, leading to a degenerate wave kinetic equation. For $\sigma\in(0,1)$, we prove the existence up to time scales $T_0 \sim \alpha^{-1-}$ and show that the second-order statistics are well approximated by the wave kinetic equation. In both cases, the solutions remain smooth while exhibiting smallness in suitable $L^\infty$-based norms despite having large total energy.

The proof proceeds in two main steps. First, we establish the
propagation of randomness for a suitably truncated equation, which
allows us to overcome the derivative loss and recover the kinetic
description. Then, we perform deterministic high-order energy estimates
and a bootstrap argument to extend the solution up to time $T_0$. The
main difficulty  lies in controlling the high-frequency
energy growth for a genuinely quasilinear equation with a general
 dispersion relation $|\xi|^\sigma$. In particular, different
regimes of $\sigma$ lead to fundamentally different four-wave resonance
structures, requiring separate treatments of trivial and nontrivial
resonant interactions. 

\end{abstract}


%\subjclass[2010]{60H15; 35R60; 35Q30}
\keywords{Majda–McLaughlin–Tabak equation, wave kinetic equation, weak turbulence, resonant interactions, Feynman diagram}

\date{\today}

\maketitle


\tableofcontents 


 \section{Introduction}
 
 
The kinetic theory of nonlinear wave systems serves as the formal foundation for the nonequilibrium statistical physics of such systems, and can be viewed as a natural extension of Boltzmann’s kinetic framework, originally introduced in the setting of interacting particle systems. The central objective of this theory is to extract macroscopic statistical dynamics from microscopic wave-wave interactions. To describe the effective energy spectrum, physicists formulated the Wave Kinetic Equation (WKE), whose role closely parallels that of the Boltzmann equation in classical particle kinetic theory, acting as a canonical bridge between microscopic interaction mechanisms and macroscopic kinetic-scale statistical descriptions. Then, Zakharov and his collaborators discovered a close connection to  hydrodynamic
 turbulence, whence the names “wave turbulence theory” and "Kolmogorov-Zakharov spectra". 

From a mathematical perspective, a rigorous derivation of the wave kinetic equation starting from dispersive equations constitutes a fundamental problem, and can be viewed as a  generalized form of Hilbert’s Sixth Problem for  statistical theory of wave systems.
Such a rigorous mathematical derivation has been developed in recent years. Through the precise mathematical analysis of Feynman diagram expansions, major breakthroughs have been achieved in the rigorous derivation of the WKE for NLS, or other semilinear models. We refer to Section \ref{sec:liter} for further relevant literature. However, for a large class of quasilinear equations, the situation is totally different, as solutions cannot be constructed through iterative applications of the Duhamel formula due to the unavoidable derivative loss.  In particular, the solution cannot be represented through iterative applications of  Feynman trees, which constitutes the core analytical strategy for semilinear equations.  This difficulty was recently overcome in the quasilinear setting by
Deng, Ionescu, and Pusateri \cite{DIP25a,DIP25} for the two-dimensional
gravity water-wave system, whose dispersion relation
corresponds to the case $\sigma=\frac12$.

Among the quasilinear models, a notable example is the toy model introduced by Majda, McLaughlin, and Tabak to conduct  numerical investigations of wave turbulence \cite{MMT97}. This equation is analytically tractable and straightforward to simulate, yet it still displays an exceptionally rich and diverse range of phenomena.
The universality of the weak turbulence scenario was called into question after studies of the MMT model revealed spectral behaviors that, in some cases, deviate from the classical Kolmogorov-Zakharov spectra predicted by wave turbulence theory. The introduction of the MMT model thus marked an important step toward a deeper understanding of wave turbulence. Nevertheless, a rigorous derivation of the quasilinear MMT equation remains an open problem.  

In this paper, we study the wave kinetic theory of the quasilinear MMT equation, which reads as:
  \begin{align*}
 \tag{MMT} \label{MMT}
\begin{cases}
i \partial_t u +2\pi|\nabla|^\sigma u \pm \lambda^{2}|\nabla|^\beta\left[ \left||\nabla|^\beta u\right|^{2} |\nabla|^\beta u\right]=0,  \quad x\in \mathbb{T}_L := [0,L],   \\
u(0,x) = u_{\textrm{in}}(x).
\end{cases}
\end{align*}
Here the sign of the nonlinearity has no effect on the kinetic description and we consider the defocusing case in this paper.
  The parameter $\sigma\in(0,1)\cup (1,2]$ controls the dispersion.  The parameter $\beta\in(0,1)$ quantifies the "degree of quasilinearity" of the equation. The restriction $0<\beta<1$ guarantees that the equation is indeed quasilinear. The case $\sigma=2,\beta=0$ corresponds to the cubic  Schr\"odinger equation (NLS).
 We rescale   the operators     $|\nabla|^\sigma$ and $|\nabla|^\beta$   into
\begin{align*} 2\pi|\nabla|^\sigma e^{ 2\pi ikx} =2\pi |k|^\sigma  e^{2\pi ik x},\ \ \
 |\nabla|^\beta e^{2\pi ikx} = |k|^\beta  e^{2\pi ik x}.
\end{align*} 

     The corresponding  wave kinetic equation reads as:
\begin{align*}
\tag{WKE}\label{WKE}
\begin{split}
\partial_t n(t, \xi) =&\mathcal K\left(n(t,\cdot)\right)(\xi), \\
\mathcal K(\phi)(\xi):=& \int_{\substack{(\xi_1, \xi_2, \xi_3)\in \R^{3}\\\xi_1-\xi_2+\xi_3=\xi,\\ |\xi_1|^\sigma-|\xi_2|^\sigma+|\xi_3|^\sigma=|\xi|^\sigma}}|\xi|^{2\beta}|\xi_1|^{2\beta}|\xi_2|^{2\beta}|\xi_3|^{2\beta} \phi \phi_1 \phi_2 \phi_3\left(\frac{1}{\phi_1}-\frac{1}{\phi_2}+\frac{1}{\phi_3}-\frac{1}{\phi}\right)\, \dif \xi_1 \dif \xi_2 \dif \xi_3,
\end{split}
\end{align*}
where we used the shorthand notations $\phi:=\phi(\xi),$ and $\phi_j:=\phi(\xi_j)$ for $i=1,2,3$. We note that \eqref{WKE} does not depend upon the sign of nonlinearity in \eqref{MMT}. In particular, as discussed in \cite{MMT97}, in the case $\sigma\in(1,2]$, the collision operator on the right hand of \eqref{WKE} is trivial, since the   dispersion relation $$\xi_1-\xi_2+\xi_3=\xi,\ \ |\xi_1|^\sigma-|\xi_2|^\sigma+|\xi_3|^\sigma=|\xi|^\sigma$$ only admits trivial solution  $\{\xi,\xi_2\}=\{\xi_1,\xi_3\}$.  In any case, the  main part of the integrand equals to 0.  The question now becomes what
 is the appropriate kinetic theory. However, in the case $\sigma \in (0,1)$, the resonant four-wave interactions admit additional nontrivial solutions. For instance, for the water-wave dispersion law $\sigma = \tfrac{1}{2}$ as in \cite{DIP25a,DIP25}, there exist the so-called Benjamin–Feir resonances (taking symmetry into account):
\begin{align*}
( \xi,\xi_1,\xi_2,\xi_3)=  \left( \frac{(a+b+\sqrt{ab})^2}{(\sqrt{a}+\sqrt{b})^2},\ a,\  -\frac{ab}{(\sqrt{a}+\sqrt{b})^2},\ b\right).
\end{align*}
However, for a general dispersive relation, the resonance structure usually does not admit an explicit characterization.


The kinetic  description  is  akin to a law of large numbers,  therefore
we start with a random distribution of the initial data. The conjecture   states that  in the limit of large torus size $L$ and vanishing strength of the nonlinearity 
\begin{align}
  \alpha:=\lambda^2L^{-1},\label{def:alpha}  
\end{align}
the effective dynamics of the Fourier-space mass density $\mathbf E |\widehat u(t, k)|^2$ ($k \in \mathbb{Z}_L:=L^{-1}\mathbb{Z}$) is well approximated  by
\begin{equation}\label{approximation}
\mathbf E |\widehat u(t, k)|^2 \approx n\left(\frac{t}{T_{\mathrm{kin}}}, k\right),
\end{equation}
 where $T_{\mathrm{kin}}\sim \alpha^{-2}=\frac{L^{2}}{\lambda^4}$ is the  kinetic timescale. While the discussion and derivation under the quasi-Gaussian random phase hypothesis are established in \cite{MMT97}, a rigorous derivation remains an open problem.



\subsection{Statement of the results}  

The spatial Fourier series of a function $u: \mathbb{T}_L \to \mathbb C$ is defined on $\mathbb{Z}_L:=L^{-1}\mathbb{Z}$ by
\begin{equation*}
\widehat{u}(k)=\int_{\mathbb{T}_L} u(x) e^{-2\pi i kx}\dif x,\quad \mathrm{\; such \ that \;}\quad u(x)=\frac{1}{L}\sum_{k \in \mathbb{Z}_L} \widehat{u}(k) \,e^{2\pi i kx}. 
\end{equation*}
 To match the kinetic approximation at time $t=0$, we consider 
well-prepared random homogeneous initial data of the form
 \begin{equation}\label{wellprepared}
 u_{\rm in}(x)=\frac{1}{L}\sum_{k \in \mathbb{Z}_L}\sqrt{n_{\textrm{in}}(k)} \, \eta_{k}(\omega)\,e^{2\pi i kx},
\end{equation}
where $n_{\mathrm{in}}$ is a Schwartz function,
and $\{\eta_k\}_{k\in\mathbb{Z}_L}$ are independent, identically distributed 
complex Gaussian random variables.

Our goal in this paper is to rigorously investigate wave turbulence
and to derive the wave kinetic equation  for the MMT model. Our main results are twofold. 
First, we establish long-time regularity for solutions to the MMT equation
on a large torus $\mathbb{T}_L$,  where the total energy is large,
but the local energy remains small.  Second, we show that as $L \to \infty$ and $\alpha := L^{-\gamma} \to 0$
for some small $\gamma \in (0,1)$, the second moment of the solution
$\mathbf{E}|\widehat{u}(t,k)|^2$ can be approximated by the power-series
expansion associated with the wave kinetic equation.   
\begin{theorem} \label{thm1}
Let $\sigma \in (0,1)\cup(1,2]$ and $\beta \in \left(0,\frac{\sigma}{4}\right)$. 
Assume that $\alpha$ satisfies the scaling law $\alpha = L^{-\gamma}$ for some $\gamma \in (0,1)$. 
Let $n_{\mathrm{in}} \in \mathcal{S}(\mathbb{R})$ be a Schwartz function, and let 
$\{\eta_k(\omega)\}_{k\in\mathbb{Z}_L}$ be i.i.d. complex normalized Gaussian random variables. Let $0<\delta \ll 1$, and define
 \begin{align}\label{def:T}
 T:= \begin{cases}
      L^{-10\delta}\min\{L,L^{\frac{(1+\gamma)(1+2\beta)}{2}} ,L^{\frac{5\gamma}{4}},L^{\frac{\gamma}{1-\beta}}\},\ &{\rm if}\ \sigma\in(1,2]; \\
       L^{-10\delta}\min\{ L^{\frac{1+2\beta}{2+2\beta-\sigma}},L^{\frac{5\gamma}{4}},L^{\frac{\gamma}{1-\beta}}\},\ &{\rm if}\ \sigma\in(0,1).
 \end{cases}
 \end{align} 
Let  $N_0:= 10 + \Bigl\lfloor \frac{100\beta^2}{\delta^2} \Bigr\rfloor$, and let
\begin{align*}
    T_0:=\begin{cases}
        T,\ \ &{\rm if}\  \sigma\in  (1,2];\\
        \min\{T,  L^{\gamma+\theta}\},\ \ &{\rm if}\  \sigma\in (0,1),
    \end{cases}
\end{align*}
 where $0<\theta\ll1$ is another constant depending on $N_0$.  

Then, with probability $\geq 1 - L^{-40}$,   the equation \eqref{MMT} admits a smooth solution on $[0,T_0]$ satisfying
\begin{align*}
  \sup_{t\in[0,T_0]} \|u(t)\|_{H^{N_0}} &\leq L^{\theta},  \ \
  \sup_{t\in[0,T_0]} \sum_{0\leq i\leq 2} \|\partial_x^i u(t)\|_{L^\infty} \leq L^{-1/2+\theta}.
\end{align*}

 Moreover, we have 
\begin{equation}\label{app}
\mathbf E |\widehat u(t, k)|^2 =n_{\mathrm{in}}(k)+\frac{t}{T_{\mathrm{kin}}}\mathcal K(n_{\mathrm{in}})(k)
+o_{l^\infty_k}\left(\frac{t}{T_{\mathrm {kin}}}\right)_{L \to \infty}
\end{equation}
for all $L^{0+} \leq t \leq T_0$, where $T_{\mathrm{kin}}=\alpha^{-2}/2$, and $o_{l^\infty_k}\left(\frac{t}{T_{\mathrm {kin}}}\right)_{L \to \infty}$ is a quantity that is bounded in $l^\infty_k$ by $L^{-\delta'}\cdot \frac{t}{T_{\mathrm {kin}}}$ for some $\delta'>0$.

\end{theorem}

 

\begin{remark}
\begin{enumerate}
   
    \item  In Theorem~\ref{thm1}, we also establish long-time regularity for solutions to \eqref{MMT} with high probability. 
The expectation $\mathbf{E}$ in \eqref{app} is taken over the event where a smooth solution exists on $[0,T_0]$, 
which has probability at least $1 - L^{-40}$. 
Outside this event, the corresponding quantity is defined to be zero.
    \item As in the quasilinear water-wave setting \cite{DIP25a,DIP25}, 
our solutions are assumed to be small in certain $L^\infty$-based norms, 
while allowing for large total energy as the size of the torus increases. 
The smallness of the $L^\infty$ norm also reflects the smallness of local energy per unit volume. 
This feature plays an essential role in the energy estimates and in controlling the quasilinear interactions over long time intervals. 
  
    \item The well-/ill-posedness theory for \eqref{MMT} itself remains a difficult and interesting problem. 
In the semilinear case $\beta=0$, this question has been extensively studied in the literature; 
see, for instance, \cite{CHHO13,CHHO14,CHKL14,CHKL15,HS15,BLLZ23} and the references therein. 
In contrast, the quasilinear regime $\beta>0$ is substantially more challenging. 
Indeed, the derivative-type nonlinear interactions amplify high-frequency components and lead to a derivative-loss mechanism, making the equation significantly harder to analyze.  
We refer to the recent work \cite{PPW26} for progress on the well-/ill-posedness theory in the quasilinear setting with $\sigma\in(1,2]$. 
    \item  As mentioned earlier, in the case $\sigma \in (1,2]$, the first-order term $\mathcal{K}(n_{\mathrm{in}})$ vanishes. In particular, we show that there is no nontrivial evolution of the second moment up to time $T_0$. This reflects the strong non-resonant structure of the equation in this regime. We also note that even the local well-posedness of \eqref{WKE} remains a highly nontrivial problem, especially in the regime $\beta > 0$, where the collision kernel becomes singular near low frequencies. Recent progress on the well-/ill-posedness theory of the WKE can be found in \cite{GIK20,GLZ25,AL25}. Our analysis does not rely on the well-posedness theory of the WKE itself; instead, we establish convergence only at the level of the first-order kinetic expansion.
 
\end{enumerate}
\end{remark}



\br We briefly comment on the time scale obtained in the present work. 
In principle, for cubic problems one expects to control solutions up to the kinetic time scale 
$T_{\mathrm{kin}} \sim \alpha^{-2}$, where wave turbulence phenomena emerge and the WKE becomes relevant. 
In the NLS case with $d \geq 2$, this was achieved in the series of works \cite{DH23a,DH23b,DH23c}. 
In the semilinear case $\beta = 0$, Vassilev \cite{Vas24} obtained a lifespan up to 
$T \sim \alpha^{-5/4+}$ in one dimension. 
However, in our setting, the one-dimensional nature of the problem combined with its quasilinear structure 
makes the analysis significantly more difficult. 
For $\sigma \in (1,2]$, we obtain a lifespan $T _0\sim \alpha^{-5/4+} \wedge \alpha^{-1/(1-\beta)+},$
while for $\sigma \in (0,1)$, we obtain $T_0 \sim \alpha^{-1-}$,
due to limitations of the energy estimates. 
We refer to Section~\ref{sec:ideaproof} for further discussion. Compared with \cite{Vas24}, our analysis requires the introduction of a renormalization procedure in order to eliminate certain resonant interactions. 
As a consequence, the resonant terms appearing in the counting argument acquire additional lower-order correction terms. 
This modification leads to the extra restriction $\alpha T^{1-\beta}<1,$ which is needed to ensure that the renormalized phase corrections remain perturbative in the counting estimates.  

Deng, Ionescu, and Pusateri \cite{DIP25} treated the two-dimensional gravity water-wave system, whose interface is one-dimensional, and established the existence of random solutions up to times of order $T \sim\alpha^{-4/3+}$. More recently, Vassilev and Wu \cite{VW26} obtained the analogous scale $ \alpha^{-4/3+}$ for the full $\beta$-FPUT system by means of a refined diagrammatic and counting procedure.
By contrast, the exponent $\alpha^{-5/4+}$ obtained here for $\sigma\in(1,2]$ results from our direct use of the molecule algorithm developed in \cite{Vas24}. We expect that, by replacing this step with the refined strategy developed in \cite{DIP25,VW26}, the propagation-of-randomness argument could also be extended to the scale $\alpha^{-4/3+}$ at least in the regime $\sigma\in(1,2]$. We do not pursue this optimization here, the main purpose of the present work is to address the the additional difficulties created by the general dispersion relation $|\xi|^\sigma$.  For $\sigma\in(0,1)$, moreover, the final lifespan is already constrained by the deterministic high-frequency energy estimate, so an improvement of the combinatorial counting alone would not extend the time scale in the main theorem.
\er
  
The derivation of the wave kinetic equation has be achieved for the cubic NLS \cite{DH23b,DH23c} and the MMT  with no derivatives in the nonlinearity  \cite{Vas24}.  
 However, under the circumstances we are considering,  the MMT equations \eqref{MMT} are quasilinear and solutions cannot
 be constructed by iteration of the Duhamel formula due to unavoidable derivative loss. Following the method  developed for 2D gravity water waves \cite{DIP25a,DIP25}, we establish our main result, Theorem \ref{thm1}, via a two-step strategy: propagation of randomness, and energy estimates.

 \textbf{Propagation of randomness.}
We begin by introducing a Fourier truncation of \eqref{MMT}. 
More precisely, we define the truncation operator
\[
\widehat{P_{\leq K}f}
=
\varphi_{\leq K}\widehat{f},
\]
where the cutoff function $\varphi_{\leq K}(x)=\varphi(x/2^K)$ is defined in Section~\ref{sec:pre}. 
The truncation parameter $K$ is chosen as
\begin{align}
    K:=\Bigl[\log_2\bigl(L^{\frac{\delta}{10\beta}}\bigr)\Bigr],
    \label{def:K}
\end{align}
where $\delta>0$ is the small constant appearing in Theorem~\ref{thm1}, 
and $[\cdot]$ denotes the floor function.

We then consider the truncated MMT equation on the time interval $[0,T]$:
\begin{align*}
\tag{Tr-MMT}\label{TrMMT}
\begin{cases}
i \partial_t u_{\rm tr}
+2\pi |\nabla|^\sigma u_{\rm tr}
+\lambda^{2}P_{\leq K}
|\nabla|^\beta
\Bigl[
\bigl||\nabla|^\beta u_{\rm tr}\bigr|^{2}
|\nabla|^\beta u_{\rm tr}
\Bigr]
=0,\\
u_{\rm tr}(0,x)=P_{\leq K}u_{\rm in}(x).
\end{cases}
\end{align*}

As a first step, we study the truncated equation, for which the iterated Duhamel expansion 
(and the associated Feynman tree representation) can be applied. 
Our goal is to establish the well-posedness of the solution, 
and to prove the convergence of solutions to the truncated MMT equation 
toward a corresponding truncated wave kinetic equation.


\begin{theorem}\label{thm2}
Let $\sigma\in (0,1)\cup (1,2],\beta\in(0,\frac\sigma4)$. We assume  $\alpha$ has the scaling law $\alpha=L^{-\gamma}$ with $\gamma\in(0,1)$. Let $n_{\mathrm{in}} \in \mathcal S$ be a Schwartz function, and $\eta_{k}(\omega)$ are i.i.d. complex normalized   Gaussian random variables. We assume $0<\delta\ll1$, and   define $T$ as in \eqref{def:T}. 
Let the truncation parameter $K$ be defined as in \eqref{def:K}.  Then, with probability $\geq 1 - L^{-40}$, the truncated MMT equation \eqref{TrMMT}  has a smooth function on $[0,T]$, and
\begin{equation}\label{approx2}
\mathbf E |\widehat u_{\rm tr}(t, k)|^2 =\varphi_{\leq K}^2(k)n_{\mathrm{in}}(k)+\frac{t}{T_{\mathrm{kin}}}\mathcal K_{\rm tr}(n_{\mathrm{in}})(k)+o_{l^\infty_k}\left(\frac{t}{T_{\mathrm {kin}}}\right)_{L \to \infty}
\end{equation}
for all $L^{0+} \leq t \leq T$, where $T_{\mathrm{kin}}=\alpha^{-2}/2$. 
Here, the truncated collision operator $\mathcal K_{\rm tr}$ is defined as 
\begin{align}
    \mathcal K_{\rm tr}(\phi)(\xi):=& \int_{\mathcal{E}}\varphi_{\leq K}^2(\xi_1)\varphi_{\leq K}^2(\xi_2)\varphi_{\leq K}^2(\xi_3)\varphi_{\leq K}^2(\xi)|\xi|^{2\beta}|\xi_1|^{2\beta}|\xi_2|^{2\beta}|\xi_3|^{2\beta}\notag\\
    &\quad\quad\quad\quad\times \phi \phi_1 \phi_2 \phi_3\left(\frac{1}{\phi_1}-\frac{1}{\phi_2}+\frac{1}{\phi_3}-\frac{1}{\phi}\right)\, \dif \xi_1 \dif \xi_2 \dif \xi_3,\label{bd:ktr}
\end{align}
where the integration domain  is defined as \begin{align}
    \mathcal{E}:=\{(\xi_1, \xi_2, \xi_3)\in \R^{3}:&\ \xi_1-\xi_2+\xi_3=\xi,\notag\\ &\ \quad\quad |\xi_1|^\sigma+\Gamma_0(\xi_1)-|\xi_2|^\sigma-\Gamma_0(\xi_2)+|\xi_3|^\sigma+\Gamma_0(\xi_3)=|\xi|^\sigma+\Gamma_0(\xi)\},\notag%\label{def:domain:trknie}
\end{align}
where $\Gamma_0(\xi)$ will be defined in \eqref{def:gamma1} below satisfying $$\sup_{\xi\in\mathbb{R}}|\Gamma_0(\xi)|\lesssim   L^{-\delta}.$$% $$ \Gamma_0(k)= C_K \cdot \alpha \varphi_{\leq K}(k)|k|^{2\beta}$$  for some constant $C_K>0$ satisfying $\sup_{K}C_K\lesssim \|n_{\mathrm{in}}(x)|x|^{2\beta}\|_{L^1}$.
\end{theorem}

Compared with the original collision kernel in \eqref{WKE}, 
the truncated collision kernel differs in two aspects. 
First, it contains the truncation factors $\varphi_{\leq K}$, 
which arise from the Fourier truncation procedure. 
Second, the integration domain contains the additional phase correction term $\Gamma_0$, 
which originates from the renormalization of trivial pairs.  


\textbf{Energy estimates.} In the second step, we estimate the energy increment of higher-order derivatives of the solution. 
The goal is to establish a deterministic bound showing that the growth of the energy can be controlled by an $L^\infty$-based norm.
\begin{theorem}\label{thm:energy inequality}
    Let $\sigma\in(0,1)\cup (1,2]$ and $\beta\in(0,\frac{\sigma}{4})$. Let $N_0\geq10$ and $T_1<T_2\in\mathbb{R}$. Assume that $u$ is a solution to \eqref{MMT} satisfying for $t'\in[T_1,T_2]$ \begin{align}
   \|u(t')\|_{H^{N_0}}\leq A,\ \  \|u(t')\|_{ {W}^{2,0}}\leq \epsilon(t'),\label{bd:thm9.1}
\end{align}
where $A>1$, and $\epsilon:[T_1,T_2]\to[0,\infty)$ is a continuous function with $\epsilon(t)\ll \lambda^{-1}$.
Then   for any $t_1\leq t_2\in[T_1,T_2]$, it holds that
\begin{align*}
 \|u(t_2)\|_{H^{N_0}}^2\lesssim  \|u(t_1)\|_{H^{N_0}}^2 +\lambda^4A^2\int_{t_1}^{t_2}\epsilon^4(s)\dif s+1_{\{\sigma\in(0,1)\}}\lambda^2A^{2-\frac{1}{4N_0}}\int_{t_1}^{t_2}\epsilon^{2+\frac{1}{4N_0}}(s)\dif s.%+1_{\{\sigma\in(0,1)\}}\lambda^2A^2\int_{t_1}^{t_2}\epsilon^2(s)\dif s.
\end{align*}
\end{theorem}
We refer to Section~\ref{sec:pre} for the definition of the $L^\infty$-based space $W^{2,0}$. Compared with \cite{DIP25a}, which focused on the gravity
water-wave dispersion relation $|\xi|^{1/2}$, we establish high-frequency
energy estimates for general power-law dispersion relations
$|\xi|^\sigma$.

Subsequently, we carry out the main bootstrap argument to establish the existence of solutions to \eqref{MMT} up to time $T_0$, and to prove that each truncated term in \eqref{approx2} converges to the corresponding term in \eqref{app}. This completes the proof of Theorem~\ref{thm1}.

\subsection{Further relevant literature}\label{sec:liter}
The study of wave kinetic theory can be traced back to Peierls’ pioneering work \cite{Pei29} on anharmonic crystals. In that study, he first introduced  the phonon Boltzmann equation, which established a statistical physics framework and later inspired extensive generalizations and developments (see subsequent works in \cite{Spo06,Spo08}). Shortly afterwards, Nordheim \cite{Nor28} proposed the quantum Boltzmann equation to describe the dynamics of quantum interacting gases (see subsequent works in \cite{UU33,Hug83,ESY03}). Then, this framework had been extended to a wide class of wave systems \cite{Dav72,ZS67,Has62,BN69}. The MMT equation is introduced by Majda, McLaughlin and Tabak \cite{MMT97}  as a model for assessing the validity of weak turbulence theory for random waves 
in an unambiguous and transparent fashion.  Under certain experimental settings, their observations appear to contradict the classical theory, revealing a new spectral behavior known as the so-called “MMT spectrum” \cite{CMMT99,CMMT01}.  For further development, we refer to \cite{ZDP04,ZGPD01, DB23}.

The main difficulty lies in the rigorous justification of the convergence of the Feynman diagram expansions. In the linear case, it is firstly derived by   Spohn \cite{Spo77} for short times, and then by  Erd\"os, Salmhofer and Yau \cite{ESY03b,EY00} for a global result.  
Subsequently, the rigorous derivation of this problem has seen substantial progress, driven primarily by the work of  \cite{BGHS18,BGHS21,CG20,CG25,DH21,DH23a,DH23b}. In particular, Deng and Hani \cite{DH23c}  extended the  justification of NLS to
 arbitrarily long times that cover the full life span of the WKE. It is worth mentioning that this approach has already been widely applied to other dispersive equations, such as the Wick  NLS \cite{HSZ24} and the NLS with additive stochastic forcing \cite{GH24}. We also refer to \cite{ST21,HRST22} for the wave turbulence theory for   KdV
 type equations.
However, in one spatial dimension, the derivation of the WKE encounters new obstacles due to the difficulty of  divergent contributions in the Feynman expansion. Notably, even at the subcritical scale $T \lesssim   \alpha^{-2}$, \cite{DIP25} identified explicit pairs of coupled Feynman diagrams that yield divergent contributions which appear to lack mutual cancellation. This divergence is an pure one-dimensional phenomenon. 
 Despite these challenges, significant progress has been made in the rigorous derivation of kinetic and long-time regularity results. In particular, Jin and  Ma~\cite{JM21} first established long-time dynamics for NLS via normal form analysis in the regime $ \gamma > 1 $. For the MMT equation with $ \beta = 0 $, Vassilev~\cite{Vas24} derived the WKE up to timescales $ T \sim \alpha^{-5/4} $.
 
However, for quasilinear equations, there are only few results on rigorous justification, due to the unavoidable
derivative loss. Recently, Deng, Ionescu and Pusateri~\cite{DIP25a,DIP25} initiated the study of long-time regularity for quasilinear water-wave systems. Wu \cite{Wu25} and Vassilev--Wu \cite{VW26} established the WKE for the $\beta$-FPUT system, first for a reduced evolution equation and subsequently for the full system, using refined diagrammatic expansions and counting arguments.
 

\subsection{Ideas of proof}\label{sec:ideaproof}
In what follows, we present an overview of the main steps of the proof, identify the new challenges that arise in this context, and describe the techniques developed to overcome them.

First, for the propagation of randomness of the truncated equation, by using the truncation function to control the  derivative, we could write the solution as a sum of Feynman trees using the iterated Duhamel formula:
\begin{align*} 
     c_k(t)=c^{\leq N}_k (t)+\mathcal{R}^{N+1}_k,
\end{align*}
where $  c^{\leq N}_k (t) =\sum_{i=1}^N c^{i}_k (t)$ are the sum of the ternary trees of order no more than $N$, and $\mathcal{R}^{N+1}_k$ is the higher order remainder term. The constant $N$ is chosen large enough but independent of $L$. The main components and  ideas of the  analysis are the following:
\begin{enumerate}
   \item  \textbf{Renormalization.} 
To formulate the expansion in terms of ternary trees, it is necessary to remove the nonlinear resonant contributions, which would otherwise lead to unfavorable counting estimates (see Figure~\ref{fig:1loop}). 
In the semilinear MMT case ($\beta=0$), these resonances can be eliminated by introducing a linear oscillatory phase depending on the conserved mass of the solution. 
In the present quasilinear setting, however, such a simple renormalization is no longer sufficient. 
Instead, we introduce an implicit deterministic renormalization factor depending on both time and frequency in order to cancel part of the resonant interactions. 
This renormalization plays a crucial role throughout the counting argument and in the derivation of the effective kinetic dynamics.

\item \textbf{Counting estimates with renormalization.} 
In the two-/three-vector counting estimates, the renormalization of trivial couples introduces additional phase correction terms, which must be incorporated into the counting analysis. 
These corrections contain factors of the form $|k|^\beta$, whose singular behavior near $k=0$ creates substantial difficulties. 
As a consequence, the low-frequency regime requires a separate treatment. 
This modification leads to the additional restriction
\[
\alpha T^{1-\beta}<1,
\]
which ensures that the renormalized phase corrections remain perturbative within the counting estimates.  This also shows that the lifespan of solutions becomes shorter as $\beta$ decreases, which further highlights the fundamental difference between the quasilinear and semilinear regimes.
  \item \textbf{Algorithm.} 
In this paper, we reduce the moment estimates for the coefficients $c_k^i$ to a combinatorial counting problem on molecules. 
To carry out this analysis, we apply the algorithm introduced in \cite{Vas24}. 
A new feature in the present setting is the presence of resonant (degenerate) nodes arising from the quasilinear interactions. 
However, in the three-vector counting step, such resonant nodes contribute only the trivial bound $L$, which is acceptable since
\[
L <L^2 T^{-1}
\]
under the condition $T<L$. Therefore, degenerate nodes do not introduce significant additional difficulties into the counting procedure. 
Moreover, the renormalization procedure eliminates the potentially problematic counting configurations. 
In particular, the bad loop-type structures responsible for unfavorable counting estimates are ruled out after renormalization.  

 \item  \textbf{Convergence of low-order ternary trees.} 
In the proof of Theorem~\ref{thm2}, the higher-order ternary trees are treated as error terms, 
so it suffices to analyze the contributions of couples of orders $0$, $1$, and $2$. 
For these low-order terms, we compute explicit formulas and identify the cancellations arising from the renormalized resonant interactions. 
Based on these cancellations, we establish the asymptotic expansion stated in Proposition~\ref{asymptotic}. 
A new difficulty in the quasilinear setting is that the resulting integrands are no longer smooth and exhibit singular behavior near $k_i=0$. 
To overcome this issue, we introduce an additional truncation away from the low-frequency region and apply a stationary phase argument. 
Consequently, the regime where some $k_i$ are close to zero requires a separate analysis.
\end{enumerate}

To control the high-frequency component, we establish deterministic estimates for the energy increment of the solution. 
In the quasilinear setting, this appears to be the only available mechanism for propagating high-regularity bounds over long time intervals. 
Due to the cubic nonlinearity, one may initially expect an energy estimate of quartic type:
\begin{align*}
 \|u(t_2)\|_{H^{N_0}}^2
 \lesssim
 \|u(t_1)\|_{H^{N_0}}^2
 +\lambda^2A^2\int_{t_1}^{t_2}\epsilon^2(s)\dif s.
\end{align*}
However, with high probability, the $L^\infty$ norm satisfies $\epsilon(t)\sim L^{-1/2},$
which implies $ \lambda^2\epsilon^2\sim \alpha^{-1}.$ 
As a consequence, such a quartic estimate would only provide a lifespan up to
\[
T\sim \alpha^{-1}.
\]
At this timescale, the low-frequency counting estimates can already be handled by relatively crude arguments, and no genuinely kinetic behavior can be observed. 
The main difficulty is therefore to improve the energy estimate beyond the quartic level.  The work of Deng, Ionescu, and Pusateri \cite{DIP25a,DIP25}
established a framework for quasilinear wave turbulence  corresponding to
the dispersion relation $|\xi|^{1/2}$. A crucial feature of this case is
the special algebraic structure of the phase function, which admits a
favorable factorization on the resonant set.  For a general power-law dispersion
$|\xi|^\sigma$, such a factorization is no longer available in general.
The key ingredients and ideas leading to this improvement are summarized as follows:  
\begin{enumerate}
\item \textbf{Symmetrization and factorization of symbols.} 
Working in Fourier space, we study the quartic energy increment
\begin{align*}
Q(u,\overline u,u,\overline u)
:=
\frac{1}{L^{3}}
\sum_{\xi_1-\xi_2+\xi_3-\xi_4=0}
\widehat{u}_{\xi_1}\overline{\widehat{u}}_{\xi_2}
\widehat{u}_{\xi_3}\overline{\widehat{u}}_{\xi_4}
M(\xi),
\end{align*}
where $M(\xi)
:=
|\xi_1|^\beta|\xi_2|^\beta|\xi_3|^\beta|\xi_4|^\beta
\bigl(|\xi_3|^{2N_0}-|\xi_4|^{2N_0}\bigr).$ The first step is to exploit symmetry and rewrite the symbol as
\[
M(\xi)
=
\frac12
|\xi_1|^\beta|\xi_2|^\beta|\xi_3|^\beta|\xi_4|^\beta
\bigl(
|\xi_1|^{2N_0}-|\xi_2|^{2N_0}
+
|\xi_3|^{2N_0}-|\xi_4|^{2N_0}
\bigr).
\]
Next, after an integration by parts in time, the analysis reduces to estimating
\[
M(\xi)\,[\Omega^{(\sigma)}(\xi)]^{-1},
\]
where $\Omega^{(\sigma)}(\xi)
:=
|\xi_1|^\sigma-|\xi_2|^\sigma
+
|\xi_3|^\sigma-|\xi_4|^\sigma.$ 
This leads naturally to a small-divisor problem associated with the resonant phase function.
To overcome this difficulty, we exploit the resonance relation
\[
\xi_1-\xi_2+\xi_3-\xi_4=0
\]
and factorize the symbol:
 \begin{align*}
    |\xi_1|^{2N_0}&-|\xi_2|^{2N_0}+|\xi_3|^{2N_0}-|\xi_4|^{2N_0}\\
    &=2N_0(2N_0-1)
     (\xi_1-\xi_2)(\xi_2-\xi_3)\\
     &\quad\quad \times \int_0^1\int_0^1((1-\theta)\eta\xi_2+(1-\theta)(1-\eta)\xi_3+\eta\theta\xi_1+\theta(1-\eta)\xi_4)^{2N_0-2}\dif \eta\dif\theta.
\end{align*}
The key point is that the integral term retains sufficient regularity for large $N_0$, while the additional factor $(\xi_1-\xi_2)(\xi_2-\xi_3)$
provides cancellation against the small divisor arising from $\Omega^{(\sigma)}$.  Here, we have to analyze the phase function
$\Omega^{(\sigma)}$ for the full range of power-law dispersions
$|\xi|^\sigma$. Unlike the special cases $\sigma=\frac12$ and
$\sigma=2$, where the phase function admits explicit algebraic
factorizations, no universal factorization is available for a general
value of $\sigma$. Therefore, we need to perform a case-by-case analysis
of the resonance geometry and establish suitable factorizations.
Such structures naturally appear in normal form transformations, where division by resonant phases may otherwise lead to a loss of derivatives.

Moreover, expressing the symbol in this form allows us to obtain robust derivative bounds needed for the application of Lemma~\ref{lem:sinfty}. 
A crucial observation is that for any multi-index $\alpha$, any $a,b>0$, and any $\gamma\in\mathbb R$,
\begin{align*}
    |\xi^\alpha\partial_\xi^\alpha[(a\xi_1+b\xi_2)^\gamma]|
    \lesssim
    (a\xi_1+b\xi_2)^\gamma.
\end{align*}
This property guarantees that the differentiated symbols preserve the same scaling behavior as the original phase factors. Subsequently, we rewrite all relevant expressions in terms of combinations of homogeneous power-law factors, which allows us to recover the appropriate scaling properties and establish uniform symbol estimates.
 
 \item \textbf{Cubic resonances.} 
To improve the quartic energy estimate, it is necessary to analyze carefully the resonant four-wave interactions, characterized by the system
\begin{align*}
    \xi_1 - \xi_2 + \xi_3 - \xi_4 = 0,
    \qquad
    |\xi_1|^{\sigma} - |\xi_2|^{\sigma}
    + |\xi_3|^{\sigma} - |\xi_4|^{\sigma} = 0.
\end{align*}
For every $\sigma\in(0,1)\cup(1,2]$, this system always admits the trivial resonant solutions $\{\xi_1,\xi_3\}=\{\xi_2,\xi_4\}.$
When $\sigma\in(1,2]$, these are in fact the only resonances. 
In contrast, for $\sigma\in(0,1)$, the system also admits a family of genuinely nontrivial resonant solutions, which we denote by $(\xi_i^*)_{1\le i\le4}$.

In the regime $\sigma\in(1,2]$, only trivial resonances occur. 
These interactions lead to integrable contributions and therefore do not produce any nontrivial growth of the energy. 
To control the associated small divisors, we perform a detailed case-by-case analysis according to the signs and relative sizes of the frequencies $\xi_i$. 
Exploiting the symmetry of the interaction, we reduce the argument to three representative configurations. 
In each case, the resonance relation admits a suitable factorization structure, allowing us to control the small divisors and derive robust derivative estimates for the corresponding symbols. 
Combined with a normal form transformation, this yields the desired sextic energy increment.

The situation is substantially more delicate when $\sigma\in(0,1)$. 
In this regime, the nontrivial resonances generate genuine singularities, since
\[
M(\xi)[\Omega^{(\sigma)}(\xi)]^{-1}
\]
blows up near the resonant set $\xi=\xi^*$. 
For the gravity water-wave equation \cite{DIP25a} (corresponding to $\sigma=\frac12$), the cubic symbol vanishes exactly on the nontrivial resonant set, so the singularity cancels:
\[
\frac{M(\xi)}{\Omega^{(\sigma)}(\xi)}
=
\frac{M(\xi)-M(\xi^*)}
{\Omega^{(\sigma)}(\xi)-\Omega^{(\sigma)}(\xi^*)}.
\]
However, the MMT model does not enjoy such a favorable algebraic structure. 
Instead, we perform a similar renormalized decomposition, while treating separately the contribution of the resonant component $M(\xi^*)$. 

A further difficulty is that, unlike the case $\sigma=\frac12$, the nontrivial resonant solutions do not admit explicit formulas. 
As a result, the analysis relies on an implicit-function argument together with a careful study of the geometry of the resonant set. 
Through this refined analysis, we ultimately recover a quartic-order contribution, leading to a slight improvement in the lifespan of the solution. 

 
\end{enumerate}

 
\subsection{Organization of the paper} In Section \ref{sec:pre}, we gather the preliminary results  that will be used throughout the text.
Then, the proof is carried out in two main steps. In Sections \ref{trees and LTE}-\ref{sec:proof of theorem 2}, we establish the propagation of randomness for the truncated equation \eqref{TrMMT}. Section \ref{trees and LTE} introduces the basic notions of trees, couples, and molecules, and expresses the solution in terms of ternary trees. In Section \ref{sec:mainest}, we gather some estimates for couples, including both counting estimates and integral bounds. Section \ref{sec:irregular} is devoted to the analysis of irregular chains and the demonstration of their cancellation. Section \ref{sec:Algorithm} presents our main algorithm for molecules. The treatment of the remainder term is given in Section \ref{bd:remainder}. Finally, in Section \ref{sec:proof of theorem 2}, we complete the proof of Theorem \ref{thm2}, which involves delicate calculations for   the low-order ternary trees. Then, in Sections~\ref{sec:acceptineq}–\ref{sec:proofthm1}, we establish the energy estimate stated in Theorem~\ref{thm:energy inequality} and present the proof of the main result, Theorem~\ref{thm1}, respectively. Finally, in Appendix~\ref{sec:app}, we include the algorithm introduced in~\cite{Vas24} for the reader’s convenience.


 \section{Preliminaries}\label{sec:pre}
 \subsection{The Littlewood-Paley projections and associated operators}
 For any $k\in\mathbb{Z}$, we define  an even smooth function $\varphi:\mathbb{R}\to[0,1]$ with  support in $[-\frac85, \frac85]$ and satisfying $\varphi=1$ in $[-\frac54, \frac54]$. We define for $k\in\mathbb{Z}$
\begin{align*}
    \varphi_k(x):=\varphi(x/2^{k})-\varphi(x/2^{k-1}),
\end{align*}
and
\begin{align*}
     \varphi_{\leq k}(x):=\varphi(x/2^{k}),\ \  \varphi_{\geq k}(x):=1-\varphi(x/2^{k-1}).
\end{align*}
Let $P_k,P_{\geq k}, P_{\leq k }$ be the  operators on $\mathbb{T}_L$ defined by the Fourier multipliers $\varphi_k, \varphi_{\leq k}, \varphi_{\geq k}$ respectively.  Moreover, we write
\begin{align*}
    P_k'=\sum_{a=-2,-1,0,1,2}P_{k+a},\ \  \varphi_k'=\sum_{a=-2,-1,0,1,2}\varphi_{k+a}.
\end{align*}
We also denote $\overline{\mathbb{Z}}:=\mathbb{Z}\cup\{-\infty\}$, and  define the operator $P_{-\infty}$ as
\begin{align*}
    P_{-\infty}f:=\frac{1}{L}\int_{\mathbb{T}_L}f\dif x.
\end{align*}
The operators $P_k$ are given by convolution in the physical space. More precisely, for $f\in L^1(\mathbb{T}_L)$ and $k\in\mathbb{Z}$, we have
\begin{align*}
    P_kf(x):&=\int_{\mathbb{T}_L}f(y)K(x-y)\dif y,\\
    K(x):&=\frac{1}{L}\sum_{\xi\in\mathbb{Z}_L}\varphi_k(\xi)e^{2\pi i\xi x}.
\end{align*}

\subsection{Multipliers }
We will often work with multipliers $m:\mathbb{R}^n\to\mathbb{C}$ operators defined by such multipliers. % We define the class of symbols as
%\begin{equation}
% S^\infty := \left\{ m : (L^{-1}\mathbb{Z})^n \to \mathbb{C} : \|m\|_{S^{\infty}} := \| \mathcal{F}^{-1} (m) \|_{L^1(\mathbb{T}_L^n)} < \infty \right\}.
%\end{equation}
For $n\geq2$, we define
 
\begin{align*}   
   \mathbb{H}_{L}^n& := \{ (\xi_1, \ldots, \xi_{n+1}) \in (L^{-1}\mathbb{Z})^{n+1} : \xi_1 + \ldots + \xi_{n+1} = 0 \},\\
    \mathbb{H}_{\infty}^n &:= \{ (\xi_1, \ldots, \xi_{n+1}) \in \mathbb{R}^{n+1} : \xi_1 + \ldots + \xi_{n+1} = 0 \},
\end{align*} 
and the associated class of symbols
\begin{align*}
    \tilde S^\infty:=\tilde S^\infty(\mathbb{H}_L^n),
\end{align*}
where for $m:\mathbb{H}_L^n\to\mathbb{C}$,
\begin{equation}
   \| m \|_{\tilde S^\infty} := \bigg{\|}\frac{1}{L^n} \sum_{\xi_1, \ldots, \xi_{n} \in L^{-1} \mathbb{Z}} m(\xi_1, \ldots, \xi_{n},-\xi_1-\xi_2...-\xi_{n}) e^{2\pi i(x_1\xi_1+x_2\xi_2+...+x_n\xi_n)}\bigg{\|}_{L^1_x(\mathbb{T}_L^n)} < \infty.
\end{equation} 
Then we have the following property on  symbols:  \begin{lemma}$($\cite[Lemma 2.1]{DIP25a}$)$\label{lem:holder}
     Let $p_1,p_2,...,p_n\in[1,\infty],n\geq3$ satisfying $\frac1{p_1}+\frac1{p_2}+...+\frac1{p_n}=1$ and $m\in\tilde{S}^\infty(\mathbb{H}_L^{n-1})$. Then it holds that
     \begin{align*}
        \left| \frac{1}{L^{n-1}}\sum_{(\xi_1,\xi_2,...,\xi_n)\in\mathbb{H}_L^{n-1}}\widehat{f}_1(\xi_1)...\widehat{f}_n(\xi_n)m(\xi_1,\xi_2,...,\xi_n)\right|\lesssim \|f_1\|_{L^{p_1}} \|f_2\|_{L^{p_2}}... \|f_n\|_{L^{p_n}}\|m\|_{\tilde{S}^\infty}.
     \end{align*}
 \end{lemma}
We also need to the following estimates on  the  multipliers.  
 \begin{lemma}$($\cite[Lemma 2.2]{DIP25a}$)$\label{lem:sinfty}
 \begin{enumerate}
     \item % If $m,m'\in S^\infty$, then $m\cdot m'\in S^\infty$ and it holds that 
     %\begin{align*}
    %     \|m\cdot m'\|_{S^\infty}\lesssim   \|m \|_{S^\infty}  \|  m'\|_{S^\infty}.
    % \end{align*}
    % Similarly,\
    If $m,m'\in \tilde{S}^\infty$, then $m\cdot m'\in \tilde{S}^\infty$ and it holds that 
     \begin{align*}
         \|m\cdot m'\|_{\tilde{S}^\infty}\lesssim   \|m \|_{\tilde{S}^\infty}  \|  m'\|_{\tilde{S}^\infty}.
     \end{align*}
   \item 
Assume that $n\geq1$ and $m:\mathbb{R}^n\to\mathbb{C}$ is a continuous compactly supported function and $\check{m}\in L^1(\mathbb{R}^n)$ is the  inverse Fourier transform of $m$. Then
\begin{align}
    \bigg \|\frac{1}{L^{n}}\sum_{\xi_1,\xi_2,...,\xi_n\in L^{-1} \mathbb{Z}}m(\xi_1,\xi_2,...,\xi_n)e^{2\pi i(x_1\xi_1+x_2\xi_2+...+x_n\xi_n)} \bigg \|_{L^1_x(\mathbb{R}^n)}\lesssim\|\check{m}\|_{L^1(\mathbb{R}^n)}.
\end{align}
\item  Assume that $\underline{k}=(k_1,...,k_{d+1})\in\mathbb{Z}^{d+1},\Lambda>0$,  and the multiplier $m:\mathbb{H}_L^d\to\mathbb{C}$  admits an extension $\tilde{m}:\mathbb{R}^{d+1}\to\mathbb{C}$, such that  $m(\xi)\varphi_{\underline{k}}(\xi)=\tilde{m}(\xi)\varphi_{\underline{k}}(\xi)$ for all $\mathbb{H}_L^d$. Here we introduce the notation $$\varphi_{\underline{k}}(\xi):=\varphi_{k_1}(\xi_1)\varphi_{k_2}(\xi_2)\cdots\varphi_{k_{d+1}}(\xi_{d+1}).$$  If it  satisfies the differential inequalities
\begin{align*}
    |2^{\alpha\cdot \underline{k}}\partial_{\xi}^\alpha[\tilde{m}(\xi_1,...,\xi_{d+1})\varphi_{\underline{k}}(\xi)]|\lesssim \Lambda,
\end{align*}
for any multi-index $\alpha=(\alpha_1,...,\alpha_{d+1})$ with $|\alpha|\leq \frac{d+2}{2}$ and any $\xi\in\mathbb{H}_{\infty}^d$, then
\begin{align*}
    \|m(\xi)\varphi_{\underline{k}}(\xi)\|_{\tilde{S}^\infty}\lesssim \Lambda.
\end{align*}
 \end{enumerate}

 \end{lemma}

\subsection{Function spaces}
In this paper, we denote $H^N/\dot{H}^N$ by the  standard inhomogeneous/ homogeneous  Sobolev spaces on $\mathbb{T}_L$ respectively.  We define the $L^\infty$ based Sobolev space $W^{N,0}$ on $\mathbb{T}_L$ by
\begin{align*}
    \|f\|_{W^{N,0}}:=\|P_{-\infty}f\|_{L^\infty}+\sum_{k\in\mathbb{Z}}\|P_kf\|_{L^\infty}2^{Nk}.
\end{align*}
We define an additional $L^\infty$ based Sobolev space $\widetilde{W}^N$ by
\begin{align*}\|f\|_{\widetilde{W}^N}:=\|f\|_{L^\infty}+\sum_{k\in\mathbb{Z}^+}\|P_kf\|_{L^\infty}2^{Nk}.
\end{align*}
In particular, for $N\geq1$, we have
\begin{align}
    \|f\|_{\widetilde{W}^N}\lesssim  \|f\|_{W^{N,0}}.
\end{align}

  \section{Tree Expansions}\label{trees and LTE}
  From this section up to Section~\ref{sec:proof of theorem 2}, we focus on the low-frequency component, namely the truncated equation~\eqref{TrMMT}, and aim to prove Theorem~\ref{thm2}.
Our analysis is based on a Feynman diagram expansion. In this section, we show how the solution can be represented in terms of ternary trees (Feynman diagrams) and formulate the main estimates associated with these trees.  
 \subsection{First reductions}\label{reductions}
We consider the truncated MMT equation as introduced in \eqref{TrMMT}.
 Let $(\widehat{u}_{\rm tr})_k(t)$ denotes the Fourier coefficients of $u_{\rm tr}(t,x)$. First, we define $$a_k(t):= e^{-2\pi i|k|^\sigma t} (\widehat u_{\rm tr})_k(t),\ \ t\in[0,T],$$ and obtain the following equation for the Fourier modes:
 \begin{equation*}%\label{akeqn}
\begin{cases}
 i \partial_t{a_k} =- \lambda^2 L^{-2} \sum\limits_{\substack{(k_1,k_2,k_{3}) \in (\mathbb{Z}_L)^3 \\ k - k_1 + k_2 -k_3 = 0}} a_{k_1}\overline{a_{k_2}}  a_{k_3} e^{2\pi i \Omega^{(\sigma)}(k_1,k_2,k_3,k)t}B_{k_1,k_2,k_3,k}\varphi_{\leq K}(k), \\ 
a_k(0) =(a_k)_{\rm in},
\end{cases}
\end{equation*}
where we define 
\begin{align*}
\Omega^{(\sigma)}(k_1,k_2,k_3,k) &:=|k_1|^\sigma-|k_2|^\sigma+|k_3|^\sigma-|k|^\sigma,\\ B_{k_1,k_2,k_3,k}&:=|k_1|^\beta|k_2|^\beta|k_3|^\beta|k|^\beta.
\end{align*}
Here and in the following, we write $\Omega^{(\sigma)}(k_1,k_2,k_3,k)=\Omega(k_1,k_2,k_3,k)$ when there is no ambiguity.
Then, we rescale the time by setting
$$
b_k(t)=a_k(Tt),\ \ t\in[0,1],
$$ 
which solves 
\begin{equation*}%\label{bkeqn}
\begin{cases}
 i \partial_t{b_k} =- \lambda^2 L^{-2} T\sum\limits_{k - k_1 + k_2 -k_3 = 0} b_{k_1}\overline{b_{k_2}}  b_{k_3} e^{2\pi i \Omega(k_1,k_2,k_3,k)Tt}B_{k_1,k_2,k_3,k}\varphi_{\leq K}(k),\\
b_k(0) =(a_k)_{\rm in}.
\end{cases}
\end{equation*}
Note that the positions of $k_1,k_3$ are symmetric above, so the sum above can be written as 
$$
\sum\limits_{ k - k_1 + k_2 -k_3 = 0}= \sum_{k_1, k_3\neq k}+2\sum_{k_1=k,k_3\neq k}+\sum_{k_1=k_2=k_3=k},
$$
which allows us to write, by introducing the renormalized variable
$$c_k=b_ke^{-2\pi i \Gamma(k,t)},\ \ t\in[0,1],$$
as
\begin{align}\label{ipartck}
 i \partial_t{c_k}& = - \lambda^2 L^{-2} T\sum\limits_{k - k_1 + k_2 -k_3 = 0} c_{k_1}\overline{c_{k_2}}  c_{k_3} e^{2\pi i \Theta(k_1,k_2,k_3,k,t)}B_{k_1,k_2,k_3,k}\varphi_{\leq K}(k) +2\pi\partial_t\Gamma(k,t)c_k\notag\\ 
& = -\alpha L^{-1} T
\bigg(   \sum\limits_{(k_1,k_2, k_{3})}^{\times} c_{k_1}\overline{c_{k_2}}  c_{k_3} e^{2\pi i \Theta(k_1,k_2,k_3,k,t)}B_{k_1,k_2,k_3,k}\varphi_{\leq K}(k) +|c_k|^2c_kB_{k,k,k,k}\varphi_{\leq K}(k) \bigg)\notag\\
 &\qquad- 2\alpha L^{-1} T\sum_{k_1}|c_{k_1}|^2B_{k_1,k_1,k,k}\varphi_{\leq K}(k)c_k+2\pi\partial_t\Gamma(k,t)c_k,
\end{align}
where we denote
$$\Theta(k_1,k_2,k_3,k,t):=\Omega(k_1,k_2,k_3,k)Tt+ \Gamma(k_1,t)- \Gamma(k_2,t)+ \Gamma(k_3,t)- \Gamma(k,t).$$ Here, we still do not emphasize $\sigma$ if it does not cause ambiguity.
We recall the relation $\alpha =\lambda^2L^{-1}$.  Here and below, the notation $\sum^{\times}_{(k_1,k_2, k_{3})}$ denotes summation under the constraints 
$k_j\in\mathbb{Z}_L$, $k_1-k_2+k_3=k$, and $k\not\in\{k_1,k_3\}$. 
The function $\Gamma(k,t)$ is introduced to handle resonant contributions via renormalization; 
its precise choice will be specified in the sequel. We additionally ask that $\Gamma(k,0)=0$.
 

Finally, we recall the well-prepared initial data (\ref{wellprepared}), which yields
\begin{equation}\label{initial0}
c_k(0)=(a_k)_{\mathrm{in}}=\varphi_{\leq K}(k)\sqrt{n_{\mathrm{in}}(k)} \cdot\eta_{k}(\omega).
\end{equation} 
 In the rest of this paper, we focus on the system \eqref{ipartck} on the time interval $[0,1]$ together with the  well-prepared  initial condition \eqref{initial0}.

\subsection{The tree expansion}
Now, we try to express the solution to \eqref{ipartck} in terms of certain ternary tree expansions, following the ideas in \cite{DH21,DH23a,Vas24,Wu25}. 
This approach provides a systematic way to represent the iterated Duhamel expansions of the solution in a combinatorial structure. 
More precisely, each node of the ternary tree corresponds to one trilinear interaction generated by the nonlinearity, while the leaves represent the initial data.  
In particular, this ternary tree formulation can be viewed as an alternative of the classical Feynman diagram expansion, 
which is often used in physics and wave turbulence theory.


\begin{definition}\label{tree1} 
\begin{enumerate}
\item Let $\mathcal{T}$ be a ternary tree. We use $\mathcal{L}$ to denote the set of leaves, use $\mathcal{N}=\mathcal{T}\backslash\mathcal{L}$  to denote the set of branching nodes, denote $n$ by the  number of branching nodes, and $\mathfrak{r} \in \mathcal N$ by the root node. Then the number of leaves is $l=2n+1$. The scale of a ternary tree $\mathcal T$ is defined as $|\mathcal T|=n$,   the number of branching nodes. In particular, the scale of a single node $\mathcal{T}=\bullet$ is 0.   The three children subtrees of $\mathcal{T}$ are defined by $\mathcal{T}_1,\mathcal{T}_2$
 and $\mathcal{T}_3$ from left to right.

\item   For each node $\mathfrak{n}\in\mathcal{N}$, let its children from left to right be $\mathfrak{n}_1$, $\mathfrak{n}_2$, $\mathfrak{n}_3$. We refer to them as  siblings of each other. We fix the sign $\zeta_{\mathfrak{n}}\in\{\pm\}$ as follows: first we fix $\zeta_{\mathfrak{r}}=\{\pm\}$, then for any node $\mathfrak{n}\in \mathcal{N}$, define $\zeta_{\mathfrak{n}_1}=-\zeta_{\mathfrak{n}_2}=\zeta_{\mathfrak{n}_3}=\zeta_{\mathfrak{n}}$. The sign of the root $\mathfrak{r}$ of $\mathcal{T}$ is denoted by the sign of $\mathcal{T}$.  The sign of subtrees $\mathcal{T}_1,\mathcal{T}_3$ are the same as $\mathcal{T}$, while the sign of $\mathcal{T}_2$ is opposite.  In the following, we typically denote by $\mathcal{T}^\zeta$ a tree with sign  $\zeta=+/-$. We define the conjugate
 $\overline{\mathcal{T}}$ of a tree $\mathcal{T}$ to be the same tree but with opposite sign.
\end{enumerate}
\end{definition}
 
 Then, we present a tree expression as a  approximation of  $c_k(t)$, which is defined  as follows: let $N$ be a integer to be determined later (independent of $L$), and let
\begin{align}
     c_k(t):=c^{\leq N}_k (t)+\mathcal{R}^{N+1}_k:=\sum_{i=1}^N c^{i}_k (t)+\mathcal{R}^{N+1}_k,\ \  c^{i}_k(t):=\sum_{|\mathcal{T}^+|=i}
 c^{\mathcal{T}^+}_k (t),\label{def:ck}
\end{align}
where the sum in the last term is taken over all the trees with sign $+$ and scale $i$.
 Here, $\mathcal{R}^{N+1}_k$ is denoted by the remainder term. 

For any tree $\mathcal{T}$ (which may have sign $+$ or $-$), the coefficient $c^{\mathcal{T}}_k(t)$ is defined inductively as:
\begin{align}
    c_k^{\bullet}(t)&=\varphi_{\leq K}(k)\sqrt{n_{\mathrm{in}}(k)} \cdot\eta_{k}^{\zeta_\mathfrak{r}}(\omega),\notag\\
     c_k^{\mathcal{T}}(t)&=\int_0^t\mathcal{W}(c^{\mathcal{T}_1}, c^{\mathcal{T}_2} , c^{\mathcal{T}_3})_k(t')\,\mathrm{d}t',\label{def:ct:int}
\end{align}
where we write $z^+=z$ and $z^-=\overline{z}$. We recall that $\mathcal{T}_1,\mathcal{T}_2,\mathcal{T}_3$ denote the three subtrees of the root $\mathfrak{r}$,  from left to right, and $\mathcal{W}$ is the trilinear operator defined  by
\begin{align}
    \label{Duhamel0}
\mathcal{W}(c^{\mathcal{T}_1}, c^{\mathcal{T}_2} , c^{\mathcal{T}_3})_k(t):%&=\frac{i\alpha T}{L}  \sum\limits_{(k_1,k_2, k_{3})}^{\times} c^{\mathcal{T}_1}_{k_1}{c^{\mathcal{T}_2}_{k_2}}  c^{\mathcal{T}_3}_{k_3} e^{2\pi i \Theta(k_1,k_2,k_3,k,t)}B_{k_1,k_2,k_3,k}\varphi_{\leq K}(k)\notag\\
%&\quad-\frac{i\alpha T}{L}c^{\mathcal{T}_1}_{k}{c^{\mathcal{T}_2}_{k}}c^{\mathcal{T}_3}_{k}B_{k,k,k,k}\varphi_{\leq K}(k)\notag\\
%&\quad+\frac{i\alpha T}{L}\sum_{k_1}\left(c^{\mathcal{T}_1}_{k_1}{c^{\mathcal{T}_2}_{k_1}} -\mathbf{E}[c^{\mathcal{T}_1}_{k_1}{c^{\mathcal{T}_2}_{k_1}} ] \right) B_{k_1,k_1,k,k}\varphi_{\leq K}(k)c^{\mathcal{T}_3}_{k}\notag\\
%&\quad+\frac{i\alpha T}{L}\sum_{k_3}\left({c^{\mathcal{T}_2}_{k_3}}c^{\mathcal{T}_3}_{k_3} -\mathbf{E}[{c^{\mathcal{T}_2}_{k_3}}c^{\mathcal{T}_3}_{k_3} ]\right) B_{k,k_3,k_3,k}\varphi_{\leq K}(k)c^{\mathcal{T}_1}_{k}\notag\\
&=\frac{i\alpha T}{L}  \sum\limits_{(k_1,k_2, k_{3})}^{\times} c^{\mathcal{T}_1}_{k_1}{c^{\mathcal{T}_2}_{k_2}}  c^{\mathcal{T}_3}_{k_3} e^{2\pi i \Theta(k_1,k_2,k_3,k,t)}B_{k_1,k_2,k_3,k}\varphi_{\leq K}(k)\notag\\
&\quad+\frac{i\alpha T}{L}\left(c^{\mathcal{T}_1}_{k}{c^{\mathcal{T}_2}_{k}}c^{\mathcal{T}_3}_{k}-c^{\mathcal{T}_1}_{k}\mathbf{E}[{c^{\mathcal{T}_2}_{k}}c^{\mathcal{T}_3}_{k}]-\mathbf{E}[{c^{\mathcal{T}_1}_{k}}c^{\mathcal{T}_2}_{k}]c^{\mathcal{T}_3}_{k}\right)B_{k,k,k,k}\varphi_{\leq K}(k)\notag\\
&\quad+\frac{i\alpha T}{L}\sum_{k_1\neq k}\left(c^{\mathcal{T}_1}_{k_1}{c^{\mathcal{T}_2}_{k_1}} -\mathbf{E}[c^{\mathcal{T}_1}_{k_1}{c^{\mathcal{T}_2}_{k_1}} ] \right) B_{k_1,k_1,k,k}\varphi_{\leq K}(k)c^{\mathcal{T}_3}_{k}\notag\\
&\quad+\frac{i\alpha T}{L}\sum_{k_3\neq k}\left({c^{\mathcal{T}_2}_{k_3}}c^{\mathcal{T}_3}_{k_3} -\mathbf{E}[{c^{\mathcal{T}_2}_{k_3}}c^{\mathcal{T}_3}_{k_3} ]\right) B_{k,k_3,k_3,k}\varphi_{\leq K}(k)c^{\mathcal{T}_1}_{k}.
\end{align}
Here, we apply a renormalization at the level of the product, rather than to the individual factors. This approach partially eliminates the resonant contributions.
 
Comparing \eqref{Duhamel0} with \eqref{ipartck}, we choose the renormalization function so that \eqref{Duhamel0} provides a good approximation to \eqref{ipartck}:
\begin{align}
    \partial_t\Gamma(k,t):&=\frac{\alpha T}{\pi L}\sum_{ |\mathcal{T}^+|\leq N,\ |\mathcal{T}^-|\leq N }\sum_{k_1}\mathbf{E}[c^{\mathcal{T}^+}_{k_1} {c^{\mathcal{T}^-}_{k_1}} ] B_{k_1,k_1,k,k}\varphi_{\leq K}(k),\label{partgamma}\\
    \Gamma(k,0)&=0,
\end{align}
where the first sum is taken over all the positive and negative trees with scale no more than $N$. 
Here we use the fact that $B_{k_1,k_1,k,k}=B_{k,k_1,k_1,k}$. We also notice that the expression $c_k^{\mathcal{T}} = c_k^{\mathcal{T}}(\Gamma)$ depends on $\Gamma$, so that \eqref{partgamma} can be regarded as an ODE system. The choice of $\Gamma$ is then obtained by solving this ODE system via a fixed point argument.  

Compared with the semi-linear case studied in \cite{DH21,DH23a,Vas24}, resonant contributions still persist in the present setting. We therefore introduce the notion of enhanced trees to  capture these residual resonances.

 \begin{definition} \label{def:enhancedtress}
\begin{enumerate}
\item  (Enhanced Trees). An enhanced tree is a ternary tree $\mathcal{T}$ equipped with an additional set $\mathcal{N}_D\subset \mathcal{N}$ of degenerate branching nodes. Here, for each degenerate node $\mathfrak{n}\in \mathcal{N}_D$ and its three children $\{\mathfrak{n}_1,\mathfrak{n}_2,\mathfrak{n}_3\}$, we choose a marked child $\mathfrak{n}_c$ from $\{\mathfrak{n}_1,\mathfrak{n}_3\}$, and label
 the remaining child   as $\mathfrak{n}_{nc}$. The same degenerate node, with different choices of marked child, is regarded as distinct types of degenerate nodes. In particular, a node $\mathfrak{n}\in\mathcal{N}_D$ may be fully degenerate, meaning that both $\mathfrak{n}_1$ and $\mathfrak{n}_3$ can be chosen as marked children.
\item (Admissible Decoration) We assign to each $\mathfrak{n}\in\mathcal{T}$ an element $k_{\mathfrak{n}}\in\mathbb{Z}_L$. We say such a decoration $(\mathcal{N}_D,k[\mathcal{T}])$ is admissible if for any $\mathfrak{n}\in\mathcal{N}$ we have $k_{\mathfrak{n}}=k_{\mathfrak{n}_1}-k_{\mathfrak{n}_2}+k_{\mathfrak{n}_3}.$
Moreover, for any $\mathfrak{n}\in\mathcal{N}/\mathcal{N}_D$, it holds that  $k_{\mathfrak{n}}\not\in\{k_{\mathfrak{n}_1},k_{\mathfrak{n}_3}\}$. For any $\mathfrak{n}\in \mathcal{N}_D$, it holds that    $k_{\mathfrak{n}}=k_{\mathfrak{n}_c},\ k_{\mathfrak{n}_2}=k_{\mathfrak{n}_{nc}}.$ In particular, we say that a node is fully degenerate if $k_{\mathfrak{n}}=k_{\mathfrak{n}_1}=k_{\mathfrak{n}_2}=k_{\mathfrak{n}_3}$, which is consistent with the choice of marked child described above. 
 Moreover, an admissible decoration is uniquely determined by the values of $k_{\mathfrak{l}}$ for $\mathfrak{l}\in\mathcal{L}$.  
\item  For any $n\in\mathcal{N}$ we denote $$\Omega_{\mathfrak{n}}:=\Omega(k_{\mathfrak{n}_1},k_{\mathfrak{n}_2},k_{\mathfrak{n}_3},k_{\mathfrak{n}}),\  \ \sum\Gamma(\mathfrak{n},t)= \Gamma(k_{\mathfrak{n}_1},t)- \Gamma(k_{\mathfrak{n}_2},t)+ \Gamma(k_{\mathfrak{n}_3},t)- \Gamma(k_{\mathfrak{n}},t),$$ and $$\Theta_{\mathfrak{n}}(t):=\Theta(k_{\mathfrak{n}_1},k_{\mathfrak{n}_2},k_{\mathfrak{n}_3},k_{\mathfrak{n}},t):=\Omega_{\mathfrak{n}}Tt+\sum\Gamma(\mathfrak{n},t).$$ %We define 
%\begin{align*}
%\epsilon_{k_{\mathfrak{n}_1},k_{\mathfrak{n}_2},k_{\mathfrak{n}_3}} = 
%\begin{cases}
%-1, & \text{if } k_{\mathfrak{n}} = k_{\mathfrak{n}_1} = k_{\mathfrak{n}_2} = k_{\mathfrak{n}_3} \text{ and } \mathfrak{n} \notin \mathcal{N}_D, \\
%1, & \text{otherwise}.
%\end{cases}
%\end{align*}

\end{enumerate}
\end{definition}

From now on, we will refer to an enhanced tree simply as a “tree” for convenience.
Note that the same underlying tree equipped with different decorations is regarded as distinct.

In the expression given by \eqref{Duhamel0}, the four terms on the right-hand side correspond to four distinct configurations: the root $\mathfrak{r}$ is non-degenerate and the decorations $k_i$ are distinct; the root $\mathfrak{r}$ is fully degenerate; the child $\mathfrak{n}_3$ of $\mathfrak{r}$ is degenerate; and the child $\mathfrak{n}_1$ of $\mathfrak{r}$ is degenerate.  %Furthermore, we notice that we allow $k_{\mathfrak{n}}=k_{\mathfrak{n}_c}=k_{\mathfrak{n}_2}=k_{\mathfrak{n}_{nc}}$ for any degenerate node. 

Then, we give the explicit expression for the tree representation $c^{\mathcal{T}}_{k}$.
\begin{proposition}\label{treedefs} 
For any enhanced tree $\mathcal{T}$ with scale $n$,  we have
\begin{align}
   \label{def:ct}
c^{\mathcal{T}}_{k}(t,\omega)&=\bigg(\frac{\alpha T}{L}\bigg)^n\sum_{\mathcal{D}}\prod_{\mathfrak{n}\in\mathcal{N}}(i\zeta_{\mathfrak{n}}%\epsilon_{k_{\mathfrak{n}_1},k_{\mathfrak{n}_2},k_{\mathfrak{n}_3}}
)\prod_{\mathfrak{n}\in\mathcal{N}}B_{k_{\mathfrak{n}_1},k_{\mathfrak{n}_2},k_{\mathfrak{n}_3},k_{\mathfrak{n}}}\prod_{\mathfrak{n}\in\mathcal{T}}\varphi_{\leq K}(k_\mathfrak{n})\notag\\ 
&\quad\quad \quad\quad \times \mathcal{A}_{\mathcal{T}}(t,\Theta [t,\mathcal{N}])\prod_{\mathfrak{l}\in\mathcal{L}}\sqrt{n_{\mathrm{in}}(k_{\mathfrak{l}})}\mathcal{B}_{\mathcal{T}}(\omega),
\end{align}where the sum on $\mathcal{D}$ is taken over all admissible decorations $(\mathcal{N}_D,k[\mathcal{T}])$ such that $k_{\mathfrak{r}}=k$, and the functions $\mathcal{A}_{\mathcal{T}}(t,\Theta [t,\mathcal{N}])$ are defined iteratively as
\begin{align}\label{def:a}
    \mathcal{A}_{\bullet}(t,\Theta [\mathcal{N}])&:=1,\notag\\
     \mathcal{A}_{\mathcal{T}}(t,\Theta [\mathcal{N}])&:=\int_0^t e^{2\pi i \zeta_{\mathfrak{r}}\Theta_{\mathfrak{r}}(t')}\prod_{i=1}^3 \mathcal{A}_{\mathcal{T}_i}(t',\Theta [\mathcal{N}_i])\dif t',
\end{align}
where $\mathcal{T}_i$ are three subtrees on $\mathcal{T}$ from left to right, and $\mathcal{N}_i$ is the set of branching nodes in  $\mathcal{T}_i$, with $\mathcal{N}=\mathcal{N}_1\cup \mathcal{N}_2\cup \mathcal{N}_3\cup \{\mathfrak{r}\}$. 

The noise terms are defined iteratively as:
\begin{align}
    \mathcal{B}_{\bullet}(\omega)&:=\eta_{k}(\omega)^{\zeta_{\mathfrak{r}}},\notag\\
    \mathcal{B}_{\mathcal{T}}(\omega)&:=\prod_{i=1}^3  \mathcal{B}_{\mathcal{T}_i}(\omega)-1_{\mathfrak{r}\in\mathcal{N}_D}  \mathcal{B}_{\mathcal{T}_{c}}(\omega)\mathbf{E}[  \mathcal{B}_{\mathcal{T}_{2}}(\omega)  \mathcal{B}_{\mathcal{T}_{{nc}}}(\omega)],\label{def:Bomega}
\end{align}
where, if $\mathfrak{r}\in\mathcal{N}_D$, we denote by $\mathcal{T}_{c}, \mathcal{T}_{2}, \mathcal{T}_{nc}$ the subtrees of $\mathfrak{r}$ rooted at the nodes $\mathfrak{n}_{c}, \mathfrak{n}_{2}, \mathfrak{n}_{nc}$, respectively. Here, the indicator $1_{\mathfrak{r}\in\mathcal{N}_D}$ counts twice (corresponding to $\mathfrak{n}_1$ and $\mathfrak{n}_3$) in the fully degenerate case. 


\end{proposition}
\begin{proof} We proceed by induction. First, \eqref{def:ct} holds for $|\mathcal{T}|=0$ and $1$ by definition and by \eqref{def:ct:int}.  Now, assume that \eqref{def:ct} holds for all trees $\mathcal{T}$ with scale at most $n$. Let $\mathcal{T}$ be a tree of scale $n+1$ (with sign $+$, without loss of generality). Its three children $\mathcal{T}_i$, $i=1,2,3$, have scales at most $n$, so that \eqref{def:ct} holds for each $\mathcal{T}_i$ by the induction hypothesis. Therefore, by \eqref{def:ct:int}, we obtain  
\begin{align}
c^{\mathcal{T}}_{k}(t,\omega)&=\bigg(\frac{\alpha T}{L}\bigg)^{n+1}\int_0^t \Bigg{\{}%1_{\mathfrak{r}\notin\mathcal{N}_D}
\sum_{k=k_1-k_2+k_3}\prod_{i=1}^3\Big{(}\sum_{\mathcal{D}_i}\prod_{\mathfrak{n}\in\mathcal{N}_i}(i\zeta_{\mathfrak{n}}%\epsilon_{k_{\mathfrak{n}_1},k_{\mathfrak{n}_2},k_{\mathfrak{n}_3}}
)\prod_{\mathfrak{n}\in\mathcal{N}_i}B_{k_{\mathfrak{n}_1},k_{\mathfrak{n}_2},k_{\mathfrak{n}_3},k_{\mathfrak{n}}}\notag\\
&\times\prod_{\mathfrak{n}\in\mathcal{T}_i}\varphi_{\leq K}(k_\mathfrak{n})\mathcal{A}_{\mathcal{T}_i}(s,\Theta [s,\mathcal{N}])\prod_{\mathfrak{l}\in\mathcal{L}_i}\sqrt{n_{\mathrm{in}}(k_{\mathfrak{l}})}\mathcal{B}_{\mathcal{T}_i}(\omega)\Big{)}e^{2\pi i \Theta_{\mathfrak{r}}(s)}i%\epsilon_{k_{1},k_{2},k_{3}}B_{k_{1},k_{2},k_{3},k}
B_{k_1 ,k_2 ,k_{3},k}\varphi_{\leq K}(k )\notag\\
&-1_{\mathfrak{r}\in\mathcal{N}_D}\sum_{k_2}
\prod_{i=1}^3\Big{(}\sum_{\mathcal{D}_i}\prod_{\mathfrak{n}\in\mathcal{N}_i}(i\zeta_{\mathfrak{n}}%\epsilon_{k_{\mathfrak{n}_1},k_{\mathfrak{n}_2},k_{\mathfrak{n}_3}}
)\prod_{\mathfrak{n}\in\mathcal{N}_i}B_{k_{\mathfrak{n}_1},k_{\mathfrak{n}_2},k_{\mathfrak{n}_3},k_{\mathfrak{n}}}\prod_{\mathfrak{n}\in\mathcal{T}_i}\varphi_{\leq K}(k_\mathfrak{n})\notag\\
&\times\mathcal{A}_{\mathcal{T}_i}(s,\Theta [s,\mathcal{N}])\prod_{\mathfrak{l}\in\mathcal{L}_i}\sqrt{n_{\mathrm{in}}(k_{\mathfrak{l}})}\mathcal{B}_{\mathcal{T}_{c}}(\omega)\mathbf{E}[\mathcal{B}_{\mathcal{T}_{2}}\mathcal{B}_{\mathcal{T}_{nc}}]\Big{)} iB_{k ,k ,k_{2},k_{2}}\varphi_{\leq K}(k )\Bigg{\}}\dif s.\notag
\end{align}
Here, in the first line on the right-hand side, the sum over $\mathcal{D}_i$ is taken over all admissible decorations of $\mathcal{T}i$ such that $k_{\mathfrak{n}_i}=k_i$.
In the second line, the sum over $\mathcal{D}_i$ is taken over all admissible decorations of $\mathcal{T}_i$ satisfying $k_{\mathfrak{n}_c}=k$ and $k_{\mathfrak{n}_2}=k_{\mathfrak{n}_{nc}}=k_2$.
We denote by $\mathcal{L}_i$ the set of leaves of $\mathcal{T}_i$.
We also remark that, when $\mathfrak{r}\in\mathcal{N}_D$, there are two possible choices of the marked child, corresponding to the two terms in \eqref{Duhamel0}.
Then, \eqref{def:ct} follows directly from the definitions \eqref{def:a} and \eqref{def:Bomega}.

\end{proof}


From now on, we write $B_\mathfrak{n}=B_{k_{\mathfrak{n}_1},k_{\mathfrak{n}_2},k_{\mathfrak{n}_3},k_{\mathfrak{n}}},\ \varphi_{\mathfrak{n}}=\varphi_{\leq K}(k_\mathfrak{n})$ for simplicity.

In order to analyze the second-order statistics of the solution, we need to consider products of tree expansions and compute their expectations. This naturally leads to pairing structures between leaves, arising from the independence of the Gaussian random variables.
 In addition, due to the presence of renormalization and degenerate nodes, it is necessary to impose further constraints on admissible pairings. 
\begin{definition}\label{def:couple}
\begin{enumerate} \item  (Couple)
    A couple $\mathcal{Q}$ consists of two trees $\mathcal{T}^+$ and $\mathcal{T}^-$, each labeled with
 opposite signs, together with the choice of degenerate nodes $\mathcal{N}_D=(\mathcal{N}_D)_+\cup( \mathcal{N}_D)_-$, and a partition $\mathcal{P}$ of the set of leaves $\mathcal{L}=\mathcal{L}_+ \cup \mathcal{L}_-$ into $(n+1)$ disjoint
 two-element subsets, where $n = n(\mathcal{T}^+) + n(\mathcal{T}^-)$ is called the order of the couple. The
 partition $\mathcal{P}$ must satisfy the condition $\zeta_l=-\zeta_{l'}$ for every pair $\{l,l'\} \in\mathcal{ P}$. For a couple
 $\mathcal{Q}=\{\mathcal{T}^+,\mathcal{T}^-,\mathcal{P},\mathcal{N}_D\}$, we denote by $\mathcal{N} = \mathcal{N}_+\cup \mathcal{N}_-$ the branching nodes. We define $$
 \zeta(\mathcal{Q}) =\prod_{\mathfrak{n}\in\mathcal{N}} (i\zeta_{\mathfrak{n}}).$$
 The trivial couple is given by two trivial trees whose roots are paired.
 A decoration $\mathcal{E}$  is obtained by decorating each tree $\mathcal{T}^+$ and $\mathcal{T}^-$, with  further requirement that $k_l = k_{l'}$ for every pair $\{l,l'\} \in\mathcal{ P}$. %We
% define
% \begin{align}
%     \epsilon_{\mathcal{E}}=\prod_{\mathfrak{n}\in\mathcal{N}} \epsilon_{k_{\mathfrak{n}_1},k_{\mathfrak{n}_2},k_{\mathfrak{n}_3}}.
% \end{align}
 A decoration  is called a $k$-decoration if $k_{\mathfrak{r}_+} = k_{\mathfrak{r}_-} = k$.
\item  (Enhanced Couples) We call couple $\mathcal{Q}=\{\mathcal{T}^+,\mathcal{T}^-,\mathcal{P},\mathcal{N}_D\}$  an enhanced couple if  the
 partition $\mathcal{P}$ further satisfies that for any $\mathfrak{n}\in\mathcal{N}_D$, the leaves   from the subtree with root $\{\mathfrak{n}_2,\mathfrak{n}_{nc}\}$  are not completely paired together. In the fully degenerate case, this condition is imposed for each choice of marked child. We refer to such a partition $\mathcal{P}$ as an enhanced pairing.
 \end{enumerate}
\end{definition}

\begin{remark}
  In this paper, the definition of enhanced couples given in Definition~\ref{def:couple} is in fact equivalent to that in \cite[Definition 3.3]{DIP25}. Note that if $k_{\mathfrak{n}_1}\neq k_{\mathfrak{n}_2}$,  then it is impossible for the subtrees rooted at $\mathfrak{n}_1$ and $\mathfrak{n}_2$ to be paired.   In the case $k_{\mathfrak{n}_1}=k_{\mathfrak{n}_2}=k_{\mathfrak{n}_1}=k_{\mathfrak{n}}$, from \eqref{Duhamel0}, we have the expression
    \begin{align*}
       c^{\mathcal{T}_1}_{k}{c^{\mathcal{T}_2}_{k}}c^{\mathcal{T}_3}_{k} -\mathbf{E}[c^{\mathcal{T}_1}_{k}{c^{\mathcal{T}_2}_{k}} ]c^{\mathcal{T}_3}_{k} -c^{\mathcal{T}_1}_{k}\mathbf{E}[c^{\mathcal{T}_2}_{k}{c^{\mathcal{T}_3}_{k}} ] ,
    \end{align*}
   which implies that neither $\mathcal{T}_1$ nor $\mathcal{T}_3$ can be paired with $\mathcal{T}_2$.
This is consistent with the definition of enhanced pairing in \cite[Definition 3.3]{DIP25}.
\end{remark}

We introduce the following Complex Isserlis' Theorem.
\begin{lemma}
\label{thm:Isserlis}\cite[Lemma A.2]{DH23a}
Let  $k_j \in \mathbb{Z}_L$ and signs $\{\zeta_j\}\subset \{+,-\}$ for $1 \leq j \leq n$ be fixed. It holds that
\begin{align*}
\mathbf{E}\left[\prod_{j=1}^n \eta^{\zeta_j}_{k_j}(\omega)\right] = 
\sum_{\mathcal{P}} 
\prod_{\{j,j'\}\in\mathcal{P} } {1}_{k_j = k_{j'}}
\end{align*}
where the summation ranges over all partitions $\mathcal{P}$ of $\{1,2,...,n\}$ into pairs $\{j,j'\}$ with  $\zeta_j = -\zeta_{j'}$.

\end{lemma}

\begin{lemma}[Enhanced Tree Pairing]

For enhanced trees $\mathcal{T}^+$ and $\mathcal{T}^-$, the expectation satisfies
\begin{align*}
\mathbf{E}\left[\mathcal{B}_{\mathcal{T}^+}\mathcal{B}_{\mathcal{T}^-}\right] = 
\sum_{ \mathcal{P}}1.\label{lem:enhanced_pairing}
\end{align*}
where the summation ranges over all enhanced pairings between $\mathcal{T}^+$ and $\mathcal{T}^-$.
\end{lemma}
The proof of this lemma is a direct consequence of the Complex Isserlis' Theorem introduced in Lemma~\ref{thm:Isserlis}. We refer to \cite[Lemma 5.14]{Wu25} for a detailed proof.


\begin{definition}[Enhanced Couple Operator]
\label{def:enhanced_couple}
For an enhanced couple  $\mathcal{Q}=\{\mathcal{T}^+,\mathcal{T}^-,\mathcal{P},\mathcal{N}_D\}$ of order $n$, we define  
\begin{align*}
(\mathcal{K}_\mathcal{Q})(t,s,k) := \bigg(\frac{\alpha T}{L}\bigg)^n \zeta(\mathcal{Q}) \sum_{\mathcal{E}}  %\epsilon_{\mathcal{E}}
\prod_{\mathfrak{n}\in\mathcal{N}}B_\mathfrak{n}\prod_{\mathfrak{n}\in\mathcal{Q}}\varphi_\mathfrak{n}
\int_{\mathcal{E}_\mathfrak{n}}\prod_{\mathfrak{n}\in\mathcal{N}}  e^{\zeta_\mathfrak{n}2\pi i\Theta_\mathfrak{n}(t_\mathfrak{n}) } \dif t_\mathfrak{n}
\prod_{\mathfrak{l} \in \mathcal{L}}^+ n_{\rm in}(k_l),
%\times \mathbf{E}\left[B_{T^+}(\eta_{T^+}(\varrho))B^*_{T^-}(\eta_{T^-}(\varrho))\right],
\end{align*}
where  the sum on $\mathcal{E}$ is taken over all $k$-decorations on $\mathcal{Q}$.
 The product $\prod_{\mathfrak{l} \in \mathcal{L}}^+$ is taken over all leaves with $+$ signs. 
The time integration domain is
\begin{align*}
\mathcal{E}_\mathfrak{n} = \{t{[\mathcal{N}]} : 0 < t_{\mathfrak{n}'} < t_\mathfrak{n}  \text{ whenever } \mathfrak{n}' \text{ is a child of } \mathfrak{n} ; 
t_\mathfrak{n}  < t \text{ for } \mathfrak{n}  \in \mathcal{N}^+, \text{ and } t_\mathfrak{n}  < s \text{ for } \mathfrak{n}  \in \mathcal{N}^-\}.
\end{align*}
\end{definition}

Then, as a direct result of \eqref{def:ct}, it holds that 
 \begin{align}
     \mathbf{E}[ c^{\mathcal{T}^+}_k (t)c^{\mathcal{T}^-}_k(s)] =\sum_{\mathcal{P}} (\mathcal{K}_\mathcal{Q})(t,s,k),
 \end{align}
where the  sum is taken over  all enhanced pairings $\mathcal{P}$ on the couple $(\mathcal{T}^+,\mathcal{T}^-)$.
In the following of this paper, we aim to establish the bounds on $(\mathcal{K}_\mathcal{Q})(t,s,k)$, for any fixed  enhanced pairing $\mathcal{P}$.

Going back to \eqref{partgamma}, we rewrite it in the form \begin{align}
    \partial_t\Gamma(k,t)&=\frac{\alpha T}{\pi L}\varphi_{\leq K}(k)|k|^{2\beta}\sum_{k_1}\varphi_{\leq K}^2(k_1)n_{\mathrm{in}}(k_1)|k_1|^{2\beta}\notag\\
    &\quad+\frac{\alpha T}{\pi L}\varphi_{\leq K}(k)|k|^{2\beta}\sum_{k_1}\sum_{\mathcal{Q}_{\leq N}}(\mathcal{K}_\mathcal{Q})(t,t,k_1)|k_1|^{2\beta},\notag
\end{align}
where the summation $\mathcal{Q}_{\leq N}$ is taken over all nontrivial admissible couples with the rank of each tree at most $N$.
We then decompose $$\Gamma(k,t)=\Gamma_0(k)Tt+\Gamma_1(k,t),$$ where 
 \begin{align}
     \Gamma_0(k)&:= \alpha  \varphi_{\leq K}(k)|k|^{2\beta}C_K,\notag\\
     \partial_t  \Gamma_1(k,t)&:=\frac{\alpha T}{\pi L}\varphi_{\leq K}(k)|k|^{2\beta}\sum_{k_1}\sum_{\mathcal{Q}_{\leq N}}(\mathcal{K}_\mathcal{Q})(t,t,k_1)|k_1|^{2\beta},\ \ \Gamma_1(k,0)=0.\label{def:gamma1}
 \end{align}
 Here  we define $$C_K:=\frac{1}{\pi L}\sum_{k_1\in\mathbb{Z}_L}\varphi_{\leq K}^2(k_1)n_{\mathrm{in}}(k_1)|k_1|^{2\beta}  \lesssim  \|n_{\mathrm{in}}(x)|x|^{2\beta}\|_{L^1}.$$
 By the choice of truncation function $\varphi_{\leq K}$, we obtain the bound
\begin{align*}
   \sup_{\xi\in\mathbb{R}} |\Gamma_0(\xi)|\lesssim  \alpha  L^\delta\lesssim L^{-\delta}.
\end{align*}
We remark that the contribution of the trivial couples, encoded in $\Gamma_0(k)$, is the dominant one among all renormalization terms. Indeed, it satisfies  
\[
\Gamma_0(k) T \sim \alpha TL^{\delta},
\]
which would impose the restrictive condition $\alpha T \lesssim 1$ if treated directly.
For this reason, we separate this term and treat it independently in the subsequent analysis.  

We now turn to the analysis of $\Gamma_1$. We observe that \eqref{def:gamma1} defines an ODE in $\Gamma_1$, since   the quantities $c_k^{\mathcal{T}} = c_k^{\mathcal{T}}(\Gamma)$ depend on $\Gamma$, and hence the kernels $\mathcal{K}_{\mathcal{Q}}$ depend implicitly on $\Gamma_1$. Therefore, \eqref{def:gamma1} should be regarded as a nonlinear integral equation for $\Gamma_1$.
 
\subsection{The  remainder term}\label{sec:rem}
In this section, we write down the equation for the remainder term $\mathcal{R}^{N+1}$. Our goal is to derive a precise equation for $\mathcal{R}^{N+1}$ and to exploit its structure in order to control it in the subsequent analysis.
 We recall that by definition,
\begin{align}\label{def:cleqn}
     c_k(t)=c^{\leq N}_k (t)+\mathcal{R}^{N+1}_k.
\end{align}
Then, combining \eqref{ipartck} and \eqref{def:ct:int} with the definition of the truncated expansion, we deduce that $\mathcal{R}^{N+1}$ satisfies the following equation:
\begin{equation}\mathcal{R}_k^{N+1}=(c_{\sim N})_k+\mathscr{L}_k(\mathcal{R}^{N+1}) +\mathscr{Q}_k(\mathcal{R}^{N+1})+\mathscr{C}_k(\mathcal{R}^{N+1}).\label{eqnr}
\end{equation} Here, the operators $\mathscr{L}$, $\mathscr{Q}$, and $\mathscr{C}$ represent, respectively, the linear, quadratic, and cubic contributions involving the remainder $\mathcal{R}^{N+1}$, while $(c_{\sim N})_k$ collects the source terms coming from trees of scale comparable to $N$.
More precisely, the relevant terms are defined as follows: 
\begin{align}
(c_{\sim N})_k&:=\sum_{\substack{n_1, n_2, n_3\leq N\\  n_1+n_2+n_3\geq N}}\int_0^t \mathcal{W}(c^{n_1},\overline {c^{n_2}},c^{n_3})_k\dif s,\notag
\end{align}
where we recall that $c^{i}_k(t):=\sum_{|\mathcal{T}^+|=i}
 c^{\mathcal{T}^+}_k (t),$ and $\mathcal{W}$ is defined in \eqref{Duhamel0}. Note that, in the second term of the expression, we take the complex conjugate, since $c^{i}$ is summed only over trees with sign $+$.
 
To define the remaining three operators, we introduce the trilinear form
\begin{align}
  \mathcal{W}_1(u,v,w)_k(t):&=\frac{i\alpha T}{L}\sum\limits_{(k_1,k_2, k_{3})}^{\times} u_{k_1}{v_{k_2}}  w_{k_3} e^{2\pi i \Theta(k_1,k_2,k_3,k,t)}B_{k_1,k_2,k_3,k}\varphi_{\leq K}(k)\notag\\
&\quad -\frac{i\alpha T}{L}u_{k}{v_{k}}w_{k}B_{k,k,k,k}\varphi_{\leq K}(k)\notag.    
\end{align}
Then we have
\begin{align}
\mathscr{L}(v)&:=\int_0^t \big(2\mathcal{W}_1(c^{\leq N}, \overline{c^{\leq N}}, v)+\mathcal{W}_1(c^{\leq N},\overline{ v}, c^{\leq N})+\mathscr{L}_1[c^{\leq N}](v)\big)\dif s, \notag \\
\mathscr{Q}(v)&:=\int_0^t \big(2\mathcal{W}_1(v,\overline{v}, c^{\leq N})+\mathcal{W}_1(v, \overline{c^{\leq N}}, v)+\mathscr{Q}_1[c^{\leq N}](v)\big)\dif s,\notag\\
\mathscr{C}(v)&:=\int_0^t  \big(\mathcal{W}_1(v, \overline{v}, v)+\mathscr{C}_1(v) \big)\dif s,\label{def:l}
\end{align} 
    where
\begin{align}
(\mathscr{L}_1)[c^{\leq N}]_k(v)&:=2\frac{i\alpha T}{L}\sum_{k_1}\left(|c^{\leq N}_{k_1}|^2 -\mathbf{E}[|c^{\leq N}_{k_1}|^2] \right) B_{k_1,k_1,k,k}\varphi_{\leq K}(k)v_k\notag\\
&\quad +2\frac{i\alpha T}{L}\sum_{k_1}\left(c_{k_1}^{\leq N}\overline{v_{k_1}}+v_{k_1}\overline{c_{k_1}^{\leq N}}\right) B_{k,k_1,k_1,k}\varphi_{\leq K}(k)c_k^{\leq N},\notag\\
(\mathscr{Q}_1)[c^{\leq N}]_k(v)&:=2\frac{i\alpha T}{L}\sum_{k_1}|v_{k_1}|^2 B_{k_1,k_1,k,k}\varphi_{\leq K}(k)c_k^{\leq N}\notag\\
    &\quad +2\frac{i\alpha T}{L}\sum_{k_1}\left(c_{k_1}^{\leq N}\overline{v_{k_1}}+v_{k_1}\overline{c_{k_1}^{\leq N}}\right) B_{k,k_1,k_1,k}\varphi_{\leq K}(k)v_k,\notag\\
(\mathscr{C}_1)_k(v)&:=2\frac{i\alpha T}{L}\sum_{k_1}|v_{k_1}|^2 B_{k_1,k_1,k,k}\varphi_{\leq K}(k)v_k.\notag
\end{align}
Here we used the symmetry of $\mathcal{W}_1$ and the symmetry of the coefficient $B_{k_1,k_2,k_3,k}$. We also re-decompose the summation as
 $$
\sum\limits_{ k - k_1 + k_2 -k_3 = 0}= \sum_{k_1, k_3\neq k}+2\sum_{k_1=k}-\sum_{k_1=k_2=k_3=k}.
$$

Moreover, by solving \eqref{eqnr} and using the definition of $\varphi_{\leq K}$, we deduce that $\mathcal{R}_k^{N+1}=0$ for all $|k|\geq \frac85\cdot 2^K$.  Let $\tilde\varphi_{\leq K}(x):=\varphi_{\leq K}(x/2)$ and  define the Fourier multiplier $\tilde\varphi_{\leq K}u=(\tilde\varphi_{\leq K}(k)u_k)_{k\in\mathbb{Z}_L}$. Then we
have  $$\tilde\varphi_{\leq K} (k)\mathcal{R}_k^{N+1}=\mathcal{R}_k^{N+1}.$$  
In particular, the remainder $\mathcal{R}^{N+1}$ is localized in frequencies where $\tilde\varphi_{\leq K}\equiv 1$, which allows us to insert this cutoff into the nonlinear terms.
From \eqref{eqnr}, it follows that
\begin{equation}\mathcal{R}_k^{N+1}=(c_{\sim N})_k+\mathscr{L}(\mathcal{R}_k^{N+1}) +\mathscr{Q}(\tilde\varphi_{\leq K}(k)\mathcal{R}_k^{N+1})+\mathscr{C}(\tilde\varphi_{\leq K}(k)\mathcal{R}_k^{N+1}).\notag
\end{equation}  
\begin{remark}
We remark that the function $\tilde\varphi_{\leq K}$ is not inserted into the linear operator $\mathscr{L}$, since it will be needed again later in the proof of Theorem~\ref{thm1}.
In particular, when treating the high-frequency component, it is not appropriate to include $\tilde\varphi_{\leq K}$ at this stage.

On the other hand, the inclusion of $\tilde\varphi_{\leq K}$ in the bilinear and trilinear operators is essential, as it ensures control of the growth of the factor $|k|^\beta$ in the nonlinear interactions.
We refer the reader to Section~\ref{sec:analysisL} for further details.
\end{remark}

\subsection{Statement of main estimates}\label{state} 
We define the $Z$ norm  
  for $\boldsymbol a(t)=(a_k(t))_{k \in \mathbb{Z}_L}$ by
\begin{equation*}
\|\boldsymbol a\|_{Z}^2 = \sup_{t\in[0,1]} L^{-1}\sum_{k \in \mathbb{Z}_L}\langle k\rangle^{10}|{a}_k(t)|^2 .\notag
\end{equation*}

  We recall that $N>0$ is a sufficiently large parameter to be fixed later.
To prove the derivation of the truncated equation, we will require estimates on the contributions of couples, as well as bounds for the linear operator $\mathscr{L}$.

\begin{proposition}[Bounds of tree couples]\label{prop:couple}
Let $1 \leq n \leq N^3$ and $t \in [0,1]$. For any function $\Gamma_1(k,t)$ satisfying
\[
\sup_k \sup_{t \in [0,1]} |\partial_t \Gamma_1(k,t)| \leq 1,
\]
we can define $c_{k}^{\mathcal{T}}=c_{k}^{\mathcal{T}}( \Gamma_1)$ as above, and it holds that
\begin{align}
\left|\sum_{|\mathcal{T}_{+}| + |\mathcal{T}_{-}| = n} \mathbf{E}\left[c_{k}^{\mathcal{T}_{+}}(t) c_{k}^{\mathcal{T}_{-}}(t)\right]\right| 
&\leq C \langle k \rangle^{-20} (\alpha T^{4/5} L^{2\delta})^n T^{-3/5}. \label{bd:couple}
\end{align}

Moreover, the right-hand side of \eqref{def:gamma1} defines a functional of $\Gamma_1$,  which we denote by $\mathscr{P}(\Gamma_1)$. For any 
\[
\sup_k \sup_{t \in [0,1]} |\partial_t \Gamma_1(k,t)| + |\partial_t {\Gamma}'_1(k,t)| \leq 1,
\]
it holds that
\begin{align*}
\sup_k \sup_{t \in [0,1]} |\mathscr{P}(\Gamma_1)| &\leq L^{-3\delta},\\
\sup_k \sup_{t \in [0,1]} |\mathscr{P}(\Gamma_1) - \mathscr{P}({\Gamma}'_1)| &\leq L^{-3\delta} \, \sup_k \sup_{t \in [0,1]} |\partial_t \Gamma_1 - \partial_t {\Gamma}'_1|.
\end{align*}
\end{proposition}

Then, by applying a fixed point argument, we obtain the existence of $\partial_t \Gamma_1$. Moreover, it holds that
\begin{align}
   \sup_k\sup_{t\in[0,1]}|   \partial_t\Gamma_1(k,t)|\leq L^{-3\delta}.\label{bd:gamma'}
\end{align}

\begin{proposition}[Bounds of linear operator $\mathscr{L}$]\label{prop:l}
  With probability $\geq 1-L^{-40}$, it holds that for $0\leq n\leq N$
\begin{equation}\label{defop} \|\mathscr{L}^n\|_{Z\to Z} \lesssim(\alpha  T^{4/5}  L^{2\delta})^{n/2} L^{50}. 
\end{equation}
\end{proposition}

These two propositions constitute the main results of the first part. The proof of Proposition~\ref{prop:couple} is presented in Sections~\ref{sec:irregular}--\ref{sec:Algorithm}, while the proof of Proposition~\ref{prop:l} is given in Section~\ref{bd:remainder}.

\subsection{ The molecules}\label{sec:molecule}
In order to establish the main estimates, we introduce the notion of molecules. Most of the notation follows \cite{DH23a,Vas24}; however, particular attention is required to handle the degenerate nodes, as in \cite{Wu25}.


Compared with the semilinear setting, the presence of degenerate nodes introduces additional constraints on the frequency assignments and pairing structures, which must be carefully incorporated into the definition of molecules. In particular, these nodes affect both the counting arguments and the available cancellations, and therefore require a refined treatment.

\begin{definition}\label{def-molecules}
\begin{enumerate}
\item (Molecules)
A  molecule $\mathbb{M}$ is a directed graph, where we call vertices as  atoms and edges as bonds. Multiple bonds between the same pair of atoms and loops are also allowed. Each atom has out-degree at most 2 and in-degree at most 2. The degree is the sum of out-degree and in-degree.
We write $v \in \mathbb{M}$ for an atom $v$ in $\mathbb{M}$, and $l \in \mathbb{M}$ for a bond $l$ in $\mathbb{M}$. We write $l \sim v$ If $v$ is an endpoint of $l$. 
We further require that $\mathbb{M}$ does not have any connected components consisting only of atoms each having degree 4 (where the degree is considered in the undirected sense).\\
For a molecule $\mathbb{M}$, we define the quantity
\begin{align*}
\chi := E - V + F,
\end{align*}
where $E$ is the number of bonds, $V$ is the number of atoms, and $F$ is the number of connected components.

\item (Atomic group) An atomic group in $\mathbb{M}$ is a subset of atoms together with some bonds connecting those atoms. 

\item (Bridge)
A single bond $l$ is called a bridge if removing it increases the number of connected components by one.
\item (Loop) A self-connecting bond is called a loop.
\end{enumerate}
\end{definition}


\begin{definition}[Molecules of couples]\label{couple-molecule}
Let $\mathcal{Q}$ be a nontrivial couple. We define the corresponding molecule $\mathbb{M} $ as follows. The atoms of $\mathbb{M}$ correspond to the branching nodes $\mathfrak{n} \in \mathcal{N}$ of $\mathcal{Q}$. For any two branching nodes $\mathfrak{n}_1, \mathfrak{n}_2$, we draw a bond between the corresponding atoms $v_1, v_2$ if one of the two case happens:
\begin{enumerate}
\item ($PC$) One of $\mathfrak{n}_1, \mathfrak{n}_2$ is a parent of the other, and if the child atom corresponds to a node with sign $-$, then the bond is directed  from the parent atom to the child atom. Otherwise, the direction is reversed;
\item ($LP$) A leaf child of $\mathfrak{n}_1$ is paired with a leaf child of $\mathfrak{n}_2$, the bond is directed from the atom whose corresponding child   has sign $-$  to  the other one with $+$.
%\item  Finally, we add one bond from the root with sign $+$ to the root with sign $-$; if one of the roots is a leaf then this bond should be between the other root and the parent node of the root paired with the leaf.
\end{enumerate}
 

%For any atom $v \in \mathbb{M}(\mathcal{Q})$, let $\mathfrak{n} = \mathfrak{n}(v)$ be its corresponding branching node in $\mathcal{Q}$. For any bond $l \sim v$, define $\mathfrak{m} = \mathfrak{m}(v,l)$ by:
%\begin{enumerate}
%\item If $l$ is PC, then $\mathfrak{m}$ is defined as the node of child.
%\item If $l$ is LP, then $\mathfrak{m}$ is the leaf that is a child of $\mathfrak{n}$.
%\end{enumerate}
\end{definition}

%Here we remark that, compared with \cite{DH23a,Vas24}, establishing this connection between the roots facilitates the discussion of the algorithm for degenerate atoms later.

\begin{definition}[Enhanced Molecule]
Let $\mathcal{Q}$ be an enhanced couple of order $n$. We construct its  enhanced molecule $\mathbb{M}(\mathcal{Q})$ in the same manner as Definition \ref{couple-molecule} with more restrictions:
\begin{enumerate}
    \item For any degenerate branching node $\mathfrak{n}\in\mathcal{N}_D$, there exists a corresponding atom $v_{\mathfrak{n}}$ in  $\mathbb{M}(\mathcal{Q})$.
    \item There are three atoms $v_c,v_2$, and $v_{nc}$ corresponding to the three children $\{\mathfrak{n}_c,\mathfrak{n}_2, \mathfrak{n}_{nc}\}$ of  $\mathfrak{n}$ (or their corresponding paired  nodes).
\end{enumerate}
\end{definition}
Then for any nontrivial couple $\mathcal{Q}$ of order $n$, the corresponding molecule $\mathbb{M}(\mathcal{Q})$ is connected.  It either has two degree 3 atoms or one degree 2
 atom, while all other atoms have degree 4. %  has   $n$ atoms and $2n$ bonds. Each atom has in-degree and out-degree 2.

The following lemma shows the correspondence between molecules and couples. The proof can be found in \cite[Proposition 9.6]{DH23a}, or \cite[Proposition 2.12]{Wu25}.
\begin{lemma}
\label{prop:molecules-couples}
 For any molecule $\mathbb{M}$ with order $n$, there are at most $C^n$ different couples $\mathcal{Q}$  such that the corresponding  molecule are exactly $\mathbb{M}$.
\end{lemma}



\begin{definition}[Decorations of Molecules]
\label{def:decorations-molecules}
Let $T>0$. Let $\mathbb{M}$ be a molecule. For each atom $v \in \mathbb{M}$, we fix $c_v \in \mathbb{Z}_L$,  and $C_v \in \mathbb{R}$.
A  $(c_v,C_v)$-decoration of $\mathbb{M}$ is a set of $k_l\in\mathbb{Z}_L $ to each  $l \in \mathbb{M}$, such  that for each atom $v$,
\begin{align*}
\sum_{l \sim v} \zeta_{v,l} \,k_l \;=\; c_v,
\quad\text{and}\quad
\bigl|\sum_{l: l\sim v} \zeta_{v,l} (|k_l|^\sigma+\Gamma_0(k_l)) - C_v\bigr|\;<\;T^{-1}.
\end{align*}
\end{definition}
 Here we define $\zeta_{v,l}=1$ if $l$ goes out from $v$, and $\zeta_{v,l}=-1$ if $l$ goes into $v$. We also recall the definition of $\Gamma_0$ given in \eqref{def:gamma1}.

In particular, for a nontrivial couple $\mathcal{Q}$, the related molecule $\mathbb{M}(\mathcal{Q})$, and a $k$-decorations of $\mathcal{Q}$, we define a $k$-decoration of $\mathbb M(\mathcal Q)$ such that $c_v=0$, and for any $l\in\mathbb{M}$,
\begin{enumerate}
\item 
 If $l$ is LP, between leaves $\mathfrak{l}$ and $\mathfrak{l}'$, then define $k_l := k_{\mathfrak{l}} = k_{\mathfrak{l}'}$.
 \item  If $l$ is PC, with child $\mathfrak{n}$, we define $k_l := k_{\mathfrak{n}}$.
% \item  If $l$ corresponds to the two roots, the define $k_l :=k $.
 \end{enumerate}
 The parameter $c_v$ in the decoration is usual taken as 
 \begin{equation*}%\label{bkeqn}
c_v=\begin{cases}
   0,& {\rm 
\ if}\ v {\rm\ has\ degree\ 2\ or\ 4},\\
 +k,&{\rm 
\ if}\ v {\rm\ has\ out-degree\ 2\ and\ in-degree\ 1},\\
 -k,& {\rm 
\ if}\ v {\rm\ has\ out-degree\ 1\ and\ in-degree\ 2}. 
\end{cases}
\end{equation*}
The choice of  $C_v$ will be decided in the following to derive different counting estimates.

At the end of this section, we introduce the notion of degenerate atoms.


\begin{definition}[Degenerate Atoms]
\label{def:degenerate-atoms}
Consider an enhanced molecule $\mathbb{M}$ equipped with a $(c_v,C_v)$-decoration as above. We say that an atom $v \in \mathbb{M}$ is degenerate if there exist two degenerate bonds $l_1, l_2 \sim v$, with opposite orientations at $v$, such that $k_{l_1}=k_{l_2}$. Furthermore, $v$ is fully degenerate if all bonds $l\sim v$ carry the same frequency $k_l$.
In particular, each degenerate branching node in the couple $\mathcal{Q}$ corresponds to a degenerate atom in the associated molecule $\mathbb{M}(\mathcal{Q})$. 
\end{definition}

\section{Collection of necessary estimates}\label{sec:mainest}
In this section, we collect several auxiliary results that will be used in the estimation of $\mathcal{K}_{\mathcal{Q}}$.
\subsection{The counting estimates}
 First, we have the following  two or three-vectors counting estimate.
\begin{lemma}\label{counting} Let $\sigma\in(0,1)\cup (1,2],\beta\in(0,\frac\sigma4)$ and  $D>1.$ Assume that $\alpha T^{1-\beta}D^{2-\sigma}<1$ and $D^2<T < L$.  Recall $\Gamma_0$ is defined in \eqref{def:gamma1}. Then, uniformly in $(k,a,b,c)\in(\mathbb{Z}_L)^4$ with $\max\{k,a,b,c\}\leq D$ and $m\in\mathbb{R}$, the sets
\begin{multline}%\label{3vectorest}
S_3=\big\{(x,y,z)\in (\mathbb{Z}_L)^3:x-y+z=k,\,\,\big||x|^\sigma+\Gamma_0(x)-|y|^\sigma-\Gamma_0(y)+|z|^\sigma+\Gamma_0(z)-m\big|\leq T^{-1},\\|x-a|\leq 1,\,\,|y-b|\leq 1,\,\,|z-c|\leq 1,\,\,\mathrm{and\ }k\not\in\{x,z\}\big\},\notag
\end{multline}
\begin{multline}%\label{2vectorest}\
S_2^+=\big\{(x,y)\in(\mathbb{Z}_L)^2:x+ y=k,\,\,\big||x|^\sigma+\Gamma_0(x)+|y|^\sigma+\Gamma_0(y)-m\big|\leq T^{-1},\\|x-a|\leq 1,\,\,|y-b|\leq 1\big\},\notag
\end{multline}
\begin{multline}%\label{2vectorest}\
S_2^-=\big\{(x,y)\in(\mathbb{Z}_L)^2:x- y=k,\,\,\big||x|^\sigma+\Gamma_0(x)-|y|^\sigma-\Gamma_0(y)-m\big|\leq T^{-1},\\\quad|x-a|\leq 1,\,\,|y-b|\leq 1\big\},\notag
\end{multline}satisfy the bounds
\begin{align*}
    \#S_3\lesssim L^{2}T^{-1}D^{2}\log L,\ \  \#S_2^+\lesssim LT^{-1/2}D,\ \  \#S_2^-\lesssim L.
    \end{align*}
    In particular, if $ |k|\geq T^{-1/2}$, it holds that
    \begin{align*}\#S_2^-\lesssim LT^{-1/2}D^{2}.
\end{align*} 
\end{lemma}

The proof follows by a similar  argument to in \cite[Section 3.8]{DIP25}.
\begin{proof}[Proof of $\#S_2^-$]
The bound $\#S_2^- \lesssim L$ follows immediately from a counting argument. Indeed, there are at most $O(L)$ possible choices of $x\in\mathbb{Z}_L$ satisfying $|x-a|\leq 1$, and once $x$ is fixed, $y = x - k$  is uniquely determined. ·

For the case $|k| \geq T^{-1/2}$, if either $|x| \leq T^{-1/2}$ or $|y| \leq T^{-1/2}$, the bound $L\cdot T^{-1/2}$ follows directly by counting the number of admissible choices of $x$ or $y$.
We therefore assume that $|x|,|y|\geq T^{-1/2}$. If $x$ and $y$ have the same sign, then by the mean value theorem, 
 \begin{align*}
        \bigg|\partial_x[|x|^{2\beta}-|x-k|^{2\beta}]\bigg|\lesssim \frac{|k|}{(\min\{|x|,|y|\})^{2-2\beta}}\lesssim |k|T^{1-\beta}.
    \end{align*} 
   If $x$ and $y$ have opposite signs (say $x>0$ without loss of generality), we observe that $$\partial_x(|x|^{2\beta}-|k-x|^{2\beta})=2\beta(x^{2\beta-1}+(k-x)^{2\beta-1})\lesssim  \frac{|k|}{(\min\{|x|,|y|\})^{2-2\beta}}.$$
    This implies that, using $D^2\leq T$,
    \begin{align*}
        \bigg|\partial_x[\varphi_{\leq K}(x)|x|^{2\beta}-\varphi_{\leq K}(x-k)|x-k|^{2\beta}]\bigg|\lesssim \frac{|k|}{(\min\{|x|,|y|\})^{2-2\beta}}\lesssim |k|T^{1-\beta}.
    \end{align*} 

   Arguing similarly for the dispersive term $|x|^\sigma$ (recalling that $\sigma\leq 2$), we obtain
    \begin{align}
       \bigg| \partial_x[||x|^\sigma+\Gamma_0(x)-|x-k|^\sigma-\Gamma_0(x-k)-m]\bigg|\geq \sigma|\sigma-1||k|D^{\sigma-2}-C\alpha |k|T^{1-\beta}\geq C|k|D^{\sigma-2},\label{bd:counting}
    \end{align}
    provided that $\alpha T^{1-\beta}D^{2-\sigma}<1$. Then, by the definition of $S_2^-$, a standard counting argument yields that the number of admissible values of $x$ is bounded by  $L\cdot T^{-1}\cdot |k|^{-1}D^{2-\sigma}\leq LT^{-1/2}D^{2-\sigma}$.
\end{proof}
\begin{proof}[Proof of $\#S_2^+$]
By the same argument as in the case of $S_2^-$, we may assume that $|x|,|y|\geq T^{-1/2}$. Then we have 
 \begin{align*}
        \bigg|\partial_x^2[\varphi_{\leq K}(x)|x|^{2\beta}+\varphi_{\leq K}(k-x)|k-x|^{2\beta}]\bigg|\lesssim \frac{1}{(\min\{|x|,|y|\})^{2-2\beta}}\lesssim T^{1-\beta},
    \end{align*}
    and for $\max\{|x|,|y|\}\leq D$, we have
    \begin{align*}
       \bigg| \partial_x^2[|x|^\sigma&+\Gamma_0(x)+|k-x|^\sigma+\Gamma_0(k-x)]\bigg|\\
       &\geq \sigma|\sigma-1|( |x|^{\sigma-2}+|k-x|^{\sigma-2})-C\alpha T^{1-\beta}\geq C\cdot D^{\sigma-2},
    \end{align*}
 provided that $\alpha T^{1-\beta}D^{2-\sigma}<1$.  Therefore, by a standard second-derivative counting argument, the number of admissible choices of $x$ is bounded by $L \cdot D^{1-\sigma/2} \cdot T^{-1/2}$.
\end{proof}
\begin{proof}[Proof of $\#S_3$] First, we consider the choice of $y$. If $|y|\leq T^{-1/2}$, then there are at most $L T^{-1/2}$ possible choices for $y$. For each such fixed $y$, by the bound on $\#S_2^+$, the number of pairs $(x,z)$ satisfying $x+z=k+y$ is at most $L T^{-1/2} D$. This yields the desired bound $L T^{-1/2}\cdot L T^{-1/2} D$.

Next, assume that $|y|\geq T^{-1/2}$. If both $|x|,|z|\leq T^{-1/2}$, then there are at most $L T^{-1/2}\cdot L T^{-1/2}$ possible choices.  Now assume that $|x|\geq T^{-1/2}$. For each fixed $z$, we consider the two-variable constraint   $x-y=k-z$ with $m$ replaced by $m-(|z|^\sigma+\Gamma_0(z))$.  Applying an argument analogous to \eqref{bd:counting}, we obtain that the number of admissible pairs $(x,y)$ is bounded by   $LT^{-1}|k-z|^{-1}D^{2-\sigma}$. Summing over all $z$, we obtain
$$\sum_{z}LT^{-1}|k-z|^{-1}D^{2-\sigma}\lesssim L^{2}T^{-1}D^{2-\sigma}\log L.$$ Here, the logarithmic factor $\log L$ arises from summing the weight $|k-z|^{-1}$ over $z\in\mathbb{Z}_L$.
\end{proof}


We also introduce the following convergence of iterates.

\begin{proposition}
    \label{asymptotic}
Let $T  < L^{\frac{(1+\gamma)(1+2\beta)}{2}-10\delta}$ if $\sigma\in(1,2]$, $T  <L^{\frac{1+2\beta}{2+2\beta-\sigma}-10\delta} $ if $\sigma\in(0,1)$, and $\phi \in \mathcal{S}(\mathbb{R})$. 
For any $t \in (0,T)$, let   \begin{align*}
 \mathscr S_t(\phi):&=\sum\limits_{\substack{ k-k_1+k_2-k_3=0,\\ \overline \Omega_k\neq0}}\phi_k \phi_{k_1} \phi_{k_2} \phi_{k_3}  \left[ \frac{1}{\phi_k} - \frac{1}{\phi_{k_1}} + \frac{1}{\phi_{k_2}} - \frac{1}{\phi_{k_3}} \right]\left| \frac{\sin(\pi t\overline \Omega_k)}{\pi t \overline \Omega_k} \right|^2\notag\\
  &\quad\quad\quad\times
|k_1|^{2\beta}|k_2|^{2\beta}|k_3|^{2\beta}|k|^{2\beta}\varphi_{\leq K}^2(k_1)\varphi_{\leq K}^2(k_2)\varphi_{\leq K}^2(k_3)\varphi_{\leq K}^2(k).\\
\mathscr K_t(\phi):&=L^{2} \int_{\xi_1-\xi_2+\xi_3=\xi }\phi(\xi) \phi(\xi_1) \phi(\xi_2) \phi(\xi_3)  \left[ \frac{1}{\phi(\xi)} - \frac{1}{\phi(\xi_1)} + \frac{1}{\phi(\xi_2)} - \frac{1}{\phi(\xi_3)} \right]    \left| \frac{\sin(\pi t\overline \Omega( \xi))}{\pi t\overline\Omega( \xi)} \right|^2 \notag\\ &\quad\quad\times |\xi_1|^{2\beta}|\xi_2|^{2\beta}|\xi_3|^{2\beta}|\xi|^{2\beta}\varphi_{\leq K}^2(\xi_1)\varphi_{\leq K}^2(\xi_2)\varphi_{\leq K}^2(\xi_3)\varphi_{\leq K}^2(\xi)   \dif \xi_1  \dif\xi_2 \, \dif\xi_3.
\end{align*}
 Then
$$
\mathscr S_t=\mathscr K_t +O(L^{2-\delta} t^{-1}).
$$
\end{proposition}

\begin{proof}
First, fix $k$ and define \begin{align*}
    g(x):&=\left|\frac{\sin \pi x}{\pi x}\right|^2\leq 1,\\
    F(k_1,k_3):&=\phi_k \phi_{k_1} \phi_{k_1+k_3-k} \phi_{k_3}  \left[ \frac{1}{\phi_k} - \frac{1}{\phi_{k_1}} + \frac{1}{\phi_{k_1+k_3-k}} - \frac{1}{\phi_{k_3}} \right]\\
    &\quad\quad\times
 |k_1|^{2\beta}|k_1+k_3-k|^{2\beta}|k_3|^{2\beta}|k|^{2\beta}\varphi_{\leq K}^2(k_1)\varphi_{\leq K}^2(k_1+k_3-k)\varphi_{\leq K}^2(k_3)\varphi_{\leq K}^2(k)\\
 &\quad\quad \times \chi(T^{\frac{1}{1+2\beta}}L^{2\delta}|k_1|)\chi(T^{\frac{1}{1+2\beta}}L^{2\delta}|k_1+k_3-k|)\chi(T^{\frac{1}{1+2\beta}}L^{2\delta}|k_3|),
\end{align*}
where   $\chi$ is a smooth cutoff satisfying $\chi(x)=0$ for $|x|\leq \frac12$ and $\chi(x)=1$ for $|x|\geq 1$. These truncations remove the singularity near $|k_i|=0$. 

Define the truncated quantities
\begin{align*}
   (\mathscr S_{\rm tr})_t&:=\sum_{(k_1, k_3)\in \mathbb{Z}_L^{2},\overline\Omega(k_1,k_3)\neq 0}F(k_1,k_3) g\left(t\overline \Omega(k_1,k_3)\right),\\
 (\mathscr K_{\rm tr})_t&:=L^{2} \int_{\xi_1,\xi_3}F(\xi_1,\xi_3) g\left(t\overline \Omega(\xi_1,\xi_3)\right)\dif \xi_1\dif \xi_3,
\end{align*}
where
$$\overline\Omega(x,y) =|x|^\sigma-|x+y-k|^\sigma+|y|^\sigma-|k|^\sigma+\Gamma_0(x)-\Gamma_0(x+y-k)+\Gamma_0(y)-\Gamma_0(k),$$
and
$ \Gamma_0(x)=C_K\frac{\alpha }{\pi }\varphi_{\leq K}(x)|x|^{2\beta}.$

In the first step,  for  the truncated asymptotic functions, we show that $$
( \mathscr S_{\rm tr})_t= (\mathscr K_{\rm tr})_t +O(L^{2-\delta} t^{-1}).
$$
Since $\widehat g(\tau)=1-|\tau|$ on $[-1,1]$ and vanishes outside, we write $e(x):=e^{2\pi i x}$. Applying Poisson summation, we obtain 
\begin{align*}
 (  \mathscr S_{\rm tr})_t&=\sum_{(k_1, k_3)}F(k_1,k_3) g\left(t\overline \Omega(k_1,k_3)\right)=\sum_{(k_1, k_3)}F(k_1,k_3)\int_{-\infty}^{\infty}\hat{g}(\tau)e(\tau t\overline \Omega(k_1,k_3))\dif \tau\\
 &=\frac1t \sum_{(k_1, k_3)}F(k_1,k_3)\int_{-\infty}^{\infty}\hat{g}(\frac\tau t)e(\tau \overline \Omega(k_1,k_3))\dif \tau=\frac1t \int_{-\infty}^{\infty}\hat{g}(\frac\tau t)\sum_{(k_1, k_3)}F(k_1,k_3)e(\tau \overline \Omega(k_1,k_3))\dif \tau\\
 &=\frac1t \int_{-\infty}^{\infty}\hat{g}(\frac\tau t)\sum_{(k_1, k_3)\in\mathbb{Z}^{2}}F(\frac{k_1}{L},\frac{k_3}L)e(\tau \overline \Omega(\frac{k_1}L,\frac{k_3}L))\dif \tau\\
 &=\frac1t \int_{-\infty}^{\infty}\hat{g}(\frac\tau t) \sum_{(c,d)\in\mathbb{Z}^{2}}\int_{(x,y)\in\mathbb{R}^{2}}F(\frac{x}{L},\frac{y}L)e(\tau \overline \Omega(\frac{x}L,\frac{y}L))e(-c\cdot x-d\cdot y) \dif x\dif y\dif \tau\\
  &=\frac1t \int_{-\infty}^{\infty}\hat{g}(\frac\tau t) \int_{(x,y)\in\mathbb{R}^{2}}F(\frac{x}{L},\frac{y}L)e(\tau \overline \Omega(\frac{x}L,\frac{y}L))\dif x\dif y\dif \tau=:  (\mathscr S_{\rm tr})_1\\
   &\quad+\frac{L^{2}}t \int_{-\infty}^{\infty}\hat{g}(\frac\tau t) \sum_{(c,d)\in\mathbb{Z}^{2}/\{0,0\}}\int_{(x,y)\in\mathbb{R}^{2}}F(x,y)e(\tau \overline \Omega(x,y)-c\cdot Lx-d\cdot Ly) \dif x\dif y\dif \tau=:  (\mathscr S_{\rm tr})_2.
\end{align*}

For the main term, a change of variables yields
\begin{align*}
  ( \mathscr S_{\rm tr})_1=\frac{L^{2}}t \int_{-\infty}^{\infty}\hat{g}(\frac\tau t) \int_{(x,y)\in\mathbb{R}^{2}}F(x,y)e(\tau \overline \Omega(x,y))\dif x\dif y\dif \tau= (\mathscr K_{\rm tr})_t.
\end{align*}

For the second term $(\mathscr S_{\rm tr})_2$, in the case $\sigma\in(1,2]$, it holds that by the cut off function $\varphi$ and $\chi$,
\begin{align*}
   |\tau  \nabla  \overline \Omega(x,y)|\lesssim T(L^{\delta(\sigma-1)}+\alpha (T^{\frac{1}{1+2\beta}}L^{2\delta})^{1-2\beta})\leq L^{1-\delta/2},
\end{align*}
where we used $|\tau|\leq T < \min\{L^{1-10\delta},L^{\frac{(1+\gamma)(1+2\beta)}{2}-10\delta}\}$.

While in the case $\sigma\in(0,1)$, we have 
\begin{align*}
   |\tau  \nabla  \overline \Omega(x,y)|\lesssim T((T^{\frac{1}{1+2\beta}}L^{2\delta})^{1-\sigma}+\alpha (T^{\frac{1}{1+2\beta}}L^{2\delta})^{1-2\beta})\leq L^{1-\delta/2},
\end{align*}
where we used $|\tau|\leq T < L^{\frac{1+2\beta}{2+2\beta-\sigma}-10\delta}$ and $2\beta<\sigma$. In summary, we have 
\begin{align*}
   |\tau  \nabla  \overline \Omega(x,y)| \leq L^{1-\delta/2}.
\end{align*}

Then for $c\neq 0$ we have
\begin{align*}
|   \nabla_x[ \tau \overline\Omega(x,y)-c\cdot Lx-d\cdot Ly]|\geq L|c|-CL^{1-\delta/2}\geq \frac12|c|L,
\end{align*}
as $L$ goes large enough.


Now, from the definition of $F(x,y)$,  once we apply integration by parts, 
we notice that by the cut off function $\varphi$ and $\chi$,  $$|\partial_{x_i}|x|^{2\beta}|\leq 2\beta|x|^{2\beta-1}\lesssim (TL^{2\delta})^{1-2\beta},$$ 
we obtain a power of $TL^{2\delta}\cdot (|c|L)^{-1}$ from $F(x,y)$. By  integrating enough times, we have a gain of $L^{-2\delta}$,  then we obtain $(\mathscr S_{\rm tr})_2 =O(L^{2-\delta} t^{-1}).$

It remains to estimate the difference between truncated and original quantities. By the support of $1-\chi$, at least one of $|k_1|,|k_1+k_3-k|,|k_3|$ is $\lesssim T^{-{\frac{1}{1+2\beta}}}L^{-2\delta}$, and hence
\begin{align*}
 | ( \mathscr K_{\rm tr})_t- \mathscr K_t|&\leq   L^{2} \int_{\xi_1,\xi_3}\phi_\xi \phi_{\xi_1} \phi_{\xi_1+\xi_3-\xi} \phi_{\xi_3}  \left[ \frac{1}{\phi_\xi} - \frac{1}{\phi_{\xi_1}} + \frac{1}{\phi_{\xi_1+\xi_3-\xi}} - \frac{1}{\phi_{\xi_3}} \right](T^{-{\frac{1}{1+2\beta}}}L^{-2\delta} L^{\delta})^{1+2\beta}\dif \xi_1\dif \xi_3\\
 &=O(L^2(T^{-{\frac{1}{1+2\beta}}}L^{-\delta})^{1+2\beta})\leq O(L^{2-\delta}T^{-1})\leq O(L^{2-\delta}t^{-1}).
\end{align*} The same argument yields
\begin{align*}
 | (\mathscr S_{\rm tr})_t- \mathscr S_t|&\leq  O(L^{2-\delta}t^{-1}).
\end{align*}Combining the above estimates, we conclude the proof.
\end{proof}


\subsection{The integral estimates on $\mathcal{A}_{\mathcal{T}}$}
In this section, we estimate the time integral $\mathcal{A}_{\mathcal{T}}$ defined in \eqref{def:a}. We rewrite it in the form 
\begin{align*}
 \mathcal{A}_{\mathcal{T}}=   \int_{\mathcal{D}}\prod_{\mathfrak{n}\in\mathcal{N}}  e^{\zeta_\mathfrak{n}2\pi i[\Omega_\mathfrak{n}Tt_\mathfrak{n}+\sum\Gamma(\mathfrak{n},t_\mathfrak{n})] } \dif t_\mathfrak{n},
\end{align*}
where the time integration domain is
\begin{align*}
\mathcal{D} := \{t{[\mathcal{N}]} : 0 < t_{\mathfrak{n}'} < t_\mathfrak{n}<t  \text{ whenever } \mathfrak{n}' \text{ is a child of } \mathfrak{n} \}.
\end{align*}
We write $$\overline\Omega_\mathfrak{n}:=\Omega_\mathfrak{n}+\sum\Gamma_0(k): =\Omega_\mathfrak{n}+\Gamma_0(k_1)-\Gamma_0(k_2)+\Gamma_0(k_3)-\Gamma_0(k),$$
and  then have 
$$\Omega_\mathfrak{n}Tt_\mathfrak{n}+\sum\Gamma(\mathfrak{n},t_\mathfrak{n})=\overline\Omega_\mathfrak{n}Tt_\mathfrak{n}+\sum\Gamma_1(\mathfrak{n},t_\mathfrak{n}).$$
We  emphasize that the notation $\overline{\Omega}_{\mathfrak{n}}$ does not denote complex conjugation, but rather the renormalized dispersion relation defined above.

Then we have the following representation:
\begin{align*}
 \mathcal{A}_{\mathcal{T}}(t,\Theta[\mathcal{N}])=   \int_{\mathcal{D}}\prod_{\mathfrak{n}\in\mathcal{N}}  e^{\zeta_\mathfrak{n}2\pi i[\overline\Omega_\mathfrak{n}Tt_\mathfrak{n}+\sum\Gamma_1(\mathfrak{n},t_\mathfrak{n})] } \dif t_\mathfrak{n}.
\end{align*}
\begin{definition}
     We  fix $d_{\mathfrak{n}}\in\{0,1\}$ for each $\mathfrak{n}\in\mathcal{N}$, we can define $q_{\mathfrak{n}}$ for each $\mathfrak{n}\in\mathcal{T}$ inductively by
\begin{equation}%\label{defqn}
q_{\mathfrak{n}}=0\mathrm{\ if\ }\mathfrak{n}\in\mathcal{L};\quad q_{\mathfrak{n}}=d_{\mathfrak{n}_1}q_{\mathfrak{n}_1}-d_{\mathfrak{n}_2}q_{\mathfrak{n}_2}+d_{\mathfrak{n}_3}q_{\mathfrak{n}_3}+\overline\Omega_{\mathfrak{n}}\mathrm{\ if\ }\mathfrak{n}\in\mathcal{N}.\end{equation}
\end{definition}
\begin{proposition}\label{prop:est:af}
Let $\mathcal{T}$ be a nontrivial tree, and we assume there are some functions  $f_{\mathfrak{n}}:[0,1]\to \mathbb{C}$ satisfying $$\sup_{s\in[0,1]}(|f_{\mathfrak{n}}(s)|+|(\partial_tf_{\mathfrak{n}})(s)|)\leq C_{\mathfrak{n}}.$$
We define
\begin{align}
 \mathcal{A}^f_{\mathcal{T}}(t,\overline\Omega[\mathcal{N}],\Gamma_1[\mathcal{N}]):=   \int_{\mathcal{D}}\prod_{\mathfrak{n}\in\mathcal{N}}  f_{\mathfrak{n}}(t_{\mathfrak{n}})e^{\zeta_\mathfrak{n}2\pi i[\overline\Omega_\mathfrak{n}Tt_\mathfrak{n}+\sum\Gamma_1(\mathfrak{n},t_\mathfrak{n})] } \dif t_\mathfrak{n}.\label{def:af}
\end{align}
 
Then for any function $\Gamma_1(k,t)$ satisfying $$\sup_k\sup_{t\in[0,1]}|\partial_t\Gamma_1(k,t)|\leq 1,$$
it holds that
\begin{align*}
    |\mathcal{A}^f_{\mathcal{T}}(t,\overline\Omega[\mathcal{N}],\Gamma_1[\mathcal{N}])|\lesssim \sum_{d_\mathfrak{n}\in\{0,1\}}\prod_{\mathfrak{n}\in\mathcal{N}}C_\mathfrak{n}\langle Tq_\mathfrak{n}\rangle^{-1}.
\end{align*}
Here the universal constant is universal and independent of choice of $\Gamma_1$.
\end{proposition}

\begin{proof}
  We prove this result by induction.
For $|\mathcal{T}|=1$, we may assume without loss of generality that $\zeta_\mathfrak{r}=+$. Then, for the case $|T\overline\Omega_k|\geq 1$, we have
   \begin{align*}
        \int_0^t &  f_{\mathfrak{r}}(s)e^{ 2\pi i[\overline\Omega_kTs+\sum\Gamma_1(k,s)] } \dif s\\
        &= \frac{f_{\mathfrak{r}}(t)e^{ 2\pi i[\overline\Omega_kTt+\sum\Gamma_1(k,t)] }- f_{\mathfrak{r}}(0)}{2\pi i\overline\Omega_kT}\\
        &\quad - \int_0^t \frac{f'_{\mathfrak{r}}(s)e^{ 2\pi i[\overline\Omega_kTs+\sum\Gamma_1(k,s)] }+f_{\mathfrak{r}}(s)e^{ 2\pi i[\overline\Omega_kTs+\sum\Gamma_1(k,s)] }2\pi i\sum \Gamma_1'(k,s)}{2\pi i\overline\Omega_kT}   \dif s.
   \end{align*}
  Using the assumptions on $f_{\mathfrak{r}}$ and $\Gamma_1$, we obtain 
  \begin{align*}
    |\mathcal{A}^f_{\mathcal{T}}(t,\overline\Omega[\mathcal{N}],\Gamma_1[\mathcal{N}])|\lesssim C_\mathfrak{n}| T\overline\Omega_k|^{-1}\lesssim C_\mathfrak{n}\langle T\overline\Omega_k\rangle^{-1}= C_\mathfrak{n}\langle Tq_\mathfrak{n}\rangle^{-1}.
\end{align*}
Moreover, in the case $|T\overline\Omega_k|\leq 1$, the bound follows directly from
   \begin{align*}
      \left|  \int_0^t   f_{\mathfrak{r}}(s)e^{ 2\pi i \sum\Gamma_1(k,s) } \dif s\right|\leq C_{\mathfrak{r}}.
   \end{align*}

From now on, we assume that the estimate holds for all trees of order at most $n$.
Let $\mathcal{T}$ be a tree with $|\mathcal{T}|=n+1$. Then there exists a branching node $\mathfrak{n}_0$ whose three children are leaves. Moreover, we can choose two sibling nodes (denoted by $\mathfrak{n}_1$ and $\mathfrak{n}_2$) such that each of them is either a leaf or has only leaves as children. We denote the parent of $\mathfrak{n}_0$ by $\mathfrak{n}_p$. The subtrees rooted at $\mathfrak{n}_i$, $i=0,1,2$, are denoted by $\mathcal{T}_i$, respectively. With this decomposition, we can write
   \begin{align*}
       \mathcal{A}^f_{\mathcal{T}}=   \int_{\mathcal{D}_1}\prod_{\mathfrak{n}\in\mathcal{N}_1}  f_{\mathfrak{n}}(t_{\mathfrak{n}})e^{\zeta_\mathfrak{n}2\pi i[\overline\Omega_\mathfrak{n}Tt_\mathfrak{n}+\sum\Gamma_1(\mathfrak{n},t_\mathfrak{n})] } \dif t_\mathfrak{n}%\int_0^{t_{\mathfrak{n}_p}}f_{\mathfrak{n}_0}(t_{\mathfrak{n}_0})e^{\zeta_{\mathfrak{n}_0}2\pi i[\Omega_{\mathfrak{n}_0}Tt_{\mathfrak{n}_0}+\Gamma(\mathfrak{n}_0,t_{\mathfrak{n}_0})] } dt_{\mathfrak{n}_0} 
       \mathcal{A}^f_{\mathcal{T}_0}(t_{\mathfrak{n}_p})\mathcal{A}^f_{\mathcal{T}_1}(t_{\mathfrak{n}_p})\mathcal{A}^f_{\mathcal{T}_2}(t_{\mathfrak{n}_p}) ,
   \end{align*}
   where $\mathcal{N}_1$ are the set of branching nodes without the branching nodes in $\mathcal{T}_i,i=0,1,2$ and  \begin{align*}
\mathcal{D}_1 = \{t{[\mathcal{N}_1]} : 0 < t_{\mathfrak{n}'} < t_\mathfrak{n}<t  \text{ whenever } \mathfrak{n}' \text{ is a child of } \mathfrak{n} \}.
\end{align*}

First, we assume that $|T\overline\Omega_{\mathfrak{n}_i}|\geq 1$ for $i=0,1,2$ (if $\mathcal{T}_i$ is nontrivial). Then, for each subtree $\mathcal{T}_i$, the quantity   $\mathcal{A}^f_{\mathcal{T}_i}$ is either equal to $1=\langle Tq_{\mathfrak{n}_i}\rangle^{-1}$ in the trivial case $|\mathcal{T}_i|=0,q_{\mathfrak{n}_i}=0$, or can be written in the form:
\begin{align*}
    \frac{f_{\mathfrak{n}_i}(t_{\mathfrak{n}_p})e^{ \zeta_{\mathfrak{n}_i}2\pi i[\overline\Omega_{\mathfrak{n}_i}Tt_{\mathfrak{n}_p}+\sum\Gamma_1(\mathfrak{n}_i,t_{\mathfrak{n}_p})] }}{2\pi i\overline\Omega_{\mathfrak{n}_i}T}+\frac{\tilde{f}_{\mathfrak{n}_i}(t_{\mathfrak{n}_p})}{2\pi i\overline\Omega_{\mathfrak{n}_i}T},
\end{align*}
where the function $\tilde{f}_{\mathfrak{n}_i}$ arises from integration by parts and satisfies $$\sup_{s\in[0,1]}(|\tilde{f}_{\mathfrak{n}_i}(s)|+|\tilde{f}'_{\mathfrak{n}_i}(s)|)\lesssim |f_{\mathfrak{n}_i}(0)| +\sup_{s\in[0,1]}|f'_{\mathfrak{n}_i}(s)| +\sup_{s\in[0,1]}|f_{\mathfrak{n}_i}(s)\sum\Gamma_1'({\mathfrak{n}_i},s)|\lesssim C_{\mathfrak{n}_i}.$$
Then we  deduce that \begin{align*}
       \mathcal{A}^f_{\mathcal{T}}=\prod_{i:|\mathcal{T}_i|=1}\frac{1}{2\pi i\overline\Omega_{\mathfrak{n}_i}T}\sum_{d_i,\overline  f_{\mathfrak{n}_p}}   \int_{\mathcal{D}_1}\prod_{\mathfrak{n}\in\mathcal{N}_1}  \overline f_{\mathfrak{n}}(t_{\mathfrak{n}})e^{\zeta_\mathfrak{n}2\pi i[\tilde\Omega_\mathfrak{n}Tt_\mathfrak{n}+\sum \tilde \Gamma_1(\mathfrak{n},t_\mathfrak{n})] } \dif t_\mathfrak{n},
   \end{align*}
where the sum is taken over all choices $d_i\in\{0,1\}$ for those indices $i$ with $|\mathcal{T}_i|= 1$, corresponding to whether integration by parts is applied to $\mathcal{T}_i$. The quantities $\tilde\Omega_{\mathfrak{n}}$ and $\tilde\Gamma_1$ are defined as before, except at the node $\mathfrak{n}_p$, where  $$\tilde \Omega_{\mathfrak{n}_p}=q_{\mathfrak{n}_p},\  \sum \tilde{\Gamma}_1=\sum \Gamma_1(k_{\mathfrak{n}_p})+\sum_{i}d_i\sum \Gamma_1(k_{\mathfrak{n}_i}).$$
   The functions $\overline f_{\mathfrak{n}}$ are defined as before, with the exception of the node $\mathfrak{n}_p$, where $\overline f_{\mathfrak{n}_p}$ is given by a product of functions among $f_{\mathfrak{n}_p}$, $f_{\mathfrak{n}_i}$, and $\tilde{f}_{\mathfrak{n}_i}$, depending on the choice of $d_i$.
 
  Since the new tree $\mathcal{T}'$, obtained by removing the nodes in $\mathcal{T}_i$, $i=0,1,2$, satisfies $|\mathcal{T}'|\leq n$, we may apply the induction hypothesis to obtain
\begin{align*}
    |\mathcal{A}^f_{\mathcal{T}}|\lesssim\prod_{i} \frac{1}{\langle Tq_i\rangle}\sum_{d_\mathfrak{n}\in\{0,1\}}\prod_{\mathfrak{n}\in\mathcal{N}_1}C_\mathfrak{n}\langle T\tilde q_\mathfrak{n}\rangle^{-1}\lesssim \sum_{d_\mathfrak{n}\in\{0,1\}}\prod_{\mathfrak{n}\in\mathcal{N}}C_\mathfrak{n}\langle Tq_\mathfrak{n}\rangle^{-1}.
\end{align*}
Here, $\tilde q_\mathfrak{n}$ is defined analogously to $q_\mathfrak{n}$, with $\overline{\Omega}_{\mathfrak{n}}$ replaced by $\tilde{\Omega}_{\mathfrak{n}}$.
We also note that in the case $|\mathcal{T}_i|=0$, one has $q_{\mathfrak{n}_i}=0$, and hence $\langle Tq_{\mathfrak{n}_i}\rangle^{-1}=1$, so that the above estimate remains valid.
 
Finally, if there exists some $ |T\overline{\Omega}_{\mathfrak{n}_i}| \leq 1 $ for $ i = 0,1,2 $,
the proof proceeds in the same manner with even simpler estimates.
In this case, we simply bound the corresponding factor involving $T\overline{\Omega}_{\mathfrak{n}_i}$ by $1$, and the rest of the argument remains unchanged.
\end{proof}

In particular, by taking $f_{\mathfrak{n}}\equiv 1$ for all $\mathfrak{n}\in\mathcal{N}$, we obtain the following estimate for the integral in \eqref{def:a}:
\begin{proposition}\label{prop:est:a}
Let $\mathcal{T}$ be a nontrivial tree. 
Then for any function $\Gamma_1(k,t)$ satisfying $$\sup_k\sup_{t\in[0,1]}|\partial_t\Gamma_1(k,t)|\leq 1,$$
it holds that
\begin{align*}
  |\mathcal{A}_{\mathcal{T}}(t,\Theta[\mathcal{N}])|\lesssim \sum_{d_\mathfrak{n}\in\{0,1\}}\prod_{\mathfrak{n}\in\mathcal{N}}\langle Tq_\mathfrak{n}\rangle^{-1}.
\end{align*}
Here the universal constant is universal and independent of choice of $\Gamma_1(k,t)$.
\end{proposition} 



\section{The Irregular Chains}\label{sec:irregular}
In this section, we introduce the notion of irregular chains, first introduced in \cite{DH23a}.
These chains give rise to the worst two-vector estimates, and our goal is to exploit a delicate cancellation mechanism among them.
Compared with previous works, an additional difficulty here comes from the symbols generated by derivatives, which we control by the
truncation functions.

\subsection{Cancellation of irregular chains}\label{irregular}


\begin{definition} [Irregular Chains]\label{def:chain-couple}
 Consider an enhanced couple $\mathcal{Q}=\{\mathcal{T}^+,\mathcal{T}^-,\mathcal{P},\mathcal{N}_D\}$. A sequence of nodes $(\mathfrak{n}_0, \ldots, \mathfrak{n}_q)$ is called a  irregular chain if, for each $0 \leq j \leq q-1$, 
\begin{enumerate}
    \item $\mathfrak{n}_{j+1}$ is a child of $\mathfrak{n}_{j}$, and the other two children of $\mathfrak{n}_{j}$ are both leaves,
    \item $\mathfrak{n}_{j}$ has a child $\mathfrak{m}_{j}$ paired with a child $\mathfrak{p}_{j+1}$ of $\mathfrak{n}_{j+1}$, and $\mathfrak m_i$ has opposite sign as $\mathfrak n_{i+1}$,
    \item  we denote $\mathfrak p_0$ to be the remaining child of $\mathfrak n_0$, and 
   $\mathfrak e$ and $\mathfrak f$ to be the remaining two children of $\mathfrak n_q$ with sign $+,-$ respectively.

\end{enumerate}
\end{definition}See Figure \ref{fig:irregular} for examples of two irregular chains. In this configuration, we note that degenerate nodes may also be present.

Let $\mathcal{H}=(\mathfrak{n}_0, \ldots, \mathfrak{n}_q)$ be an irregular chain. By the  definition of irregular chain, $\mathfrak{p}_i$ has the same sign as $\mathfrak{n}_i$ so there are only two choices (the left node or the right node). We denote $\zeta_{\mathfrak{n}_i}=\zeta_i$. We fix the degeneracy of each atoms $\mathfrak{n}_0, \ldots, \mathfrak{n}_q$. Then if we fix  these relative positions and $\zeta_0$, the structure of $\mathcal{H}$ is uniquely
 determined by the signs of $\zeta_i,1\leq i\leq q$.
 
\begin{definition}\label{def:twist}

\begin{enumerate}
\item (Congruent relation) 
 We define two irregular chains to be congruent (denoted $\mathcal{H}'\equiv \mathcal{H}$), if the sign $\zeta_0$ are the same and  the relative
 position of each $\mathfrak{p}_j$  are the same (both left or right) for these two irregular chains. Then the  the irregular chains $\mathcal{H}$  in this congruence class is uniquely determined
 by the signs $\zeta_{i},1\leq i\leq q$.
  We  call any two congruent irregular chains in a congruence class as twists of each  other.  See Figure \ref{fig:irregular} for an example of two irregular chains as twists.
 \item (Congruent couple) We can consider an irregular chain $\mathcal{H}$ in a couple $\mathcal{Q}$ and all its
 twists. Note that if we twist $\mathcal{H}$, then this twist can be trivially extended to $\mathcal{Q}$, as long as we
 keep the remaining part of the couple $\mathcal{Q}$ unchanged (note that the positions of $\mathfrak{e}$ and $\mathfrak{f}$ may be
 switched, but the sub-trees attached at $\mathfrak{e}$ and $\mathfrak{f}$ are unchanged).  For a given couple $\mathcal{Q}$ and we  suppose that $\mathcal{H}_i,1\leq i\leq r$ are disjoint irregular chains in $\mathcal{Q}$. Then we
  perform twists separately and arbitrarily to  each $\mathcal{H}_i$, starting from $\mathcal{Q}$, and obtain a family of couple $\mathcal{Q}'$, and call them congruent couples, with notation $\mathcal{Q}\equiv\mathcal{Q}'$.  
  \item (Decoration relation) We fix one irregular chain $\mathcal{H}$ in a couple $\mathcal{Q}$, and consider its decoration projected on $\mathcal{H}$.  Then the
  decoration on $\mathcal{H}$ can be uniquely determined by the vectors
 $(k_i, l_i), 0 \leq i\leq q +1$, which are defined as follows:\\ for   $0 \leq i\leq q $ we have $(k_i,l_i)=(k_{\mathfrak{n}_i},k_{\mathfrak{p}_i})$ if $\zeta_i=+$, and $(k_i,l_i)=(k_{\mathfrak{p}_i},k_{\mathfrak{n}_i})$ if $\zeta_i=-$; for $j=q+1$ we denote $(k_{q+1},l_{q+1}) = (k_{\mathfrak{e}},k_{\mathfrak{f}})$.  See Figure \ref{fig:decoration} for the corresponding decorations of irregular chains in Figure \ref{fig:irregular}.
 \item (Gaps) By direct observation,   we have  \begin{align}
     k_0-l_0=k_1-l_1=\cdots=k_{q+1}-l_{q+1}:=h.
 \end{align} 
 We call the irregular chain a large gap (LG) or small gap (SG) provided $h> T^{-1/2}$ or $h\leq T^{-1/2}$ respectively.
 \end{enumerate}
\end{definition}




\begin{figure}
\begin{subfigure}
\centering
\begin{tikzpicture}[level distance=1.2cm,
  level 1/.style={sibling distance=1.2cm},
  level 2/.style={sibling distance=1.2cm}]
\tikzstyle{hollow node}=[circle,draw,inner sep=1.6]
\tikzstyle{solid node}=[circle,draw,inner sep=1.6,fill=black]
\tikzset{
red node/.style = {circle,draw=black,fill=red,inner sep=1.6},
blue node/.style= {circle,draw = black, fill= blue,inner sep=1.6}, 
purple node/.style= {circle,draw = black, fill= purple,inner sep=1.6}, 
orange node/.style= {circle,draw = black, fill= orange,inner sep=1.6},
yellow node/.style= {circle,draw = black, fill= yellow,inner sep=1.6},
green node/.style = {circle,draw=black,fill=green,inner sep=1.6}}
\node[solid node, label = right:{$\mathfrak{n}_0(+)$}] at (-5.5,.5){}
    child{node[solid node, label = right: {$\mathfrak{n}_1(+)$}]{}
        child{node[solid node, label=below: $\mathfrak{n}_2(+)$]{}
             child{node[solid node, label=below: $\mathfrak{e}(+)$]{}}
             child{node[solid node, label=below: $\mathfrak{f}(-)$]{}}
             child{node[green node, label=below: $\mathfrak{p}_2(+)$]{}}
        }
        child{node[green node, label=below: $\mathfrak{m}_1(-)$]{}}
        child{node[red node, label=below: $\mathfrak{p}_1(+)$]{}}
    }
    child{node[red node, label=below: $\mathfrak{m}_0(-)$]{}}
    child{node[solid node, label=below: $\mathfrak{p}_0(+)$]{}}
    
;


\node[solid node, label = right: {$\mathfrak{n}_0(+)$}] at (1,0.5){}
    child{node[red node, label=below: $\mathfrak{m}_0(+)$]{}}
    child{node[solid node, label = right: {$\mathfrak{n}_1(-)$}]{}
        child{node[green node, label=below: $\mathfrak{m}_1(-)$]{}}
        child{node[solid node, label=below: $\mathfrak{n}_2(+)$]{}
             child{node[solid node, label=below: $\mathfrak{e}(+)$]{}}
             child{node[solid node, label=below: $\mathfrak{f}(-)$]{}}
             child{node[green node, label=below: $\mathfrak{p}_2(+)$]{}}
             }
        child{node[red node, label=below: $\mathfrak{p}_1(-)$]{}}
    }
    child{node[solid node, label=below: $\mathfrak{p}_0(+)$]{}}
    
;
\end{tikzpicture}
\end{subfigure}
\caption{Two examples of congruent chains. Points with the same color are paired.}
\label{fig:irregular}
\end{figure}



\begin{figure}
\begin{subfigure}%[t]{1\linewidth}
\centering
\begin{tikzpicture}[level distance=1.2cm,
  level 1/.style={sibling distance=1.2cm},
  level 2/.style={sibling distance=1.2cm}]
\tikzstyle{hollow node}=[circle,draw,inner sep=1.6]
\tikzstyle{solid node}=[circle,draw,inner sep=1.6,fill=black]
\tikzset{
red node/.style = {circle,draw=black,fill=red,inner sep=1.6},
blue node/.style= {circle,draw = black, fill= blue,inner sep=1.6}, 
purple node/.style= {circle,draw = black, fill= purple,inner sep=1.6}, 
orange node/.style= {circle,draw = black, fill= orange,inner sep=1.6},
yellow node/.style= {circle,draw = black, fill= yellow,inner sep=1.6},
green node/.style = {circle,draw=black,fill=green,inner sep=1.6}}
\node[solid node, label = right:{$k_0(+)$}] at (-5.5,.5){}
    child{node[solid node, label = right: {$k_1(+)$}]{}
        child{node[solid node, label=below: $k_2(+)$]{}
             child{node[solid node, label=below: $k_3(+)$]{}}
             child{node[solid node, label=below: $l_3(-)$]{}}
             child{node[green node, label=below: $l_2(+)$]{}}
        }
        child{node[green node, label=below: $l_2(-)$]{}}
        child{node[red node, label=below: $l_1(+)$]{}}
    }
    child{node[red node, label=below: $l_1(-)$]{}}
    child{node[solid node, label=below: $l_0(+)$]{}}
    
;


\node[solid node, label = right: {$k_0(+)$}] at (1,0.5){}
    child{node[red node, label=below: $k_1(+)$]{}}
    child{node[solid node, label = right: {$l_1(-)$}]{}
        child{node[green node, label=below: $l_2(-)$]{}}
        child{node[solid  node, label=below: $k_2(+)$]{}
             child{node[solid node, label=below: $k_3(+)$]{}}
             child{node[solid node, label=below: $l_3(-)$]{}}
             child{node[green node, label=below: $l_2(+)$]{}}
             }
        child{node[red node, label=below: $k_1(-)$]{}}
    }
    child{node[solid node, label=below: $l_0(+)$]{}}
    
;
\end{tikzpicture}
\end{subfigure}
\caption{The corresponding decorated irregular chains are illustrated in Figure \ref{fig:irregular}.}\label{fig:decoration}
\end{figure}



In particular, for any congruent chain, we adopt the shorthand notation $\Omega_i=\Omega_{\mathfrak{n}_i}, \ \sum\Gamma(i,t)=\sum\Gamma(\mathfrak{n}_i,t)$  and $\Theta(i,t)=\Theta(\mathfrak{n}_i,t)$ for any congruent chains  for short. Moreover,  the expressions of   $\zeta_i\Omega_{ i}$ and $\zeta_i\sum\Gamma(i,t)$ are independent of the choice of $\zeta_i$: $$\zeta_i\Omega_i=|k_i|^\sigma- |k_i -h|^\sigma + |k_{i+1}- h|^\sigma-|k_{i+1}|^\sigma,$$
and
$$ \zeta_i\sum\Gamma(i,t)=\Gamma(k_i,t)- \Gamma(k_i -h,t)+ \Gamma(k_{i+1}- h,t)- \Gamma(k_{i+1},t).$$

 \begin{definition}
      (Splicing) Consider a couple $\mathcal{Q}$ and an irregular chain  $\mathcal{H}=(\mathfrak{n}_0,...,\mathfrak{n}_q)$ in it. We apply Splicing to this chain is to remove those nodes $(\mathfrak{n}_1,...,\mathfrak{n}_q)$ and their leaf
children $(\mathfrak{m}_0,...,\mathfrak{m}_{q-1},\mathfrak{p}_0,...,\mathfrak{p}_{q-1})$ from $\mathcal{Q}$ and redefining the children of $\mathfrak{n}_0$ by $\mathfrak{p}_0,\mathfrak{e}$ and $\mathfrak{f}$
  with their position determined by  their
 relative positions.  See Figure \ref{fig:splicing} for an example. We denote the rest couple by $\mathcal{Q}^{sp}$.
 \end{definition}

\begin{figure}
\begin{subfigure}%[t]{1\linewidth}
\centering
\begin{tikzpicture}[level distance=1.2cm,
  level 1/.style={sibling distance=1.2cm},
  level 2/.style={sibling distance=1.2cm}]
\tikzstyle{hollow node}=[circle,draw,inner sep=1.6]
\tikzstyle{solid node}=[circle,draw,inner sep=1.6,fill=black]
\tikzset{
red node/.style = {circle,draw=black,fill=red,inner sep=1.6},
blue node/.style= {circle,draw = black, fill= blue,inner sep=1.6}, 
purple node/.style= {circle,draw = black, fill= purple,inner sep=1.6}, 
orange node/.style= {circle,draw = black, fill= orange,inner sep=1.6},
yellow node/.style= {circle,draw = black, fill= yellow,inner sep=1.6},
green node/.style = {circle,draw=black,fill=green,inner sep=1.6}}
\node[solid node, label = right:{$\mathfrak{n}_0(+)$}] at (-5.5,.5){}
    child{node[solid node, label=below: $\mathfrak{e}(+)$]{}}
    child{node[solid node, label=below: $\mathfrak{f}(-)$]{}}
    child{node[solid node, label=below: $\mathfrak{p}_0(+)$]{}}
;


\node[solid node, label = right: {$k_0(+)$}] at (1,0.5){}
    child{node[solid node, label=below: $k_3(+)$]{}}
    child{node[solid node, label=below: $l_3(-)$]{}}
    child{node[solid node, label=below: $l_0(+)$]{}}   
;
\end{tikzpicture}
\end{subfigure}
\caption{The congruent couple in Figure \ref{fig:irregular} after splicing}
\label{fig:splicing}
\end{figure}


\begin{remark}
 For a non-degenerate atom $\mathfrak{n}_0$, after splicing, it may become degenerate. In this case, the two subtrees rooted at $\mathfrak{n}_2$ and $\mathfrak{n}_{nc}$ may become completely paired, so that the resulting couple is no longer enhanced in the sense of Definition \ref{def:couple}.
For example, in Figure \ref{fig:irregular} (or Figure \ref{fig:splicing}), if $k_0=k_3$ and $l_0=l_3$, then the node $\mathfrak{n}_0$ becomes degenerate, while the leaves in the subtree rooted at ${\mathfrak{f}_2, \mathfrak{p}_{0}}$ may be completely paired.

Similarly, for an originally degenerate node, the splicing procedure may also destroy the enhanced structure. We refer to such nodes as bad degenerate nodes.

\end{remark}  

 It is easy to see that for any $\mathcal{Q}'\equiv\mathcal{Q}$, one has $\mathcal{Q}^{sp}=\mathcal{Q}'^{sp}$.
We then define 
 \begin{align}
     \mathcal{G}_\mathcal{Q}(t,s,k):=\sum_{\mathcal{Q}'\equiv\mathcal{Q}}\mathcal{K}_\mathcal{Q}(t,s,k),\label{def:gq}
 \end{align}
 where the sum is taken over all congruent couples $\mathcal{Q}\equiv\mathcal{Q}'$.  
\begin{lemma}\label{lem-splice}
For an enhanced couple $\mathcal Q$ and one SG irregular chain $(\mathfrak n_0, \ldots, \mathfrak n_{q})$ in $\mathcal{Q}$, 
\begin{align*}
    \mathcal{G}_\mathcal{Q}(t,s,k) :&= \bigg(\frac{\alpha T}{L}\bigg)^{n_{sp}} \zeta(\mathcal{Q}^{sp}) \sum_{\mathcal{E}^{sp}}  %\tilde{\epsilon}_{\mathcal{E}^{sp}}
    \prod_{\mathfrak{n}\in\mathcal{N}^{sp}}B_\mathfrak{n}\prod_{\mathfrak{n}\in\mathcal{Q}^{sp}}\varphi_\mathfrak{n}
\int_0^1\int_{\mathcal{E}^{sp}_\tau} P_q(\tau, t_{\mathfrak n_0}, k[\mathcal N^{sp}])\\
&\quad\quad\quad\quad\quad\quad\quad\quad\times\prod_{\mathfrak{n}\in\mathcal{N}^{sp}}  e^{\zeta_\mathfrak{n}2\pi i\Theta_\mathfrak{n}(t_\mathfrak{n}) } \dif t_\mathfrak{n}\dif \tau
\prod_{\mathfrak{l} \in \mathcal{L}^{sp}}^+ n_{\rm in}(k_l),
\end{align*}
 for some function  $P_q$ which satisfies 
\begin{align}
\sup_{|k_{\mathfrak n_0} - k_{\mathfrak p_0}| \leq T^{-1/2}} \left|\left|P_q(\tau,t_0,  k[\mathcal N^{sp}])\right|\right|_{l^\infty} &\lesssim \left(\alpha T^{1/2}L^\delta\right)^q.\label{eq-pestimate}
%\sup_{|k_{\mathfrak n_0} - k_{\mathfrak p_0}| \leq T^{-1/2}} \left|\left|\frac{d}{dt_0}P_q(\tau,t_0,  k[\mathcal N^{sp}])\right|\right|_{l^\infty} &\lesssim \left(\alpha  T^{1/2}\right)^q.\label{eq-dpestimate}
\end{align}

In the above expression,  $\mathcal{Q}^{sp}$ is denoted 
the remained couple   by splicing out its single irregular chain (we abuse this notation a little). Its corresponding decorations, nodes, leaves are defined by a similar manner. The time internal is defined as 
\begin{equation}\label{eq-sig-setE}
\mathcal E^{sp}_{\tau} = \mathcal E^{sp} \cap \{ t_{\mathfrak n_0} \geq (t_{\mathfrak e}\vee t_{\mathfrak f} )+ \tau \}.
\end{equation}
Furthermore, $%\tilde{\epsilon}_
{\mathcal{E}^{sp}}$ permits degeneracies at $\mathfrak{n}_0$, while the subtrees permit complete pairing.
\end{lemma}
\begin{proof}
Without loss of generality, we assume that $\mathcal{Q}$ is the element satisfying  $\zeta_{i}=+$ for any  $0\leq i\leq q$. Then we have $k_{\mathfrak n_i} = k_i$, $k_{\mathfrak p_i} = k_{\mathfrak m_{i-1}} = l_i$, and $t_i = t_{\mathfrak n_i},0\leq i\leq q$,  $t_{q+1} = t_{\mathfrak e}\vee t_{\mathfrak f} $.  
  Then we define $\Theta(t)=\Omega Tt+\Gamma(t)$, where 
    \begin{align*}
    \Omega = \sum_{j = 0}^q \Omega_{j} = |k_{0}|^\sigma- |k_{0} -h|^\sigma + |k_{q+1}- h|^\sigma-|k_{q+1}|^\sigma.\\
   \Gamma(t)=\Gamma(k_0,t)- \Gamma(k_0 -h,t)+ \Gamma(k_{q+1}- h,t)- \Gamma(k_{q+1},t).
    \end{align*}
   
    For all $\mathcal Q'\equiv \mathcal{Q}$, the difference between  $\mathcal K_{\mathcal Q}$ and $\mathcal K_{\mathcal Q'}$ (run over all $\mathcal{Q}$) can be identified as:
 \begin{align*}
    \bigg(\frac{i\alpha T}{L}\bigg)^q&\sum\limits_{k_i \in \mathbb{Z}_L}  %\prod_{i=0}^q \epsilon_{k_{i+1}l_{i+1}l_j}
    \prod_{i=0}^q |k_i|^\beta|k_{i+1}|^\beta|l_i|^\beta|l_{i+1}|^\beta\ \prod_{i=1}^q \varphi_{\leq K}({l_i})\varphi_{\leq K}({k_i})\notag\\ 
    &\prod_{i = 1}^q [\varphi_{\leq K}(l_i)n_{\rm in}(k_i)-\varphi_{\leq K}(k_i)n_{\rm in}(l_i)]
\int_{t_{q+1} < t_q < \ldots < t_1 < t_0}\prod_{i=0}^q  e^{\zeta_{i}2\pi i\Theta_i(t_i) } \dif t_1\dif t_2...\dif t_q,
 \end{align*}
    The resulting expressions are derived as follows:
    \begin{align*} 
    & %\epsilon_{k_{q+1}l_{q+1}l_{0}}
    |k_0|^\beta|k_{q+1}|^\beta|l_0|^\beta|l_{q+1}|^\beta e^{ 2\pi i \Theta(t_0)} \bigg(\frac{i\alpha T}{L}\bigg)^q\sum\limits_{k_i \in \mathbb{Z}_L}% \epsilon_{\mathcal{H}}
    \prod_{i=1}^q |k_i|^{2\beta}|l_i|^{2\beta}\prod_{i=1}^q \varphi_{\leq K}({l_i})\varphi_{\leq K}({k_i}) \\
     &\times\prod_{i = 1}^q [\varphi_{\leq K}(l_i)n_{\rm in}(k_i)-\varphi_{\leq K}(k_i)n_{\rm in}(l_i)]
\int_{t_{q+1} < t_q < \ldots < t_1 < t_0} \prod_{j = 1}^q e^{2\pi i( \Theta_{j}(t_{j})-\Theta_j(t_{0}))} \dif t_1\dif t_2...\dif t_q\\
    & = %\epsilon_{k_{q+1}l_{q+1}l_{0}}
    |k_0|^\beta|k_{q+1}|^\beta|l_0|^\beta|l_{q+1}|^\beta e^{ 2\pi i \Theta(t_0)}  \int_{\tau\in [0,t_0-t_{q+1}]}P_q(\tau, t_0, k[\mathcal N^{sp}]) \dif \tau,
    \end{align*}
  %  where we define $\epsilon_\mathcal{H}:=\prod_{i=0}^q \epsilon_{k_{i+1}l_{i+1}l_j} \cdot \epsilon_{k_{q+1}l_{q+1}l_{0}}^{-1}=\pm1$, 
  and
    \begin{align}\label{bd:irre:cancel}
    P_q&(\tau, t_0, k[\mathcal N^{sp}]) :=\bigg(\frac{i\alpha T}{L}\bigg)^q\sum\limits_{k_i \in \mathbb{Z}_L} %\epsilon_{\mathcal{H}}
    \prod_{i=1}^q |k_i|^{2\beta}|l_i|^{2\beta}\prod_{i=1}^q \varphi_{\leq K}({l_i})\varphi_{\leq K}({k_i}) \notag\\
     &\times\prod_{i = 1}^q [\varphi(l_i)n_{\rm in}(k_i)-\varphi(k_i)n_{\rm in}(l_i)]
\int_{0<s_0 <s_1< \ldots < s_q < \tau} \prod_{j = 1}^q e^{2\pi i(  \Theta_{j}(t_0-s_{j})-\Theta_j(t_{0}))} \dif s_1\dif s_2...\dif s_q.
    \end{align}
    Using the  decaying  of $  n_{\mathrm{in}}$,   we have for the SG case where $|h|\leq T^{-1/2}$,
    \begin{align*}
    \frac{\alpha T}{L} \sum_{k \in \mathbb{Z}_L} \left| \varphi(k)n_{\mathrm{in}}(k - h)-\varphi(k-h)n_{\mathrm{in}}(k)\right| &\lesssim \alpha  T |h| \lesssim \alpha T^{\frac{1}{2}}.
    \end{align*}
 Together with the fact that
    \begin{align*}
        |k_i|^{2\beta}|l_i|^{2\beta} \varphi_{\leq K}({l_i})\varphi_{\leq K}({k_i})\lesssim L^{\delta}, 
    \end{align*} we imply \eqref{eq-pestimate}. 
   Now we could write $ \mathcal{G}_\mathcal{Q}(t,s,k)$ as desired expression by noticing that $\tau=t_{\mathfrak{n}_0}-(t_{\mathfrak e}\vee t_{\mathfrak f} )\in(0,1]$. 
    \end{proof}



\subsection{The choice of irregular chains}\label{sec:choicechain}
The remaining question is how to select such irregular chains and apply the splicing procedure to them. In general, we cannot simply splice all chains, as this may generate new irregular chains or additional terms, which in turn lead to unfavorable counting estimates. 
\begin{definition}\label{def:doublebond}
Consider  an enhanced couple $\mathcal{Q}$, and the corresponding molecule $\mathbb{M}(\mathcal{Q})$, we assume that there are two atoms, $v_1$ and $v_2$, connected by  two bonds. Then we denote by $\mathfrak{n}_j = \mathfrak{n}(v_j)$ for $j=1,2$.  Then we call the double bound is 

\begin{enumerate} 
\item type CL if: two  bonds is composed of one LP bond and one PC bond, and all other leaf-children  remain unpaired.
\item type CN if: two bonds are  both   LP, and the left two leaves are not paired.
\end{enumerate}
According to different choice of paired leaves between $\mathfrak{n}_j$, the bonds can be the opposite directions or same directions.
\end{definition}

\begin{definition}
\begin{enumerate}
\item (Double Chains)\label{def:chain-molecule}
Consider an enhanced couple $\mathcal{Q}$ and the associated molecule $\mathbb{M}(\mathcal{Q})$. A sequence of nodes $\mathbb{H}=(v_0,\ldots,v_q)$ is called a double chain if, for each $0\leq i\leq q-1$, the pair $(v_i,v_{i+1})$ is connected by two bonds.
Given a molecule $\mathbb{M}(\mathcal{Q})$ and a chain $\mathbb{H}$, there is at most one CN double bond; see, for instance, \cite[Proposition 5.7]{DH23a}.
We further call $\mathbb{H}$ a hyperchain if $v_0$ and $v_q$ are connected by a single bond, and a  pseudo-hyperchain if $v_0$ and $v_q$ are connected through another common atom $v\notin\{v_0,\ldots,v_q\}$.
 
\item  (Irregular chain) A double chain is called an irregular chain if all its bonds are CL and are oriented in opposite directions. Similarly, we define irregular hyperchains and irregular pseudo-hyperchains in the same way.
%\item  (Wide Ladder) A  wide ladder is a collection of double chains $\mathbb{H}_i,1\leq i\leq m0,$ where $\mathbb{H}_i=(v_0^{(i)}, \ldots, v_{q^{(i)}}^{(i)}),$ each chain is called a rung. For each $1 \leq i \leq m - 1$, there is a bond connecting the nodes $v_0^{(i)}$ and $v_0^{(i+1)}$, as well as a bond connecting the nodes $v_{q^{(i)}}^{(i)}$ and $v_{q^{(i+1)}}^{(i+1)}$.
\end{enumerate}
\end{definition}
In the following, we collectively refer to these structures as (irregular) chain-like objects. For clarity, the term “chain” will be reserved for chains in the strict sense, excluding hyperchains and pseudo-hyperchains.

In Definition \ref{def:chain-couple}, we also introduced irregular chains for an enhanced couple. In fact, given a couple $\mathcal{Q}$ and the corresponding molecule $\mathbb{M}(\mathcal{Q})$, each irregular chain in $\mathcal{Q}$ corresponds to an irregular chain in $\mathbb{M}(\mathcal{Q})$, and conversely, each irregular chain in $\mathbb{M}(\mathcal{Q})$ corresponds to a unique irregular chain in $\mathcal{Q}$.
Moreover, for any irregular chain $\mathcal{H}$ and the associated double-bond chain $\mathbb{H}$ in $\mathbb{M}$, if $\mathcal{Q}^{sp}$ denotes the couple after splicing, then the corresponding molecule $\mathbb{M}^{sp}=\mathbb{M}(\mathcal{Q}^{sp})$ is obtained from $\mathbb{M}$ by merging all atoms in $\mathbb{H}$ into a single atom. We refer to this procedure as splicing at the level of molecules.

 Consider an enhanced couple $\mathcal Q$ and the associated enhanced molecule $\mathbb{M}(\mathcal{Q})$. Let $\mathscr C$ denote the collection of all maximal small-gap (SG) irregular chain-like objects. Here, “maximal” means that the chain cannot be extended further by adding adjacent nodes. In particular, these objects are pairwise disjoint.
 
We then apply the splicing procedure to the irregular chains in the following manner:

\medskip

\boxed{\rm Operation\ Splicing}
We consider each $\mathcal{C}\in \mathscr C$ and proceed as follows:
\begin{enumerate}
\item If $\mathcal{C}$ is a chain, we apply the splicing procedure directly to the corresponding irregular chain.
\item If $\mathcal{C}$ is a hyperchain or a pseudo-hyperchain, we select a middle double bond. This decomposes $\mathcal{C}$ into two chains (possibly consisting of a single node) on each side of the middle node, and we apply splicing to each of these two chains.

After this procedure, in the case of a hyperchain, the remaining connection becomes a triple bond, while in the case of a pseudo-hyperchain, it reduces to a single double-bond pseudo-hyperchain. 
\end{enumerate}

\medskip
We note that no new double bonds are created during the splicing procedure.

\begin{remark}
  We remark that when applying the splicing procedure to a chain, its type should be determined with respect to the current molecule, rather than the original one, since splicing one chain may change the type of other chains. We refer to \cite{Vas24,Wu25} for a detailed discussion.
In particular, wide ladder structures may arise in the molecule; see, for example, Figure \ref{fig:ladder}. 
\begin{figure}
\begin{tikzpicture}[scale = .9]
    \tikzstyle{every node} = [circle, scale = .5,  draw = black]
    \node (1) at (-5,0) {};
    \node (2) at (-4,0) {};
    \node (3) at (-3,0) {};
    \node (4) at (-2,0) {};
    \node (5) at (-1,0) {};
    \node (6) at (-5,1) {};
    \node (7) at (-4,1) {};
    \node (8) at (-3,1) {};
    \node (9) at (-2,1) {};
    \node (10) at (-1,1) {};
    
  

    \draw[-,  thick] (1.25) -- (2.155);
    \draw[-, thick] (2.205) -- (1.335);
    \draw[-,  thick] (2.25) -- (3.155);
    \draw[-, thick] (3.205) -- (2.335);
    \draw[-,  thick] (3.25) -- (4.155);
    \draw[-, thick] (4.205) -- (3.335);
    \draw[-,  thick] (4.25) -- (5.155);
    \draw[-, thick] (5.205) -- (4.335);

    \draw[-,  thick] (1) -- (6);
    \draw[-,  thick] (5) -- (10);
    
     \draw[-,  thick] (6.25) -- (7.155);
    \draw[-, thick] (7.205) -- (6.335);
    \draw[-,  thick] (7.25) -- (8.155);
    \draw[-, thick] (8.205) -- (7.335);
    \draw[-,  thick] (8.25) -- (9.155);
    \draw[-, thick] (9.205) -- (8.335);
    \draw[-,  thick] (9.25) -- (10.155);
    \draw[-, thick] (10.205) -- (9.335);
  
    
\end{tikzpicture}
\hspace{1cm}
\begin{tikzpicture}[scale = .9]
    \tikzstyle{every node} = [circle, scale = .5,  draw = black]
    \node (1) at (-5,0) {};
    \node (2) at (-4,0) {};
    \node (3) at (-3,0) {};
    \node (4) at (-2,0) {};
    \node (5) at (-1,0) {};
    \node (6) at (-3,1) {};
   
    
  

    \draw[-,  thick] (1.25) -- (2.155);
    \draw[-, thick] (2.205) -- (1.335);
    \draw[-,  thick] (2.25) -- (3.155);
    \draw[-, thick] (3.205) -- (2.335);
    \draw[-,  thick] (3.25) -- (4.155);
    \draw[-, thick] (4.205) -- (3.335);
    \draw[-,  thick] (4.25) -- (5.155);
    \draw[-, thick] (5.205) -- (4.335);

    \draw[-,  thick] (1) -- (6);
    \draw[-,  thick] (6) -- (5);
   
    
\end{tikzpicture}
\caption{In such a wide ladder configuration, applying the splicing operation to the first chain may cause the second chain to transform from an ordinary chain into a pseudo-hyperchain.}
\label{fig:ladder}
\end{figure}
\end{remark}
After the splicing operation, the molecule $\mathbb{M}(\mathcal{Q}^{sp})$ remains connected. By the construction of the irregular chains, the only small-gap (SG) irregular chains that remain have length one and occur only within pseudo-hyperchains.
We also note that there may exist SG CN double chains with opposite orientations; however, these are not irregular chains. 

 
\begin{lemma} \label{cor-M}
For an enhanced couple $\mathcal{Q}$, we apply splicing to all the  chain-like objects as above, and  consider the resulting couple $\mathcal Q^{sp}$. Then we have
\begin{align*}
    \mathcal{G}_\mathcal{Q}(t,s,k) :&= \bigg(\frac{\alpha T}{L}\bigg)^{n_{sp}} \zeta(\mathcal{Q}^{sp}) \sum_{\mathcal{E}^{sp}}  %\tilde{\epsilon}_{\mathcal{E}^{sp}}
    \prod_{\mathfrak{n}\in\mathcal{N}^{sp}}B_\mathfrak{n}\prod_{\mathfrak{n}\in\mathcal{Q}^{sp}}\varphi_\mathfrak{n}
\int_{[0,1]^{\mathcal{N}_0^{sp}}}\int_{\mathcal{E}^{sp}_\tau} \prod_{\mathfrak{n}_0\in \mathcal{N}_0^{sp}}P_{q_{\mathfrak{n}_0}}(\tau_{\mathfrak{n}_0}, t_{\mathfrak n_0}, k[\mathcal N^{sp}])\\
&\quad\quad\quad\quad\quad\quad\quad\quad\times\prod_{\mathfrak{n}\in\mathcal{N}^{sp}}  e^{\zeta_\mathfrak{n}2\pi i\Theta_\mathfrak{n}(t_\mathfrak{n}) } \dif t_\mathfrak{n}\prod_{\mathfrak{n}_0\in \mathcal{N}_0^{sp}}\dif \tau_{\mathfrak{n}_0}
\prod_{\mathfrak{l} \in \mathcal{L}^{sp}}^+ n_{\rm in}(k_l),
\end{align*}

In the above expression,  $\mathcal N_0^{sp}$ are the nodes at which an irregular chain was spliced out and $P_{q_{\mathfrak n_0}}$ is given in Lemma \ref{lem-splice} with $q$ replaced by $q_{\mathfrak n_0}$, which is the length or irregular chain spliced out below $\mathfrak n_0$.    The time internal is defined as 
  $$\mathcal{E}^{sp}_{\pmb{\tau}} = \mathcal E \cap \{t_{\mathfrak n_0} \geq t_{\mathfrak n_0^{2}} + \tau_{\mathfrak n_0}, t_{\mathfrak n_0^{3}} + \tau_{\mathfrak n_0}\}_{\mathfrak n_0 \in \mathcal N_0^{sp}}.$$
  Furthermore, $\tilde{\epsilon}_{\mathcal{E}^{sp}}$ allows degeneracies at each $\mathfrak{n}_0\in\mathcal{N}_0^{sp}$, while the corresponding subtrees may admit complete pairing. We refer to such nodes as bad degenerate nodes.
\end{lemma}
\begin{proof}
   This Lemma is  a direct consequence of Lemma  \ref{lem-splice} by taking all the chains into consideration.
\end{proof}


\section{The Algorithm on Molecules } \label{sec:Algorithm}
In this section, we consider the couple $\mathcal{Q}$ and the associated molecules $\mathbb{M}(\mathcal{Q})$. 
We first apply the splicing operation to the irregular chains, as in Lemma~\ref{cor-M}, 
to obtain $\mathcal{Q}^{sp}$. 
We then implement the algorithm introduced in \cite{Vas24} on the molecule $\mathbb{M}(\mathcal{Q}^{sp})$.
This procedure allows us to control the number of decorations of a molecule in a step-by-step manner. 
Suppose that at some stage of the algorithm we have a molecule $\mathbb{M}^{pre}$, 
and after performing one operation we obtain $\mathbb{M}^{post}$. 
Let $\mathfrak{D}^{pre}$ and $\mathfrak{D}^{post}$ denote the corresponding numbers of decorations. 
We associate to this operation a counting factor $\mathfrak{C}$ such that
\begin{align*}
    \mathfrak{D}^{pre} \leq \mathfrak{C}\,\mathfrak{D}^{post}.
\end{align*}

Compared to \cite{Vas24}, the main additional difficulty arises from possible resonances among atoms. 
We will show that the renormalization of the equation eliminates certain resonant interactions, 
so that these degenerate atoms do not lead to any additional difficulties. 
 

 


\subsection{The algorithm}\label{sec:algo:dege}
First, we perform a pre-processing step: we apply the splicing operation to all  double chains, 
not only to the negative chains or the LG chains, in the same manner as in Operation Splicing before.

We next show that the resulting molecule contains no loops. 
Indeed, the splicing operation (and also the pre-processing step) cannot generate loops. 
Therefore, a loop could only arise if there exists a degenerate node $\mathfrak{n}$ 
whose two children are paired with each other. 
However, this configuration is excluded by the enhanced assumption in Definition \ref{def:couple}. 
 

We now claim that the molecule contains two atoms of degree $3$, 
with all remaining atoms of degree $4$, rather than a single atom of degree $2$. 
We argue by contradiction. There are two possible cases.
In the first case, one tree in the couple is trivial and is paired with a child 
(say $\mathfrak{n}_1$) of the other root $\mathfrak{r}$. 
Then the two subtrees rooted at $\mathfrak{n}_2$ and $\mathfrak{n}_3$ must be fully paired, 
which contradicts the enhanced assumption in Definition \ref{def:couple}. 
In the second case, suppose that the splicing operation (or the pre-processing step) 
produces a molecule with a single degree $2$ atom instead of two degree $3$ atoms. 
In this situation, we may instead splice all but one of the double bonds, 
and then begin the algorithm by removing the two degree $3$ atoms. 
The resulting molecule still contains two atoms of degree $3$.

In Figure~\ref{fig:1loop}, we present two counterexamples that lead to undesired estimates 
in the absence of the enhanced assumption.
 
  
\begin{figure}
\begin{tikzpicture}[level distance=.8cm,
  level 1/.style={sibling distance=.8cm},
  level 2/.style={sibling distance=.8cm}]
\tikzstyle{hollow node}=[circle,draw,inner sep=1.6]
\tikzstyle{solid node}=[circle,draw,inner sep=1.6,fill=black]
\tikzset{
red node/.style = {circle,draw=black,fill=red,inner sep=1.6},
blue node/.style= {circle,draw = black, fill= blue,inner sep=1.6}, 
purple node/.style= {circle,draw = black, fill= purple,inner sep=1.6}, 
orange node/.style= {circle,draw = black, fill= orange,inner sep=1.6},
yellow node/.style= {circle,draw = black, fill= yellow,inner sep=1.6},
green node/.style = {circle,draw=black,fill=green,inner sep=1.6}}


\node[solid node, label = right: $k$ ] at (-5.5,.5){}
    child{node[solid node, label = left: $k$ ]{}
        child{node[solid node, label=left: $k$]{}
             child{node[solid node, label=left: $k$]{}
                  child{node[purple node, label=left: $k$]{}}
                  child{node[orange node, label=below: $k_4$]{}}
                  child{node[orange node, label=below: $k_4$]{}}
             }
             child{node[blue node, label=below: $k_3$]{}}
             child{node[blue node, label=below: $k_3$]{}}
        }
        child{node[green node, label=below: $k_2$]{}}
        child{node[green node, label=below: $k_2$]{}}
    }
    child{node[red node, label=below: $k_1$]{}}
    child{node[red node, label=below: $k_1$]{}}
;
\node[purple node, label=right: $k$] at (-3.5,.5){};


\end{tikzpicture}
\hspace{1cm}
\begin{tikzpicture}[level distance=.8cm,
  level 1/.style={sibling distance=2.3cm},
  level 2/.style={sibling distance=.8cm}]
\tikzstyle{hollow node}=[circle,draw,inner sep=1.6]
\tikzstyle{solid node}=[circle,draw,inner sep=1.6,fill=black]
\tikzset{
red node/.style = {circle,draw=black,fill=red,inner sep=1.6},
blue node/.style= {circle,draw = black, fill= blue,inner sep=1.6}, 
purple node/.style= {circle,draw = black, fill= purple,inner sep=1.6}, 
orange node/.style= {circle,draw = black, fill= orange,inner sep=1.6},
yellow node/.style= {circle,draw = black, fill= yellow,inner sep=1.6},
green node/.style = {circle,draw=black,fill=green,inner sep=1.6}}


\node[solid node, label = right: $k$ ] at (-5.5,.5){}
    child{node[solid node, label = left: $k_1$ ]{}
        child{node[purple node, label=below: $k_1$]{}}
        child{node[green node, label=below: $k_2$]{}}
        child{node[orange node, label=below: $k_2$]{}}
    }
    child{node[solid node, label=right: $k_1$]{}
        child{node[purple node, label=below: $k_1$]{}}
        child{node[green node, label=below: $k_2$]{}}
        child{node[orange node, label=below: $k_2$]{}}}
    child{node[blue node, label=below: $k$]{}}
;

\node[blue node, label=right: $k$] at (-3.5,.5){};

\end{tikzpicture}
\caption{If the enhanced assumption is not imposed, then in the case $n = 4$ shown on the left, 
all sibling children are mutually paired (as indicated by nodes of the same color). 
In this situation, the parameters $k_i$, $1 \leq i \leq 4$, can be chosen independently, 
leading to the worst-case counting estimate of order $L^{4}$. 
The corresponding molecule structure contains four loops.\\
The example on the right shows that problematic counting estimates may also arise 
from configurations that do not contain loops. 
In this case, nodes of the same color are mutually paired, 
and the resulting counting estimate for the parameters $k_1$ and $k_2$ is of order $L^{2}$, 
which may  not be bounded by $(LT^{-1/5})^3 T^{-3/5}$. 
Under the enhanced assumption, such configurations are excluded. }
\label{fig:1loop}
\end{figure}

 In the following, we apply the algorithm introduced in \cite[Section 7]{Vas24} 
to remove all remaining atoms. 
Compared with the assumptions in \cite[Section 7.1]{Vas24}, 
degenerate atoms may still be present in our setting. 
However, as we shall see, under the assumption that no loops are present, 
these degenerate atoms do not introduce additional difficulties.

For the reader’s convenience, we provide a complete description of the algorithm 
in Appendix~\ref{sec:app}. 
Here we summarize the types of operations according to the values of 
$\mathfrak{C}$ and $\Delta \chi$:
\begin{itemize}
    \item \textbf{Bridge Operations} (including Operation 1). 
    In this case, $\Delta \chi = 0$. 
    By \cite[Lemma 9.14]{DH23a}, for a bridge $l$, the parameter $k_l$ is fixed, 
    hence $\mathfrak{C} = 1$. 
    Let $m_0$ denote the total number of such operations.
    
    \item \textbf{Sole Atom Operations} (including Operation 9). 
    Here $\Delta \chi = 0$ and $\mathfrak{C} = 1$. 
    Let $m_1$ denote their total number.
    
    \item \textbf{Two-Vector Counting Operations} (including Operation 0 and Operations 5--8). 
    In this case, $\Delta \chi = -1$ (since no bridge is present), 
    and by Lemma~\ref{counting}, $\mathfrak{C} = L$. 
    Let $m_2$ denote the number of such operations.
    
    \item \textbf{Three-Vector Counting Operations} (including Operations 2--4). 
    Here $\Delta \chi = -2$ (as neither bridges nor loops are present). 
    If the atom is non-degenerate, then by Lemma~\ref{counting},
    \[
    \mathfrak{C} = L^{2}T^{-1}D^2 \log L.
    \]
    If the atom is degenerate, then trivially $\mathfrak{C} = L$, 
    which is bounded by $L^{2}T^{-1}D^2 \log L$ since $T < L$. 
    We assume throughout that the decoration parameter $D$ remains bounded. 
    Let $m_3$ denote the total number of such operations.
\end{itemize}


Then,  we need the following key bound  as in \cite[Proposition 7.6]{Vas24}:

\begin{lemma}
  For any molecule $\mathbb{M}(\mathcal{Q})$ obtained after applying  Operation Splicing 
and the  Pre-processing Step,  we have
  \begin{align}\label{prop-op-bound}
    m_2 &\leq 3m_3 -3.
\end{align}
\end{lemma}In fact, the proof does not depend on the decorations of the molecule, 
and hence degenerate atoms do not affect the result.

\begin{proposition}\label{prop-rigidity:nodege}
 We consider a enhanced couple $\mathcal{Q}^{sp}$ as in Proposition \ref{cor-M}, and its corresponding molecule $\mathbb M = \mathbb M(\mathcal Q^{sp})$, with scale $n$. For a $(k_v,C_v)$-decoration as in Definition \ref{def:decorations-molecules} for some fixed $C_v$, we additionally that there are $k_v^0\in\mathbb{Z}_L$ satisfying $|k_v^0|\leq D$ for some $D>1$, and for any $v\in\mathbb{M}$, $$|k_v-k_v^0|\leq 1.$$
Then, the number $\mathfrak D$ of such  decorations is bounded by 
\begin{equation} 
\mathfrak D \leq C^n L^{n}T^{-\frac{n}{5} - \frac{3}{5}}(\log L)^nD^{2n}.
\end{equation}
Here the universal constant $C$ is uniform in the decoration and the choice of $k_v^0$ and $C_v$.
\end{proposition}
\begin{proof}
For such a molecule $\mathbb{M}$, we first consider the double bonds, 
namely those for which the two bonds are either oriented in the same direction 
or have LG-opposite directions (recall Definition \ref{def:twist} for LG). 
These bonds are not included in $m_2$. 
Once such a bond is removed, we strictly eliminate one double bond, 
and the corresponding change is
$\Delta \chi = -2 + 1 + 0 = -1.$
By Lemma~\ref{counting}, we may then apply the advantageous two-vector counting estimate
$L T^{-\frac{1}{2}} D^2$
either for LG-irregular chains or for same-directional chains. 
In summary, we remove all such bonds and have
$\mathfrak{C} = L T^{-\frac{1}{2}} D^2$ and $\Delta \chi = -1.$
These two bounds are comparable   to the three-vector counting estimate 
that will be used in the subsequent analysis.

 
 Then  we apply the algorithm above to the remaining molecule and cumulate the count estimate for each instance. By counting on $\chi$ we have $2m_3 + m_2 = n$, together with \eqref{prop-op-bound}, it holds that $5m_3 - 3 \geq n$, which implies that 
\begin{align*}
    \mathfrak D &\leq C^n (L^{2}T^{-1}D^2 \log L)^{m_3}L^{dm_2}= C^n L^{n}( \log L)^{m_3}  T^{-m_3}D^{2m_3} \\
    & \leq C^n L^{n}  T^{-\frac{n}{5} - \frac{3}{5}} (\log L)^nD^{2n}.
    \end{align*}
\end{proof} 


 


\subsection{Bounds on couples} \label{sec-boundoncouple}
In previous sections, we have establish the estimates on the counting problem. In this section, we are in position to prove  Proposition \ref{prop:couple}.
\begin{proof}[Proof of Proposition \ref{prop:couple}]
First, we consider to prove \eqref{bd:couple}, the left hand side equals to
     \begin{align}
     \sum_{\mathcal{Q}_{n}}\mathbf{E}[ c^{\mathcal{T}^+}_k (s)c^{\mathcal{T}^-}_k(s)] = \sum_{|\mathcal{Q}|=n}(\mathcal{K}_\mathcal{Q})(t,s,k).
 \end{align}
 where the sum on the right hand side is taken over all the enhanced couples with $n$ nodes. Now since $n\leq N^3$, we know that there are at most  $O(C^{N^3}(N^3)!)$ such kinds of couples. By taking into account the different ways of choosing degenerate nodes and pairing, there are also at most $O(C^{N^3}(N^3)!)$ ways.  Then by classify the couple by the congruent relation, we also have at most $O(C^{N^3}(N^3)!)$ ways. Since this number is independent of $L$, in the following we could fix one single couple $\mathcal{Q}$ with fixed decoration, and   then  consider its  congruent class and the expression $\mathcal{G}_{\mathcal{Q}}$, which is defined in Lemma \ref{lem-splice}.

In the expression in    Lemma \ref{lem-splice}, by the choice of $K$, the symbols of derivatives  are bounded by \begin{align*}
\prod_{\mathfrak{n}\in\mathcal{N}^{sp}}B_\mathfrak{n}\prod_{\mathfrak{n}\in\mathcal{Q}^{sp}}\varphi_\mathfrak{n}\leq L^{n_{sp}\delta}.
    \end{align*}
    Then, we aim to bound the oscillatory integral for the fixed $\tau_{\mathfrak{n}_0}$:
    \begin{align}
      \int_{\mathcal{E}^{sp}_\tau} \prod_{\mathfrak{n}_0\in \mathcal{N}_0^{sp}}P_{q_{\mathfrak{n}_0}}(\tau_{\mathfrak{n}_0}, t_{\mathfrak n_0}, k[\mathcal N^{sp}])\prod_{\mathfrak{n}\in\mathcal{N}^{sp}}  e^{\zeta_\mathfrak{n}2\pi i\Theta_\mathfrak{n}(t_\mathfrak{n}) } \dif t_\mathfrak{n}\label{eq-osc}.
    \end{align}
Since the condition $\sup_k\sup_{t\in[0,1]}|\partial_t\Gamma(k,t)|\leq 1$ asked    in Proposition \ref{prop:est:af} is satisfied.
    We apply    Proposition \ref{prop:est:af} with \begin{align*}f_{\mathfrak{n}_0}(t_{\mathfrak n_0})=\left(\alpha T^{1/2}L^\delta\right)^{-q_{\mathfrak n_0}}P_{q_{\mathfrak n_0}}(\tau_{\mathfrak n_0}, t_{\mathfrak n_0}, k[\mathcal N^{sp}]),\end{align*} which satisfies the assumptions by \eqref{eq-pestimate}. Then we obtain
    \begin{align}
    &\left| \sum_{\mathcal{Q}_{n}}\mathbf{E}[ c^{\mathcal{T}^+}_k (t)c^{\mathcal{T}^-}_k(s)]\right| \notag\\
    &\lesssim 
      C^{n}\left( \frac{\alpha T}{L}\right)^{n_{{{sp}}}}(\alpha  T^{1/2}L^\delta )^{\sum q_{\mathfrak{n}_0}} \sum_{\mathcal{E}^{sp}}  %\tilde{\epsilon}_{\mathcal{E}^{sp}}
      L^{n_{sp}\delta}\notag\\
&\quad\quad\quad  \times
\sup_{\pmb \tau \in [0,1]^{\mathcal N_0^{sp}}}\bigg{|}\int_{\mathcal{E}^{sp}_\tau} \prod_{\mathfrak{n}_0\in \mathcal{N}_0^{sp}}f_{\mathfrak{n}_0}(t_{\mathfrak n_0})\prod_{\mathfrak{n}\in\mathcal{N}^{sp}}  e^{\zeta_\mathfrak{n}2\pi i[\overline\Omega_\mathfrak{n}Tt_\mathfrak{n}+\sum\Gamma_1(\mathfrak{n},t_\mathfrak{n})]  } \dif t_\mathfrak{n}\bigg{|}
\prod_{\mathfrak{l} \in \mathcal{L}^{sp}}^+ n_{\rm in}(k_l).\end{align}
Then using the decay of the initial data $n_{\rm in}$ and the truncation functions $\prod_{\mathfrak{n}\in\mathcal{Q}^{sp}}\varphi_\mathfrak{n}$, we could assume that $|k_l-k_l^0|\leq1$ for $k_l^0\in \mathbb{Z}^d$, and only consider the expression under the restrictions $|k_l^0|\leq L^{\delta}$. 
 Moreover, by rewrite $n_{\rm in}(k_l)=\langle k_l\rangle^{-50}\tilde n_{\rm in}(k_l)$, we have the above expression (under the restriction  $|k_l-k_l^0|\leq1$ ) is bounded by

\begin{align}
 &\prod_{l\in\mathcal{L}^{sp}}\langle k_l\rangle^{-30} C^{n}\left( \frac{\alpha T}{L}\right)^{n_{{sp}}}(\alpha  T^{1/2} L^\delta)^{\sum q_{\mathfrak{n}_0}} \sum_{{\lambda[\mathcal N^{sp}]}}\sum_{\mathcal{E}_\lambda^{sp}}  %\tilde{\epsilon}_{\mathcal{E}^{sp}_\lambda}
 L^{n_{sp}\delta} \notag\\
      & \quad \quad\times 
\sup_{\pmb \tau \in [0,1]^{\mathcal N_0^{sp}}} \left|\int_{\mathcal{E}^{sp}_{\pmb{\tau}}}\prod_{\mathfrak{n}_0 \in \mathcal{N}^{sp}_0}f_{\mathfrak{n}_0}(t_{\mathfrak n_0})\prod_{\mathfrak{n} \in \mathcal{N}^{sp}}e^{\zeta_\mathfrak{n}2\pi i[\overline\Omega_\mathfrak{n}Tt_\mathfrak{n}+\sum\Gamma_1(\mathfrak{n},t_\mathfrak{n})]  } \dif  t_{\mathfrak{n}}\right|\notag\\
      &\lesssim \langle k\rangle^{-20} C^{n}\left( \frac{\alpha T}{L}\right)^{n_{{{sp}}}}(\alpha  T^{1/2} L^\delta)^{\sum q_{\mathfrak{n}_0}} \sup_{\lambda[\mathcal N^{{sp}}]}\sum_{\mathcal E_{\lambda }^{{sp}}}%\tilde \epsilon_{\tilde{\mathcal E}_{\lambda }^{{sp}}}
      L^{n_{sp}\delta}\notag\\
      & \quad \quad\times\sum_{{\lambda[\mathcal N^{sp}]}}\sup_{\pmb \tau \in [0,1]^{\mathcal N_0^{sp}}}\left|\int_{\mathcal{E}^{sp}_{\pmb{\tau}}}\prod_{\mathfrak{n}_0 \in \mathcal{N}^{sp}_0}f_{\mathfrak{n}_0}(t_{\mathfrak n_0})\prod_{\mathfrak{n} \in \mathcal{N}^{sp}}e^{\zeta_\mathfrak{n}2\pi i[\overline\Omega_\mathfrak{n}Tt_\mathfrak{n}+\sum\Gamma_1(\mathfrak{n},t_\mathfrak{n})]  } \dif t_{\mathfrak{n}}\right|,\label{eq-sp1}
    \end{align}
    where $\lambda[\mathcal N^{sp}] \in \mathbb{Z}^{|\mathcal N^{sp}|}$ and ${\mathcal E}_{\lambda}^{sp}$ is taken over all the decorations of the couple satisfying $|T\overline\Omega_{\mathfrak n} - \lambda_{\mathfrak n}| \leq 1$ for each $\mathfrak n \in \mathcal N^{sp}$. Moreover, it holds that $\sum q_{\mathfrak{n}_0}+ n_{sp} = n$. We apply Proposition \ref{prop:est:af} to bound the last line above by
    \begin{align}
        &\sum_{{\lambda[\mathcal N^{sp}]}}\sup_{\pmb \tau \in [0,1]^{\mathcal N_0^{sp}}}\sum_{d_\mathfrak{n}\in\{0,1\}}\prod_{\mathfrak{n}\in\mathcal{N}^{sp}}\langle Tq_\mathfrak{n}\rangle^{-1}
        %&\quad\quad \lesssim C^{n_{sp}}\sum_{\lambda[\mathcal N]} \sum_{d_{\mathfrak n} \in \{0,1\}}\frac{1}{\langle T{q}_{\mathfrak n}\rangle} 
        \lesssim C^{n_{sp}} (\log L)^{n_{sp}},\notag
    \end{align}
    where we used the fact that $d_\mathfrak{n}$ only has two choices, and since $k_l\in\mathbb{Z}_L$ with $|k_l-k_l^0|\leq1$, there are at most $L^{4}$ choices of $\lambda_\mathfrak{n}$.

    To bound the  term 
      $  \sup_{\lambda[\mathcal N^{{sp}}]}\sum_{\mathcal E_{\lambda }^{{sp}}}%\tilde \epsilon_{\tilde{\mathcal E}_{\lambda }^{{sp}}}
     1 ,$
    we transfer the couples into molecule as in Section \ref{sec:Algorithm}.  By choosing $\delta>0$ small enough such that $\alpha T^{1-\beta}L^{2\delta}<1$, it suffices to apply Lemma \ref{counting} with $D=L^{\delta}$, we then obtain that \eqref{eq-sp1} is bounded by
    \begin{align}
     \langle k\rangle^{-20} &C^{n}\left( \alpha T \right)^{n_{{{sp}}}}(\alpha  T^{1/2}L^\delta )^{\sum q_{\mathfrak{n}_0}}L^{n_{sp}\delta}T^{-\frac{n_{sp}}{5} - \frac{3}{5}}  (\log L)^{2n_{sp}}D^{2n_{sp}}\notag\\
     &\leq  \langle k\rangle^{-20} C^{n}(\alpha  T^{4/5} L^{4\delta})^nT^{- \frac{3}{5}},\notag
    \end{align}
    which implies \eqref{bd:couple}.
     
   Then we prove the next statement.  We recall that by \eqref{def:gamma1}, \eqref{bd:couple}
 \begin{align*}
   | \mathscr{P}( \Gamma_1)|&=\left|\frac{\alpha T}{\pi L}\sum_{k_1}\sum_{\mathcal{Q}_{\leq N}}\mathbf{E}[ c^{\mathcal{T}^+}_{k_1} (t)c^{\mathcal{T}^-}_{k_1}(t)] ( \Gamma_1)B_{k_1,k_1,k,k}\varphi_{\leq K}(k)\right|\\
   &\lesssim \alpha T^{\frac25}L^\delta \cdot L^{-1}\sum_{k_1\in\mathbb{Z}_L}\langle k_1\rangle^{-20} \cdot \frac{\alpha L^{4\delta}  T^{4/5}}{1-\alpha L^{4\delta} T^{4/5}} \leq L^{-3\delta},
 \end{align*}
 where we used $\alpha L^{4\delta} T^{4/5}\leq L^{-2\delta}$.  Here we also use the fact that the sum over $k_1$ is restricted to $|k_1| \lesssim L^\delta$, so the symbol term is bounded together with the
  truncation function $\varphi_{\leq K}(k)$.



In the end,   it remains to  bound the last term, namely,  $\mathscr{P}[\Gamma_1]-\mathscr{P}[{\Gamma}'_1]$, under the assumption on $\Gamma_1$ and $\Gamma_1'$. By a similar argument as before, we only need to handle the term 
\begin{align*}
    \mathcal{G}_Q[\Gamma_1](t,k)-    \mathcal{G}_Q[{\Gamma}'_1](t,k),
\end{align*}
 where $ \mathcal{G}_Q$ is defined in \eqref{def:gq}. From the definition, the only terms that involve the  $\Gamma_1$  are in the time integrals, so we
 only need to estimate the term 
 \begin{align}
    \int_{\mathcal{D}}\prod_{\mathfrak{n}\in\mathcal{N}}  e^{\zeta_\mathfrak{n}2\pi i\overline\Omega_\mathfrak{n}Tt_\mathfrak{n}}e^{\zeta_\mathfrak{n}2\pi i\sum\Gamma_1(\mathfrak{n},t_\mathfrak{n}) }-\prod_{\mathfrak{n}\in\mathcal{N}}  e^{\zeta_\mathfrak{n}2\pi i\overline\Omega_\mathfrak{n}Tt_\mathfrak{n}}e^{\zeta_\mathfrak{n}2\pi i\sum\Gamma'_1(\mathfrak{n},t_\mathfrak{n}) } \dif t_\mathfrak{n}.\label{bd:inge_diff}
 \end{align}
  By taking difference, it suffices to estimate the following expression: for $|\overline\Omega_kT|\geq 1$,
 
  \begin{align*}
        \int_0^t &  e^{ 2\pi i[\overline\Omega_kTs+\sum\Gamma_1(k,s)] } - e^{ 2\pi i[\overline\Omega_kTs+\sum \Gamma'_1(k,s)] } \dif s\\
        &= \frac{ e^{ 2\pi i[\overline\Omega_kTt+\sum\Gamma_1(k,t)] }-e^{ 2\pi i[\overline\Omega_kTt+\sum  \Gamma'_1(k,t)] }}{2\pi i\overline\Omega_kT}\\
        &\quad - \int_0^t \frac{  e^{ 2\pi i[\overline\Omega_kTs+\sum\Gamma_1(k,s)] }2\pi i\sum \partial_s\Gamma_1(k,s)- e^{ 2\pi i[\overline\Omega_kTs+\sum \Gamma'_1(k,s)] }2\pi i\sum  \partial_s\Gamma_1'(k,s)}{2\pi i\overline\Omega_kT}   \dif s.
   \end{align*}

Now we go back to \eqref{bd:inge_diff}, by taking difference, at a loss of $|\mathcal{Q}|=n\leq N^3$, we only to deal with a similar expression as \eqref{def:af} with $f=1$, and there is a node $\mathfrak{n}$, such that the  term in the product $e^{ 2\pi i \sum\Gamma_1(k,s) } $ is replaced $ e^{ 2\pi i \sum\Gamma_1(k,s)} - e^{ 2\pi i\sum \Gamma'_1(k,s)} $, and for the remaining node, this term is $\Gamma_1$ or $ \Gamma'_1$. Then since $$\sup_k\sup_{t\in[0,1]}|\partial_t\Gamma_1(k,t)|+|\partial_t\Gamma'_1(k,t)|\leq 1,$$ we  could estimate this expression by a similar way as Lemma \ref{prop:est:af} and we only need to deal with the integral at node $\mathfrak{n}$. From the above expression,  the first term equals to
   \begin{align*}
       \frac{ e^{ 2\pi i\overline\Omega_kTt}}{2\pi i\overline\Omega_kT}(e^{2\pi i\sum\Gamma_1(k,t) }-e^{ 2\pi i \sum  \Gamma'_1(k,t) }).
   \end{align*}  
   Then it contribution is $\frac{ e^{ 2\pi i\overline\Omega_kTt}}{2\pi i\overline\Omega_kT}$ with an extra term $(e^{2\pi i\sum\Gamma_1(k,t) }-e^{ 2\pi i \sum  \Gamma'_1(k,t) })$.  Then for this term, we obtain the same bound as in Lemma \ref{prop:est:af}, with extra  product of $|\partial_t\Gamma_1-\partial_t \Gamma'_1|$.
For the second term, it is bounded by 
   \begin{align*}
  \frac{1}{ \overline\Omega_kT}   \left(  \left | \sum \partial\Gamma_1-\sum\partial \Gamma_1' \right|+ \left|e^{ 2\pi i \sum\Gamma_1} -e^{ 2\pi i \sum \Gamma_1'}  \right|\right)\lesssim \frac{1}{ \langle\overline\Omega_kT\rangle}|\partial_t\Gamma_1-\partial_t\Gamma_1'|.
   \end{align*}
   For the case $|\overline\Omega_kT|< 1$, we could directly consider the difference $$|  e^{ 2\pi i\sum\Gamma_1(k,s) } - e^{ 2\pi i\sum \Gamma'_1(k,s) }|  \lesssim\frac{1}{ \langle\overline\Omega_kT\rangle} |  e^{ 2\pi i\sum\Gamma_1(k,s) } - e^{ 2\pi i\sum \Gamma'_1(k,s) }|,$$ which is bounded by the same argument as before.
   Then for this term, we also obtain the same bound as in Lemma \ref{prop:est:af}, with extra  product of $|\partial_t\Gamma_1-\partial_t \Gamma'_1|$. 
   Then  in summary we obtain that \eqref{bd:inge_diff} is bounded by 
   \begin{align*}
       N^3\cdot\sum_{d_\mathfrak{n}\in\{0,1\}}\prod_{\mathfrak{n}\in\mathcal{N}}C_\mathfrak{n}\langle Tq_\mathfrak{n}\rangle^{-1} \sup_k\sup_{t\in[0,1]}|\partial_t\Gamma_1-\partial_t \Gamma'_1|.
   \end{align*}
   Then by following the same calculation as before, the term $   \mathcal{G}_Q[\Gamma_1](t,k)-    \mathcal{G}_Q[{\Gamma}'_1](t,k)$ is bounded by  $$ \langle k\rangle^{-20} C^{n}(\alpha  T^{4/5} L^{4\delta})^nT^{- \frac{2}{5}}\sup_k\sup_{t\in[0,1]}|\partial_t\Gamma_1-\partial_t \Gamma'_1|.$$
   Taking sum on $n$ and $k$ we obtain the desired bound by a similar argument as before.
 
\end{proof}


\section{The Remainder Term}\label{bd:remainder}
In this section, we study the remainder term $\mathcal{R}^{N+1}$, which solves the following fixed point problem:
\begin{equation}\mathcal{R}^{N+1}=(1-\mathscr{L})^{-1}[c_{\sim N}+\mathscr{Q}(\tilde\varphi_{\leq K}\mathcal{R}^{N+1})+\mathscr{C}(\tilde\varphi_{\leq K}\mathcal{R}^{N+1})].\label{def:rkn:fix}
\end{equation}
We refer to Section \ref{sec:rem} for the notation. Since $\mathscr{Q}$ and $\mathscr{C}$ are bilinear and trilinear operators, respectively, it suffices to estimate the linear operator $\mathscr{L}$.
In what follows, we show that, with a high probability, the operators $\mathscr{L}^n$ are bounded on $Z$ for all $n$, as stated in Proposition \ref{prop:l}. Consequently, the bound for $(1-\mathscr{L})^{-1}$ follows directly.

\subsection{The analysis of linear operator $\mathscr{L}$}\label{sec:analysisL}
To estimate the linear operator $\mathscr{L}$, we further analyze its structure through the tree expansion developed in the previous sections.
\begin{proposition} \label{prop-Lm}
Let $\mathscr{ L}$ be defined as in \eqref{def:l}. Note that $\mathscr L^{n}$ is a  linear operator for $n \geq 0$. Define its kernels $(\mathscr L^{n})_{k l}^\zeta(t,s)$ for $\zeta \in \{\pm\}$ by 
\begin{equation*}
(\mathscr L^{n} b)_k (t) = \sum_{\zeta \in \{\pm\}} \sum_{l} \int_{\R} (\mathscr L^{n})_{kl}^{\zeta} (t,s) b_{l}(s)^{\zeta}\dif  s.
\end{equation*}
Then, for each $1 \leq n \leq N$ and $\zeta \in \{\pm\}$, we can decompose
\begin{equation}
(\mathscr L^{n})^\zeta_{k, l} = \sum_{n \leq m \leq N^3} (\mathscr 
L^{n})_{k, l}^{m, \zeta},\notag
\end{equation}
such that for any $n \leq m \leq N^3$ and $k,l \in \mathbb{Z}_L$ and $t,s \in [0,1]$ with $t > s$, we have 
\begin{equation} \label{eq-Lnm}
\mathbf{E} |(\mathscr L^{n})_{k, l}^{m, \zeta}(t,s)|^2 \lesssim \langle k-\zeta l\rangle^{-20}  L^{40} (\alpha  T^{4/5} L^{2\delta})^m.
\end{equation}
    \end{proposition}
First, to analyze the expression of $(\mathscr L^{n})_{k, l}^{m, \zeta}$, we introduce the notion of flower trees, as in \cite[Section 11]{DH23a}.
\begin{definition}\label{def:flowtree}
A flower tree is a tree $\mathcal{T}$ with a   particular  leaf $\mathfrak{f}$  with is  called the  flower. Choosing a different leaf $\mathfrak{f}$ for the same tree $\mathcal{T}$ yields a different flower tree. We call the unique path connecting the root $\mathfrak{r}$ and the flower $\mathfrak{f}$  as the stem. A flower couple is a pair of flower trees whose flowers are coupled.

The  height  of a flower tree $\mathcal{T}_{\mathfrak{f}}$ is defined as the length of stem, i,e, number of branching nodes along its stem. A flower tree of height $n$ could be obtained by adding two sub-trees on each nodes on the stem by $n$ times. A flower tree is called an  admissible  tree if each of the added sub-trees has scale at most $N$.
\end{definition}

We then complete the proof of Proposition~\ref{prop-Lm} by expressing $\mathscr{L}^n$ as a sum over flower trees.
\begin{proof}[Proof of Proposition \ref{prop-Lm}]
  By the definition of the operator $\mathscr{L}$ in \eqref{def:l}, we observe that applying $\mathscr{L}$ corresponds to attaching two subtrees $\mathcal{T}_1$ and $\mathcal{T}_2$ to a single node, where the scales of these subtrees are at most $N$. Hence, the resulting object is an admissible tree in the sense of Definition \ref{def:flowtree}.   In particular, we remark that the absence of terms such as $\mathbf{E}[v_{k_1}\overline {c_{k_1}^{\leq N}}]$ reflects the fact that such contributions cannot be fully paired, in accordance with the definition of flower trees.  
Iterating this construction $n$ times, the operator $\mathscr{L}^n$ corresponds to building an admissible tree of height $n$. Therefore,
    \begin{align*}
        (\mathscr L^{n})^\zeta_{k, l} = \sum_{n \leq m \leq N^3} (\mathscr 
L^{n})_{k, l}^{m, \zeta}=\sum_{n \leq m \leq N^3} \sum_{\mathcal{T}}\tilde{c}_{kl}^{\mathcal{T}}(t,s), 
\end{align*}
where the sum over $\mathcal{T}$ runs over all flower trees with height $n$, scale $m$, and satisfying $\zeta_{\mathfrak{r}}=+$ and $\zeta_{\mathfrak{f}}=\zeta$.  The coefficient $\tilde{c}_{kl}^{\mathcal{T}}(t,s)$ is given by
\begin{align}
    \tilde{c}_{kl}^{\mathcal{T}}(t,s):=\bigg(\frac{\alpha T}{L}\bigg)^m\sum_{\mathcal{D}}\prod_{\mathfrak{n}\in\mathcal{N}}(i\zeta_{\mathfrak{n}}%\epsilon_{k_{\mathfrak{n}_1},k_{\mathfrak{n}_2},k_{\mathfrak{n}_3}}
    )\prod_{\mathfrak{n}\in\mathcal{N}}B_{k_{\mathfrak{n}_1},k_{\mathfrak{n}_2},k_{\mathfrak{n}_3},k_{\mathfrak{n}}}\prod_{\mathfrak{f}\neq \mathfrak{n}\in\mathcal{T}}\varphi_{\leq K}(k_\mathfrak{n})\notag\\
    \times \int_{\mathcal{D}_\mathfrak{n}}\prod_{\mathfrak{n}\in\mathcal{N}}  e^{\zeta_\mathfrak{n}2\pi i \Theta_\mathfrak{n}(t) } \dif t_\mathfrak{n} \delta({t_{\mathfrak{f}^p}-s)}\prod_{\mathfrak{f}\neq \mathfrak{l}\in\mathcal{L}}\sqrt{n_{\mathrm{in}}(k_{\mathfrak{l}})}\tilde{\mathcal{B}}_{\mathcal{T}_\mathfrak{f}}(\omega)1_{k_\mathfrak{f}=l},\label{def:tildect}
\end{align}
where the sum over $\mathcal{D}$ runs over all $k$-decorations of $\mathcal{T}_{\mathfrak{f}}$, and $\mathfrak{f}^p$ denotes the parent of $\mathfrak{f}$. The term $\tilde{\mathcal{B}}_{\mathcal{T}_\mathfrak{f}}$ is defined analogously to \eqref{def:Bomega}, with $\eta_{\mathfrak{f}}(\omega)$ replaced by $1$.  

Next, we consider flower couples $\mathcal{Q}$ whose subtrees all have height $n$ and scale $m$. Then
 \begin{align*}
     \mathbf{E} |(\mathscr L^{n})_{k, l}^{m, \zeta}(t,s)|^2 =\sum_{\mathcal{Q}} (\tilde{\mathcal{K}}_\mathcal{Q})(t,s,k), 
 \end{align*}
 where the sum on $\mathcal{Q}$ runs over all the flower couples $\mathcal{Q} $ such that each substree is with  height $n$ and scale $m$. And 
 \begin{align*}
   (\tilde{\mathcal{K}}_\mathcal{Q})(t,s,k):&= \bigg(\frac{\alpha T}{L}\bigg)^{2m} \zeta(\mathcal{Q}) \sum_{\mathcal{E}}  %\epsilon_{\mathcal{E}}
   \prod_{\mathfrak{n}\in\mathcal{N}}B_\mathfrak{n}\prod_{\mathfrak{f}\neq \mathfrak{n}\in\mathcal{Q}}\varphi_\mathfrak{n}\\
  &\quad\quad \times 
\int_{\mathcal{E}_\mathfrak{n}}\prod_{\mathfrak{n}\in\mathcal{N}}  e^{\zeta_\mathfrak{n}2\pi i\Theta_\mathfrak{n}(t_\mathfrak{n}) } \dif t_\mathfrak{n}\prod\delta({t_{\mathfrak{f}^p}-s)}
\prod_{\mathfrak{f}\neq \mathfrak{l} \in \mathcal{L}}^+ n_{\rm in}(k_{\mathfrak{l}})1_{k_\mathfrak{f}=l}.
 \end{align*}
 
To prove \eqref{eq-Lnm}, we follow the same strategy as in Section~\ref{sec-boundoncouple}, with the only difference being the presence of flower structures.
 
First, for any flower couple $\mathcal{Q}_{\mathfrak{f}}$, we consider its congruence class $\mathcal{Q}_{\mathfrak{f}}\equiv \mathcal{Q}'_{\mathfrak{f}'}$, where $\mathfrak{f}'$ is the image of $\mathfrak{f}$. As in Section~\ref{sec:choicechain}, we select all irregular chains and apply splicing. If the flower lies within an irregular chain, we skip its parent and instead consider the two resulting subchains. This introduces at most two double bonds, resulting in a loss of at most $L^{2}$.

 
We then proceed as in Proposition \ref{prop:couple}. By the choice of truncation functions,
 $$\prod_{\mathfrak{n}\in\mathcal{N}}B_\mathfrak{n}\prod_{\mathfrak{f}\neq \mathfrak{n}\in\mathcal{Q}}\varphi_\mathfrak{n}\lesssim L^{m\delta}.$$
 Most terms can be handled as before. For the special term $B_{\mathfrak{f}^p}$, let $\mathfrak{f}^{s,1},\mathfrak{f}^{s,2}$ denote the siblings of $\mathfrak{f}$. Using
 $k_{\mathfrak{f}^{p}}=k_{\mathfrak{f}}-k_{\mathfrak{f}^{s,1}}+k_{\mathfrak{f}^{s,2}}$, we obtain
\begin{align}
B_{\mathfrak{f}^p}&\varphi_{\leq K}(k_{\mathfrak{f}^{p}})\varphi_{\leq K}(k_{\mathfrak{f}^{s,1}})\varphi_{\leq K}(k_{\mathfrak{f}^{s,2}})\notag\\
    &\leq |k_{\mathfrak{f}^p}|^{\beta}|k_{\mathfrak{f}^{s,1}}|^{\beta}|k_{\mathfrak{f}^{s,2}}|^{\beta}(|k_{\mathfrak{f}^p}|+|k_{\mathfrak{f}^{s,1}}|+|k_{\mathfrak{f}^{s,2}}|)^{\beta}\varphi_{\leq K}(k_{\mathfrak{f}^{p}})\varphi_{\leq K}(k_{\mathfrak{f}^{s,1}})\varphi_{\leq K}(k_{\mathfrak{f}^{s,2}})\lesssim L^{\delta/2}.\label{bd:flowerB}
\end{align}

 The decay $\langle k-\zeta l\rangle^{-20}$ follows from the indicator $1_{k_\mathfrak{f}=l}$ replacing $n_{\rm in}(k_{\mathfrak{f}})$. Indeed, $k-\zeta l$ is a linear combination of $k_{\mathfrak{l}}$ for $\mathfrak{f}\neq \mathfrak{l}\in\mathcal{L}$, and the same argument as in \eqref{eq-sp1} applies. For the Dirac term $\delta(t_{\mathfrak{f}^p}-s)$, we return to \eqref{eq-sp1} and observe that omitting the integration over two flower nodes results in a loss of at most $L^{20}$. Finally, by the truncation condition and the relation
 $k_{\mathfrak{f}^{p}}=k_{\mathfrak{f}}-k_{\mathfrak{f}^{s,1}}+k_{\mathfrak{f}^{s,2}}$, we obtain $|k_\mathfrak{f}|\lesssim L^\delta$. Thus, as in Section~\ref{sec-boundoncouple}, we may assume $|k_l|\lesssim L^\delta$ for all $l\in\mathcal{L}$. The counting estimates from Section~\ref{sec:mainest} then remain valid.
\
\end{proof}


To prove Proposition~\ref{prop:l}, we apply large deviation estimates to obtain pathwise bounds with a high probability. We begin with the following result from \cite[Lemma 6.3]{Wu25}:


\begin{lemma}[Gaussian Hypercontractivity]\label{lem:Gaussian Hypercontractivity}
     Let $\{\eta_k\}$ be i.i.d. Gaussian variables.  Given $\zeta_j\in\{+,-\}$ and define 
 $$X=\sum_{k_i,1\leq i\leq n}a_{k_1,...,k_n}\prod^n_{j=1}\eta^{\zeta_j}_{k_j}(\omega),$$
 where $a_{k_1,...,k_n}$ are constants. Then, for $q\geq2$, it holds
$$ \mathbf{E}[|X|^q]\leq (q-1)^{nq/2}  \mathbf{E}[|X|^2]^{q/2}.$$
\end{lemma}

We are now in a position to prove Proposition~\ref{prop:l}.
\begin{proof}[Proof of Proposition \ref{prop:l}]
    By definition, we have
\begin{align*}
\left\|\mathscr{L}^{n}(a)\right\|_{Z}^{2}= & \sup _{0 \leq t \leq 1} L^{-1} \sum_{k \in \mathbb{Z}_{L}}\langle k\rangle^{10}\left|\sum_{\zeta \in\{ \pm\}} \sum_{l} \int_{0 \leq s \leq t} \sum_{n \leq m \leq N^{3}}\left(\mathscr{L}^{n}\right)_{k, l}^{m, \zeta} a_{l}^{\zeta}(s) \mathrm{d} s\right|^{2} \\
\lesssim & \sup _{0 \leq s \leq t \leq 1} \sup _{\zeta} \sup _{m} L^{-1} N^{3} \sum_{k \in \mathbb{Z}_{L}}\langle k\rangle^{10}\left| \sum_{l \in \mathbb{Z}_{L}}\langle l\rangle^{-5}\left(\mathscr{L}^{n}\right)_{k, l}^{m, \zeta}(t, s)\langle l\rangle^{5} a_{l}^{\zeta}(s)\right|^{2} \\
\lesssim & \|a\|_{Z}^{2} N^{3} \sup _{0 \leq s \leq t \leq 1} \sup _{\zeta} \sup _{m} \sum_{k \in \mathbb{Z}_{L}}\sum_{l \in \mathbb{Z}_{L}}\langle k\rangle^{10}\langle l\rangle^{-10}\left|\left(\mathscr{L}^{n}\right)_{k, l}^{m, \zeta}(t, s)\right|^2 \\
\lesssim & \|a\|_{Z}^{2} N^{3} \sup _{0 \leq s \leq t \leq 1} \sup _{\zeta} \sup _{m} \sup _{k, l}\left(\langle k-\zeta  l\rangle^{9}\left|\left(\mathscr{L}^{n}\right)_{k, l}^{m, \zeta}(t, s)\right|\right)^{2}\\
& \quad\times\sum_{k \in \mathbb{Z}_{L}}\sum_{l \in \mathbb{Z}_{L}}\langle k\rangle^{10}\langle l\rangle^{-10} \langle k-\zeta  l\rangle^{-18}\\
\lesssim & \|a\|_{Z}^{2}  L^2\sup _{0 \leq s \leq t \leq 1} \sup _{\zeta} \sup _{m} \sup _{k, l}\left(\langle k-\zeta  l\rangle^{9}\left|\left(\mathscr{L}^{n}\right)_{k, l}^{m, \zeta}(t, s)\right|\right)^{2}.
\end{align*}

It now suffices to bound the above expression with probability  $\geq 1 - L^{-40}$. From the definition of $\left(\mathscr{L}^{n}\right)_{k, l}^{m, \zeta}$ in \eqref{def:tildect}, we observe that there are at most $L^{1+\delta}\cdot L^{1+\delta}$ possible choices for $(k,l)$ (by an argument similar to \eqref{bd:flowerB}).
Moreover, the number of choices for $(\zeta,m)$ is finite and bounded by $2\cdot N^3$. Hence, it suffices to establish the desired bound for fixed $(\zeta,m,k,l)$.

Then, by \eqref{eq-Lnm} and an argument analogous to \cite[Proposition 12.1]{DH23a}, 
combining the Gaussian hypercontractivity (Lemma~\ref{lem:Gaussian Hypercontractivity}) 
with the Gagliardo--Nirenberg inequality, we obtain
\begin{align*}
  \mathbf{E}\left|  \sup _{0 \leq s \leq t \leq 1}   \sup _{k, l} \langle k-\zeta  l\rangle^{9}\left|\left(\mathscr{L}^{n}\right)_{k, l}^{m, \zeta}(t, s)\right|\right|^p\lesssim (p L^{40} (\alpha  T^{4/5} L^{2\delta})^{m})^{p/2}.
\end{align*}
Consequently, by Chebyshev’s inequality, with probability  $\geq 1 - L^{-40}$, we have

\begin{align*}
   \sup _{0 \leq s \leq t \leq 1}   \sup _{k, l}  \langle k-\zeta  l\rangle^{9}\left|\left(\mathscr{L}^{n}\right)_{k, l}^{m, \zeta}(t, s)\right|\lesssim L^{45} (\alpha  T^{4/5} L^{2\delta})^{n/2}.
\end{align*} 


\end{proof}
 

\subsection{The estimate on the remainder}

\begin{proposition}\label{prop:ckn}
      With probability $\geq 1-L^{-40}$, it holds that 
      \begin{align*}
        \langle k\rangle^{9}| c_k^{n}(t)|\leq  (p^2\alpha  T^{4/5}  L^{2\delta} )^{n}L,\ \ \langle k\rangle^{9}|(c_{\sim N})_k(t)|\leq (p^2\alpha  T^{4/5}  L^{2\delta} )^{N}L,
      \end{align*}
     for any $k \in\mathbb{Z}_L$, $t \in[0,1]$, and all $0 \leq  n \leq N^3$. Here $p>1$ is a universal constant. We recall \eqref{def:ck} and \eqref{eqnr} for the notations. 
  \end{proposition}
      
      Moreover, by the choice of $\alpha$ and $T$, we have
      \begin{align*}
          (p^2\alpha  T^{4/5}  L^{2\delta} )^{1/2}\leq L^{-\delta},
      \end{align*}
      which implies 
      \begin{align*}
          \|c^n\|_{Z}\lesssim L^{1-n\delta},\ \ \|c_{\sim N}\|_{Z}\lesssim L^{1-N\delta}.
      \end{align*}


\begin{proof}
   By Proposition~\ref{prop:couple}, we obtain
\begin{align*}
 \mathbf{E}| c_k^{n}(t)|^2
&\leq C \langle k \rangle^{-20} (\alpha T^{4/5} L^{2\delta})^{2n} T^{-3/5},\\
 \mathbf{E}|(c_{\sim N})_k(t)|^2
&\leq C \langle k \rangle^{-20} (\alpha T^{4/5} L^{2\delta})^{2N} T^{-3/5}.
\end{align*}
The estimate for $c_{\sim N}$ is obtained in the same way as for $c^n$, with $n$ replaced by $N$. Indeed, $c_{\sim N}$ is defined as the sum of $c_k^{\mathcal{T}^+}$ over all trees $\mathcal{T}^+$ of scale greater than $N$ whose three subtrees have scale at most $N$.
 Then, following the same argument as in \cite[Proposition 12.1]{DH23a}, we apply Gaussian hypercontractivity together with the Gagliardo–Nirenberg inequality to upgrade the above moment bounds to the desired   estimates. 
\end{proof}

The following proposition shows the existence of the remainder term $\mathcal{R}^{N+1}$.

\begin{proposition}\label{prop:rnsmall}
    With probability $\geq 1-L^{-40}$,  the mapping defined on the right hand  side of \eqref{def:rkn:fix} is a contraction mapping from the set $\{b : \|b\|_{ Z} \leq  L^{-500}\}$ to itself.
\end{proposition}
\begin{proof}
We first restrict to the complement of an exceptional set of probability at most $L^{-40}$.  
We then estimate the operators $\mathscr{Q}$ and $\mathscr{C}$ defined in \eqref{def:l}. By definition, we have
\begin{align}
  \mathcal{W}_1&(\tilde\varphi_{\leq K}u,\tilde\varphi_{\leq K}v,\tilde\varphi_{\leq K}w)_k(t)\notag\\
  &=\frac{i\alpha T}{L}  \sum\limits_{(k_1,k_2, k_{3})}^{\times} u_{k_1}{v_{k_2}}  w_{k_3} e^{2\pi i \Theta(k_1,k_2,k_3,k,t)}B_{k_1,k_2,k_3,k}\tilde\varphi_{\leq K}(k_1)\tilde\varphi_{\leq K}(k_2)\tilde\varphi_{\leq K}(k_3)\varphi_{\leq K}(k)\notag\\
&\quad -\frac{i\alpha T}{L}u_{k}{v_{k}}w_{k}B_{k,k,k,k}\varphi_{\leq K}(k)\tilde\varphi_{\leq K}(k)^3\notag.    
\end{align}
Here we recall that $\tilde\varphi_{\leq K}u=(\tilde\varphi_{\leq K}(k)u_k)_{k\in\mathbb{Z}_L}$ (see Section~\ref{sec:rem}). By the truncation properties of $\tilde\varphi_{\leq K}$, we obtain
 
\begin{align*}
B_{k_1,k_2,k_3,k}\tilde\varphi_{\leq K}(k_1)\tilde\varphi_{\leq K}(k_2)\tilde\varphi_{\leq K}(k_3)\varphi_{\leq K}(k)\leq L^\delta,\\
B_{k,k,k,k}\varphi_{\leq K}(k)\tilde\varphi_{\leq K}(k)^3\leq L^\delta.
\end{align*}
Consequently,
\begin{align}\label{bd:est:w1}
  &\bigg\| \int_0^t\mathcal{W}_1(\tilde\varphi_{\leq K}u,\tilde\varphi_{\leq K}v,\tilde\varphi_{\leq K}w)(s)\dif s\bigg\|_{Z}\notag\\
  & \leq  \int_0^t\big\|\mathcal{W}_1(\tilde\varphi_{\leq K}u,\tilde\varphi_{\leq K}v,\tilde\varphi_{\leq K}w)(s)\big\|_{Z}\dif s\leq L^{21}\|u\|_Z\|v\|_Z\|w\|_Z,
   \end{align}
   and similarly,
   \begin{align*}
    \bigg\| \int_0^t(\mathscr{Q}_1)_k[\tilde\varphi_{\leq K}u](\tilde\varphi_{\leq K}v)\dif s\bigg\|_{Z}&\leq L^{21}\|u\|_Z\|v\|_Z^2,\\
    \bigg\| \int_0^t  (\mathscr{C}_1)_k(\tilde\varphi_{\leq K}v)\dif s\bigg\|_{Z}&\leq L^{21}\|v\|_Z^3.
\end{align*}
Here we recall the definition of $\mathscr{Q}_1$ and $\mathscr{C}_1$ in Section \ref{sec:rem}.

 Using that $\tilde\varphi_{\leq K}c^{\leq N}=c^{\leq N}$, we rewrite \begin{align*} 
\mathscr{Q}(\tilde\varphi_{\leq K}v)&:=\int_0^t \big(2\mathcal{W}_1(\tilde\varphi_{\leq K}v,\tilde\varphi_{\leq K}\overline{v}, \tilde\varphi_{\leq K}c^{\leq N})\\
&\quad\quad\quad +\mathcal{W}_1(\tilde\varphi_{\leq K}v, \tilde\varphi_{\leq K}\overline{c^{\leq N}},\tilde\varphi_{\leq K} v)+\mathscr{Q}_1[\tilde\varphi_{\leq K}c^{\leq N}](\tilde\varphi_{\leq K}v)\big)\dif s,\\
\mathscr{C}(\tilde\varphi_{\leq K}v)&:=\int_0^t \mathcal{W}_1(\tilde\varphi_{\leq K}v, \tilde\varphi_{\leq K}\overline{v}, \tilde\varphi_{\leq K}v)+\mathscr{C}_1(\tilde\varphi_{\leq K}v)\dif s.
\end{align*} 

Let $\|b\|_{Z}\leq L^{-500}$. By Proposition~\ref{prop:ckn} (for sufficiently large $N$), we obtain, with probability   $\geq 1-L^{-40}$,
\begin{align}
   \|c_{\sim N}& +\mathscr{Q}(\tilde\varphi_{\leq K}b)+\mathscr{C}(\tilde\varphi_{\leq K}b)\|_Z\notag\\
   &\leq \|c_{\sim N}\|_Z + L^{21}\|b\|_Z^2(\sum_{n\leq N}\|c^n\|_Z)+ L^{21}\|b\|_Z^3\leq  L^{-600}.\label{bd:jlc}
\end{align}
Then we write  
\begin{align*}
    (1 -\mathscr{L})^{-1} =(1-\mathscr{L}^N)^{-1}(1+\mathscr{L} +...+\mathscr{L}^{N-1}),
\end{align*}
which is a map on $Z$. Here we understand  $(1-\mathscr{L}^N)^{-1}$  by the Neumann series. By applying Proposition \ref{prop:l}, we have 
\begin{align}\label{bd:1-l}
    \| (1 -\mathscr{L})^{-1} \|_{Z\to Z} \lesssim L^{70}.
\end{align}
Together with \eqref{bd:jlc} we finish the proof.
\end{proof}
Then, by a fixed point argument together with Proposition~\ref{prop:rnsmall}, we obtain that, with probability   $\geq 1 - L^{-40}$, there exists a solution $\mathcal{R}^{N+1}$ to \eqref{def:rkn:fix} satisfying  \begin{align}
     \|\mathcal{R}^{N+1}\|_{ Z} \leq  L^{-500}.\label{bd:mathcalr}
\end{align} 


\section{Proof of Theorem \ref{thm2}}\label{sec:proof of theorem 2} 


\begin{proof}[Proof of Theorem \ref{thm2}]
We denote by $E$ the complement of all exceptional sets appearing in Proposition~\ref{prop:l} and Proposition~\ref{prop:rnsmall}, so that
     $$\mathbf{P}(E) \geq  1-L^{-40}.$$
 By the reductions in Section~\ref{reductions}, we have
 $$\mathbf{E}|\widehat{u}_{\rm tr}(k,t)|^2=\mathbf{E}|c_k(s)|^2,\ \ s = \frac tT.$$
 We further denote $\mathbf{E}_E[\ \cdot\ ]:=\mathbf{E}[ 1_E\cdot\ ]$. Then
 \begin{align*}
  \mathbf E|\widehat u_{\rm tr}(k, t)|^2&=   \mathbf{E}_E |c_k(s)|^2+O(L^{-40}) \\
  &=\mathbf{E}_E| c^{\leq N}_k (s)|^2 +2\mathrm{Re}\mathbf{E} _E[c^{\leq N}_ k (s)\overline{\mathcal{R}^{N+1}_k (s)}] +\mathbf{E}_E |\mathcal{R}^{N+1}_k (s)|^2+O(L^{-40}) .
 \end{align*}
 For the terms involving the remainder, by Proposition~\ref{prop:ckn} and \eqref{bd:mathcalr}, we obtain
 \begin{align*}
     2\mathrm{Re}\mathbf{E} _E[c^{\leq N}_ k (s)\overline{\mathcal{R}^{N+1}_k (s)}] +\mathbf{E}_E |\mathcal{R}^{N+1}_k (s)|^2=O(L^{-10}). 
 \end{align*}
Next, using Proposition~\ref{prop:couple}, we estimate
\begin{align*}
    \mathbf{E}_E| c^{\leq N}_k (s)|^2%&= \sum_{n=0}^{2N}\sum_{|\mathcal{T}^+|+|\mathcal{T}^-|=n}\mathbf{E}_E[c^{\mathcal{T}^+}_kc^{\mathcal{T}^-}_k(s)]\notag\\ 
    &\leq \sum_{n=0}^{2}\sum_{|\mathcal{T}^+|+|\mathcal{T}^-|=n}\mathbf{E}_E[c^{\mathcal{T}^+}_kc^{\mathcal{T}^-}_k(s)]+\sum_{n=3}^{2N} C \langle k\rangle^{-20}  (\alpha  T^{4/5}  L^{2\delta})^nT^{- \frac{3}{5}}\\
    &\leq \sum_{n=0}^{2}\sum_{|\mathcal{T}^+|+|\mathcal{T}^-|=n}\mathbf{E}_E[c^{\mathcal{T}^+}_kc^{\mathcal{T}^-}_k(s)]+C\frac{T}{T_{\mathrm{kin}}} \cdot\frac{\alpha  T^{4/5}  L^{6\delta}}{1-\alpha  T^{4/5}  L^{2\delta}}\\
    &=\sum_{n=0}^{2}\sum_{|\mathcal{T}^+|+|\mathcal{T}^-|=n}\mathbf{E}_E[c^{\mathcal{T}^+}_kc^{\mathcal{T}^-}_k(s)]+O(L^{-\delta}\frac{T}{T_{\mathrm{kin}}}).
\end{align*}
Here we used the fact that $\alpha  T^{4/5}  L^{6\delta}\leq L^{-\delta}.$

It remains to estimate the lower-order terms appearing in the above expression.
For the case $n=0$, by definition,   $c_k^{\bullet}(s)= \varphi_{\leq K}(k)\sqrt{n_{\mathrm{in}}(k)}\cdot\eta_k(\omega)^\zeta$, where $\zeta=\pm$ depends on the sign of the root.
Taking expectation, we obtain that the zeroth-order contribution is  $\varphi_{\leq K}^2(k)\, n_{\mathrm{in}}(k)$.

 For the case $n=1$, we have $\mathbf{E}_E[c^{\mathcal{T}^+}_kc^{\mathcal{T}^-}_k(s)]=0$,  since there is no admissible  decorations in this case.

 It remains to consider the case $n=2$. In this case, we need to estimate the terms $\mathbf{E}|c^1_k(s)|^2$ and $2\mathrm{Re}\mathbf{E}\left[\overline{c^0_k(s)}\, c^2_k(s)\right]$. For the former, by definition, we have
 
\begin{align}
c_{k}^{1}(s)&=\frac{i\alpha T}{L}  \sum\limits_{(k_1,k_2, k_{3})}^{\times} (a_{\rm in})_{k_1}\overline {(a_{\rm in})_{k_2}}  (a_{\rm in})_{k_3} \int_0^se^{2\pi i[\overline\Omega_kTs+\sum\Gamma_1(k,s)]}\dif sB_{k_1,k_2,k_3,k}\varphi_{\leq K}(k)\notag\\
&\quad+\frac{i\alpha Ts}{L} \left(|(a_{\rm in})_{k}|^2 -2\mathbf{E}|(a_{\rm in})_{k}|^2  \right) B_{k,k,k,k}\varphi_{\leq K}(k) (a_{\rm in})_{k}\notag\\
&\quad+2\frac{i\alpha Ts}{L}\sum_{k_1\neq k}\left(|(a_{\rm in})_{k_1}|^2 -\mathbf{E}|(a_{\rm in})_{k_1}|^2  \right) B_{k_1,k_1,k,k}\varphi_{\leq K}(k) (a_{\rm in})_{k}\notag\\
&=\frac{i\alpha T}{L} \sum\limits_{(k_1,k_2, k_{3}),\overline \Omega_k\neq 0}^{\times} (a_{\rm in})_{k_1}\overline {(a_{\rm in})_{k_2}}  (a_{\rm in})_{k_3} \frac{e^{2\pi i\overline \Omega_kTs}-1}{2\pi i\overline\Omega_kT}B_{k_1,k_2,k_3,k}\varphi_{\leq K}(k)\notag\\
&\quad+\frac{i\alpha T}{L}  \sum\limits_{(k_1,k_2, k_{3}),\overline \Omega_k\neq 0}^{\times} (a_{\rm in})_{k_1}\overline {(a_{\rm in})_{k_2}}  (a_{\rm in})_{k_3}\frac{Er_1(s)}{T\overline\Omega_k}B_{k_1,k_2,k_3,k}\varphi_{\leq K}(k)\notag\\
&\quad+\frac{i\alpha T}{L} \sum\limits_{(k_1,k_2, k_{3}),\overline \Omega_k= 0}^{\times} (a_{\rm in})_{k_1}\overline {(a_{\rm in})_{k_2}}  (a_{\rm in})_{k_3} \int_0^se^{2\pi i\sum\Gamma_1(k,u)}\dif uB_{k_1,k_2,k_3,k}\varphi_{\leq K}(k)\notag\\
&\quad+\frac{i\alpha Ts}{L} \left(|(a_{\rm in})_{k}|^2 -2\mathbf{E}|(a_{\rm in})_{k}|^2  \right) B_{k,k,k,k}\varphi_{\leq K}(k) (a_{\rm in})_{k}\notag\\
&\quad+2\frac{i\alpha Ts}{L}\sum_{k_1\neq k}\left(|(a_{\rm in})_{k_1}|^2 -\mathbf{E}|(a_{\rm in})_{k_1}|^2  \right) B_{k_1,k_1,k,k}\varphi_{\leq K}(k) (a_{\rm in})_{k},\label{bd:ck1}
\end{align}
where we define the error term by
\begin{align}
    Er_1(s):= -\int_0^se^{2\pi i[ \overline\Omega_kTu+\sum\Gamma_1(k,u)]}\sum\partial_u\Gamma_1(k,u)\dif u+ \frac{e^{2\pi i[ \overline\Omega_kTs+\sum\Gamma_1(k,s)]}-e^{2\pi i \overline\Omega_kTs}}{2\pi i}.\notag
\end{align}
By \eqref{bd:gamma'}, we obtain
\begin{align}
   | Er_1(s)|&\lesssim s\sup_{u\in[0,s]}\sup_k|\partial_t\Gamma_1(k,u)|+|e^{2\pi i \sum\Gamma_1(k,s)}-1|\notag\\
   &\lesssim s \|\partial_t\Gamma_1(k,s)\|_{L_k^\infty C^0_s}+|  \sum\Gamma_1(k,s)|\lesssim s \|\partial_t\Gamma_1(k,s)\|_{L_k^\infty C^0_s}\lesssim sL^{-3\delta}.\label{def:er}
\end{align}

We denote each term on the right-hand side of \eqref{bd:ck1} by $C_i(s)$ for $1 \leq i \leq 5$, and estimate them in turn.   By Lemma~\ref{thm:Isserlis} together with a direct computation, we obtain

\begin{align*}
\mathbf{E}|C_1(s)|^2&=\frac{2\alpha^2T^2s^2
}{L^{2}}\sum_{(k_1, k_2, k_3); \overline \Omega_k\neq 0}^\times n_{\mathrm{in}}(k_1)n_{\mathrm{in}}(k_2)n_{\mathrm{in}}(k_3)\left|\frac{\sin \pi\overline \Omega _kTs}{\pi \overline\Omega _kTs}\right|^2\\
&\quad\quad\times B^2_{k_1,k_2,k_3,k}\varphi_{\leq K}^2(k_1)\varphi_{\leq K}^2(k_2)\varphi_{\leq K}^2(k_3)\varphi_{\leq K}^2(k),\\
\mathbf{E}|C_4(s)|^2&=\frac{\alpha^2 T^2s^2}{L^{2}}\mathbf{E}[(|\eta_{k}(\omega)|^2 -2 )^2|\eta_{k}(\omega)|^2]B^2_{k,k,k,k}\varphi_{\leq K}^8(k)n_{\mathrm{in}}(k)|n_{\mathrm{in}}(k)|^2\\
&=\frac{2\alpha^2 T^2s^2}{L^{2}}B^2_{k,k,k,k}\varphi_{\leq K}^8(k)n_{\mathrm{in}}^3(k),\\
    \mathbf{E}|C_5(s)|^2&=\frac{4\alpha^2 T^2s^2}{L^{2}}\sum_{k_1\neq k}\mathbf{E}(|\eta_{k_1}(\omega)|^2 -1 )^2 \mathbf{E}|\eta_{k}(\omega)|^2 B^2_{k_1,k_1,k,k}\varphi_{\leq K}^4(k_1)\varphi_{\leq K}^4(k)n_{\mathrm{in}}(k)|n_{\mathrm{in}}(k_1)|^2\\
    &=\frac{4\alpha^2 T^2s^2}{L^{2}}\sum_{k_1\neq k}B^2_{k_1,k_1,k,k}\varphi_{\leq K}^4(k_1)\varphi_{\leq K}^4(k)n_{\mathrm{in}}(k)n_{\mathrm{in}}^2(k_1).
\end{align*}
Moreover, by Lemma~\ref{thm:Isserlis} and the resonance relations, we have for $1\leq i\leq 4$, \begin{align*}
   \mathrm{Re} \mathbf{E}[C_i(s)\overline{C_5(s)}]=0,
\end{align*} 
and   \begin{align*}
   \mathrm{Re} \mathbf{E}[(C_1(s)+C_2(s)+C_3(s))\overline{C_4(s)}]+\mathrm{Re} \mathbf{E}[(C_1(s)+C_2(s))\overline{C_3(s)}]=0.
\end{align*}  

For the term $C_3$, since $\overline \Omega_k=0$, we apply the 3-counting estimate in Lemma~\ref{counting} with $T$ replaced by $T L^{4\delta}$ and choose $D=L^\delta$. Then
\begin{align*}
   \mathbf{E}|C_3(s)|^2 &\lesssim \frac{\alpha^2 T^2s^2}{L^{2}} \sum\limits_{(k_1,k_2, k_{3}),\overline \Omega_k= 0}^{\times}  n_{\mathrm{in}}(k_1)n_{\mathrm{in}}(k_2)n_{\mathrm{in}}(k_3)
B^2_{k_1,k_2,k_3,k}\varphi_{\leq K}^2(k_1)\varphi_{\leq K}^2(k_2)\varphi_{\leq K}^2(k_3)\varphi_{\leq K}^2(k)\notag\\
   &\lesssim \frac{\alpha^2 T^2s^2}{L^{2}}\cdot  L^{2-4\delta}T^{-1}\log L\cdot L^\delta=O(\frac{t}{T_{\rm kin}}L^{-\delta}).
\end{align*}
%By a similar argument as before, for the case $|k_i|\geq D$, we use the decay of $n_{\mathrm{in}}(k_i)$ to achieve the desired result.


Similarly, applying the 3-counting estimate in Lemma~\ref{counting} again, we obtain
\begin{align*}
\mathbf{E}&[C_1\overline{C_2}(s)+C_2\overline{C_1}(s)+|C_2(s)|^2]\\
&\lesssim\frac{\alpha^2T^2}{L^{2}}\sum_{(k_1, k_2, k_3); \overline \Omega_k\neq 0}^\times n_{\mathrm{in}}(k_1)n_{\mathrm{in}}(k_2)n_{\mathrm{in}}(k_3)\frac{|Er_1|^2+|Er_1|}{|T\overline \Omega_k|^2} \notag\\
&\quad \quad \times B^2_{k_1,k_2,k_3,k}\varphi_{\leq K}^2(k_1)\varphi_{\leq K}^2(k_2)\varphi_{\leq K}^2(k_3)\varphi_{\leq K}^2(k)\\
&\lesssim\frac{\alpha^2T^2}{L^{2}} \cdot L^{2}T^{-1}\log L\cdot \sum_{j\neq 0}\frac{sL^{-3\delta}}{j^2}\cdot L^\delta=O(\frac{t}{T_{\rm kin}}L^{-\delta}).
\end{align*}

 
In summary, we obtain
\begin{align*}
     \mathbf{E}|c^1_k(s)|^2=\mathbf{E}|C_1(s)|^2+\mathbf{E}|C_4(s)|^2+\mathbf{E}|C_5(s)|^2+O(\frac{t}{T_{\rm kin}}L^{-\delta}).
\end{align*}


For the term $2\mathbf{E}\mathrm{Re}\overline{c^0_k(s)}c^2_k(s)$, we assume that the root and the other branching node are decorated by $k$ and $k'$, respectively. A direct inspection of the admissible decorations shows that the two branching nodes must be degenerate simultaneously. As a consequence, only the configurations illustrated in Figure~\ref{fig:type2tree} are possible for the corresponding enhanced couples $\mathcal{Q}$.
Moreover, in each of these cases, the resonance relations imply that  $\overline \Omega_k+\overline \Omega_{k'}=0$ and $\sum\Gamma_1(k,s)+\sum\Gamma_1(k',s)=0$.

\begin{figure}
\begin{subfigure}%[t]{1\linewidth}
\centering
\begin{tikzpicture}[level distance=.8cm,
  level 1/.style={sibling distance=.8cm},
  level 2/.style={sibling distance=.8cm}]
\tikzstyle{hollow node}=[circle,draw,inner sep=1.6]
\tikzstyle{solid node}=[circle,draw,inner sep=1.6,fill=black]
\tikzset{
red node/.style = {circle,draw=black,fill=red,inner sep=1.6},
blue node/.style= {circle,draw = black, fill= blue,inner sep=1.6}, 
purple node/.style= {circle,draw = black, fill= purple,inner sep=1.6}, 
orange node/.style= {circle,draw = black, fill= orange,inner sep=1.6},
yellow node/.style= {circle,draw = black, fill= yellow,inner sep=1.6},
green node/.style = {circle,draw=black,fill=green,inner sep=1.6}}


\node[solid node, label = right: $k$ ] at (-5.5,.5){}
    child{node[solid node, label = left: $k_1$ ]{}
        child{node[blue node, label=below: $k_2$]{}}
        child{node[red node, label=below: $k_3$]{}}
        child{node[green node, label=below: $k$]{}}
    }
    child{node[blue node, label=below: $k_2$]{}}
    child{node[red node, label=below: $k_3$]{}}
;
\node[green node, label=right: $k$] at (-3.5,.5){};



\node[solid node, label = right: $k$ ] at (.5,.5){}
    child{node[solid node, label = left: $k$ ]{}
        child{node[green node, label=below: $k$]{}}
        child{node[green node, label=below: $k$]{}}
        child{node[green node, label=below: $k$]{}}
    }
    child{node[green node, label=below: $k$]{}}
    child{node[green node, label=below: $k$]{}}
;
\node[green node, label=right: $k$] at (3,.5){};



\node[solid node, label = right:{$k$}] at (-5.5,-2.7){}
    child{node[solid node, label = left: $k_1$ ]{}
        child{node[blue node, label=below: $k_1$]{}}
        child{node[red node, label=below: $k$]{}}
        child{node[green node, label=below: $k$]{}}
    }
    child{node[blue node, label=below: $k_1$]{}}
    child{node[red node, label=below: $k$]{}}
;
\node[green node, label=right: $k$] at (-3,-2.7){};

 

\node[solid node, label = right:{$k$}] at (.5,-2.7){}
    child{node[solid node, label = left: $k$ ]{}
        child{node[green node, label=below: $k$]{}}
        child{node[red node, label=below: $k_3$]{}}
        child{node[blue node, label=below: $k_3$]{}}
    }
    child{node[blue node, label=below: $k_3$]{}}
    child{node[red node, label=below: $k_3$]{}}
;
\node[green node, label=right: $k$] at (3,-2.7){};

\end{tikzpicture}
\end{subfigure}
\caption{All possible types of couples $\mathcal{Q}$ of order $n=2$, where the second branching node is the left child of the root. Symmetry should also be taken into account.
Nodes with the same color are paired with each other. The first configuration in the first row correspond to the non-degenerate case, while the two in the second row correspond to the degenerate case.}
\label{fig:type2tree}
\end{figure}

Now, similarly to \eqref{bd:ck1}, we distinguish four cases: the non-degenerate case with $\overline \Omega_k,\overline \Omega_{k'}\neq 0$; the non-degenerate case with $\overline \Omega_k=\overline \Omega_{k'}=0$; the fully degenerate case; and the partially degenerate case with $k\neq k_1$.

We begin with the non-degenerate case with $\overline \Omega_k,\overline \Omega_{k'}\neq 0$. 
In this case, the time integral $\mathcal{A}_{\mathcal{T}}$ takes the form
 \begin{align*}
     &\int_0^se^{2\pi i[\overline \Omega_kTs_1+\sum\Gamma_1(k,s_1)]}\int_0^{s_1}e^{2\pi i[\overline \Omega_{k'}Ts_2+\sum\Gamma_1(k',s_2)]}\dif s_2\dif s_1\\
     &=  \int_0^se^{2\pi i[\overline \Omega_kTs_1+\sum\Gamma_1(k,s_1)]}\notag\\ 
     &\quad\quad\quad\quad\times\Big[\frac{e^{2\pi i[\overline \Omega_{k'}Ts_1+\sum\Gamma_1(k',s_1)]}-1}{2\pi i\overline \Omega_{k'}T}-\int_0^{s_1}\frac{\partial_t\sum\Gamma_1(k',s_2)}{\overline \Omega_{k'}T}e^{2\pi i[\overline \Omega_{k'}Ts_2+\sum\Gamma_1(k',s_2)]}\dif s_2\Big]\dif s_1\\
      &=\frac1{2\pi i\overline \Omega_{k'}T}  \int_0^s1-e^{2\pi i[\overline \Omega_kTs_1+\sum\Gamma_1(k,s_1)]}\dif s_1+\widetilde {Er}_2=\frac{-s}{2\pi i\overline \Omega_{k}T}  -\frac{  e^{2\pi i\overline \Omega_kTs}-1}{4\pi^2|\overline \Omega_{k}T|^2}  +Er_2.
 \end{align*}
where we used the fact that $\overline \Omega_k+\overline \Omega_{k'}=0$ and $\sum\Gamma_1(k,s)+\sum\Gamma_1(k',s)=0$ among all the possible pairs. Here the error term  $Er_2$  is defined as 
 \begin{align*}
     Er_2&:=-\frac{  e^{2\pi i[\overline \Omega_kTs+\sum\Gamma_1(k,s)]}- e^{2\pi i\overline \Omega_kTs}}{4\pi^2|\overline \Omega_{k}T|^2}\\
     &\quad-\int_0^se^{2\pi i[\overline \Omega_kTs_1+\sum\Gamma_1(k,s_1)]} \int_0^{s_1}\frac{\partial_t\sum\Gamma_1(k',s_2)}{\overline \Omega_{k'}T}e^{2\pi i[\overline \Omega_{k'}Ts_2+\sum\Gamma_1(k',s_2)]}\dif s_2 \dif s_1\\
     &\quad+\int_0^s \frac1{2\pi i\overline \Omega_{k'}T}  \frac{\partial_t\sum\Gamma_1(k,s_1)}{\overline \Omega_{k}T}   e^{2\pi i[\overline \Omega_kTs_1+\sum\Gamma_1(k,s_1)]}\dif s_1.
 \end{align*}
We now estimate $Er_2$. For the second term above, we apply Proposition~\ref{prop:est:af} with
  $$f=\int_0^{s_1}\partial_t\sum\Gamma_1(k',s_2)e^{2\pi i[\overline \Omega_{k'}Ts_2+\sum\Gamma_1(k',s_2)]}\dif s_2.$$By \eqref{bd:gamma'}, we have
 \begin{align*}
     |f|+|\partial_tf|\lesssim \sup_k \|\partial_t\Gamma_1(k,s)\|_{C^0_s}\leq L^{-3\delta}.
 \end{align*}
 Hence, Proposition~\ref{prop:est:af} yields
\begin{align*}
   \left| \int_0^se^{2\pi i[\overline \Omega_kTs_1+\sum\Gamma_1(k,s_1)]} \int_0^{s_1}\frac{\partial_t\sum\Gamma_1(k',s_2)}{\overline \Omega_{k'}T}e^{2\pi i[\overline \Omega_{k'}Ts_2+\sum\Gamma_1(k',s_2)]}\dif s_2 \dif s_1\right|\lesssim s\frac{\sup_k \|\partial_t\Gamma_1(k,s)\|_{C^0_s}}{|T\overline\Omega_k|^2}..
\end{align*}
Consequently,
 \begin{align*}
     | Er_2|\lesssim \sup_k\sup_{s\in[0,1]}\frac{|\sum\Gamma_1(k,s)|+s\|\partial_t\Gamma_1(k,s)\|_{C^0_s}}{|\overline \Omega_{k}T|^2}\lesssim \frac{ s \|\partial_t\Gamma_1(k,s)\|_{C^0_sL_k^\infty}}{|\overline \Omega_{k}T|^2}\lesssim \frac{ sL^{-3\delta}}{|\overline \Omega_{k}T|^2}.
 \end{align*}
 
The contribution of this case to $2\,\mathbf{E}\mathrm{Re}\,\overline{c^0_k(s)}c^2_k(s)$ is
 

\begin{align*}
2\bigg(\frac{\alpha T}{L}\bigg)^2&\sum\limits_{(k_1,k_2, k_{3}),\overline \Omega_k\neq  0}^{\times}n_{\mathrm{in}}(k_1)n_{\mathrm{in}}(k_2)n_{\mathrm{in}}(k_3)n_{\mathrm{in}}(k)\left(\frac{1}{n_{\mathrm{in}}(k_1)}-\frac{1}{n_{\mathrm{in}}(k_2)}+\frac{1}{n_{\mathrm{in}}(k_3)}\right)\\
&\times \Big(\left|\frac{\sin \pi \overline \Omega _kTs}{\pi\overline  \Omega _kT}\right|^2+2\mathrm{Re}(Er_2) \Big)
B^2_{k_1,k_2,k_3,k}\varphi_{\leq K}^2(k_1)\varphi_{\leq K}^2(k_2)\varphi_{\leq K}^2(k_3)\varphi_{\leq K}^2(k),
\end{align*}
where the term with $\mathrm{Re}(Er_2) $ is bounded by using 3 number counting estimate in Lemma \ref{counting} by:
\begin{align*}
    \frac{\alpha^2 T^2}{L^{2}}\cdot L^{2}T^{-1}\log L\cdot sL^{-3\delta}\cdot\sum_{j\neq 0}\frac{1}{j^2}\cdot L^\delta=O(\frac{t}{T_{\rm kin}}L^{-\delta}).
\end{align*}

Second, we consider the non-degenerate case with $\overline \Omega_k=\overline \Omega_{k'}=0$ and $k\neq k_1,k_3$. Its contribution is bounded by
\begin{align*}
\bigg(\frac{\alpha T}{L}\bigg)^2&\sum\limits_{(k_1,k_2, k_{3}),\overline \Omega_k= 0}^{\times}n_{\mathrm{in}}(k_1)n_{\mathrm{in}}(k_2)n_{\mathrm{in}}(k_3)n_{\mathrm{in}}(k)\left(\frac{1}{n_{\mathrm{in}}(k_1)}+\frac{1}{n_{\mathrm{in}}(k_2)}+\frac{1}{n_{\mathrm{in}}(k_3)}\right)\\
&\times 
B^2_{k_1,k_2,k_3,k}\varphi_{\leq K}^2(k_1)\varphi_{\leq K}^2(k_2)\varphi_{\leq K}^2(k_3)\varphi_{\leq K}^2(k),
\end{align*}
which is bounded by $O\!\left(\frac{t}{T_{\rm kin}}L^{-\delta}\right)$ by the same argument as for $C_3$. Here all terms are nonnegative, as we take absolute values. 

Third, we consider the fully degenerate case with  $k= k_1=k_3=k_2$, its contribution is 
\begin{align*}
   2i^2(1-1+1)\frac{\alpha^2 T^2s^2}{L^{2}}B^2_{k,k,k,k}\varphi_{\leq K}^8(k)n_{\mathrm{in}}(k)|n_{\mathrm{in}}(k)|^2=-\mathbf{E}|C_4(s)|^2,
\end{align*}
where the factor $1-1+1$ corresponds to the three possible positions (left, middle, or right) of the second branching node relative to the root. The prefactor $2$ reflects that there are exactly two admissible pairings, since the root node can only pair with two leaves carrying the $+$ sign (see Figure~\ref{fig:type2tree}).
 

Finally, we consider the partially degenerate case with $k_1\neq k$ or $k_3\neq k$. Its contribution is
\begin{align*}
    4i^2(1-1+1)\frac{\alpha^2 T^2s^2}{L^{2}}\sum_{k_1\neq k}B^2_{k_1,k_1,k,k}\varphi_{\leq K}^4(k_1)\varphi_{\leq K}^4(k)n_{\mathrm{in}}(k)|n_{\mathrm{in}}(k_1)|^2=-\mathbf{E}|C_5(s)|^2,
\end{align*}
where the factor $4$ comes from the four admissible pairings, corresponding to permutations of the two $+$ leaves in Figure~\ref{fig:type2tree}.
 




Combining the above estimates, we obtain
\begin{align*}\mathbf{E}|c^1_k(t)|^2+2\mathrm{Re}\mathbf{E}[\overline{c^0_k(t)}c^2_k(t)]=\frac{2\alpha^2t^2}{L^{2}}\cdot\mathscr{S}_{t}(n_{\mathrm{in}}) +O(\frac{t}{T_{\rm kin}}L^{-\delta}),\end{align*} where we recall  \begin{align*}
 \mathscr S_t(\phi):&=\sum\limits_{\substack{ k-k_1+k_2-k_3=0,\\ \overline \Omega_k\neq0}}\phi_k \phi_{k_1} \phi_{k_2} \phi_{k_3}  \left[ \frac{1}{\phi_k} - \frac{1}{\phi_{k_1}} + \frac{1}{\phi_{k_2}} - \frac{1}{\phi_{k_3}} \right]\left| \frac{\sin(\pi t\overline \Omega_k)}{\pi t \overline \Omega_k} \right|^2\notag\\
  &\quad\quad\quad\times
|k_1|^{2\beta}|k_2|^{2\beta}|k_3|^{2\beta}|k|^{2\beta}\varphi_{\leq K}^2(k_1)\varphi_{\leq K}^2(k_2)\varphi_{\leq K}^2(k_3)\varphi_{\leq K}^2(k).
\end{align*} 
Consequently,
\begin{align*}\mathbf{E}|\widehat{u}_{\rm tr}(k,t)|^2=\varphi_{\leq K}^2(k)n_{\mathrm{in}}(k)+\frac{2\alpha^2t^2}{L^{2}}\mathscr{S}_{t}(n_{\mathrm{in}})(k)+O(\frac{t}{T_{\rm kin}}L^{-\delta }).\end{align*}

By the asymptotic estimate in Proposition~\ref{asymptotic}, we further obtain \begin{align*}\mathbf{E}|\widehat{u}_{\rm tr}(k,t)|^2=\varphi_{\leq K}^2(k)n_{\mathrm{in}}(k)+\frac{2\alpha^2t^2}{L^{2}}\mathscr{K}_{t}(n_{\mathrm{in}})(k)+O(\frac{t}{T_{\rm kin}}L^{-\delta}),\end{align*}
 where we recall $\mathscr K_t$ is defined by 

\begin{align*}
\mathscr K_t(\phi):&=L^{2} \int_{\xi_1-\xi_2+\xi_3=\xi }\phi(\xi) \phi(\xi_1) \phi(\xi_2) \phi(\xi_3)  \left[ \frac{1}{\phi(\xi)} - \frac{1}{\phi(\xi_1)} + \frac{1}{\phi(\xi_2)} - \frac{1}{\phi(\xi_3)} \right]    \left| \frac{\sin(\pi t\overline \Omega( \xi))}{\pi t\overline\Omega( \xi)} \right|^2 \notag\\ &\quad\quad\times |\xi_1|^{2\beta}|\xi_2|^{2\beta}|\xi_3|^{2\beta}|\xi|^{2\beta}\varphi_{\leq K}^2(\xi_1)\varphi_{\leq K}^2(\xi_2)\varphi_{\leq K}^2(\xi_3)\varphi_{\leq K}^2(\xi)   \dif \xi_1  \dif\xi_2 \, \dif\xi_3.
\end{align*}

Finally, the proof follows from the identity that, for any smooth function $f$,
\begin{align*}
t \int \left| \frac{\sin(\pi t x)}{\pi t x} \right|^2 f(x)\dif x =  f(0) + O(t^{-1}),
\end{align*}
which completes the proof.
\end{proof}




\section{An acceptable energy inequality}\label{sec:acceptineq}
In the previous sections, we derived the truncated wave kinetic equation   from the truncated MMT equation. 
From the present section on, we shift our focus to the high-frequency component of the full MMT equation.
The analysis is based on a suitable energy inequality, as stated in Theorem~\ref{thm:energy inequality}. 
The key idea is to show that the increment of the high-frequency energy can be controlled by an $L^\infty$-type norm of the solution. 
This control will play a crucial role when dealing with solutions arising from random initial data.
 

 
Throughout this section and the next section, we denote the nonlinear term  by 
\begin{align*}
    \mathscr{B}(u,v,w):= -\lambda^2 |\nabla|^\beta\left[|\nabla|^\beta u \cdot|\nabla|^\beta\overline v   \cdot|\nabla|^\beta w\right].
\end{align*}

Let $N_0\geq 10$. We define the operator $J^{N_0}$ via the Fourier multiplier
\begin{align*}
    m_{N_0}(\xi):=\varphi_{\geq 0}(\xi)\langle\xi\rangle^{N_0}.
\end{align*}
We then introduce the high-frequency energy profile on the time interval $[T_1,T_2]$:
\begin{align*}
    A(t):=\int |J^{N_0}u(t)|^2\dif x\leq \|u(t)\|_{H^{N_0}}^2\leq A^2,
\end{align*}
where the last inequality holds under the assumption of 
Theorem~\ref{thm:energy inequality}.  

Our goal is to derive the following energy increment estimate for $A(t)$:
 \begin{proposition}\label{prop:highenergy} 
Under the assumption of Theorem \ref{thm:energy inequality}, it holds for $t_1\leq t_2\in[T_1,T_2]$ that: 
     \begin{align*}
  |   A(t_2)-A(t_1)|&\lesssim\lambda^2[\|u(t_1)\|_{H^{N_0}}\epsilon(t_1)]^2+\lambda^2[\|u(t_2)\|_{H^{N_0}}\epsilon(t_2)]^2\\
  &\quad+\lambda^4A^2\int_{t_1}^{t_2}\epsilon^4(s)\dif s+1_{\{\sigma\in(0,1)\}}\lambda^2A^{2-\frac{1}{4N_0}}\int_{t_1}^{t_2}\epsilon^{2+\frac{1}{4N_0}}(s)\dif s.
\end{align*}
 \end{proposition}

The proof of Proposition~\ref{prop:highenergy}, which forms the core of this section, will be presented later. 
Assuming this result, Theorem~\ref{thm:energy inequality} follows directly by combining it with the conservation of mass. 



\begin{proof}[Proof of Theorem \ref{thm:energy inequality}]
In Proposition~\ref{prop:highenergy}, we derived an energy increment estimate 
for the high-frequency component. 
In the case $\sigma\in(1,2]$, using the smallness condition 
$\lambda^2 \epsilon^2(t)\ll 1$, we obtain 
 \begin{align*}
    A(t_2)&\leq A(t_1)+ \|u(t_1)\|_{H^{N_0}}^2+\frac12\|u(t_2)\|_{H^{N_0}}^2+\lambda^4A^2\int_{t_1}^{t_2}\epsilon^4(s)\dif s\\
    &\leq  2\|u(t_1)\|_{H^{N_0}}^2+\frac12\|u(t_2)\|_{H^{N_0}}^2+\lambda^4A^2\int_{t_1}^{t_2}\epsilon^4(s)\dif s.
\end{align*}

For the low-frequency part, we use the conservation of mass to deduce
\begin{align*}
    \|P_{\leq 4}u(t_2)\|_{H^{N_0}}^2&\lesssim \|u(t_2)\|_{L^2}^2=\|u(t_1)\|_{L^2}^2 \leq \|u(t_1)\|_{H^{N_0}}^2.
\end{align*}
 
Combining the high- and low-frequency estimates, we obtain
    \begin{align*}
        \|u(t_2)\|_{H^{N_0}}^2&\leq  A(t_2)+\|P_{\leq 4}u(t_2)\|_{H^{N_0}}^2\\
        &\leq   C\|u(t_1)\|_{H^{N_0}}^2+\frac12\|u(t_2)\|_{H^{N_0}}^2+\lambda^4A^2\int_{t_1}^{t_2}\epsilon^4(s)\dif s.
    \end{align*}
  Absorbing the term $\tfrac12\|u(t_2)\|_{H^{N_0}}^2$ into the left-hand side yields the desired estimate.

The case $\sigma\in(0,1)$ follows by the same argument.
\end{proof}
 

\subsection{ Proof of  Proposition \ref{prop:highenergy}}
 
Since $u$ solves \eqref{MMT}, a standard energy computation yields, for $t\in[T_1,T_2]$,
\begin{align*}
   i( A(t_2)-A(t_1))&=\int_{t_1}^{t_2}\langle i\partial_tJ^{N_0}u,J^{N_0}u\rangle+\langle J^{N_0}u, -i\partial_tJ^{N_0} u\rangle\dif t\\
    &=\int_{t_1}^{t_2} \langle \mathscr{B}(u,u,u),J^{2N_0}{u} \rangle- \langle J^{2N_0}u,\mathscr{B}({u},{u},u) \rangle\dif t.
\end{align*}
Here $\langle\cdot,\cdot\rangle$ denotes the standard complex $L^2$ inner product on $\mathbb{T}_L$. Passing to Fourier variables, we obtain
\begin{align*}
  |   A(t_2)-A(t_1)|&\lesssim\lambda^2\left|\int_{t_1}^{t_2} Q(u,\overline u,u,\overline u)\dif t\right|,
\end{align*}
where \begin{align*}
Q(u,\overline u,u,\overline u):=\frac{1}{L^{3}}\sum_{\xi_1-\xi_2+\xi_3-\xi_4=0}\widehat{u}_{\xi_1}\overline{\widehat{u}}_{\xi_2}\widehat{u}_{\xi_3}\overline{\widehat{u}}_{\xi_4}B_{\xi_1,\xi_2,\xi_3,\xi_4}(m_{N_0}(\xi_3)^2-m_{N_0}(\xi_4)^2).
\end{align*}
Here we recall  $B_{\xi_1,\xi_2,\xi_3,\xi_4}=|\xi_1|^\beta|\xi_2|^\beta|\xi_3|^\beta|\xi_4|^\beta.$ By symmetry under permutations of the indices, we rewrite $Q$ as
 \begin{align*}
Q(u,\overline u,u,\overline u)&=\frac{1}{2L^{3}}\sum_{\xi_1-\xi_2+\xi_3-\xi_4=0}\widehat{u}_{\xi_1}\overline{\widehat{u}}_{\xi_2}\widehat{u}_{\xi_3}\overline{\widehat{u}}_{\xi_4}B_{\xi_1,\xi_2,\xi_3,\xi_4}\\
&\quad\quad\quad\times(m_{N_0}(\xi_1)^2-m_{N_0}(\xi_2)^2+m_{N_0}(\xi_3)^2-m_{N_0}(\xi_4)^2).
\end{align*}

Then it suffices to establish an estimate for $Q$. 
We introduce the notation $$\varphi_{\underline{k}}(\xi):=\varphi_{k_1}(\xi_1)\varphi_{k_2}(\xi_2)\varphi_{k_3}(\xi_3)\varphi_{k_4}(\xi_4)$$ and let $\tilde{k}_1 \ge \tilde{k}_2 \ge \tilde{k}_3 \ge \tilde{k}_4$ 
denote the decreasing rearrangement of the set 
$\{k_1,k_2,k_3,k_4\} \subset \mathbb{Z}$.

Let $\epsilon' > 0$ be a small constant. 
For any $s \in [t_1,t_2]$ and $k \in \mathbb{Z}$, define
    \begin{align*}
        A_k(s):=\sum_{l\in\mathbb{Z}}2^{-\epsilon'|l-k| }\|P_lu(s)\|_{H^{N_0}}\sim \sum_{l\in\mathbb{Z}}2^{-\epsilon'|l-k|}\cdot 2^{\max\{l, 0\}N_0}\|P_lu(s)\|_{L^2}.
    \end{align*}
   By Young's inequality, for any $s\in[t_1,t_2]$,
    \begin{align*}
        \sum_{k\in\mathbb{Z}} |A_k(s)|^2\lesssim\|u(s)\|_{H^{N_0}}^2\lesssim A^2 .
    \end{align*}
In particular, for $\sigma \in (0,1)$, define
    \begin{align*}
        \tilde{A}_k(s):&=\sum_{l\in\mathbb{Z}}2^{-\epsilon'|l-k| }\|P_lP_{\geq-10}u(s)\|_{H^{N_0-\frac12+\beta}}\\
        &\sim \sum_{l\in\mathbb{Z}}2^{-\epsilon'|l-k|}\cdot 2^{\max\{l, 0\}\cdot (N_0-\frac12+\beta)}\|P_lP_{\geq-10}u(s)\|_{L^2}.
    \end{align*}
   By Young's inequality and the Gagliardo--Nirenberg inequality, for any $s\in[t_1,t_2]$ and $\beta \in (0,\tfrac14)$,
    \begin{align*}
        \sum_{k\in\mathbb{Z}} |\tilde{A}_k(s)|^2&\lesssim\|P_{\geq-10}u(s)\|_{H^{N_0-\frac12+\beta}}^2\lesssim  \| u(s)\|_{\dot{H}^{N_0-\frac14}}^2\\
        &\lesssim \| u(s)\|_{{H}^{N_0 }}^{2(1-\frac1{4N_0-2})} \| u(s)\|_{L^\infty}^{ 2\frac1{4N_0-2}}\lesssim A^{2(1-\frac1{4N_0})}\epsilon^{2\frac1{4N_0}},
    \end{align*}
where the implicit constant is independent of the size of the torus $L$.

  By the same argument as \cite[Lemma 5.2, (5.20)]{DIP25a}, 
to obtain the desired bound in Proposition~\ref{prop:highenergy}, 
it suffices to prove that \begin{align}
2^{\epsilon'(\max\{\tilde k_3,0\}-\min\{\tilde k_4,0\})}&\left|\int_{t_1}^{t_2} Q(P_{k_1}u,P_{k_2}\overline u,P_{k_3}u,P_{k_4}\overline u)\dif t\right|\notag\\
&\lesssim [A_l(t_1)\epsilon(t_1)]^2+[A_l(t_2)\epsilon(t_2)]^2+\lambda^2\int_{t_1}^{t_2}A_l^2(s)\epsilon^4(s)\dif s\notag\\
&\quad+1_{\{\sigma\in(0,1)\}}\int_{t_1}^{t_2}\tilde{A}_l(s)A_l(s)\epsilon^2(s)\dif s,\label{bd:9.2mid}
\end{align}
for any $(k_1,k_2,k_3,k_4)\in\mathbb{Z}_l:=\{(k_1,k_2,k_3,k_4)\in\mathbb{Z}^4,l=\max\{k_1,k_2,k_3,k_4\} \}$.  Here we recall that $P_k$ are the  Littlewood-Paley projections.

In what follows, we use the shorthand

\begin{align}
   \Psi(\xi)&:=  \Psi(\xi_1,\xi_2,\xi_3,\xi_4)=m_{N_0}(\xi_1)^2-m_{N_0}(\xi_2)^2+m_{N_0}(\xi_3)^2-m_{N_0}(\xi_4)^2,\notag\\
   \Omega^{(\sigma)}(\xi)&:= \Omega^{(\sigma)}(\xi_1,\xi_2,\xi_3,\xi_4)=|\xi_1|^\sigma-|\xi_2|^\sigma+|\xi_3|^\sigma-|\xi_4|^\sigma,\label{def:psiphi}
\end{align}
and let \begin{align*}
Q^{(\sigma)}(u,\overline u,u,\overline u):=\frac{1}{L^{3}}\sum_{\xi_1-\xi_2+\xi_3-\xi_4=0}\widehat{u}_{\xi_1}\overline{\widehat{u}}_{\xi_2}\widehat{u}_{\xi_3}\overline{\widehat{u}}_{\xi_4}B_{\xi_1,\xi_2,\xi_3,\xi_4}\frac{\Psi(\xi)}{\Omega^{(\sigma)}(\xi)}.
\end{align*}
By writing 
\begin{align*}
    Q^{(\sigma)}&(P_{k_1}u,P_{k_2}\overline u,P_{k_3}u,P_{k_4}\overline u)(t_2)- Q^{(\sigma)}(P_{k_1}u,P_{k_2}\overline u,P_{k_3}u,P_{k_4}\overline u)(t_1)\\
    &=\int_{t_1}^{t_2}\partial_tQ^{(\sigma)}(P_{k_1}u,P_{k_2}\overline u,P_{k_3}u,P_{k_4}\overline u)(s)\dif s,
\end{align*}
together with the fact that $u$  is the
solution to 
\eqref{MMT}, we obtain 
\begin{align}
& \left| \int_{t_1}^{t_2}Q (P_{k_1}u,P_{k_2}\overline u,P_{k_3}u,P_{k_4}\overline u)(s)\dif s\right|\notag\\
  &\leq | Q^{(\sigma)}(P_{k_1}u,P_{k_2}\overline u,P_{k_3}u,P_{k_4}\overline u)(t_2)- Q^{(\sigma)}(P_{k_1}u,P_{k_2}\overline u,P_{k_3}u,P_{k_4}\overline u)(t_1)|\notag\\
 &+  \bigg{|}\int_{t_1}^{t_2}Q^{(\sigma)}  (P_{k_1}\mathscr{B}(u,u,u),P_{k_2}\overline u,P_{k_3}u,P_{k_4}\overline u)(s)-Q^{(\sigma)}  (P_{k_1}u,P_{k_2}\mathscr{B}(\overline u,\overline u,\overline u),P_{k_3}u,P_{k_4}\overline u)(s)\notag\\
 &\quad +Q^{(\sigma)}  (P_{k_1} u,P_{k_2}\overline u,P_{k_3}\mathscr{B}(u,u,u),P_{k_4}\overline u)(s)-Q^{(\sigma)}  (P_{k_1} u ,P_{k_2}\overline u,P_{k_3}u,P_{k_4}\mathscr{B}(\overline u,\overline u,\overline u))(s)\dif s\bigg{|}.\label{bd:exp:Q}
\end{align}

We note that similar expressions appear in \cite[Section 5]{DIP25a} 
and \cite[Section 3.1]{MMT97}. 
In what follows, we estimate each term in \eqref{bd:exp:Q}, 
distinguishing between the cases $\sigma \in (0,1)$ and $\sigma \in (1,2]$. 


\subsubsection{Case 1: $\sigma=2$.}

Although this case is closely related to the range $\sigma\in(1,2)$
considered below, we present a separate and more detailed argument for
clarity. Under the constraint
\[
    \xi_1-\xi_2+\xi_3-\xi_4=0,
\]
we have
\begin{align*}
    \Omega^{(2)}(\xi)
    =2(\xi_3-\xi_4)(\xi_4-\xi_1).
\end{align*}

We first consider the boundary terms in \eqref{bd:exp:Q}, which are of
the form
\begin{align*}
&Q^{(2)}
  (P_{k_1}u,P_{k_2}\overline u,P_{k_3}u,P_{k_4}\overline u)
\\
&\qquad
=\frac{1}{L^{3}}
\sum_{\xi_1-\xi_2+\xi_3-\xi_4=0}
\widehat{u}_{\xi_1}
\overline{\widehat{u}}_{\xi_2}
\widehat{u}_{\xi_3}
\overline{\widehat{u}}_{\xi_4}
B_{\xi_1,\xi_2,\xi_3,\xi_4}
\frac{\Psi(\xi)}{\Omega^{(2)}(\xi)}
\varphi_{\underline k}(\xi).
\end{align*}
The quotient $\Psi/\Omega^{(2)}$ is understood through the smooth
extension constructed below.

Using the constraint
$\xi_1-\xi_2+\xi_3-\xi_4=0$ and applying the mean value theorem twice,
we obtain, for $N_0\geq10$,
\begin{align}
\Psi(\xi)
&=(\xi_1-\xi_2)
  \int_0^1
  (m_{N_0}^2)'
  \bigl(\theta\xi_1+(1-\theta)\xi_2\bigr)
  \dif\theta 
-(\xi_4-\xi_3)
  \int_0^1
  (m_{N_0}^2)'
  \bigl(\theta\xi_4+(1-\theta)\xi_3\bigr)
  \dif\theta
\notag\\
&=(\xi_1-\xi_2)
  \int_0^1
  \Big[
      (m_{N_0}^2)'
      \bigl(\xi_2+\theta(\xi_1-\xi_2)\bigr) 
      -(m_{N_0}^2)'
      \bigl(\xi_3+\theta(\xi_1-\xi_2)\bigr)
  \Big]
  \dif\theta
\notag\\
&=(\xi_1-\xi_2)(\xi_2-\xi_3)
  \int_0^1\int_0^1
  (m_{N_0}^2)''
  \bigl(
      \eta\xi_2+(1-\eta)\xi_3
      +\theta(\xi_1-\xi_2)
  \bigr)
  \dif\eta\dif\theta
\notag\\
&=(\xi_1-\xi_2)(\xi_2-\xi_3)
  \int_0^1\int_0^1
  (m_{N_0}^2)''
  \Big(
      (1-\theta)\eta\xi_2
      +(1-\theta)(1-\eta)\xi_3
\notag\\
&\hspace{5.5cm}
      +\eta\theta\xi_1
      +\theta(1-\eta)\xi_4
  \Big)
  \dif\eta\dif\theta.
\label{bd:phi=}
\end{align}
 Hence, 
\begin{align*}
\frac{\Psi(\xi)}{\Omega^{(2)}(\xi)}
&=
\frac12
\int_0^1\int_0^1
(m_{N_0}^2)''
\Big(
    (1-\theta)\eta\xi_2
    +(1-\theta)(1-\eta)\xi_3
\\
&\hspace{4.6cm}
    +\eta\theta\xi_1
    +\theta(1-\eta)\xi_4
\Big)
\dif\eta\dif\theta.
\end{align*}
The right-hand side provides a smooth extension of
$\Psi/\Omega^{(2)}$ across the resonant set.

Together with the bound
\[
    |(m_{N_0}^2)''(\xi)|
    \lesssim |\xi|^{2N_0-2},
\]
this gives 
\begin{align}
\left|
    \varphi_{\underline{k}}(\xi)
    \frac{\Psi(\xi)}{\Omega^{(2)}(\xi)}
\right|
\lesssim
2^{(2N_0-2)\widetilde{k}_1}.
\notag
\end{align}
It follows that
\begin{align}
&\left|
B_{\xi_1,\xi_2,\xi_3,\xi_4}
\frac{\Psi(\xi)}{\Omega^{(2)}(\xi)}
\varphi_{\underline k}(\xi)
\right| 
\lesssim
2^{(2N_0-2)\widetilde{k}_1}
2^{\beta(
    \widetilde{k}_1+\widetilde{k}_2
    +\widetilde{k}_3+\widetilde{k}_4)}
\lesssim
2^{(2N_0-2+2\beta)\widetilde{k}_1}
2^{\beta(\widetilde{k}_3+\widetilde{k}_4)}.
\notag
\end{align}

We next verify the corresponding derivative bounds. Since
$(m_{N_0}^2)''$ has polynomial growth, for every derivative order
needed below,
\[
    \big|
        \xi^\alpha
        \partial_\xi^\alpha
        (m_{N_0}^2)''
    \big|
    \lesssim
    |\xi|^{2N_0-2}.
\]
Denoting the affine expression in \eqref{bd:phi=} by
$X_{\eta,\theta}(\xi)$, the chain rule gives
\begin{align*}
\left|
\xi_1\partial_{\xi_1}
\left[
    \frac{\Psi}{\Omega^{(2)}}
\right]
\right|
&\lesssim
\int_0^1\int_0^1
|\eta\theta\xi_1|\,
\left|
    (m_{N_0}^2)'''
    \bigl(X_{\eta,\theta}(\xi)\bigr)
\right|
\dif\eta\dif\theta
\\
&\lesssim
\int_0^1\int_0^1
|\eta\theta\xi_1|\,
|X_{\eta,\theta}(\xi)|^{2N_0-3}
\dif\eta\dif\theta \lesssim
2^{(2N_0-2)\widetilde{k}_1}.
\end{align*}
By repeated application of the chain rule, the same argument yields,
for every relevant multi-index $\alpha$,
\begin{align}
\left|
\xi^\alpha\partial_{\xi}^\alpha
\left[
    \frac{\Psi}{\Omega^{(2)}}
\right]
\right|
\lesssim
2^{(2N_0-2)\widetilde{k}_1}.
\label{bd:example}
\end{align}

Each localized factor
$|\xi_j|^\beta\varphi_{k_j}(\xi_j)$ satisfies the corresponding
symbol estimates. Hence, for every relevant multi-index $\alpha$,
\begin{align*}
&\left|
\xi^\alpha\partial_{\xi}^\alpha
\left(
B_{\xi_1,\xi_2,\xi_3,\xi_4}
\frac{\Psi}{\Omega^{(2)}}
\varphi_{\underline k}(\xi)
\right)
\right| 
\lesssim
2^{(2N_0-2+2\beta)\widetilde{k}_1}
2^{\beta(\widetilde{k}_3+\widetilde{k}_4)}.
\end{align*}
By Lemma~\ref{lem:sinfty}, it follows that
\begin{align}
\left\|
B_{\xi_1,\xi_2,\xi_3,\xi_4}
\frac{\Psi}{\Omega^{(2)}}
\varphi_{\underline k}(\xi)
\right\|_{\widetilde S^\infty}
\lesssim
2^{(2N_0-2+2\beta)\widetilde{k}_1}
2^{\beta(\widetilde{k}_3+\widetilde{k}_4)}.
\label{bd:bpsi/phi2}
\end{align}

By Lemma~\ref{lem:holder}, for $\epsilon'>0$ sufficiently small,
we obtain
\begin{align*}
&2^{\epsilon'
   (\max\{\widetilde k_3,0\}
    -\min\{\widetilde k_4,0\})}
\left|
Q^{(2)}
(P_{k_1}u,P_{k_2}\overline u,
 P_{k_3}u,P_{k_4}\overline u)(s)
\right|
\\
&\quad
\lesssim
A_{\widetilde k_1}(s)
A_{\widetilde k_2}(s)
\epsilon^2(s)
2^{\beta\widetilde{k}_3
   -(2-\epsilon')\max\{\widetilde{k}_3,0\}} 
2^{\beta\widetilde{k}_4
   -2\max\{\widetilde{k}_4,0\}
   -\epsilon'\min\{\widetilde k_4,0\}}.
\end{align*}
Here we used the fact that the frequency constraint implies $ 2^{\widetilde k_1}
    \sim
    2^{\widetilde k_2}.$ 
Consequently, with
$l=\widetilde k_1=\max\{k_1,k_2,k_3,k_4\}$, the slow variation of the
frequency envelope gives
\[
    A_{\widetilde k_2}(s)\lesssim A_l(s).
\]
Choosing $0<\epsilon'<\beta$ sufficiently small, all the remaining
dyadic factors are bounded. Evaluating the above estimate at
$s=t_1$ and $s=t_2$, we obtain the first two terms on the right-hand
side of \eqref{bd:9.2mid}. 

Next, we consider the space-time integrals in \eqref{bd:exp:Q}.  
By symmetry, it suffices to focus on the first term:
 \begin{align*}
     \int_{t_1}^{t_2}Q^{(2)}  (P_{k_1}\mathscr{B}(u,u,u),P_{k_2}\overline u,P_{k_3}u,P_{k_4}\overline u)\dif s.
 \end{align*}
 We denote $\mathscr{B}(u):=\mathscr{B}(u,u,u)$, and establish the following estimates:
 \begin{align}
    2^{(N_0-2\beta)\max\{k_1,0\}} \|P_{k_1}  \mathscr{B}(u)\|_{L^2}\lesssim   \lambda^2 A_{k_1}\epsilon^2,\ \
   \|P_{k_1}  \mathscr{B}(u)\|_{L^\infty}\lesssim \lambda^22^{\beta k_1}  \epsilon^3,\label{bd:pkbu}
 \end{align}
 where we suppress the time dependence and write $\epsilon(t)=\epsilon$.
 
We first consider $k \geq 0$ and $p \in \{2, \infty\}$. Then we have
 \begin{align}
     \|P_k(uv)\|_{L^p}&\lesssim \|P_k(P'_kuP_{\leq k-4}v)\|_{L^p}+ \|P_k(P_{\leq k-4}uP'_kv)\|_{L^p}+\sum_{a,b\geq k-4,|a-b|\leq 8}\|P_k(P_auP_bv)\|_{L^p}\notag\\
     &\lesssim \sum_{a\geq k-4}\|P_au\|_{L^p}\|v\|_{\widetilde{W}^0}+\sum_{a\geq k-4}\|P_av\|_{L^p}\|u\|_{\widetilde{W}^0},\label{bd:pkuv}
 \end{align}
 where we place the low frequency terms  in $L^\infty$ and place the high frequency terms in $L^p$. We recall the definition of $\widetilde{W}^{0}$ in Section \ref{sec:pre}.
In particular, by \eqref{bd:thm9.1} we have 
 \begin{align*}
      \|P_k(|\nabla|^\beta u\cdot |\nabla|^\beta \overline u)\|_{L^p}\lesssim \sum_{a\geq k-4}\|P_a|\nabla|^\beta u\|_{L^p}\epsilon.
 \end{align*}
 
 This implies that
 \begin{align*}
    \|\ ||\nabla|^\beta u|^2\|_{\widetilde{W}^1}&\lesssim \epsilon^2+\sum_{k\geq0}2^k\sum_{a\geq k-4}\|P_a|\nabla|^\beta u\|_{L^\infty}\epsilon\\
    &\lesssim \epsilon^2+\sum_{a\geq-4}\sum_{0\leq k\leq a+4} 2^k \|P_a|\nabla|^\beta u\|_{L^\infty}\epsilon\lesssim \epsilon^2.
 \end{align*}

Then, returning to \eqref{bd:pkuv}, we obtain
\begin{align}
     \|P_k(||\nabla|^\beta u|^2|\nabla|^\beta u)\|_{L^p}
     &\lesssim \sum_{a\geq k-4}\|P_a||\nabla|^\beta u|^2\|_{L^p}\epsilon+\sum_{a\geq k-4}\|P_a|\nabla|^\beta u\|_{L^p}\epsilon^2 \notag\\
     &\lesssim \sum_{a\geq k-4}\sum_{b\geq a-4}\|P_b|\nabla|^\beta u\|_{L^p}\epsilon^2.\notag
 \end{align}
 
This implies that
\begin{align*}
    2^{(N_0-\beta)k}\|P_{k}(||\nabla|^\beta u|^2|\nabla|^\beta u)\|_{L^2}&\lesssim 2^{(N_0-\beta)k}\sum_{a\geq k-4}\sum_{b\geq a-4}2^{-b(N_0-\beta)}\|P_b  u\|_{H^{N_0}}\epsilon^2\\
    &\lesssim \sum_{b\geq k-8}\sum_{b+4\geq a\geq k-4}2^{(N_0-\beta)(k-b)}\|P_b  u\|_{H^{N_0}}\epsilon^2\\
    &\lesssim \sum_{b\geq k-8}2^{(N_0-\beta-\epsilon'')(k-b)}\|P_b  u\|_{H^{N_0}}\epsilon^2\lesssim A_k\epsilon^2,\\
     \|P_{k}(||\nabla|^\beta u|^2|\nabla|^\beta u)\|_{L^\infty}&\lesssim \sum_{a\geq k-4}\sum_{b\geq a-4}2^{b\beta}\|P_b  u\|_{L^\infty}\epsilon^2\\
    &\lesssim \sum_{b\geq k-8}(b+8-k)2^{b\beta}\|P_b  u\|_{L^\infty}\epsilon^2\lesssim \epsilon^3.
\end{align*}

By \eqref{bd:thm9.1} and the definition of $\mathscr{B}(u)$, we have for $k_1\geq0$
\begin{align*}
2^{(N_0-2\beta)k_1} \|P_{k_1}  \mathscr{B}(u)\|_{L^2}
&\lesssim 
\lambda^2 2^{(N_0-\beta)k_1}
\|P_{k_1}(\big||\nabla|^\beta u\big|^2\,|\nabla|^\beta u)\|_{L^2}\lesssim 
\lambda^2 A_{k_1}\epsilon^2,
\end{align*}
and
  \begin{align*}
   \|P_{k_1}  \mathscr{B}(u)\|_{L^\infty}\lesssim \lambda^22^{\beta k_1}  \epsilon^3.
 \end{align*}
This proves \eqref{bd:pkbu} for $k_1\geq 0$.

For $k_1<0$, we argue similarly. For $p\in\{2,\infty\}$, we have
\begin{align*}
  \|  P_{\leq 0}(uv)\|_{L^p}\lesssim \sum_{a,b\geq0,|a-b|\leq 8}\|P_0^*(P_a^*uP_b^*v)\|_{L^p}\lesssim \sum_{a \geq0}\|P_a^*u\|_{L^p}\|v\|_{\widetilde{W}^0}+\sum_{a \geq0}\|P_a^*v\|_{L^p}\|u\|_{\widetilde{W}^0},
\end{align*}
where we define $P_0^*=P_{\leq 0},P_a^*=P_a$ for $a\geq1$. Since no frequency growth occurs in the low-frequency regime, 
the same argument as above applies and yields \eqref{bd:pkbu} for $k_1<0$.



 Then, with \eqref{bd:pkbu} in hand, for $\epsilon'>0$ sufficiently small, it follows from Lemma~\ref{lem:holder} and \eqref{bd:bpsi/phi2} that
 \begin{align*}
    & 2^{\epsilon'(\max\{\tilde k_3,0\}-\min\{\tilde k_4,0\})}|  Q^{(2)}  (P_{k_1}\mathscr{B}(u),P_{k_2}\overline u,P_{k_3}u,P_{k_4}\overline u)(s)|\\
     &\quad\lesssim \lambda^2 A_{\tilde k_1}(s)A_{\tilde k_2}(s)\epsilon^4(s)2^{\beta \tilde{k}_3-2\max\{\tilde{k}_3,0\}}\cdot 2^{\beta \tilde{k}_4-2\max\{\tilde{k}_4,0\}-\epsilon' \min\{\tilde k_4,0\}}\cdot   2^{(4\beta-2)\max\{\tilde{k}_1,0\}}.
 \end{align*}
    By using the fact that $\beta\in(0,\frac12)$, we obtain the third term on the right side of \eqref{bd:9.2mid}.




 \subsubsection{Case 2: $\sigma\in(1,2)$.}
 In this case, the estimate is similar to Case 1. It suffices to establish an analogue of \eqref{bd:bpsi/phi2} 
for the regime $\sigma \in (1,2)$. 
However, the analysis is slightly more involved and requires a case-by-case discussion.

We start from the resonance condition
$\xi_1 - \xi_2 + \xi_3 - \xi_4 = 0.$
For convenience, we perform the change of variables 
$\xi_2 \mapsto -\xi_2$ and $\xi_4 \mapsto -\xi_4$. 
Under this transformation, the expressions in \eqref{def:psiphi} remain unchanged, 
and the resonance condition becomes
\begin{align*}
\xi_1 + \xi_2 + \xi_3 + \xi_4 = 0.
\end{align*}
We note that if $\xi_i = 0$ for some $i \in \{1,2,3,4\}$, 
then $Q^{(\sigma)}$ vanishes due to the factor 
$B_{\xi_1,\xi_2,\xi_3,\xi_4} = |\xi_1|^\beta|\xi_2|^\beta|\xi_3|^\beta|\xi_4|^\beta$. 
Hence, in the sequel we assume $\xi_i \neq 0$ for all $i$.

Let $\iota_i = \mathrm{sgn}(\xi_i)$ and denote by 
$\iota = (\iota_1,\iota_2,\iota_3,\iota_4)$ 
the associated sign pattern. 
It is clear that there is no solution when 
$\iota = (+,+,+,+)$ or $\iota = (-,-,-,-)$, 
since in that case the identity 
$\xi_1+\xi_2+\xi_3+\xi_4=0$ cannot hold.

Taking symmetry into account, it remains to consider the following three cases:

\textbf{Case 2.1: $\sigma\in(1,2)$ and
$\iota=(-,-,+,+)$.}
We replace $\xi_1$ and $\xi_2$ by $-\xi_1$ and $-\xi_2$,
respectively, so that all frequencies are nonnegative.
Under this convention, the resonance condition becomes
\[
    \xi_1+\xi_2=\xi_3+\xi_4.
\]
By simultaneously exchanging $(\xi_1,\xi_4)$ and
$(\xi_3,\xi_2)$ if necessary, we may assume that $ \xi_1-\xi_4=\xi_3-\xi_2>0.$
Then
\begin{align}
|\Omega^{(\sigma)}|
&=\xi_1^\sigma-\xi_4^\sigma
  +\xi_3^\sigma-\xi_2^\sigma
\notag\\
&=(\xi_1-\xi_4)
  \int_0^1
  \sigma\bigl(\theta\xi_1+(1-\theta)\xi_4\bigr)^{\sigma-1}
  \dif\theta 
 +(\xi_3-\xi_2)
  \int_0^1
  \sigma\bigl(\theta\xi_3+(1-\theta)\xi_2\bigr)^{\sigma-1}
  \dif\theta
\notag\\
&=(\xi_3-\xi_2)
\left[
  \int_0^1
  \sigma\bigl(\theta\xi_1+(1-\theta)\xi_4\bigr)^{\sigma-1}
  \dif\theta 
 +\int_0^1
  \sigma\bigl(\theta\xi_3+(1-\theta)\xi_2\bigr)^{\sigma-1}
  \dif\theta
\right].
\label{bd:omega2.2}
\end{align}

Since all frequencies are nonnegative and $\sigma-1>0$, we have
\begin{align*}
&\int_0^1
\sigma\bigl(\theta\xi_1+(1-\theta)\xi_4\bigr)^{\sigma-1}
\dif\theta
\geq
\int_0^1\sigma(\theta\xi_1)^{\sigma-1}\dif\theta
=\xi_1^{\sigma-1},
\\
&\int_0^1
\sigma\bigl(\theta\xi_1+(1-\theta)\xi_4\bigr)^{\sigma-1}
\dif\theta
\geq
\int_0^1\sigma\bigl((1-\theta)\xi_4\bigr)^{\sigma-1}
\dif\theta
=\xi_4^{\sigma-1}.
\end{align*}
Consequently,
\begin{align*}
&\int_0^1
\sigma\bigl(\theta\xi_1+(1-\theta)\xi_4\bigr)^{\sigma-1}
\dif\theta
\geq
\frac12\bigl(
\xi_1^{\sigma-1}+\xi_4^{\sigma-1}
\bigr).
\end{align*}
Applying the same argument to the second integral in
\eqref{bd:omega2.2}, we obtain
\begin{align}
|\Omega^{(\sigma)}|
&\gtrsim_\sigma
(\xi_1^{\sigma-1}
 +\xi_2^{\sigma-1}
 +\xi_3^{\sigma-1}
 +\xi_4^{\sigma-1})
(\xi_3-\xi_2)
\notag\\
&\gtrsim_\sigma
(\xi_1+\xi_2+\xi_3+\xi_4)^{\sigma-1}
(\xi_3-\xi_2).
\label{eq:omega-lower-case21}
\end{align}
In the last inequality, we used $0<\sigma-1<1$.

After the above changes of variables, the representation
\eqref{bd:phi=} gives
\begin{align}
\left|
\varphi_{\underline{k}}(\xi)\Psi(\xi)
\right|
\lesssim
\varphi_{\underline{k}}(\xi)
|(\xi_1+\xi_2)(\xi_2-\xi_3)|
2^{(2N_0-2)\tilde{k}_1}.\notag
\end{align}  Thus,
\begin{align}
\left|
\varphi_{\underline{k}}(\xi)
\frac{\Psi(\xi)}{\Omega^{(\sigma)}(\xi)}
\right|
&\lesssim
\varphi_{\underline{k}}(\xi)
\frac{
(\xi_1+\xi_2)
2^{(2N_0-2)\tilde{k}_1}
}{
(\xi_1+\xi_2+\xi_3+\xi_4)^{\sigma-1}
} \lesssim
2^{(2N_0-\sigma)\tilde{k}_1}.\notag
\end{align}
Here we used the fact that, on the support of
$\varphi_{\underline{k}}$, $ \xi_1+\xi_2+\xi_3+\xi_4
    \sim 2^{\tilde{k}_1}.$ 

We next establish the corresponding derivative bounds. Define
\begin{align*}
G(\xi)
:=
\frac{\Omega^{(\sigma)}(\xi)}{\xi_3-\xi_2}.
\end{align*}
By \eqref{bd:omega2.2},
\begin{align*}
G(\xi)
&=
\int_0^1
\sigma\bigl(\theta\xi_1+(1-\theta)\xi_4\bigr)^{\sigma-1}
\dif\theta +
\int_0^1
\sigma\bigl(\theta\xi_3+(1-\theta)\xi_2\bigr)^{\sigma-1}
\dif\theta.
\end{align*} We know that
\begin{align}
\left|
\xi_1\partial_{\xi_1}
\bigl[
(\theta\xi_1+(1-\theta)\xi_4)^{\sigma-1}
\bigr]
\right|
&\lesssim
\theta\xi_1
(\theta\xi_1+(1-\theta)\xi_4)^{\sigma-2}
\notag\\
&\lesssim
(\theta\xi_1+(1-\theta)\xi_4)^{\sigma-1}.\notag
\end{align}
The same estimate holds for derivatives with respect to the other
frequency variables. Repeated application of the chain rule therefore
gives, for every relevant multi-index $\alpha$,
\begin{align}
\left|
\xi^\alpha\partial_\xi^\alpha G(\xi)
\right|
\lesssim_\alpha G(\xi).
\label{eq:G-symbol-case21}
\end{align}

 Then, we introduce the Fa\`a di Bruno's formula:
\begin{align*}
    \frac{d^n}{dx^n} f(g(x)) = \sum C_{m_1,m_2,...,m_n,n} f^{(m_1 + \cdots + m_n)}(g(x)) \cdot \prod_{j=1}^{n} \bigl(g^{(j)}(x)\bigr)^{m_j},
\end{align*}
where the sum is over all $n$-tuples of nonnegative integers 
$(m_{1},\ldots ,m_{n})$ satisfying the constraint
$1\cdot m_{1}+2\cdot m_{2}+3\cdot m_{3}+\cdots +n\cdot m_{n}=n.$ 
Using \eqref{eq:G-symbol-case21}, we conclude that
\begin{align}
\left|
\xi^\alpha\partial_\xi^\alpha
\left[
G(\xi)^{-1}
\right]
\right|
\lesssim_\alpha
G(\xi)^{-1}.
\label{eq:G-inverse-symbol-case21}
\end{align}

On the other hand, using \eqref{bd:phi=} after the above changes of
variables and applying the chain rule as in \eqref{bd:example}, we
obtain
\begin{align}
\left|
\xi^\alpha\partial_\xi^\alpha
\left[
\varphi_{\underline{k}}(\xi)
\frac{\Psi(\xi)}{\xi_3-\xi_2}
\right]
\right|
\lesssim
2^{(2N_0-1)\tilde{k}_1}.
\label{eq:psi-divided-symbol-case21}
\end{align}
Indeed, after cancellation of $\xi_3-\xi_2$, the remaining expression
contains one factor of size at most $2^{\tilde{k}_1}$ and one factor
of $(m_{N_0}^2)''$, whose size is bounded by
$2^{(2N_0-2)\tilde{k}_1}$.

Combining \eqref{eq:omega-lower-case21},
\eqref{eq:G-inverse-symbol-case21}, and
\eqref{eq:psi-divided-symbol-case21}, and also using the standard
derivative bounds for the Littlewood--Paley cutoffs, we conclude that
\begin{align}
\left|
\xi^\alpha\partial_\xi^\alpha
\left[
\varphi_{\underline{k}}(\xi)
\frac{\Psi(\xi)}{\Omega^{(\sigma)}(\xi)}
\right]
\right|
&=
\left|
\xi^\alpha\partial_\xi^\alpha
\left[
\varphi_{\underline{k}}(\xi)
\frac{\Psi(\xi)}{\xi_3-\xi_2}
\left(
\frac{\Omega^{(\sigma)}(\xi)}
     {\xi_3-\xi_2}
\right)^{-1}
\right]
\right| \lesssim
2^{(2N_0-\sigma)\tilde{k}_1}.\notag
\end{align}


\textbf{Case 2.2: $\sigma\in(1,2)$ and
$\iota=(-,+,+,+)$.}
We replace $\xi_1$ by $-\xi_1$, so that all frequencies are
nonnegative. Under this convention, the resonance condition becomes
\[
    \xi_1=\xi_2+\xi_3+\xi_4.
\]
We use the elementary inequality
\[
    a^\sigma-b^\sigma
    \geq a^{\sigma-1}(a-b),
    \qquad a>b>0.
\] 

By symmetry between $\xi_2$ and $\xi_4$, we may assume that
$\xi_2\geq\xi_4$. We then have
\begin{align*}
\Omega^{(\sigma)}
&=(\xi_2+\xi_3+\xi_4)^\sigma
  -\xi_2^\sigma+\xi_3^\sigma-\xi_4^\sigma
\\
&\geq
(\xi_2+\xi_3+\xi_4)^{\sigma-1}
(\xi_3+\xi_4)-\xi_4^\sigma
\\
&=
\Big(
(\xi_2+\xi_3+\xi_4)^{\sigma-1}
-\xi_4^{\sigma-1}
\Big)\xi_4
+
(\xi_2+\xi_3+\xi_4)^{\sigma-1}\xi_3.
\end{align*}
Since $\xi_2\geq\xi_4$, we have  
$
     \xi_2+\xi_3+\xi_4\geq2\xi_4,
$
and hence
$
    \xi_4^{\sigma-1}
    \leq
    2^{1-\sigma}\xi_1^{\sigma-1}.
$
Therefore, setting
$
    \epsilon_\sigma:=1-2^{1-\sigma}>0,
$
we obtain
\begin{align}
\Omega^{(\sigma)}
&\geq
\epsilon_\sigma
(\xi_2+\xi_3+\xi_4)^{\sigma-1}\xi_4
+
(\xi_2+\xi_3+\xi_4)^{\sigma-1}\xi_3 \geq
\epsilon_\sigma
\xi_1^{\sigma-1}(\xi_3+\xi_4).
\label{eq:omega-lower-case22}
\end{align}

After the above changes of variables, the representation
\eqref{bd:phi=} yields
\begin{align}
\left|
\varphi_{\underline k}(\xi)\Psi(\xi)
\right|
\lesssim
(\xi_2+\xi_3)(\xi_3+\xi_4)
2^{(2N_0-2)\tilde k_1}.\notag
\end{align}  Moreover,
\begin{align}
\left|
\varphi_{\underline{k}}(\xi)
\frac{\Psi(\xi)}{\Omega^{(\sigma)}(\xi)}
\right|
&\lesssim
\frac{
(\xi_2+\xi_3)(\xi_3+\xi_4)
2^{(2N_0-2)\tilde{k}_1}
}{
\xi_1^{\sigma-1}(\xi_3+\xi_4)
}  \lesssim
2^{(2N_0-\sigma)\tilde{k}_1}.\notag
\end{align} 

We next establish the corresponding derivative bounds. We first claim
that, for $i=2,3,4$,
\begin{align}
    \left|
    \xi_i\partial_{\xi_i}\Omega^{(\sigma)}
    \right|
    \lesssim
    \Omega^{(\sigma)}.\notag
\end{align}
Since the roles of $\xi_2$ and $\xi_4$ are symmetric, it suffices to
consider $i=2,3$. We compute
\begin{align*}
\frac1\sigma
\xi_2\partial_{\xi_2}\Omega^{(\sigma)}
&=
\xi_2(\xi_2+\xi_3+\xi_4)^{\sigma-1}
-\xi_2^\sigma
\\
&\leq
(\xi_2+\xi_3+\xi_4)^\sigma
-\xi_2^\sigma+\xi_3^\sigma-\xi_4^\sigma,
\\
\frac1\sigma
\xi_3\partial_{\xi_3}\Omega^{(\sigma)}
&=
\xi_3(\xi_2+\xi_3+\xi_4)^{\sigma-1}
+\xi_3^\sigma
\\
&\leq
(\xi_2+\xi_3+\xi_4)^\sigma
-\xi_2^\sigma+\xi_3^\sigma-\xi_4^\sigma.
\end{align*}
For the second inequality, we used
\begin{align*}
    \xi_2^\sigma+\xi_4^\sigma
    \leq
    (\xi_2+\xi_4)
    (\xi_2+\xi_3+\xi_4)^{\sigma-1}.
\end{align*} 
The estimate for $i=4$ follows by symmetry.

For higher-order derivatives, we use the mean-value representation
\begin{align}
\frac1\sigma
\partial_{\xi_2}\Omega^{(\sigma)}
&=
(\xi_2+\xi_3+\xi_4)^{\sigma-1}
-\xi_2^{\sigma-1}
\notag\\
&=
(\xi_3+\xi_4)
\int_0^1
(\sigma-1)
\bigl(
\xi_2+\theta(\xi_3+\xi_4)
\bigr)^{\sigma-2}
\,\dif\theta.\notag
\end{align}
Since each component in the expression are of power law,  by the application of the chain rule gives, for every relevant
multi-index $\alpha$,
\begin{align}
\left|
\xi_2\xi^\alpha
\partial_\xi^\alpha
\left[
\partial_{\xi_2}\Omega^{(\sigma)}
\right]
\right|
\lesssim_\alpha
\left|
\xi_2\partial_{\xi_2}\Omega^{(\sigma)}
\right|
\lesssim
\Omega^{(\sigma)}.
\label{eq:higher-derivative-xi2-case22}
\end{align}
The same conclusion holds for $\xi_4$ by symmetry.

For the $\xi_3$ derivative, we use
\[
    \frac1\sigma
    \partial_{\xi_3}\Omega^{(\sigma)}
    =
    (\xi_2+\xi_3+\xi_4)^{\sigma-1}
    +\xi_3^{\sigma-1}.
\]
Both terms satisfy the standard homogeneous symbol bounds. Thus,
for every relevant multi-index $\alpha$,
\begin{align}
\left|
\xi_3\xi^\alpha
\partial_\xi^\alpha
\left[
\partial_{\xi_3}\Omega^{(\sigma)}
\right]
\right|
\lesssim_\alpha
\left|
\xi_3\partial_{\xi_3}\Omega^{(\sigma)}
\right|
\lesssim
\Omega^{(\sigma)}.
\label{eq:higher-derivative-xi3-case22}
\end{align}
It follows from
\eqref{eq:higher-derivative-xi2-case22}--%
\eqref{eq:higher-derivative-xi3-case22} that
\begin{align*}
\left|
\xi^\alpha\partial_\xi^\alpha
\Omega^{(\sigma)}
\right|
\lesssim_\alpha
\Omega^{(\sigma)}
\end{align*}
for every nonzero relevant multi-index $\alpha$. 

On the other hand, the representation \eqref{bd:phi=} and the chain
rule imply
\begin{align}
\left|
\xi^\alpha\partial_\xi^\alpha
\left[
\varphi_{\underline k}(\xi)\Psi(\xi)
\right]
\right|
\lesssim_\alpha
(\xi_2+\xi_3)(\xi_3+\xi_4)
2^{(2N_0-2)\tilde k_1}.
\label{eq:psi-derivative-case22}
\end{align} 

Combining \eqref{eq:omega-lower-case22}  and
\eqref{eq:psi-derivative-case22}, and applying Leibniz's rule, we
conclude that
\begin{align}
\left|
\xi^\alpha\partial_\xi^\alpha
\left[
\varphi_{\underline{k}}(\xi)
\frac{\Psi(\xi)}{\Omega^{(\sigma)}(\xi)}
\right]
\right|
&\lesssim
\frac{
(\xi_2+\xi_3)(\xi_3+\xi_4)
2^{(2N_0-2)\tilde k_1}
}{
\xi_1^{\sigma-1}(\xi_3+\xi_4)
} \lesssim
2^{(2N_0-\sigma)\tilde k_1}.\notag
\end{align}
  
\textbf{Case 2.3: $\sigma\in(1,2)$ and
$\iota=(-,+,-,+)$.}
We replace $\xi_1$ and $\xi_3$ by $-\xi_1$ and $-\xi_3$,
respectively, so that all frequencies are nonnegative.
Under this convention, the resonance condition becomes
\[
    \xi_1+\xi_3=\xi_2+\xi_4.
\]
Using the mean value theorem twice, we write
\begin{align}
\Omega^{(\sigma)}
&=\xi_1^\sigma-\xi_2^\sigma+\xi_3^\sigma-\xi_4^\sigma
\notag\\
&=(\xi_1-\xi_2)
  \int_0^1
  \sigma\bigl(\xi_2+\theta(\xi_1-\xi_2)\bigr)^{\sigma-1}
  \dif\theta 
-(\xi_4-\xi_3)
  \int_0^1
  \sigma\bigl(\xi_3+\theta(\xi_4-\xi_3)\bigr)^{\sigma-1}
  \dif\theta
\notag\\
&=(\xi_1-\xi_2)\sigma
  \int_0^1
  \Big[
      \bigl(\xi_2+\theta(\xi_1-\xi_2)\bigr)^{\sigma-1}
 -
      \bigl(\xi_3+\theta(\xi_1-\xi_2)\bigr)^{\sigma-1}
  \Big]
  \dif\theta
\notag\\
&=(\xi_1-\xi_2)(\xi_2-\xi_3)\sigma(\sigma-1)
\notag\\
&\quad\times
\int_0^1\int_0^1
\Big(
    (1-\theta)\eta\xi_2
    +(1-\theta)(1-\eta)\xi_3
    +\eta\theta\xi_1
    +\theta(1-\eta)\xi_4
\Big)^{\sigma-2}
\dif\eta\dif\theta.
\label{bd:omega2.3}
\end{align}
Here we used $ \xi_1-\xi_2=\xi_4-\xi_3.$ 
For brevity, set
\begin{align*}
X_{\eta,\theta}(\xi)
:=
(1-\theta)\eta\xi_2
+(1-\theta)(1-\eta)\xi_3
+\eta\theta\xi_1
+\theta(1-\eta)\xi_4.
\end{align*}
The coefficients in this expression are nonnegative and sum to one.
Hence, 
\[
    X_{\eta,\theta}(\xi)^{\sigma-2}
    \gtrsim
    2^{(\sigma-2)\widetilde k_1}.
\]
Therefore,
\begin{align}
\left|\Omega^{(\sigma)}(\xi)\right|
\gtrsim_\sigma
\left|
(\xi_1-\xi_2)(\xi_2-\xi_3)
\right|
2^{(\sigma-2)\widetilde k_1}.
\label{eq:omega-lower-case23}
\end{align}

By \eqref{bd:phi=}, $\Psi$ contains the same two factors
$(\xi_1-\xi_2)(\xi_2-\xi_3)$. Consequently,  
\begin{align}
\left|
\varphi_{\underline{k}}(\xi)
\frac{\Psi(\xi)}{\Omega^{(\sigma)}(\xi)}
\right|
&\lesssim
\frac{
\left|
(\xi_1-\xi_2)(\xi_2-\xi_3)
\right|
2^{(2N_0-2)\widetilde{k}_1}
}{
\left|
(\xi_1-\xi_2)(\xi_2-\xi_3)
\right|
2^{(\sigma-2)\widetilde{k}_1}
} \lesssim
2^{(2N_0-\sigma)\widetilde{k}_1}.\notag
\end{align}

We next establish the corresponding derivative estimates. Define
\begin{align*}
G_\sigma(\xi)
:=
\frac{\Omega^{(\sigma)}(\xi)}
     {(\xi_1-\xi_2)(\xi_2-\xi_3)}.
\end{align*}
By \eqref{bd:omega2.3},
\begin{align*}
G_\sigma(\xi)
=
\sigma(\sigma-1)
\int_0^1\int_0^1
X_{\eta,\theta}(\xi)^{\sigma-2}
\dif\eta\dif\theta>0.
\end{align*}
For each $i\in\{1,2,3,4\}$, we have
\[
    |\xi_i\partial_{\xi_i}X_{\eta,\theta}(\xi)|
    \leq X_{\eta,\theta}(\xi).
\] 
Repeated application of the chain rule therefore yields, for every
relevant multi-index $\alpha$,
\begin{align}
\left|
\xi^\alpha\partial_\xi^\alpha
G_\sigma(\xi)
\right|
\lesssim_\alpha
G_\sigma(\xi).\notag
\end{align}
Since $G_\sigma>0$, the multivariable Fa\`a di Bruno formula gives
\begin{align}
\left|
\xi^\alpha\partial_\xi^\alpha
\left[
G_\sigma(\xi)^{-1}
\right]
\right|
\lesssim_\alpha
G_\sigma(\xi)^{-1}.
\label{eq:G-inverse-case23}
\end{align}
Furthermore, \eqref{eq:omega-lower-case23} implies
\begin{align}
G_\sigma(\xi)^{-1}
\lesssim
2^{(2-\sigma)\widetilde k_1}.
\label{eq:G-inverse-size-case23}
\end{align}

On the other hand, using \eqref{bd:phi=} and the chain rule as in
\eqref{bd:example}, we obtain the absolute dyadic bound
\begin{align}
\left|
\xi^\alpha\partial_\xi^\alpha
\left[
\varphi_{\underline k}(\xi)
\frac{\Psi(\xi)}
     {(\xi_1-\xi_2)(\xi_2-\xi_3)}
\right]
\right|
\lesssim_\alpha
2^{(2N_0-2)\widetilde k_1}.
\label{eq:psi-divided-symbol-case23}
\end{align}
Notice that we use an absolute dyadic bound here, rather than an
estimate relative to the size of $\Psi$, since the latter may vanish.

Combining \eqref{eq:G-inverse-case23},
\eqref{eq:G-inverse-size-case23}, and
\eqref{eq:psi-divided-symbol-case23}, and applying Leibniz's rule, we
obtain
\begin{align}
\left|
\xi^\alpha\partial_\xi^\alpha
\left[
\varphi_{\underline k}(\xi)
\frac{\Psi(\xi)}{\Omega^{(\sigma)}(\xi)}
\right]
\right|
&=
\left|
\xi^\alpha\partial_\xi^\alpha
\left[
\varphi_{\underline k}(\xi)
\frac{\Psi(\xi)}
     {(\xi_1-\xi_2)(\xi_2-\xi_3)}
G_\sigma(\xi)^{-1}
\right]
\right| \lesssim
2^{(2N_0-\sigma)\widetilde k_1}.
\label{bd:case2.3}
\end{align}

Combining the estimates obtained in Cases 2.1--2.3 and applying
Lemma~\ref{lem:sinfty}, we conclude that
\begin{align}
\left\|
B_{\xi_1,\xi_2,\xi_3,\xi_4}
\frac{\Psi(\xi)}{\Omega^{(\sigma)}(\xi)}
\varphi_{\underline k}(\xi)
\right\|_{\widetilde S^\infty}
&\lesssim
2^{(2N_0-\sigma)\widetilde{k}_1}
2^{\beta(
\widetilde{k}_1+\widetilde{k}_2
+\widetilde{k}_3+\widetilde{k}_4)}
\notag\\
&\lesssim
2^{(2N_0-\sigma+2\beta)\widetilde{k}_1}
2^{\beta(\widetilde{k}_3+\widetilde{k}_4)}.
\notag
\end{align}

We are now in a position to estimate the terms in
\eqref{bd:exp:Q}. For the boundary terms, Lemma~\ref{lem:holder}
gives
\begin{align*}
&2^{\epsilon'
(\max\{\widetilde k_3,0\}
-\min\{\widetilde k_4,0\})}
\left|
Q^{(\sigma)}
(P_{k_1}u,P_{k_2}\overline u,
 P_{k_3}u,P_{k_4}\overline u)(s)
\right|
\\
&\quad\lesssim
A_{\widetilde k_1}(s)
A_{\widetilde k_2}(s)
\epsilon^2(s)
2^{\beta\widetilde{k}_3
-(2-\epsilon')\max\{\widetilde{k}_3,0\}}
\\
&\qquad\qquad\times
2^{\beta\widetilde{k}_4
-2\max\{\widetilde{k}_4,0\}
-\epsilon'\min\{\widetilde k_4,0\}}
2^{(4\beta-\sigma)\max\{\widetilde{k}_1,0\}}.
\end{align*}
Here the factor $  2^{(4\beta-\sigma)\max\{\widetilde k_1,0\}}$ 
is a harmless weakening of the sharper boundary estimate, chosen so
that the boundary and space-time contributions have the same dyadic
form.
  

For the space-time contributions, using \eqref{bd:pkbu},
Lemma~\ref{lem:holder}, and the above multiplier bound, we similarly
obtain
\begin{align*}
&2^{\epsilon'
(\max\{\widetilde k_3,0\}
-\min\{\widetilde k_4,0\})}
\left|
Q^{(\sigma)}
(P_{k_1}\mathscr B(u),P_{k_2}\overline u,
 P_{k_3}u,P_{k_4}\overline u)(s)
\right|
\\
&\quad\lesssim
\lambda^2
A_{\widetilde k_1}(s)
A_{\widetilde k_2}(s)
\epsilon^4(s)
2^{\beta\widetilde{k}_3
-(2-\epsilon')\max\{\widetilde{k}_3,0\}}
\\
&\qquad\qquad\times
2^{\beta\widetilde{k}_4
-2\max\{\widetilde{k}_4,0\}
-\epsilon'\min\{\widetilde k_4,0\}}
2^{(4\beta-\sigma)\max\{\widetilde{k}_1,0\}}.
\end{align*} 

Since $4\beta<\sigma$ and $0<\epsilon'<\beta$, the dyadic factors are
bounded by $1$. Therefore,  integration over $[t_1,t_2]$ yields $\lambda^2
    \int_{t_1}^{t_2}
    A_l(s)^2\epsilon^4(s)\dif s,$ 
which is the third term on the right-hand side of
\eqref{bd:9.2mid}. 

\subsubsection{Case 3: $\sigma\in(0,1)$.}

In this regime, the analysis becomes more delicate due to the presence
of nontrivial resonant interactions. As in Case 2, we first replace
$\xi_2$ and $\xi_4$ by $-\xi_2$ and $-\xi_4$, respectively. The
momentum constraint then takes the form
\[
    \xi_1+\xi_2+\xi_3+\xi_4=0.
\]

We decompose
\begin{align*}
&Q(P_{k_1}u,P_{k_2}\overline u,
   P_{k_3}u,P_{k_4}\overline u)
\\
&\qquad
=
Q_1(P_{k_1}u,P_{k_2}\overline u,
    P_{k_3}u,P_{k_4}\overline u)
+
Q_2(P_{k_1}u,P_{k_2}\overline u,
    P_{k_3}u,P_{k_4}\overline u),
\end{align*}
where
\begin{align*}
&Q_1(P_{k_1}u,P_{k_2}\overline u,
     P_{k_3}u,P_{k_4}\overline u)
\\
&\qquad
:=
\frac{1}{L^3}
\sum_{\mathcal Q_1}
\widehat u_{\xi_1}
\overline{\widehat u}_{-\xi_2}
\widehat u_{\xi_3}
\overline{\widehat u}_{-\xi_4}
B_{\xi_1,\xi_2,\xi_3,\xi_4}
\Psi(\xi_1^*,\xi_2,\xi_3^*,\xi_4)
\varphi_{\underline k}(\xi),
\\
&Q_2:=Q-Q_1.
\end{align*}
Here $\sum_{\mathcal Q_1}$ denotes summation over
$\xi\in(L^{-1}\mathbb Z)^4$ satisfying
\[
    \xi_1+\xi_2+\xi_3+\xi_4=0
\]
and whose sign configuration belongs to
\begin{align*}
\{
&(-,+,+,+),(+,-,+,+),(+,+,-,+),(+,+,+,-),
\\
&(+,-,-,-),(-,+,-,-),(-,-,+,-),(-,-,-,+)
\}.
\end{align*}

We next define the modified frequencies. Consider, for example, the
configuration $\iota=(-,+,+,+).$ 
We set $ \xi_1^*
    :=
    -(\xi_2+\xi_3^*+\xi_4),$ 
where $\xi_3^*$ is the unique positive solution of
\begin{align}
    (x+\xi_2+\xi_4)^\sigma
    +x^\sigma
    -\xi_2^\sigma
    -\xi_4^\sigma
    =0.
\label{eq:resonant-root}
\end{align} 
Since $0<\sigma<1$, the strict concavity of $x\mapsto x^\sigma$
implies  \eqref{eq:resonant-root} admits a unique solution
satisfying
\[
    0<\xi_3^*<\min\{\xi_2,\xi_4\}.
\] 
  

We aim to prove that
\begin{align}
&2^{\epsilon'
(\max\{\tilde k_3,0\}-\min\{\tilde k_4,0\})}
\left|
\int_{t_1}^{t_2}
Q_1(P_{k_1}u,P_{k_2}\overline u,
    P_{k_3}u,P_{k_4}\overline u)(s)
\,\dif s
\right|
\notag\\
&\qquad
\lesssim
\int_{t_1}^{t_2}
\widetilde A_l(s)A_l(s)\epsilon^2(s)
\,\dif s,
\label{bd:q1}
\end{align}
and
\begin{align}
&2^{\epsilon'
(\max\{\tilde k_3,0\}-\min\{\tilde k_4,0\})}
\left|
\int_{t_1}^{t_2}
Q_2(P_{k_1}u,P_{k_2}\overline u,
    P_{k_3}u,P_{k_4}\overline u)(s)
\,\dif s
\right|
\notag\\
&\qquad
\lesssim
[A_l(t_1)\epsilon(t_1)]^2
+
[A_l(t_2)\epsilon(t_2)]^2
+
\lambda^2
\int_{t_1}^{t_2}
A_l^2(s)\epsilon^4(s)
\,\dif s.
\label{bd:q2}
\end{align}
 

The estimate for $Q_1$ is treated in Case 3.0 below. 
For $Q_2$, the argument is similar to Case 2, and it suffices to establish a symbol bound analogous to \eqref{bd:bpsi/phi2}.
For convenience, we rewrite $\xi_2,\xi_4$ as$-\xi_2,-\xi_4$, so that
\[
\xi_1+\xi_2+\xi_3+\xi_4=0.
\]
Taking symmetry into account, we consider the sign configurations in Cases 3.1–3.3.

\textbf{Case 3.0: $\sigma\in(0,1)$ and the proof of
\eqref{bd:q1}.}
By symmetry, it suffices to consider the sign configuration
$\iota=(-,+,+,+)$. Using \eqref{bd:phi=} after the corresponding
changes of signs, we obtain
\begin{align*}
&\left|
\Psi(\xi_1^*,\xi_2,\xi_3^*,\xi_4)
\varphi_{\underline{k}}(\xi)
\right|
\lesssim
(\xi_2+\xi_3^*)
(\xi_3^*+\xi_4)
(\xi_2+\xi_3^*+\xi_4)^{2N_0-2}
\varphi_{\underline{k}}(\xi)
\\
&\qquad
\lesssim
\xi_2\xi_4
(\xi_2+\xi_4)^{2N_0-2}
\varphi_{\underline{k}}(\xi) 
\lesssim
2^{\widetilde{k}_1+\widetilde{k}_3
  +(2N_0-2)\widetilde{k}_1}.
\end{align*}
Here we used $ 0<\xi_3^*<\min\{\xi_2,\xi_4\}.$ 
Moreover,  
\[
    \max\{\xi_2,\xi_4\}\lesssim 2^{\widetilde k_1},
    \qquad
    \min\{\xi_2,\xi_4\}\lesssim 2^{\widetilde k_3}.
\]
It follows that
\begin{align*}
&\left|
B_{\xi_1,\xi_2,\xi_3,\xi_4}
\Psi(\xi_1^*,\xi_2,\xi_3^*,\xi_4)
\varphi_{\underline{k}}(\xi)
\right| 
\lesssim
2^{(2N_0-1+2\beta)\widetilde{k}_1}
2^{(1+\beta)\widetilde{k}_3}
2^{\beta\widetilde{k}_4}.
\end{align*}

We next establish the corresponding derivative estimates. Regard $ 
    \xi_3^*=\xi_3^*(\xi_2,\xi_4)$
as a function of $(\xi_2,\xi_4)$, and write
$
    S:=\xi_2+\xi_3^*(\xi_2,\xi_4)+\xi_4.
$
By definition,
\[
    S^\sigma+(\xi_3^*)^\sigma
    -\xi_2^\sigma-\xi_4^\sigma=0.
\]
Implicit differentiation with respect to $\xi_2$ gives
\begin{align}
\partial_{\xi_2}\xi_3^*
&=
\frac{
\xi_2^{\sigma-1}
-
S^{\sigma-1}
}{
S^{\sigma-1}
+
(\xi_3^*)^{\sigma-1}
} =
\frac{
(1-\sigma)(\xi_3^*+\xi_4)
\displaystyle\int_0^1
\bigl(
\xi_2+\theta(\xi_3^*+\xi_4)
\bigr)^{\sigma-2}
\,\dif\theta
}{
S^{\sigma-1}
+
(\xi_3^*)^{\sigma-1}
}.
\label{eq:derivative-xi3star}
\end{align} 
We claim that
\begin{align}
|\xi_4\partial_{\xi_4}\xi_3^*|+|\xi_2\partial_{\xi_2}\xi_3^*|\lesssim \xi_3^*  . \notag
\end{align} We verify the $\xi_2\partial_{\xi_2}$ case. It suffices to show
\begin{align*}
    \xi_2^{\sigma}-  \xi_2(\xi_2+\xi_3^*+\xi_4)^{\sigma-1}\leq  \xi_3^*(\xi_2+\xi_3^*+\xi_4)^{\sigma-1}+(\xi_3^*)^{\sigma}.
\end{align*}
By the definition of $\xi_3^*$, it suffices to check
\begin{align*}
    (\xi_2+\xi_3^*+\xi_4)^\sigma-\xi_4^\sigma-  \xi_2(\xi_2+\xi_3^*+\xi_4)^{\sigma-1}\leq  \xi_3^*(\xi_2+\xi_3^*+\xi_4)^{\sigma-1}.
\end{align*}
Since $\sigma\in(0,1)$,    we use  $(\xi_2+\xi_3^*+\xi_4)^\sigma=(\xi_2+\xi_3^*+\xi_4)\cdot (\xi_2+\xi_3^*+\xi_4)^{\sigma-1}$ to conclude the desired bound. The $\xi_4$ case follows similarly. 

For higher order derivatives, $\partial^\alpha \xi_3^*$ can be written as a ratio of expressions with homogeneous  power law structure.
More generally, for every multi-index
$\alpha=(\alpha_2,\alpha_4)$ needed below, repeated implicit
differentiation yields
\begin{align}
\left|
\xi_2^{\alpha_2}
\xi_4^{\alpha_4}
\partial_{\xi_2}^{\alpha_2}
\partial_{\xi_4}^{\alpha_4}
\xi_3^*
\right|
\lesssim_\alpha
\xi_3^*.
\label{bd:part2xi3}
\end{align}  


We now estimate derivatives of the modified symbol. Set
\[
\Psi^*(\xi_2,\xi_4)
:=
\Psi\bigl(
-\xi_2-\xi_3^*(\xi_2,\xi_4)-\xi_4,\,
\xi_2,\,
\xi_3^*(\xi_2,\xi_4),\,
\xi_4
\bigr).
\]
By the chain rule,
\begin{align*}
\xi_2\partial_{\xi_2}\Psi^*
&=
\xi_2
\left[
-(1+\partial_{\xi_2}\xi_3^*)\partial_1
+\partial_2
+(\partial_{\xi_2}\xi_3^*)\partial_3
\right]
\Psi
\bigl(
\xi_1^*,\xi_2,\xi_3^*,\xi_4
\bigr).
\end{align*}
Using   the factorized
representation of $\Psi$ above, and the symbol estimates for
$(m_{N_0}^2)''$, we obtain
\begin{align*}
\left|
\xi_2\partial_{\xi_2}\Psi^*
\right|
&\lesssim
2^{\widetilde k_1}
2^{\widetilde k_3
  +(2N_0-2)\widetilde k_1}
\lesssim
2^{\widetilde k_3
  +(2N_0-1)\widetilde k_1}.
\end{align*}
Higher-order derivatives are controlled in the same way by repeated
application of the chain rule and \eqref{bd:part2xi3}. Consequently,
for every relevant multi-index $\alpha$,
\begin{align}
\left|
\xi^\alpha\partial_\xi^\alpha
\left[
B_{\xi_1,\xi_2,\xi_3,\xi_4}
\Psi(\xi_1^*,\xi_2,\xi_3^*,\xi_4)
\varphi_{\underline{k}}(\xi)
\right]
\right|
\lesssim_\alpha
2^{(2N_0-1+2\beta)\widetilde{k}_1}
2^{(1+\beta)\widetilde{k}_3}
2^{\beta\widetilde{k}_4}.
\label{eq:q1-symbol-derivatives}
\end{align}
Here derivatives with respect to $\xi_1$ and $\xi_3$ only act on the
remaining localized factors, while the modified symbol itself depends
on $\xi_2$ and $\xi_4$.

Lemma~\ref{lem:sinfty} now gives
\begin{align}
\left\|
B_{\xi_1,\xi_2,\xi_3,\xi_4}
\Psi(\xi_1^*,\xi_2,\xi_3^*,\xi_4)
\varphi_{\underline{k}}(\xi)
\right\|_{\widetilde S^\infty}
\lesssim
2^{(2N_0-1+2\beta)\widetilde{k}_1}
2^{(1+\beta)\widetilde{k}_3}
2^{\beta\widetilde{k}_4}.
\label{eq:q1-sinfty}
\end{align}

We next return to
\[
    Q_1(P_{k_1}u,P_{k_2}\overline u,
        P_{k_3}u,P_{k_4}\overline u).
\]
We may insert an additional cutoff on the largest-frequency factor:
\begin{align*}
&Q_1(P_{k_1}u,P_{k_2}\overline u,
     P_{k_3}u,P_{k_4}\overline u)
\\
&\quad
=
\frac{1}{L^3}
\sum_{-\xi_1=\xi_2+\xi_3+\xi_4}
\varphi_{\geq-10}(\xi_1)
\widehat u_{\xi_1}
\overline{\widehat u}_{-\xi_2}
\widehat u_{\xi_3}
\overline{\widehat u}_{-\xi_4}
B_{\xi_1,\xi_2,\xi_3,\xi_4} 
\Psi(\xi_1^*,\xi_2,\xi_3^*,\xi_4)
\varphi_{\underline{k}}(\xi).
\end{align*}
Indeed, on the support of $1-\varphi_{\geq-10}(\xi_1)$, all the
actual frequencies satisfy
\[
    |\xi_j|\leq|\xi_1|,
\]
while
\[
    0<\xi_3^*<\min\{\xi_2,\xi_4\},
    \qquad
    |\xi_1^*|
    =
    \xi_2+\xi_3^*+\xi_4
    \lesssim|\xi_1|.
\]
Thus all four arguments of $\Psi$ lie in the region where
$m_{N_0}$ vanishes, and hence
\[
    \Psi(\xi_1^*,\xi_2,\xi_3^*,\xi_4)=0.
\]
 Moreover, under the constraint $-\xi_1 = \xi_2 + \xi_3 + \xi_4$, we have 
\[
2^{k_1}\sim 2^{\tilde{k}_1}.
\]
Applying Lemma \ref{lem:holder}, we obtain
\begin{align*}
 |Q_1(P_{k_1}u,P_{k_2}\overline u,P_{k_3}u,P_{k_4}\overline u)(s)|\lesssim \tilde{ A}_{\tilde k_1}(s)A_{\tilde k_2}(s)\epsilon^2(s)2^{(1+\beta) \tilde{k}_3-2\max\{\tilde{k}_3,0\}}\cdot 2^{\beta \tilde{k}_4-2\max\{\tilde{k}_4,0\}}.
\end{align*}
 This completes the proof of \eqref{bd:q1}.

\textbf{Case 3.1: $\sigma\in(0,1)$ and
$\iota=(-,+,+,+)$.}
We replace $\xi_1$ by $-\xi_1$, so that all frequencies are
nonnegative. Under this convention, the resonance condition becomes
\[
    \xi_1=\xi_2+\xi_3+\xi_4.
\]

To estimate $Q_2$, we apply a normal-form transformation analogous to
\eqref{bd:exp:Q}. We are led to study the multiplier
\begin{align}
\varphi_{\underline k}(\xi)
\frac{\Psi(\xi)-\Psi(\xi^*)}{\Omega^{(\sigma)}(\xi)}
&=
\varphi_{\underline{k}}(\xi)
\frac{\Psi(\xi)-\Psi(\xi^*)}
     {\Omega^{(\sigma)}(\xi)-\Omega^{(\sigma)}(\xi^*)},
\label{eq:case31-difference-quotient}
\end{align}
where $ \xi^*=(\xi_1^*,\xi_2,\xi_3^*,\xi_4)$ 
and, by construction,
\[
    \Omega^{(\sigma)}(\xi^*)=0.
\]

Fix $(\xi_2,\xi_4)$ and define
\begin{align*}
\widetilde\Psi(x)
&:=
\Psi\bigl(
    -\xi_2-x-\xi_4,\,
    \xi_2,\,
    x,\,
    \xi_4
\bigr),
\\
\widetilde\Omega(x)
&:=
(\xi_2+x+\xi_4)^\sigma
-\xi_2^\sigma+x^\sigma-\xi_4^\sigma.
\end{align*}
 
By the mean value theorem,
\begin{align*}
\left|
\widetilde\Psi(\xi_3)
-\widetilde\Psi(\xi_3^*)
\right|
&=
|\xi_3-\xi_3^*|
\left|
\int_0^1
\widetilde\Psi'
\bigl(
    \theta\xi_3+(1-\theta)\xi_3^*
\bigr)
\,\dif\theta
\right|
\\
&\lesssim
|\xi_3-\xi_3^*|
\bigl(
\max\{\xi_3,\xi_3^*,\xi_2,\xi_4\}
\bigr)^{2N_0-1}
\\
&\lesssim
|\xi_3-\xi_3^*|
2^{(2N_0-1)\widetilde{k}_1}.
\end{align*}
Similarly, 
\begin{align*}
\left|
\widetilde\Omega(\xi_3)
-\widetilde\Omega(\xi_3^*)
\right|
&\gtrsim_\sigma
|\xi_3-\xi_3^*|
\bigl(
\max\{\xi_3,\xi_3^*,\xi_2,\xi_4\}
\bigr)^{\sigma-1}
\\
&\gtrsim_\sigma
|\xi_3-\xi_3^*|
2^{(\sigma-1)\widetilde{k}_1}.
\end{align*}
Here we used
$
    0<\xi_3^*<\min\{\xi_2,\xi_4\}$
and
$
    \max\{\xi_3,\xi_2,\xi_4\}
    \lesssim  \xi_1
    \sim 2^{\widetilde k_1}.
$

Consequently, the quotient in
\eqref{eq:case31-difference-quotient} admits a continuous extension
across $\xi_3=\xi_3^*$ and satisfies
\begin{align}
\left|
\varphi_{\underline{k}}(\xi)
\frac{\Psi(\xi)-\Psi(\xi^*)}
     {\Omega^{(\sigma)}(\xi)}
\right|
\lesssim
2^{(2N_0-\sigma)\widetilde{k}_1}.
\label{bd:case3.1}
\end{align}

We next establish the corresponding derivative bounds. Set
\begin{align*}
z_\theta
:=
\theta\xi_3+(1-\theta)\xi_3^*,
\end{align*}
and write
\begin{align*}
D(\xi)
&:=
\frac{
\widetilde\Omega(\xi_3)
-\widetilde\Omega(\xi_3^*)
}{
\xi_3-\xi_3^*
}=
\sigma
\int_0^1
\left[
    (\xi_2+z_\theta+\xi_4)^{\sigma-1}
    +z_\theta^{\sigma-1}
\right]
\,\dif\theta.
\end{align*}
 

By \eqref{bd:part2xi3}, the dependence of $\xi_3^*$ on
$(\xi_2,\xi_4)$ satisfies the required symbol-type estimates. Hence,
for every relevant multi-index $\alpha$,
\begin{align}
\left|
\xi^\alpha\partial_\xi^\alpha
D(\xi)
\right|
\lesssim_\alpha
D(\xi).
\label{eq:D-symbol-case31}
\end{align}
For example, the integrands that arise after differentiation are
controlled by
\begin{align*}
&\left|
\xi^\alpha\partial_\xi^\alpha
\left[
(\xi_2+z_\theta+\xi_4)^{\sigma-1}
+z_\theta^{\sigma-1}
\right]
\right| 
\lesssim_\alpha
(\xi_2+z_\theta+\xi_4)^{\sigma-1}
+z_\theta^{\sigma-1}.
\end{align*}
The multivariable Fa\`a di Bruno formula therefore gives
\begin{align}
\left|
\xi^\alpha\partial_\xi^\alpha
\left[D(\xi)^{-1}\right]
\right|
\lesssim_\alpha
D(\xi)^{-1}.
\label{eq:D-inverse-case31}
\end{align}

For the numerator, define
\begin{align*}
N(\xi)
&:=
\frac{
\widetilde\Psi(\xi_3)
-\widetilde\Psi(\xi_3^*)
}{
\xi_3-\xi_3^*
} =
\int_0^1
\widetilde\Psi'(z_\theta)
\,\dif\theta.
\end{align*}
Using the symbol estimates for $m_{N_0}^2$, the chain rule, and
\eqref{bd:part2xi3}, we obtain
\begin{align}
\left|
\xi^\alpha\partial_\xi^\alpha
\left[
\varphi_{\underline k}(\xi)N(\xi)
\right]
\right|
\lesssim_\alpha
2^{(2N_0-1)\widetilde k_1}.
\label{eq:N-symbol-case31}
\end{align}
Combining  
\eqref{eq:D-inverse-case31}, and
\eqref{eq:N-symbol-case31}, we conclude that
\begin{align}
\left|
\xi^\alpha\partial_\xi^\alpha
\left[
\varphi_{\underline{k}}(\xi)
\frac{\Psi(\xi)-\Psi(\xi^*)}
     {\Omega^{(\sigma)}(\xi)}
\right]
\right|
\lesssim_\alpha
2^{(2N_0-\sigma)\widetilde{k}_1}.
\label{eq:case31-all-derivatives}
\end{align}
Thus, the bound \eqref{bd:case3.1} remains valid under all derivatives
required in Lemma~\ref{lem:sinfty}.

\medskip

\textbf{Case 3.2: $\sigma\in(0,1)$ and
$\iota=(-,-,+,+)$.}
We replace $\xi_1$ and $\xi_2$ by $-\xi_1$ and $-\xi_2$,
respectively, so that all frequencies are nonnegative. Under this
convention, the resonance condition becomes
\[
    \xi_1+\xi_2=\xi_3+\xi_4.
\]
By symmetry, we may assume that $\xi_1-\xi_4=\xi_3-\xi_2>0.$ 
By the mean value theorem,
\begin{align}
|\Omega^{(\sigma)}|
&=
\xi_1^\sigma-\xi_4^\sigma
+\xi_3^\sigma-\xi_2^\sigma
\notag\\
&=
(\xi_1-\xi_4)
\left[
\int_0^1
\sigma
\bigl(
    \theta\xi_1+(1-\theta)\xi_4
\bigr)^{\sigma-1}
\,\dif\theta
\right.
\notag\\
&\hspace{3.6cm}\left.
+
\int_0^1
\sigma
\bigl(
    \theta\xi_3+(1-\theta)\xi_2
\bigr)^{\sigma-1}
\,\dif\theta
\right].
\label{eq:omega-case32-mvt}
\end{align}
Since $\sigma-1<0$,
we have
\begin{align*}
\int_0^1
\sigma
\bigl(
    \theta\xi_1+(1-\theta)\xi_4
\bigr)^{\sigma-1}
\,\dif\theta
\geq
\sigma(\xi_1+\xi_4)^{\sigma-1}.
\end{align*} 
\begin{align*}
\int_0^1
\sigma
\bigl(
    \theta\xi_3+(1-\theta)\xi_2
\bigr)^{\sigma-1}
\,\dif\theta
\geq
\sigma(\xi_3+\xi_2)^{\sigma-1}.
\end{align*}
Consequently,
\begin{align}
|\Omega^{(\sigma)}|
&\gtrsim_\sigma
(\xi_1-\xi_4)\left[
    (\xi_1+\xi_4)^{\sigma-1}
    +(\xi_3+\xi_2)^{\sigma-1}
\right]
\notag\\
&\gtrsim_\sigma
(\xi_1-\xi_4)
(\xi_1+\xi_2+\xi_3+\xi_4)^{\sigma-1}.\notag
\end{align}

On the other hand, \eqref{bd:phi=} gives
\begin{align}
\left|
\varphi_{\underline{k}}(\xi)\Psi(\xi)
\right|
\lesssim
|(\xi_1+\xi_2)(\xi_2-\xi_3)|
2^{(2N_0-2)\widetilde{k}_1}.
\label{eq:psi-case32}
\end{align}  It follows that 
\begin{align}
\left|
\varphi_{\underline{k}}(\xi)
\frac{\Psi(\xi)}{\Omega^{(\sigma)}(\xi)}
\right|
&\lesssim
\frac{
(\xi_1+\xi_2)
2^{(2N_0-2)\widetilde{k}_1}
}{
(\xi_1+\xi_2+\xi_3+\xi_4)^{\sigma-1}
} \lesssim
2^{(2N_0-\sigma)\widetilde{k}_1}.\notag
\end{align} 

For the derivative estimates, define
\begin{align*}
G(\xi)
:=
\frac{\Omega^{(\sigma)}(\xi)}{\xi_1-\xi_4}.
\end{align*}
By \eqref{eq:omega-case32-mvt}, $G(\xi)>0$. Moreover, since each term  in the expression is of power law, repeated
application of the chain rule gives, for every relevant multi-index
$\alpha$,
\begin{align}
\left|
\xi^\alpha\partial_\xi^\alpha G(\xi)
\right|
\lesssim_\alpha
G(\xi).\notag
\end{align} 

Since $G>0$, the multivariable Fa\`a di Bruno formula yields
\begin{align}
\left|
\xi^\alpha\partial_\xi^\alpha
\left[G(\xi)^{-1}\right]
\right|
\lesssim_\alpha
G(\xi)^{-1}.
\label{eq:G-inverse-case32}
\end{align} 

For the numerator, \eqref{eq:psi-case32} and the chain rule give the
absolute dyadic estimate
\begin{align}
\left|
\xi^\alpha\partial_\xi^\alpha
\left[
\varphi_{\underline k}(\xi)
\frac{\Psi(\xi)}{\xi_1-\xi_4}
\right]
\right|
\lesssim_\alpha
2^{(2N_0-1)\widetilde k_1}.
\label{eq:psi-divided-case32}
\end{align}
Combining \eqref{eq:G-inverse-case32}  and
\eqref{eq:psi-divided-case32}, we obtain
\begin{align}
\left|
\xi^\alpha\partial_\xi^\alpha
\left[
\varphi_{\underline{k}}(\xi)
\frac{\Psi(\xi)}{\Omega^{(\sigma)}(\xi)}
\right]
\right|
\lesssim_\alpha
2^{(2N_0-\sigma)\widetilde{k}_1}.\notag
\end{align} 

\textbf{Case 3.3: $\sigma\in(0,1)$ and
$\iota=(-,+,-,+)$.}
We replace $\xi_1$ and $\xi_3$ by $-\xi_1$ and $-\xi_3$,
respectively, so that all frequencies are nonnegative.
Under this convention, the resonance condition becomes
\[
    \xi_1-\xi_2=\xi_4-\xi_3.
\]
Using \eqref{bd:omega2.3}, we write
\begin{align}
\Omega^{(\sigma)}
&=
(\xi_1-\xi_2)(\xi_2-\xi_3)\sigma(\sigma-1)
\notag\\
&\quad\times
\int_0^1\int_0^1
\Big(
    (1-\theta)\eta\xi_2
    +(1-\theta)(1-\eta)\xi_3
    +\eta\theta\xi_1
    +\theta(1-\eta)\xi_4
\Big)^{\sigma-2}
\,\dif\eta\,\dif\theta.\notag
\end{align}
For brevity, set
\[
X_{\eta,\theta}(\xi)
:=
(1-\theta)\eta\xi_2
+(1-\theta)(1-\eta)\xi_3
+\eta\theta\xi_1
+\theta(1-\eta)\xi_4.
\]
Since the coefficients in $X_{\eta,\theta}$ are nonnegative and sum
to one, on the support of $\varphi_{\underline k}$ we have 
\[
    X_{\eta,\theta}(\xi)^{\sigma-2}
    \gtrsim
    2^{(\sigma-2)\widetilde k_1}.
\]
Consequently,
\begin{align}
\left|\Omega^{(\sigma)}(\xi)\right|
\gtrsim_\sigma
\left|
(\xi_1-\xi_2)(\xi_2-\xi_3)
\right|
2^{(\sigma-2)\widetilde{k}_1}.\notag
\end{align}

On the other hand, by \eqref{bd:phi=},
\begin{align*}
\left|
\varphi_{\underline{k}}(\xi)\Psi(\xi)
\right|
\lesssim
\left|
(\xi_1-\xi_2)(\xi_2-\xi_3)
\right|
2^{(2N_0-2)\widetilde{k}_1}.
\end{align*}  Moreover,
\begin{align}
\left|
\varphi_{\underline{k}}(\xi)
\frac{\Psi(\xi)}{\Omega^{(\sigma)}(\xi)}
\right|
&\lesssim
\frac{
\left|
(\xi_1-\xi_2)(\xi_2-\xi_3)
\right|
2^{(2N_0-2)\widetilde{k}_1}
}{
\left|
(\xi_1-\xi_2)(\xi_2-\xi_3)
\right|
2^{(\sigma-2)\widetilde{k}_1}
} \lesssim
2^{(2N_0-\sigma)\widetilde{k}_1}.\notag
\end{align}

We next establish the corresponding derivative estimates. Define
\begin{align*}
G_\sigma(\xi)
:=
-\frac{\Omega^{(\sigma)}(\xi)}
{(\xi_1-\xi_2)(\xi_2-\xi_3)} 
=
\sigma(1-\sigma)
\int_0^1\int_0^1
X_{\eta,\theta}(\xi)^{\sigma-2}
\,\dif\eta\,\dif\theta
>0.
\end{align*} Repeated application of the chain rule yields, for
every relevant multi-index $\alpha$,
\begin{align}
\left|
\xi^\alpha\partial_\xi^\alpha
G_\sigma(\xi)
\right|
\lesssim_\alpha
G_\sigma(\xi).\notag
\end{align}
Indeed, each derivative produces a factor of the form
$X_{\eta,\theta}^{\sigma-2-r}$, while the corresponding scaled
frequency factors contribute at most $X_{\eta,\theta}^r$.

Since $G_\sigma>0$, the multivariable Fa\`a di Bruno formula implies
\begin{align}
\left|
\xi^\alpha\partial_\xi^\alpha
\left[
G_\sigma(\xi)^{-1}
\right]
\right|
\lesssim_\alpha
G_\sigma(\xi)^{-1}.
\label{eq:G-inverse-case33}
\end{align} 

For the numerator, \eqref{bd:phi=} and the chain rule give the absolute
dyadic estimate
\begin{align}
\left|
\xi^\alpha\partial_\xi^\alpha
\left[
\varphi_{\underline k}(\xi)
\frac{\Psi(\xi)}
{(\xi_1-\xi_2)(\xi_2-\xi_3)}
\right]
\right|
\lesssim_\alpha
2^{(2N_0-2)\widetilde k_1}.
\label{eq:psi-divided-case33}
\end{align} 

Finally,
\[
\frac{\Psi(\xi)}{\Omega^{(\sigma)}(\xi)}
=
-
\frac{\Psi(\xi)}
     {(\xi_1-\xi_2)(\xi_2-\xi_3)}
G_\sigma(\xi)^{-1}.
\]
Combining \eqref{eq:G-inverse-case33} and
\eqref{eq:psi-divided-case33}, and applying Leibniz's rule, we obtain
\begin{align}
\left|
\xi^\alpha\partial_\xi^\alpha
\left[
\varphi_{\underline{k}}(\xi)
\frac{\Psi(\xi)}{\Omega^{(\sigma)}(\xi)}
\right]
\right|
\lesssim_\alpha
2^{(2N_0-\sigma)\widetilde{k}_1}.\notag
\end{align}
 



In summary, under all circumstances (case 3.1-3.3), by Lemma~\ref{lem:sinfty}, we obtain 

\begin{align}
  \| B_{\xi_1,\xi_2,\xi_3,\xi_4}\frac{\tilde{\Psi}(\xi)}{\Omega^{(\sigma)}(\xi)}\varphi_{\underline k}(\xi)\|_{\tilde{S}^\infty}\lesssim 2^{(2N_0-\sigma)\tilde{k}_1}2^{\beta(\tilde{k}_1+\tilde{k}_2+\tilde{k}_3+\tilde{k}_4)}\lesssim2^{(2N_0-\sigma+2\beta)\tilde{k}_1}2^{\beta(\tilde{k}_3+\tilde{k}_4)}.\label{bd:9.m3}
\end{align} 
Here $\tilde{\Psi}(\xi)$ denotes either $\Psi(\xi)$ or $\Psi(\xi)-\Psi(\xi^*)$.   To estimate \eqref{bd:q2}, by the same decomposition as in \eqref{bd:exp:Q},  it suffices to bound
\begin{align*}
 | Q_2^{(\sigma)}(P_{k_1}u,P_{k_2}\overline u,P_{k_3}u,P_{k_4}\overline u)(t_i)|
 &+  \bigg{|}\int_{t_1}^{t_2}Q_2^{(\sigma)}  (P_{k_1}\mathscr{B}(u,u,u),P_{k_2}\overline u,P_{k_3}u,P_{k_4}\overline u)(s) \bigg{|}.
\end{align*}
Under the condition $\beta<\sigma/4$, 
these terms can be estimated exactly as in Case 2, 
with the bound \eqref{bd:9.m3} replacing the corresponding symbol estimate there. 
This completes the proof.  
 

\section{Proof of the Main Result: Theorem \ref{thm1}}\label{sec:proofthm1}
This section is devoted to the proof of the main result, which consists of two parts.
The first part concerns the well-posedness of equation \eqref{MMT}.
We have already established the well-posedness of the truncated equation \eqref{TrMMT} with a high probability.
Using the acceptable  energy inequality, we extend this result to \eqref{MMT} up to time $T_0$ via a bootstrap argument.
This part is presented in Section~\ref{sec:well-posed}.
The second part concerns the derivation of the wave kinetic equation stated in \eqref{app}.
Since the derivation for the truncated equation has already been obtained, it remains to justify the convergence as $K \to \infty$.
This part is addressed in Section~\ref{sec:mainresult}

 

\subsection{The well-posedness of \eqref{MMT}}\label{sec:well-posed}
In this section, we begin the proof of the main result by establishing the well-posedness of the equation \eqref{MMT}. This will be achieved by combining the energy estimates with a bootstrap argument. First, we introduce the following hyper-contractivity  estimate, the proof of which can be found in \cite[Proposition 2.2]{DIP25}.
 
\begin{proposition}[Multilinear Gaussian estimates] \label{prop:dapiancha}
Assume that $g_k : \Omega \to \mathbb{C}, k \in \{1,...,n\}$,
 are i.i.d. normalized complex Gaussian random variables on a probability space $\Omega$ and define
 \begin{align*}
 G(\omega) :=\sum_{k_1,...,k_r\in\{1,...,n\}}
 a_{k_1,...,k_r} \prod _{i=1}^rg^{\zeta_i}_{k_i}(\omega),
 \end{align*}
 where $r \geq 1, \omega \in\Omega, \zeta_j \in \{+,-\}$, and $ a_{k_1,...,k_r}\in\mathbb{C}$. Then, for any $A \geq 1$, it holds that
 \begin{align*}
     \mathbf{P}(\omega: |G(\omega)| \geq A\|G(\omega)\|_{L^2_{\omega}})
 \leq  Ce^{-A^{\frac2r}/C},
 \end{align*}
 for some constant $C \geq 1$.
\end{proposition}

Let $\delta>0$ be a fixed small constant such that  $\delta_0 := \frac{\delta}{10\beta}<\frac1{200}$. We define
$
N_0 := 10 + \Bigl\lfloor \frac{1}{\delta_0^2} \Bigr\rfloor.$
By the decay property of Schwartz functions and the definition  \eqref{wellprepared}, we have
 
\begin{align*}
   \mathbf{E}\|u_{\mathrm{in}} \|^2_{H^{N_0}}\lesssim L^{-1}\sum_{k\in\mathbb{Z}_L} \langle k\rangle^{2N_0} n_{\mathrm{in}}(k) \lesssim 1. 
\end{align*}
Then, by Proposition~\ref{prop:dapiancha}, it follows that with probability $\geq 1-L^{-40}$,
\begin{align}\label{bd:uinhn0}
   \|u_{\mathrm{in}} \|_{H^{N_0}} \lesssim L^\theta,
\end{align}
where $\theta>0$ is another small constant to be determined later.

Moreover, the truncated solution $u_{\rm tr}$ satisfies the following $L^\infty$-bound: 
With probability $\geq 1-L^{-40}$, it holds that 
    \begin{align}
  \sup_{t\in[0,T]}  \|u_{\rm tr}(t)\|_{W^{2,0}}\lesssim L^{-1/2+\theta}.\label{bd:utrw2infty}
\end{align}
    
We recall that $ c_k(s)= e^{-2\pi i[|k|^\sigma Ts+\Gamma(k,s)]} (\widehat u_{\rm tr})_k(Ts)$ and have for $t=Ts\in[0,T]$,
\begin{align}
    u_{\rm tr}(t,x)=\frac{1}{L}\sum_{k \in \mathbb{Z}_L}  c_k(s) e^{2\pi i[|k|^\sigma Ts+\Gamma(k,s)]} e^{2\pi i kx}.\notag
\end{align}
Thus,
\begin{align}
    \langle \nabla \rangle^4u_{\rm tr}(t,x)=\frac{1}{L}\sum_{k \in \mathbb{Z}_L}   \langle k \rangle^4(c^{\leq N}_k  +\mathcal{R}^{N+1}_k)e^{2\pi i[|k|^\sigma Ts+\Gamma(k,s)]} e^{2\pi i kx}.\notag
\end{align}
By the estimates on tree couples in \eqref{bd:couple} and the remainder bound \eqref{bd:mathcalr}, we obtain
\begin{align}
    \mathbf{E}|\langle \nabla \rangle^4u_{\rm tr}(t,x)|^2\lesssim \frac{1}{L^{2}}\sum_{k \in \mathbb{Z}_L}\sum_{i,j=0}^N   \langle k \rangle^8  \mathbf{E}[c^{i}_k\overline c^{j}_k]  +L^{-800}\lesssim  L^{-1}.\notag
\end{align}
Applying Proposition~\ref{prop:dapiancha}, we deduce that with probability $\geq 1-L^{-40}$,
\begin{align*}
   \sup_t\|u_{\rm tr}(t)\|_{W^{2,0}}\lesssim L^{-1/2+\theta}
\end{align*}
for small constant $\theta>0$.  
The uniformity in space-time follows from a standard discretization argument, see e.g.\ \cite{DH21,DH23a}.


We now present the key step in establishing the well-posedness of \eqref{MMT}.

\begin{proposition}[Main Bootstrap Argument]\label{prop:bootstrap}
Let $\sigma\in(0,1)\cup (1,2]$ and $\beta\in(0,\frac{\sigma}{4})$. 
There exists a sufficiently small constant $\theta>0$ such that, with probability $\geq 1-L^{-40}$, the following holds.
 Assume that $t \in [0,T_0]$ and $u$ is a solution to \eqref{MMT} satisfying
 $t \in [0,T_0]$ and $u$ is a solution to \eqref{MMT} satisfying
\begin{align*}
   \sup_{t'\in[0,t]} \|u(t')\|_{H^{N_0}}\leq L^{4\theta},\ \ \sup_{t'\in[0,t]} \|u(t')\|_{W^{2,0}}\leq L^{-1/2+4\theta}.
\end{align*}
Then we have 
\begin{align}
   \sup_{t'\in[0,t]} \|u(t')\|_{H^{N_0}}\leq L^{2\theta},\ \ \sup_{t'\in[0,t]} \|u(t')\|_{W^{2,0}}\leq L^{-1/2+2\theta}.\label{bd:bootstrap}
\end{align}
\end{proposition}

The bootstrap argument in Proposition~\ref{prop:bootstrap} follows directly from Proposition~\ref{prop:wh5d}, together with the acceptable energy inequality stated in Theorem~\ref{thm:energy inequality}.

\begin{proposition} \label{prop:wh5d}
   Under the assumption of Proposition \ref{prop:bootstrap}, we define $w:=u-u_{\rm tr}=u-P_{\leq K} u$ and have  with probability $\geq 1-L^{-40}$,
\begin{align*}
   \sup_{t'\in[0,t]} \|w(t')\|_{W^{2,0}}\lesssim   \sup_{t'\in[0,t]} \|w(t')\|_{H^{5}}\leq L^{-100}.
\end{align*}
\end{proposition}

\begin{proof}[Proof of Proposition \ref{prop:wh5d}]

We recall that $u$ and $u_{\rm tr}=P_{\leq K} u$ solve
\begin{align*}
    i \partial_t u +2\pi |\nabla|^\sigma u&= \mathscr{B}(u,u,u),\ \ u(0)=u_{\rm in},\\
     i \partial_t u_{\rm tr} +2\pi |\nabla|^\sigma u_{\rm tr}&=P_{\leq K} \mathscr{B}(u_{\rm tr},u_{\rm tr},u_{\rm tr}),\ \ u_{\rm tr}(0)=P_{\leq K}  u_{\rm in}.
\end{align*}
Taking the difference and setting $w:=u-u_{\rm tr}$, we obtain
\begin{align*}
    i \partial_t w +2\pi |\nabla|^\sigma w=P_{\leq K}\sum_{i=1}^3 \mathscr{B}^i[u_{\rm tr}](w)+\mathscr R,\ \ w(0)=P_{> K} u_{\rm in},
\end{align*}
where   $\mathscr{B}^1[u_{\rm tr}](w):=\mathscr{B}(w,u_{\rm tr},u_{\rm tr}),\mathscr{B}^2[u_{\rm tr}](w):=\mathscr{B}(u_{\rm tr},w,u_{\rm tr}),$ $\mathscr{B}^3[u_{\rm tr}](w):=\mathscr{B}(u_{\rm tr},u_{\rm tr},w)$, and the remainder $\mathscr R$ is given by 
\begin{align*}
    \mathscr R:=P_{>K} \mathscr{B}(u,u,u)+P_{\leq K}\left[\mathscr{B}(u,u,u)-\mathscr{B}(u_{\rm tr},u_{\rm tr},u_{\rm tr})- \sum_{
    i=1}^3 \mathscr{B}^i[u_{\rm tr}](w)\right].
\end{align*}
Here  $P_{>K}=1-P_{\leq K}$.

Following the same reduction as in Section~\ref{reductions}, we introduce the analogous transformation
$$d_k(s):= e^{-2\pi i[|k|^\sigma Ts+\Gamma(k,s)]} \widehat w_k(Ts),\ \ s\in[0,1],$$
which satisfy
\begin{align}
    d_k=d_k(0)+\widetilde{\mathscr{L}}( d)_k+\tilde{\mathscr{R}}_k.\label{eq:overck}
\end{align}
Here $d_k(0):=[1-\varphi_{\leq K}(k)](\widehat{u}_{\rm in})_k$, and the linear operator $\widetilde{\mathscr{L}}$ is defined as 
\begin{align}
\widetilde{\mathscr{L}}(v)&:=\int_0^t \big(2\mathcal{W}_1(c,\overline c, v)+\mathcal{W}_1(c, \overline  v, c)+\widetilde{\mathscr{L}_1}[c](v)\big)\dif s,  \notag
\end{align} 
    where
\begin{align}
\widetilde{\mathscr{L}_1}[c]_k(v):&=2\frac{i\alpha T}{L}\sum_{k_1}\left(|c_{k_1}|^2 -\mathbf{E}[|c^{\leq N}_{k_1}|^2] \right) B_{k_1,k_1,k,k}\varphi_{\leq K}(k)v_k\notag\\
&+2\frac{i\alpha T}{L}\sum_{k_1}\left(c_{k_1}\overline {v_{k_1}}+v_{k_1}\overline {c_{k_1}} \right) B_{k,k_1,k_1,k}\varphi_{\leq K}(k)c_k.\notag
\end{align}
We recall that $$ c_k(t)= e^{-2\pi i[|k|^\sigma Tt+\Gamma(k,t)]} (\widehat u_{\rm tr})_k(Tt).$$ The term $\mathbf{E}[|c^{\leq N}_{k_1}|^2]$ arises from the renormalization $\Gamma(k,t)$ and is therefore consistent with the previous definitions. 

Finally, the remainder term is given by 
\begin{align*}
  \tilde{\mathscr{R}}_k(t):=\int_0^t e^{-2\pi i[|k|^\sigma Ts+\Gamma(k,s)]}  \widehat{\mathscr{R}}_k(Ts)\dif s.
\end{align*}


We recall the linear operator $\mathscr{L}$ defined in \eqref{def:l}.  
Then, by \eqref{eq:overck}, we obtain  
\begin{align}
 d_k=(1-\mathscr{L})   ^{-1}[d_k(0)+(\widetilde{\mathscr{L}}-\mathscr{L})_k( d)+\tilde{\mathscr{R}}_k].\label{eq:dk}
\end{align}
Then, we know that 
\begin{align*}
  \|w\|_{L_{[0,t]}^\infty H^5}^2=  \|d\|_{Z'}^2:=\sup_{s\in[0,t/T]}L^{-1}\sum_{k \in \mathbb{Z}_L}\langle k\rangle^{10}|d_k(s)|^2.
\end{align*}
In the following, we slightly abuse notation and continue to denote this norm by $Z$.

By definition, we have
\begin{align*}
    \tilde\varphi_{\leq K}(k)c_k=c_k,\ \  \tilde\varphi_{\leq K}(k)c^{\leq N}_k=c^{\leq N}_k, 
\end{align*}
which implies that 
\begin{align*}
    \mathscr{L}(v)&:=\int_0^t \big(2\mathcal{W}_1( \tilde\varphi_{\leq K}c^{\leq N},  \tilde\varphi_{\leq K}\overline{c^{\leq N}}, v)+\mathcal{W}_1( \tilde\varphi_{\leq K}c^{\leq N},\overline{ v},  \tilde\varphi_{\leq K}c^{\leq N})+\mathscr{L}_1[ \tilde\varphi_{\leq K}c^{\leq N}](v)\big)\dif s,  \\
\widetilde{\mathscr{L}}(v)&:=\int_0^t \big(2\mathcal{W}_1(\tilde\varphi_{\leq K} c,\tilde\varphi_{\leq K}\overline c, v)+\mathcal{W}_1(\tilde\varphi_{\leq K}c, \overline v, \tilde\varphi_{\leq K}c)+\widetilde{\mathscr{L}_1}[\tilde\varphi_{\leq K}c](v)\big)\dif s. 
\end{align*}

Taking the difference and exploiting the properties of $\tilde{\varphi}_{\leq K}$, a calculation analogous to \eqref{bd:est:w1} yields the following bound: with probability $\geq 1 - L^{-40}$,
\begin{align*}
   \| (\widetilde{\mathscr{L}}-\mathscr{L})(v)\|_{Z}
  & \lesssim  L^{21} \|c -c^{\leq N}\|_{Z}( \|c \|_Z+ \|c^{\leq N}\|_Z)\|v\|_{Z}\\
   &\lesssim L^{21} \|\mathcal{R}^{N+1}\|_{Z}( \|\mathcal{R}^{N+1}\|_{Z}+ \|c^{\leq N}\|_{Z})\|v\|_{Z}\leq L^{-100}\|v\|_{Z}.
\end{align*}
Here we used the decomposition $c=c^{\leq N}+\mathcal{R}^{N+1}$.  The bounds for $c^{\leq N}$ and $\mathcal{R}^{N+1}$ are given in \eqref{bd:mathcalr} and Proposition~\ref{prop:ckn}, respectively. In particular, with probability $\geq 1-L^{-40}$,
\begin{align}
    \|c^{\leq N} \|_{Z}\leq \sum_{i=1}^N\|c^{i} \|_{Z}\lesssim L.\label{bd:cz}
\end{align}

By \eqref{bd:1-l}, the operator $(1-\mathscr{L})^{-1}$ is bounded on $Z$. 
Combining this with \eqref{eq:dk}, we obtain 
\begin{align}
 \|d\|_{Z}\lesssim L^{70}[\|d(0)\|_{Z}+\|\tilde{\mathscr{R}}\|_{Z}].\label{bd:overlinec:boot}
\end{align}
We now apply a bootstrap argument again. Assume that
\begin{align}
 \sup_{t'\in[0,t]} \|w(t')\|_{H^{5}}\leq   \|d\|_{Z}\leq  L^{-100}.\label{bd:boot:wh5}
\end{align}
We aim to improve this bound to  \begin{align}
 \sup_{t'\in[0,t]} \|w(t')\|_{H^{5}}\leq   \|d\|_{Z}\leq  \frac12L^{-100}.\notag
\end{align}
Since the initial data $u_{\rm in}$ is a Schwartz function, by \eqref{def:K} and \eqref{bd:uinhn0} we have
\begin{align*}
    \|d(0)\|_{Z}\leq \|P_{>K}u_{\rm in}\|_{H^{5}}\lesssim L^{-\delta_0 (N_0-5)}\|u_{\rm in}\|_{H^{N_0}}\lesssim L^{-\frac{1}{\delta_0}} \|u_{\rm in}\|_{H^{N_0}}\leq  L^{-190}.
\end{align*}

For the remainder term, by definition it suffices to estimate the following two contributions:
\begin{align*}
    \|\tilde{\mathscr{R}}\|_Z\ &\lesssim \|P_{>K} \mathscr{B}(u,u,u)\|_{L^\infty_{[0,t]}H^{5}}\notag\\
    &\quad +\bigg\|P_{\leq K}[\mathscr{B}(u,u,u)-\mathscr{B}(u_{\rm tr},u_{\rm tr},u_{\rm tr})-\sum_{i=1}^3 \mathscr{B}^i[u_{\rm tr}](w)]\bigg\|_{L^\infty_{[0,t]}H^{5}}.
\end{align*}
For the first term above, by the bootstrap assumption in Proposition~\ref{prop:bootstrap}, we have
\begin{align*}
   \|P_{>K} \mathscr{B}(u,u,u)\|_{H^{5}}& \lesssim L^{-\delta_0(N_0-6)}\| \mathscr{B}(u,u,u)\|_{H^{N_0-1}}\\
   &\lesssim L^{-\frac1{\delta_0}}\|u\|_{H^{N_0}}^3\leq L^{-200} \|u\|_{H^{N_0}}^3\leq L^{-190},
\end{align*}
where we used that $H^{N_0}$ is an algebra for $N_0>10$, and chose $\delta>0$ sufficiently small.


For the second term, by direct expansion,
\begin{align*}
  P_{\leq K}[&\mathscr{B}(u,u,u)-\mathscr{B}(u_{\rm tr},u_{\rm tr},u_{\rm tr})-\sum_{i=1}^3 \mathscr{B}^i[u_{\rm tr}](w)]\\
  &=  P_{\leq K}[ \mathscr{B}(u_{\rm tr},w,w)+\mathscr{B}(w,u_{\rm tr},w)+\mathscr{B}(w,w,u_{\rm tr})+\mathscr{B}(w,w,w)].
\end{align*}
Using the bootstrap assumption \eqref{bd:boot:wh5} together with Sobolev embedding, we obtain
\begin{align*}
   \|P_{\leq K}  \mathscr{B}(u_{\rm tr},w,w)\|_{H^{5}}\lesssim L^{20} \|u_{\rm tr}\|_{H^{5}}\|w\|_{H^{5}}^2 
\end{align*}
and similarly for the remaining terms. Hence,
\begin{align}
   &\bigg\|P_{\leq K}[\mathscr{B}(u,u,u)-\mathscr{B}(u_{\rm tr},u_{\rm tr},u_{\rm tr})-\sum_{i=1}^3 \mathscr{B}^i[u_{\rm tr}](w)]\bigg\|_{L^\infty_{[0,t]}H^{5}}\notag\\
   &\quad \lesssim L^{20} \|u_{\rm tr}\|_{L^\infty_{[0,t]}H^{5}}\|w\|_{L^\infty_{[0,t]}H^{5}}^2+L^{20} \|w\|_{L^\infty_{[0,t]}H^{5}}^3\lesssim L^{-175},\label{bd:mid6}
\end{align}
where we used \eqref{bd:cz} to deduce
\begin{align*}
    \|u_{\rm tr}\|_{L_{[0,t]}^\infty H^{5}}\lesssim \|c\|_{Z}\lesssim L.
\end{align*}

Combining the above bounds with \eqref{bd:overlinec:boot}, we obtain
\begin{align}
\sup_{t'\in[0,t]} \|w(t')\|_{H^{5}}\leq  \|d_k\|_{Z}\leq L^{70} \cdot L^{-175}\leq \frac12L^{-100}.\notag
\end{align}
 which improves the bootstrap assumption and thus closes the argument.
  
   
   
   Finally, by the Sobolev embedding estimate (\cite[(2.27)]{DIP25a})  
\begin{align*}
\|P_kf\|_{L^\infty}\lesssim (2^{k/2}+L^{-1/2})\|P_kf\|_{L^2},
\end{align*}
we obtain the embedding $H^{5} \subset W^{2,0}$,  which completes the proof.
\end{proof}


Now, we are in position to prove the main bootstrap argument in Proposition \ref{prop:bootstrap}.
\begin{proof}[Proof of Proposition \ref{prop:bootstrap}]
 Under the assumption of  Proposition \ref{prop:wh5d}, together with  \eqref{bd:utrw2infty} we have
 \begin{align}
   \sup_{t'\in[0,t]} \|u(t')\|_{W^{2,0}}\leq  \sup_{t'\in[0,t]} \|u_{\rm tr}(t')\|_{W^{2,0}}+ \sup_{t'\in[0,t]} \|w(t')\|_{W^{2,0}}\lesssim L^{-1/2+\theta}\leq  L^{-1/2+2\theta},\notag
\end{align}
which yields the second estimate in \eqref{bd:bootstrap}. Here we take $L$ sufficiently large to absorb the implicit constant.
 

For the first estimate in \eqref{bd:bootstrap}, we use the bootstrap assumptions
\begin{align*}
   \sup_{t'\in[0,t]} \|u(t')\|_{H^{N_0}}\leq L^{4\theta},\ \ \sup_{t'\in[0,t]} \|u(t')\|_{W^{2,0}}\leq L^{-1/2+4\theta}.
\end{align*}
Applying Theorem~\ref{thm:energy inequality} with 
$A=L^{4\theta},\epsilon=L^{-1/2+4\theta}$, we note that $\lambda^2\epsilon^2=\alpha L\cdot L^{-1+8\theta}\ll 1$ for $\theta>0$ sufficiently small.

For the case $\sigma\in(1,2]$, using   \eqref{bd:uinhn0}, we obtain
\begin{align}
\sup_{t'\in[0,t]}  \|u(t')\|_{H^{N_0}}^2&\lesssim  \|u_{\rm in}\|_{H^{N_0}}^2+\lambda^4A^2T\epsilon^4\notag\\
&\lesssim L^{\theta}+\alpha^2 L^{2}\cdot L^{8\theta}\cdot TL^{-2+16\theta}\leq L^{2\theta},\notag
\end{align}
where we used the fact that $\alpha^2T\ll1$, and chose $\theta>0$ small enough such that $\alpha^2TL^{25\theta}\ll1$.

Similarly, for the case $\sigma\in(0,1)$,  we have instead \begin{align}
\sup_{t'\in[0,t]}  \|u(t')\|_{H^{N_0}}^2&\lesssim  \|u_{\rm in}\|_{H^{N_0}}^2+\lambda^4A^2T\epsilon^4+\lambda^2A^{2-\frac{1}{4N_0}}T\epsilon^{2+\frac{1}{4N_0}}\notag\\ 
&\lesssim L^{\theta}+\alpha L\cdot L^{8\theta}\cdot TL^{-1-\frac{1}{8N_0}+8\theta}\leq L^{2\theta},\notag
\end{align}
where  we used $\alpha TL^{-\theta}<1$, and chose $\theta>0$ small enough such that $20\theta<\frac{1}{8N_0}$.

\end{proof}

\subsection{The derivation of WKE}\label{sec:mainresult}
The main bootstrap argument in Proposition~\ref{prop:bootstrap} yields the existence of a smooth solution to \eqref{MMT} up to time $T_0$, satisfying both the $H^{N_0}$-bound and the $L^\infty$-type bound. We now turn to the verification of \eqref{app}. 

 From Theorem~\ref{thm2}, for the truncated solution $u_{\rm tr}$ to \eqref{TrMMT}, we have
\begin{equation}
    \mathbf E |\widehat u_{\rm tr}(t, k)|^2 =\varphi^2_{\leq K}(k)n_{\mathrm{in}}(k)+\frac{t}{T_{\mathrm{kin}}}\mathcal K_{\rm tr}(n_{\mathrm{in}})(k)+o_{l^\infty_k}\left(\frac{t}{T_{\mathrm {kin}}}\right)_{L \to \infty}.\notag
\end{equation}
where $\mathcal K_{\rm tr}$ is defined in \eqref{bd:ktr}.

%\begin{proposition}\label{prop:ukconverge}
%It holds that
 %   \begin{align*}
 %  \mathbf E |\widehat u_{\rm tr}(t, k)|^2- \mathbf E |\widehat u (t, k)|^2 =o_{l^\infty_k}\left(\frac{T}{T_{\mathrm {kin}}}\right)_{L \to \infty}.
%\end{align*}
%\end{proposition}
%\begin{proof}
 %   We recall that  $ c_k(t)= e^{-2\pi i[|k|^\sigma Tt+\Gamma(k,t)]} (\widehat u_{\rm tr})_k(Tt)$, which implies that  by Proposition \ref{prop:ckn} and Proposition \ref{prop:wh5d}.
  %  \begin{align*}
 %       |\widehat u_{\rm tr}(t,k)-\widehat u(t,k)|\lesssim |\widehat w(t,k)|\lesssim L\|w\|_{L_t^\infty H^{5}}\lesssim L^{-99},\\ |\widehat u_{\rm tr}(t,k)|\lesssim |c_k(s)|\lesssim L\|c^{\leq N}\|_{Z}+L\|\mathcal{R}^{ N+1}\|_{Z}\lesssim L^{2}.
%    \end{align*}
 %   By Cauchy's inequality, it holds that
 %\begin{align*}
 %    |   \mathbf E |\widehat u_{\rm tr}(t, k)|^2- \mathbf E |\widehat u (t,k)|^2 |\leq   (  \mathbf E |\widehat u_{\rm tr}( k)-\widehat u ( k)|^2 )^{1/2}(  \mathbf E |\widehat u_{\rm tr}( k)-\widehat u ( k)|^2 +  \mathbf E |\widehat u_{\rm tr}( k) |^2)^{1/2}\lesssim L^{-10}.
%\end{align*}

%\end{proof}
 We now show that   \begin{align*}
    \mathbf E |\widehat u (t, k)|^2-\mathbf E |\widehat u_{\rm tr}(t, k)|^2 =(1-\varphi^2_{\leq K}(k))n_{\mathrm{in}}(k)+o_{l^\infty_k}\left(\frac{t}{T_{\mathrm {kin}}}\right)_{L \to \infty}.
\end{align*}
Recall that $\widehat u=\widehat u_{\rm tr}+\widehat w$, and define
\begin{align*}
    c_k(t)= e^{-2\pi i[|k|^\sigma Tt+\Gamma(k,t)]} (\widehat u_{\rm tr})_k(Tt),\ \ d_k(t):= e^{-2\pi i[|k|^\sigma Tt+\Gamma(k,t)]} \widehat w_k(Tt).
\end{align*}
It suffices to prove that for $s\in[0,1]$,
\begin{align*}
    \mathbf E&|c_k(s)+d_k(s)|^2- \mathbf E|c_k(s)|^2-\mathbf E|c_k(0)+d_k(0)|^2+\mathbf E|c_k(0)|^2\\
    &=\mathbf E[(d_k(s)+d_k(0)+2c_k(s))(d_k(s)-d_k(0))]+2\mathbf E[(c_k(s)-c_k(0))d_k(0)]
\end{align*}
is of order $o_{l^\infty_k}\left(\frac{Ts}{T_{\mathrm {kin}}}\right)_{L \to \infty}$.  Expanding the expression, we obtain  that with probability $\geq 1-L^{-40}$,
\begin{align*}
    |d_k(s)+c_k(s)|&\lesssim L^{1/2}(\|d\|_{Z}+\|c\|_{Z})\lesssim L^2,\\
    |d_k(0)|&\lesssim|(1-\varphi_{\leq K}(k))n_{\mathrm{in}}(k)|\lesssim  L^{-100}\sup_k\left||k|^{100/\delta_0}\cdot n_{\mathrm{in}}(k)\right|\lesssim L^{-100}.
\end{align*}
By \eqref{eq:overck},  \eqref{bd:cz}, 
\eqref{bd:boot:wh5},  \eqref{bd:mid6} it holds that
\begin{align*}
   | d_k(s)-d_k(0)|&\lesssim |\widetilde{\mathscr{L}}( d)_k(s)|+|\tilde{\mathscr{R}}_k(s)|\lesssim  L^{1/2}\| \widetilde{\mathscr{L}}_k(d)(s)\|_{l_L^2}+ L^{1/2}\| \tilde{\mathscr{R}}_k(s)\|_{l_L^2}\\
   &\lesssim s\cdot L^{30}\|c\|_{Z}^2\|d\|_{Z}+s\cdot L^{-100}\lesssim
    sL^{-50},
\end{align*}
where the factor $s$ comes from the time integration. 

Similarly, by \eqref{ipartck} and  \eqref{bd:cz}  
\begin{align*}
     | c_k(s)-c_k(0)|&\lesssim |\widetilde{\mathscr{L}}( c)_k(s)|\lesssim  L^{1/2}\| \widetilde{\mathscr{L}}_k(c)(s)\|_{l_L^2} \lesssim s\cdot L^{30}\|c\|_{Z}^3 \lesssim
    sL^{40}.
\end{align*}

Combining the above estimates, we obtain
\begin{align*}
    \mathbf E&[(d_k(s)+d_k(0)+2c_k(s))(d_k(s)-d_k(0))]+2\mathbf E[(c_k(s)-c_k(0))d_k(0)]\\
    &\lesssim L^2\cdot   sL^{-50}+sL^{40}\cdot L^{-100} =o_{l^\infty_k}\left(\frac{Ts}{T_{\mathrm {kin}}}\right)_{L \to \infty}.
\end{align*}
 

 Then, it suffices to prove the convergence of the collision kernel.\begin{align*} 
   \mathcal K(n_{\mathrm{in}})(k)-\mathcal K_{\rm tr}(n_{\mathrm{in}})(k)
   =o_{l^\infty_k}(1),
\end{align*}
   We need to consider the truncation function $\varphi_{\leq K}$ and  $\Gamma_0$ in the domain of integration. Under the condition $\xi=\xi_1-\xi_2+\xi_3$, it is easy to see that
\begin{align*}
   (1-\varphi_{\leq K}^2(\xi_1)) &|\xi_1|^{2\beta} |n_{\mathrm{in}}(\xi_1) |\lesssim  L^{-\delta}\left||\xi_1|^{2\beta+1}\cdot n_{\mathrm{in}}(\xi_1)\right|,\\
    (1-\varphi_{\leq K}^2(\xi_1)) &|\xi_1|^{2\beta}|n_{\mathrm{in}}(\xi_2)n_{\mathrm{in}}(\xi_3)n_{\mathrm{in}}(\xi) |\\
    &\lesssim L^{-\delta }\left||\xi_1|^{2\beta+1}\cdot n_{\mathrm{in}}(\xi_2)n_{\mathrm{in}}(\xi_3)n_{\mathrm{in}}(\xi)\right|\\
    &\lesssim L^{-\delta }\left|(|\xi_2|^{2\beta+1}+|\xi_3|^{2\beta+1}+|\xi|^{2\beta+1})\cdot n_{\mathrm{in}}(\xi_2)n_{\mathrm{in}}(\xi_3)n_{\mathrm{in}}(\xi)\right|.
\end{align*}
Then, for the Schwarz function $n_{\mathrm{in}}$, it holds that
\begin{align*}
  \mathcal K(n_{\mathrm{in}})(k)-\mathcal K _{\rm tr,1}(n_{\mathrm{in}})(k)=o_{l^\infty_k}(1),
\end{align*}
where 
\begin{align}
    \mathcal K_{\rm tr,1}(\phi)(\xi):=& \int_{\xi=\xi_1-\xi_2+\xi_3,\  {\Omega}(\xi)=0}\varphi_{\leq K}^2(\xi_1)\varphi_{\leq K}^2(\xi_2)\varphi_{\leq K}^2(\xi_3)\varphi_{\leq K}^2(\xi)|\xi|^{2\beta}|\xi_1|^{2\beta}|\xi_2|^{2\beta}|\xi_3|^{2\beta}  \notag\\
   &\ \ \  \  
   \times \phi \phi_1 \phi_2 \phi_3\left(\frac{1}{\phi_1}-\frac{1}{\phi_2}+\frac{1}{\phi_3}-\frac{1}{\phi}\right)\, \dif \xi_1 \dif \xi_2 \dif \xi_3.\notag
\end{align} 

It remains to prove that
\begin{align*}
%  \lim_{L\to\infty}  \sup_k | \mathcal K_{\rm tr,1}(n_{\mathrm{in}})(k)-\mathcal K(n_{\mathrm{in}})(k) |= 
   \mathcal K_{\rm tr,1}(n_{\mathrm{in}})(k)-\mathcal K_{\rm tr}(n_{\mathrm{in}})(k)
   =o_{l^\infty_k}(1).
\end{align*}
We recall that uniformly in $k\in\mathbb{R}$,
 \begin{align*}
     \Gamma_0(k)=C_K\frac{\alpha }{\pi }\varphi_{\leq K}(k)|k|^{2\beta}\leq C\alpha L^\delta.
 \end{align*} 

By writing $n_{\rm in }(\xi)=\langle \xi\rangle^{-100}\tilde{n}_{\rm in }(\xi)$, it suffices to consider the problem on a compact domain. 
Using the relation $\xi=\xi_1-\xi_2+\xi_3$, we rewrite the collision kernel as a function $W(\xi_1,\xi_3)$. We aim to prove
\begin{align*}
      \left|\int_{\Omega(\xi_1,\xi_3)+  \sum \Gamma_0(\alpha,\xi_1,\xi_3)=0}\chi (\xi_1,\xi_3)W (\xi_1,\xi_3)\dif \xi_1\dif \xi_3-\int_{\Omega(\xi_1,\xi_3)=0}\chi (\xi_1,\xi_3)W (\xi_1,\xi_3)\dif \xi_1\dif \xi_3\right| =o_{l^\infty_k}(1).
\end{align*}
where $\chi$ is a cutoff supported on $|\xi_1-a|\leq1$, $|\xi_3-b|\leq1$.

We decompose the domain into
\begin{align*}
  E_1:=  &\{(\xi_1,\xi_3):\min\{|\xi_1|,|\xi_3|,|\xi|\}\leq L^{-\epsilon}  \},\\
   E_2:=  &\{(\xi_1,\xi_3):\min\{|\xi_1-\xi|,|\xi_3-\xi|\}\leq L^{-3\epsilon}  \},
\end{align*}
for some $\epsilon>0$ small enough.  On the set $E_1$, we use the symbol $|\xi_1|^\beta|\xi_2|^\beta|\xi_3|^\beta|\xi|^\beta$ to gain smallness, , which yields a contribution of order $o_{l^\infty_k}(1)$.  On the set $E_2\cap (E_1)^c$, we use the cancellation property $W(\xi,\xi_3)=0$. For instance, if $|\xi_1-\xi|\leq L^{-3\epsilon}$, then
\begin{align*}
     |W (\xi_1,\xi_3)|=|W (\xi_1,\xi_3)-W (\xi,\xi_3)|\leq \|W'\|_{C^1_{(\xi_1,\xi_3)}}\cdot L^{-3\epsilon} =o_{l^\infty_k}(1).
\end{align*}

We now consider this problem on the domain $(E_1\cup E_2)^c$. On this set, all variables stay away from singular configurations. 
We further decompose the domain into finitely many subdomains where the signs of $\xi_i$ are fixed. 
On each subdomain, the functions $\Omega$, $\Gamma_0$, and $W$ are $C^1$, with derivatives bounded by $L^{C\epsilon}$. Define the perturbed resonance surface
 \[
S_\alpha := \{(\xi_1,\xi_3):\Omega(\xi_1,\xi_3)+\sum\Gamma_0(\xi_1,\xi_3)=0\}.
\]
By the non-degeneracy condition $|\nabla \Omega|\gtrsim L^{-10\epsilon}$, 
the implicit function theorem implies that $S_\alpha$ can be parametrized as
\[
\xi_1=F(\xi_3,\alpha),
\]
with $F$ smooth.

 We then compare the integrals along $S_\alpha$ and $S_0$:
\[
\Big|\int_{S_\alpha} W\,\dif S - \int_{S_0} W\,\dif S\Big| =\Big|\int_{-\delta}^\delta W(F\big(y,\alpha\big),y)-W(F\big(y,0\big),y)\dif y\Big|\leq \alpha \|W\|_{C^1_{x}}\|F\|_{C_\alpha^1},
\]
Moreover, differentiating the identity $S(F(y,\alpha),y,\alpha)=0,$
we obtain
\[
|\partial_\alpha F|
\lesssim
\frac{|\partial_\alpha S|}{|\nabla S|}
\lesssim L^{10\epsilon}.
\]
Hence,
\[
\left|
\int_{S_\alpha} W\,\dif S
-
\int_{S_0} W\,\dif S
\right|
\lesssim \alpha L^{10\epsilon},
\]
which tends to $0$ as $L\to\infty$.

Since the domain is compact, a finite covering argument completes the proof.


%On the set $E^c$, by the similar argument as in the  proof of  Proposition \ref{prop:highenergy}, we bound $\Omega^{(\sigma)}$  depending on the value of $\sigma$. In case 1, we have \begin{align*}
 %  | \Omega^{(2)}|=2|(\xi_3-\xi)(\xi-\xi_1)|\geq 2 L^{-\sigma \gamma/4}.
%\end{align*}
%In case 2.1, we have
%\begin{align*}
%   | \Omega^{(\sigma)}|& \geq  \epsilon_\sigma(\xi_2+\xi_3+\xi)^{\sigma-1} \xi\geq c L^{-\sigma \gamma/4}.
%\end{align*}
%%In case 2.2, we have
%\begin{align*}
 %  | \Omega^{(\sigma)}| 
%    & \geq c(\xi_1+\xi_2+\xi_3+\xi)^{\sigma-1} (\xi_1-\xi)\geq c L^{-\sigma \gamma/4}.
%\end{align*}
%In case 2.3, with the help of the truncation functions $\varphi_{\leq K}^2(\xi_1)\varphi_{\leq K}^2(\xi_2)\varphi_{\leq K}^2(\xi_3)\varphi_{\leq K}^2(\xi)$,  we have \begin{align*}
  %  |\Omega^{(\sigma)}| 
%    &=|(\xi_1-\xi_2)(\xi_2-\xi_3)|\sigma(\sigma-1) (\xi^*)^{\sigma-2}  \geq cL^{-\gamma/2-\delta(2-\sigma)}
%\end{align*}
%\begin{align*}
%    |\Omega^{(\sigma)}(\xi_1,\xi_3)|\geq c L^{-\sigma \gamma/4}+cL^{-\gamma/2-\delta(2-\sigma)},
%\end{align*}
%which implies that by \eqref{bd:mid7}, 
%\begin{align*}
%|\overline{\Omega}(\xi_1,\xi_3)|\geq |\Omega^{{(\sigma)}}(\xi_1,\xi_3)|-C  L^{-\gamma+\delta}>0.
%\end{align*}
%So, we know that 
%\begin{align*}
 %  \mathcal K_{\rm tr}(n_{\mathrm{in}}) (\xi_1,\xi_3)=   \int_{\{\overline{\Omega}(\xi_1,\xi_3)=0\}\cap E}W(\xi_1,\xi_3)\dif \xi_1\dif \xi_3=O(L^{-\delta'}).
%\end{align*}

\ 
\noindent{\bf Acknowledgment.}  
The author would like to thank Prof. Yu Deng for valuable discussions and
helpful suggestions. His insights and comments have been very helpful in
the development of this work.

\appendix
 \renewcommand{\appendixname}{Appendix~\Alph{section}}
  \renewcommand{\theequation}{A.\arabic{equation}}

\section{The Algorithm in \cite{Vas24}}\label{sec:app}
In this section, we recall the algorithm associated with the molecule $\mathbb{M}(\mathcal{Q}^{sp})$.  We start   the following at Operation 1, where each time we remove an atom, we also remove all  bonds connected to it:
 \begin{enumerate}
     \item[(0)] If possible, remove an atom of degree 2 with a double bond. Go to (0).
     \item  Otherwise, if possible, remove a bridge (recall Definition \ref{def-molecules}). Go to (0).
     \item  Otherwise, if possible, remove an atom of degree 3 with a triple bond. Go to (1).
     \item Otherwise, if possible, remove an atom of degree 3 with a double bond. Go to (0).
     \item  Otherwise, if possible, remove an atom of degree 3 with only single bonds. Go to (1).
     \item Otherwise, if possible, remove an atom of degree 2 with only single bonds which  satisfy at least one of the following:
\begin{itemize}
    \item  The atom is connected to two atoms which are themselves connected by a double  bond.
    \item  The atom is connected to two atoms which are themselves connected by a single  bond, which becomes a bridge if the operation were to be performed.
    \item  The atom is connected to two atoms which are themselves connected by a single  bond and there is no double bond at either. 
    \item  The atom is connected to two atoms which are themselves not connected.
\end{itemize}
 Go to (1).
 \item  Otherwise, if possible, remove an atom of degree 2 with only single bonds connected  to two atoms which are themselves connected by a single bond, and there is a double  bond at precisely one of them. Go to (1).
 \item  Otherwise, if possible, remove an atom of degree 2 with only single bonds which is  connected to two atoms which are themselves connected by a triple bond. Go to (2).
 \item  Otherwise, if possible, remove an atom of degree 2 with only single bonds connected  to two atoms which are themselves connected by a single bond, and there is a double bond at both of them. Go to (1).
 \item  Otherwise, remove a sole atom (degree 0) with no edges. Repeat.
 \end{enumerate}
 

\begin{thebibliography}{99}
\bibitem[AL25]{AL25}
I. Ampatzoglou, T. L\'eger. On the ill-posedness of kinetic wave equations. Nonlinearity, 2025, 38(11): 115004.
\bibitem[BGHS18]{BGHS18}
T. Buckmaster, P. Germain, Z. Hani, J. Shatah. Effective dynamics of the nonlinear Schrödinger equation on large domains[J]. Communications on Pure and Applied Mathematics, 2018, 71(7): 1407-1460.
\bibitem[BGHS21]{BGHS21}
T. Buckmaster, P. Germain, Z. Hani and J. Shatah. Onset of the wave turbulence description
of the longtime behavior of the nonlinear Schr\"odinger equation. Invent. Math. 225 (2021),
no. 3, 787-855.
\bibitem[BLLZ23]{BLLZ23}
E. Brun, G. Li, R. Liu, Y. Zine. Global well-posedness of one-dimensional cubic fractional nonlinear
Schr\"odinger equations in negative Sobolev spaces, arXiv:2311.13370. (2023).
\bibitem[BN69]{BN69}
D. J. Benney and A. C. Newell. Random Wave Closures. In: Studies in Applied Mathematics 48.1
 (Mar. 1969), pp. 29–53. issn: 1467-9590.
 \bibitem[CG20]{CG20}
 C. Collot and P. Germain. Derivation of the homogeneous kinetic wave equation: longer time
 scales. 2020.  
 \bibitem[CG25]{CG25}
 C. Collot and P. Germain.  On the derivation of the homogeneous kinetic wave equation.
 In: Comm. Pure Appl. Math. 78.4 (2025), pp. 856–909. issn: 0010-3640,1097-0312.  
  \bibitem[CHHO13]{CHHO13}
[8] Y. Cho, H. Hajaiej, G. Hwang T. Ozawa; On the Cauchy problem of fractional Schr\"odinger equation with
Hartree type nonlinearity, Funkcial. Ekvac., 56 (2013) 193–224.
 \bibitem[CHHO14]{CHHO14}
 Y. Cho, H. Hajaiej, G. Hwang T. Ozawa; On the orbital stability of fractional Schr\"odinger equations,
Commun. Pure Appl. Anal. 13 (2014) 1267–1282.
 \bibitem[CHKL14]{CHKL14}
[10] Y. Cho, G. Hwang, S. Kwon and S. Lee; Profle decompositions of fractional Schr\"odinger equations with
angularly regular data, J. Differential Equations, 256 (2014), 3011–3037.
 \bibitem[CHKL15]{CHKL15}
 Y. Cho, G. Hwang, S. Kwon, S. Lee; Well-posedness and ill-posedness for the cubic fractional Schr\"odinger
equations, Disc. Cont. Dyn. Systems 35 7 (2015) 2863–2880.
 \bibitem[CMMT99]{CMMT99}
 D. Cai, A. J. Majda, D. W. McLaughlin, and E. G. Tabak. Spectral bifurcations in dispersive wave
 turbulence. Proceedings of the National Academy of Sciences, 96(25):14216–14221, 1999.
  \bibitem[CMMT01]{CMMT01}
  D. Cai, A. J. Majda, D. W. McLaughlin, and E. G. Tabak. Dispersive wave turbulence in one dimension.
 Physica D: Nonlinear Phenomena, 152:551–572, 2001.
\bibitem[Dav72]{Dav72}
R. Davidson. Methods in Nonlinear Plasma Theory, 1st edn. Academic Press, San Diego (1972)
\bibitem[DB23]{DB23}
 R. S. D\'u, O. B\"uhler. The impact of frequency bandwidth on a one-dimensional model for dispersive wave turbulence[J]. Journal of Nonlinear Science, 2023, 33(5): 81.
\bibitem[DH21]{DH21}
 Y. Deng and Z. Hani. On the derivation of the wave kinetic equation for NLS.  Forum Math.  Pi 9 (2021), Paper No. e6, 37. issn: 2050-5086.
 \bibitem[DH23a]{DH23a}
 Y. Deng and Z. Hani. Full derivation of the wave kinetic equation. In: Invent. Math. 233.2 (2023), pp. 543-724. issn: 0020-9910,1432-1297. 
  \bibitem[DH23b]{DH23b}
 Y. Deng and Z. Hani. Derivation of the wave kinetic equation: full range of scaling laws. 2023. 
 \bibitem[DH23c]{DH23c}
  Y. Deng and Z. Hani.  Long time justification of wave turbulence theory. 2023. arXiv: 2311.10082.
  \bibitem[DHM24]{DHM24}
 Y. Deng, Z. Hani, and X. Ma. Long time derivation of the Boltzmann equation from hard
 sphere dynamics. 2024. arXiv: 2408.07818.
\bibitem[DIP25a]{DIP25a}
Y. Deng, A. D. Ionescu, F. Pusateri. On the wave turbulence theory of 2D gravity waves, I: deterministic energy estimates. Communications on Pure and Applied Mathematics, 2025, 78(2): 211-322.
\bibitem[DIP25b]{DIP25}
Y. Deng, A. D. Ionescu, F. Pusateri. On the wave turbulence theory of 2D gravity waves, II: propagation of randomness. arXiv:2504.14304, 2025.
\bibitem[ESY03a]{ESY03}
L. Erd\"os, M. Salmhofer,  H. T. Yau. On the quantum Boltzmann equation. Journal of statistical physics, 2004, 116(1): 367-380.
\bibitem[ESY03b]{ESY03b}\
L. Erd\"os, M. Salmhofer,  H. T. Yau. Quantum diffusion of the random Schr\"odinger
 evolution in the scaling limit. In: Acta Math. 200.2 (2008), pp. 211–277. 
 \bibitem[EY00]{EY00}
 L. Erd\"os,    H. T. Yau. Linear Boltzmann equation as the weak coupling limit of a random
 Schr\"odinger equation. In: Comm. Pure Appl. Math. 53.6 (2000), pp. 667–735. 
 \bibitem[GH24]{GH24}
  R. Grande and Z. Hani. Rigorous derivation of damped-driven wave turbulence theory. 2024.
 arXiv: 2407.10711.
\bibitem[GIK20]{GIK20}
P. Germain, A. D. Ionescu, M. B. Tran. Optimal local well-posedness theory for the kinetic wave equation. Journal of Functional Analysis, 2020, 279(4): 108570.
\bibitem[GLZ25]{GLZ25}
P. Germain, J. La, K.Z. Zhang. Local well-posedness for the kinetic mmt model[J]. Communications in Mathematical Physics, 2025, 406(1): 18.
\bibitem[Has62]{Has62}
K. Hasselmann. On the non-linear energy transfer in a gravity-wave spectrum. I. General theory. In:
 J. Fluid Mech. 12 (1962), pp. 481–500. issn: 0022-1120,1469-7645.
 \bibitem[HRST22]{HRST22}
  A. Hannani, M. Rosenzweig, G. Staffilani, and M.-B. Tran, On the Wave Turbulence Theory for a stochastic KdV
 type equation-generalization for the inhomogeneous kinetic limit, 2022, arXiv:2210.17445.
\bibitem[HS15]{HS15}
Y. Hong, Y. Sire. On fractional Schr¨odinger equations in Sobolev spaces, Commun. Pure Appl. Anal. 14
(2015) 2265–2282.
 \bibitem[HSZ24]{HSZ24}
  Z. Hani, J. Shatah, and H. Zhu. Inhomogeneous turbulence for the Wick nonlinear Schr\"odinger
 equation. In: Communications on Pure and Applied Mathematics (Apr. 2024).
\bibitem[Hug83]{Hug83}
 N. M. Hugenholtz. Derivation of the Boltzmann equation for a Fermi gas. Journal of statistical physics, 1983, 32(2): 231-254.
\bibitem[JM21]{JM21}
 K. Jin, X. Ma. Cancellations of resonances and long time dynamics of cubic Schr\"odinger equation on $\mathbb{T}$. Commun. Math. Phys. 381(3), 1309–1368 (2021).
 
   \bibitem[MMT97]{MMT97}
    A. J. Majda, D. W. McLaughlin, and E. G. Tabak. A one-dimensional model for dispersive wave
 turbulence. Journal of Nonlinear Science 7.1 (Feb. 1997), pp. 9-44. issn: 1432-1467. 
\bibitem[Nor28]{Nor28}
 L. W. Nordhiem. On the kinetic method in the new statistics and application in the electron theory of
 conductivity. In: Proceedings of the Royal Society of London. Series A 119.783 (July 1928), pp. 689
698. issn: 2053-9150.
\bibitem[Pei29]{Pei29}
 R.E. Peierls. Zur kinetischen Theorie der W\"armeleitung in Kristallen. Ann. Phys. 3, 1055–1101  (1929)
\bibitem[PPW26]{PPW26}
M. Panthee, J. Patterson, Y. Wang. 
On the well-posedness of the initial value problem for the MMT model. 
arXiv:2601.07771.
 
 \bibitem[Spo77]{Spo77}
 H. Spohn. Derivation of the transport equation for electrons moving through random impurities.
 In: J. Statist. Phys. 17.6 (1977), pp. 385–412. issn: 0022-4715,1572-9613. 
\bibitem[Spo06]{Spo06}
 H. Spohn. The phonon Boltzmann equation, properties and link to weakly anharmonic lattice
 dynamics. In: J. Stat. Phys. 124.2-4 (2006), pp. 1041–1104. issn: 0022-4715,1572-9613. 

\bibitem[Spo08]{Spo08} H. Spohn. On the Boltzmann equation for weakly nonlinear wave equations. In: Boltzmann’s
 legacy. ESI Lect. Math. Phys. Eur. Math. Soc., Z\"urich, 2008, pp. 145–159. isbn: 978-3-03719-057-9.
  \bibitem[ST21]{ST21}
   G. Staffilani and M.-B. Tran, On the wave turbulence theory for stochastic and random multidimensional KdV type
 equations, 2021, arXiv:2106.09819.
 \bibitem[UU33]{UU33}
  E. A. Uehling and G. E. Uhlenbeck. Transport Phenomena in Einstein-Bose and Fermi-Dirac Gases.
 I. In: Physical Review 43.7 (Apr. 1933), pp. 552–561.issn: 0031-899X.
\bibitem[Vas25]{Vas24}
   K. D. Vassilev. One-Dimensional Wave Kinetic Theory. Communications in Mathematical Physics, 2025, 406(12): 293.
\bibitem[VW26]{VW26}
K. D. Vassilev, and B. Wu. Rigorous Derivation of the Wave Kinetic Equation for full $\beta $-FPUT System. arXiv preprint arXiv:2605.19308 (2026).
\bibitem[Wu25]{Wu25} B. Wu. Rigorous Derivation of the Wave Kinetic Equation for $\beta $-FPUT System. arXiv:2506.02948, 2025.

\bibitem[ZDP04]{ZDP04}
V. Zakharov, F. Dias, A. Pushkarev. One-dimensional wave turbulence. Physics Reports, 2004, 398(1): 1-65.
\bibitem[ZGPD01]{ZGPD01}
V. Zakharov, P. Guyenne, A. Pushkarev, and F. Dias. Wave turbulence in one-dimensional models.
 Physica D: Nonlinear Phenomena, 152:573–619, 2001.
\bibitem[ZS67]{ZS67}
 G. M. Zaslavski and R. Z. Sagdeev. Limits of Statistical Description of a Nonlinear Wave Field. In: Soviet Physics JETP 25.4 (Oct. 1967), pp. 718–724.
\end{thebibliography}

\bigskip 




 \end{document}
